\input{preamble}

% OK, start here.
%
\begin{document}

\title{Morphismes de schémas}


\maketitle

\phantomsection
\label{section-phantom}

\tableofcontents

\section{Introduction}
\label{section-introduction}

\noindent
Dans ce chapitre, nous introduisons certains types de morphismes de schémas.
Une référence fondamentale est \cite{EGA}.




















\section{Immersions fermées}
\label{section-closed-immersions}

\noindent
Dans cette section, nous précisons certains des résultats obtenus précédemment sur les
immersions fermées de schémas. Rappelons qu'un morphisme de schémas $i : Z \to X$
est par définition une immersion fermée si (a) $i$ induit un homéomorphisme sur
une partie fermée de $X$, (b) $i^\sharp : \mathcal{O}_X \to i_*\mathcal{O}_Z$
est surjectif et (c) le noyau de $i^\sharp$ est engendré localement par des
sections, voir Schémas, Définitions \ref{schemes-definition-immersion}
et \ref{schemes-definition-closed-immersion-locally-ringed-spaces}. Il s'avère
que, puisque $Z$ et $X$ sont des schémas, il existe de nombreuses
façons de caractériser une immersion fermée.

\begin{lemma}
\label{lemma-closed-immersion}
Soit $i : Z \to X$ un morphisme de schémas.
Les conditions suivantes sont équivalentes :
\begin{enumerate}
\item Le morphisme $i$ est une immersion fermée.
\item Pour tout ouvert affine $\Spec(R) = U \subset X$,
il existe un idéal $I \subset R$ tel que
$i^{-1}(U) = \Spec(R/I)$ en tant que schémas sur $U = \Spec(R)$.
\item Il existe un recouvrement ouvert affine $X = \bigcup_{j \in J} U_j$,
$U_j = \Spec(R_j)$ et, pour tout $j \in J$, il existe
un idéal $I_j \subset R_j$ tel que
$i^{-1}(U_j) = \Spec(R_j/I_j)$ en tant que schémas sur $U_j = \Spec(R_j)$.
\item Le morphisme $i$ induit un homéomorphisme de $Z$ sur une partie fermée
de $X$ et $i^\sharp : \mathcal{O}_X \to i_*\mathcal{O}_Z$ est surjectif.
\item Le morphisme $i$ induit un homéomorphisme de $Z$ sur une partie fermée
de $X$, l'application $i^\sharp : \mathcal{O}_X \to i_*\mathcal{O}_Z$ est surjective,
et le noyau $\Ker(i^\sharp)\subset \mathcal{O}_X$ est un faisceau quasi-cohérent
d'idéaux.
\item Le morphisme $i$ induit un homéomorphisme de $Z$ sur une partie fermée
de $X$, l'application $i^\sharp : \mathcal{O}_X \to i_*\mathcal{O}_Z$ est surjective,
et le noyau $\Ker(i^\sharp)\subset \mathcal{O}_X$ est un
faisceau d'idéaux engendré localement par des sections.
\end{enumerate}
\end{lemma}

\begin{proof}
La condition (6) est notre définition d'une immersion fermée, voir Schémas,
Définitions \ref{schemes-definition-closed-immersion-locally-ringed-spaces}
et \ref{schemes-definition-immersion}.
Ainsi (6) $\Leftrightarrow$ (1). Nous avons (1) $\Rightarrow$ (2) d'après
Schémas, Lemme \ref{schemes-lemma-closed-subspace-scheme}.
Il est immédiat que (2) $\Rightarrow$ (3).

\medskip\noindent
Supposons (3). Chacun des morphismes
$\Spec(R_j/I_j) \to \Spec(R_j)$ est
une immersion fermée, voir
Schémas, Exemple \ref{schemes-example-closed-immersion-affines}.
Ainsi $i^{-1}(U_j) \to U_j$ est un homéomorphisme sur son image
et $i^\sharp|_{U_j}$ est surjectif. Ainsi $i$ est un homéomorphisme
sur son image et $i^\sharp$ est surjectif, puisque cela peut se
vérifier localement. Nous concluons que (3) $\Rightarrow$ (4).

\medskip\noindent
L'implication (4) $\Rightarrow$ (1) est
Schémas, Lemme \ref{schemes-lemma-characterize-closed-immersions}.
L'implication (5) $\Rightarrow$ (6) est immédiate.
Et l'implication (6) $\Rightarrow$ (5) découle
de Schémas, Lemme \ref{schemes-lemma-closed-subspace-scheme}.
\end{proof}

\begin{lemma}
\label{lemma-closed-immersion-ideals}
Soit $X$ un schéma. Soient $i : Z \to X$ et
$i' : Z' \to X$ des immersions fermées et considérons
les faisceaux d'idéaux $\mathcal{I} = \Ker(i^\sharp)$ et $\mathcal{I}' = \Ker((i')^\sharp)$
de $\mathcal{O}_X$.
\begin{enumerate}
\item Le morphisme $i : Z \to X$ se factorise en $Z \to Z' \to X$ pour
certains $a : Z \to Z'$ si et seulement si $\mathcal{I}' \subset \mathcal{I}$. Si cela se produit,
alors $a$ est une immersion fermée.
\item Nous avons $Z \cong Z'$ sur $X$ si et
seulement si $\mathcal{I} = \mathcal{I}'$.
\end{enumerate}
\end{lemma}

\begin{proof}
Cela résulte de notre discussion sur les sous-espaces
fermés dans Schémas, section \ref{schemes-section-closed-immersion}
en particulier
Schémas, Lemmes \ref{schemes-lemma-closed-immersion}
et \ref{schemes-lemma-characterize-closed-subspace}. Cela découle également
de manière simple de la caractérisation (3)
dans le lemme \ref{lemma-closed-immersion} ci-dessus.
\end{proof}

\begin{lemma}
\label{lemma-closed-immersion-bijection-ideals}
Soit $X$ un schéma.
Soit $\mathcal{I} \subset \mathcal{O}_X$ un faisceau d'idéaux. Les conditions suivantes
sont équivalentes :
\begin{enumerate}
\item $\mathcal{I}$ est engendré localement par des
sections comme un faisceau de $\mathcal{O}_X$-modules,
\item $\mathcal{I}$ est quasi-cohérent comme
un faisceau de $\mathcal{O}_X$-modules, et
\item existe une immersion fermée $i : Z \to X$ de schémas dont le
faisceau d'idéaux correspondant $\Ker(i^\sharp)$ est égal à $\mathcal{I}$.
\end{enumerate}
\end{lemma}

\begin{proof}
L'équivalence de (1) et (2) découle immédiatement
d'après Schémas, Lemme \ref{schemes-lemma-closed-subspace-scheme}.
Si (1) est vrai, alors il y a un sous-espace fermé
$i : Z \to X$ avec $\mathcal{I} = \Ker(i^\sharp)$
par Schémas, Définition \ref{schemes-definition-closed-subspace} et Exemple \ref{schemes-example-closed-subspace}.
Par Schémas, Lemme \ref{schemes-lemma-closed-subspace-scheme}
c'est une immersion fermée de schémas et (3) tient.
Inversement, si (3) est satisfaite, alors
(2) est vrai par Schémas, Lemme
\ref{schemes-lemma-closed-subspace-scheme} (ce qui s'applique car une immersion
fermée de schémas est a fortiori une immersion fermée
d'espaces localement annelés).
\end{proof}

\begin{lemma}
\label{lemma-base-change-closed-immersion}
Le changement de base d'une immersion fermée est une immersion fermée.
\end{lemma}

\begin{proof}
Voir Schémas, Lemme \ref{schemes-lemma-base-change-immersion}.
\end{proof}

\begin{lemma}
\label{lemma-composition-closed-immersion}
Une composition d'immersions fermées est une immersion fermée.
\end{lemma}

\begin{proof}
Nous avons vu cela
dans Schémas, Lemme \ref{schemes-lemma-composition-immersion}, mais voici une
autre démonstration.
En effet, cela découle de la caractérisation (3) des immersions
fermées dans le lemme \ref{lemma-closed-immersion}. En effet, si
$I \subset R$ est un idéal, et $\overline{J} \subset R/I$ est un idéal,
alors $\overline{J} = J/I$ pour un idéal $J \subset R$ qui contient
$I$ et $(R/I)/\overline{J} = R/J$.
\end{proof}

\begin{lemma}
\label{lemma-closed-immersion-quasi-compact}
Une immersion fermée est quasi-compacte.
\end{lemma}

\begin{proof}
Ce lemme fait double
emploi avec Schémas, Lemme \ref{schemes-lemma-closed-immersion-quasi-compact}.
\end{proof}

\begin{lemma}
\label{lemma-closed-immersion-separated}
Une immersion fermée est séparée.
\end{lemma}

\begin{proof}
Ce lemme est un cas particulier
d'après Schémas, Lemme \ref{schemes-lemma-immersions-monomorphisms}.
\end{proof}





\section{Immersions}
\label{section-immersions}

\noindent
Dans cette section, nous collectons quelques faits sur les immersions.

\begin{lemma}
\label{lemma-immersion-permanence}
Soit $Z \to Y \to X$ des morphismes de schémas.
\begin{enumerate}
\item Si $Z \to X$ est une immersion, alors $Z \to Y$ est une immersion.
\item Si $Z \to X$ est une immersion quasi-compacte et $Y \to X$ est
quasi-séparé, alors $Z \to Y$ est une immersion quasi-compacte.
\item Si $Z \to X$ est une immersion fermée et que $Y \to X$ est séparé, alors
$Z \to Y$ est une immersion fermée.
\end{enumerate}
\end{lemma}

\begin{proof}
Dans chaque cas, la démonstration consiste à considérer le diagramme commutatif
$$
\xymatrix{
Z \ar[r] \ar[rd] & Y \times_X Z \ar[r] \ar[d] & Z \ar[d] \\
& Y \ar[r] & X
}
$$
où la composition des flèches horizontales supérieures est
l’identité. Démontrons (1). La première flèche horizontale
est une coupe de $Y \times_X Z \to Z$, d'où
une immersion d'après Schémas, Lemme \ref{schemes-lemma-section-immersion}.
La flèche $Y \times_X Z \to Y$ est un changement de base de $Z \to X$
donc une immersion (Schémas, Lemme \ref{schemes-lemma-base-change-immersion}). Enfin,
une composition d'immersions est une immersion
(Schémas, Lemme \ref{schemes-lemma-composition-immersion}). Cela prouve (1). Les deux autres résultats
se prouvent exactement de la même manière.
\end{proof}

\begin{lemma}
\label{lemma-factor-quasi-compact-immersion}
Soit $h : Z \to X$ une immersion. Si
$h$ est quasi-compact, alors on
peut factoriser $h = i \circ j$ avec $j : Z \to \overline{Z}$
une immersion ouverte et $i : \overline{Z} \to X$ une immersion fermée.
\end{lemma}

\begin{proof}
Notons que $h$ est quasi-compact
et quasi-séparé (voir Schémas, Lemme \ref{schemes-lemma-immersions-monomorphisms}).
Ainsi $h_*\mathcal{O}_Z$ est un faisceau quasi-cohérent
de $\mathcal{O}_X$-modules d'après Schémas, Lemme
\ref{schemes-lemma-push-forward-quasi-coherent}.
Cela implique que $\mathcal{I} = \Ker(\mathcal{O}_X \to h_*\mathcal{O}_Z)$ est un
faisceau quasi-cohérent d'idéaux,
voir Schémas, section \ref{schemes-section-quasi-coherent}.
Soit $\overline{Z} \subset X$ le sous-schéma fermé correspondant
à $\mathcal{I}$, voir Lemme \ref{lemma-closed-immersion-bijection-ideals}. Par Schémas,
Lemme \ref{schemes-lemma-characterize-closed-subspace} le morphisme $h$
se factorise en
$h = i \circ j$ où $i : \overline{Z} \to X$ est le morphisme d'inclusion. Pour voir que $j$
est une immersion ouverte, choisissez un sous-schéma ouvert
$U \subset X$ tel que $h$ induit une immersion fermée de
$Z$ dans $U$. Il est alors clair que $\mathcal{I}|_U$
est le faisceau d'idéaux correspondant à l'immersion fermée
$Z \to U$. Nous voyons donc que $Z = \overline{Z} \cap U$.
\end{proof}

\begin{lemma}
\label{lemma-factor-reduced-immersion}
Soit $h : Z \to X$ une immersion. Si
$Z$ est réduit, alors on
peut factoriser $h = i \circ j$ avec $j : Z \to \overline{Z}$ une
immersion ouverte et $i : \overline{Z} \to X$ une immersion fermée.
\end{lemma}

\begin{proof}
Soit $\overline{Z} \subset X$ l'adhérence de $h(Z)$ munie de la structure de sous-schéma
fermé induit réduit, voir
Schémas, Définition \ref{schemes-definition-reduced-induced-scheme}. Par Schémas, Lemme \ref{schemes-lemma-map-into-reduction}
le morphisme $h$ se factorise en $h = i \circ j$
avec $i : \overline{Z} \to X$
le morphisme d'inclusion et $j : Z \to \overline{Z}$. De la définition d'une
immersion on voit qu'il existe un sous-schéma ouvert
$U \subset X$ tel que $h$ se factorise par une
immersion fermée dans $U$. Ainsi $\overline{Z} \cap U$
et $h(Z)$ sont des sous-schémas fermés réduits
de $U$ avec le même ensemble fermé sous-jacent.
Ainsi par l'unicité dans Schémas, Lemme \ref{schemes-lemma-reduced-closed-subscheme}
on voit que $h(Z) \cong \overline{Z} \cap U$. Ainsi $j$ induit
un isomorphisme de $Z$ avec $\overline{Z} \cap U$. En d’autres termes,
$j$ est une immersion ouverte.
\end{proof}



\begin{example}
\label{example-thibaut}
Voici un exemple d'immersion qui n'est pas une composition d'une
immersion ouverte suivie d'une immersion
fermée. Soit
$k$ un corps. Soit $X = \Spec(k[x_1, x_2, x_3, \ldots])$.
Soit $U = \bigcup_{n = 1}^{\infty} D(x_n)$. Alors $U \to X$ est
une immersion ouverte. Considérons
les idéaux
$$
I_n =
(x_1^n, x_2^n, \ldots, x_{n - 1}^n, x_n - 1, x_{n + 1}, x_{n + 2}, \ldots)
\subset
k[x_1, x_2, x_3, \ldots][1/x_n].
$$
Notez que $I_n k[x_1, x_2, x_3, \ldots][1/x_nx_m] = (1)$ pour tout $m \not = n$. D'où
les idéaux quasi-cohérents $\widetilde I_n$
sur $D(x_n)$ s'accordent sur $D(x_nx_m)$, à
savoir $\widetilde I_n|_{D(x_nx_m)} = \mathcal{O}_{D(x_n x_m)}$ si $n \not = m$. Ces idéaux se
collent donc à un faisceau quasi-cohérent d'idéaux $\mathcal{I} \subset \mathcal{O}_U$.
Soit $Z \subset U$ le sous-schéma
fermé correspondant à $\mathcal{I}$. Ainsi
$Z \to X$ est une immersion.

\medskip\noindent
Nous affirmons que nous ne pouvons pas
factoriser $Z \to X$ en $Z \to \overline{Z} \to X$, où $\overline{Z} \to X$
est une immersion fermée et $Z \to \overline{Z}$ une immersion ouverte. En effet, $\overline{Z}$
devrait être défini par un idéal $I \subset k[x_1, x_2, x_3, \ldots]$ tel que
$I_n = I k[x_1, x_2, x_3, \ldots][1/x_n]$. Mais le seul élément $f \in k[x_1, x_2, x_3, \ldots]$
qui se retrouve dans tous les $I_n$ est $0$
! Par conséquent, $I$ n'existe pas.
\end{example}

\begin{lemma}
\label{lemma-check-immersion}
Soit $f : Y \to X$ un morphisme de schémas. Si pour tout $y \in Y$
il existe un sous-schéma ouvert $f(y) \in U \subset X$ tel que
$f|_{f^{-1}(U)} : f^{-1}(U) \to U$ est une immersion, alors $f$ est une
immersion.
\end{lemma}

\begin{proof}
Cette affirmation découle facilement de la discussion sur les sous-schémas
fermés à la fin de Schémas, section \ref{schemes-section-immersions} mais nous
donnerons également une démonstration
détaillée. Soit $Z \subset X$ l'adhérence de $f(Y)$. Puisque la formation
de l'adhérence commute à la restriction aux ouverts, l'hypothèse montre que
$f(Y) \subset Z$ est ouvert. Par conséquent,
$Z' = Z \setminus f(Y)$ est fermé. Ainsi $X' = X \setminus Z'$ est un sous-schéma
ouvert de $X$ et $f$ se factorise en
$f : Y \to X'$ suivi de l'inclusion. Si $y \in Y$ et $U \subset X$
sont comme dans l'énoncé du lemme, alors $U' = X' \cap U$
est un voisinage ouvert de $f'(y)$ tel que $(f')^{-1}(U') \to U'$
est une immersion (Lemme \ref{lemma-immersion-permanence}) à image fermée.
Il s'agit donc d'une immersion fermée,
voir Schémas, Lemme \ref{schemes-lemma-immersion-when-closed}. Puisqu'être une
immersion fermée est locale sur la cible (par exemple
par le lemme \ref{lemma-closed-immersion}), nous concluons que
$f'$ est une immersion fermée comme voulu.
\end{proof}








\section{Immersions fermées et faisceaux quasi-cohérents}
\label{section-closed-immersions-quasi-coherent}

\noindent
Le lemme suivant fait finalement pour les faisceaux quasi-cohérents
sur les schémas ce que Modules, Lemme \ref{modules-lemma-i-star-exact} fait pour les faisceaux
abéliens. Voir également
la discussion dans Modules, section \ref{modules-section-closed-immersion}.

\begin{lemma}
\label{lemma-i-star-equivalence}
Soit $i : Z \to X$ une immersion fermée de schémas.
Soit $\mathcal{I} \subset \mathcal{O}_X$ le faisceau quasi-cohérent d'idéaux définissant
$Z$. Le foncteur
$$
i_* :
\QCoh(\mathcal{O}_Z)
\longrightarrow
\QCoh(\mathcal{O}_X)
$$
est exact, pleinement fidèle, avec image essentielle des $\mathcal{O}_X$-modules
quasi-cohérents $\mathcal{G}$ tels que $\mathcal{I}\mathcal{G} = 0$.
\end{lemma}

\begin{proof}
Une immersion fermée est quasi-compacte et
séparée, voir Lemmes \ref{lemma-closed-immersion-quasi-compact} et \ref{lemma-closed-immersion-separated}.
D'où Schémas, Lemme \ref{schemes-lemma-push-forward-quasi-coherent}
s'applique et l'image directe d'un faisceau quasi-cohérent
sur $Z$ est bien un
faisceau quasi-cohérent sur $X$.

\medskip\noindent
Par Modules, Lemme \ref{modules-lemma-i-star-equivalence} le foncteur $i_*$
est pleinement fidèle.

\medskip\noindent
Passons maintenant à la description de l'image essentielle
du foncteur $i_*$. On a $\mathcal{I}(i_*\mathcal{F}) = 0$
pour tout module $\mathcal{O}_Z$ quasi-cohérent, par exemple
par Modules, Lemme \ref{modules-lemma-i-star-equivalence}. Supposons
ensuite que $\mathcal{G}$ soit
un module $\mathcal{O}_X$ quasi-cohérent tel que
$\mathcal{I}\mathcal{G} = 0$. Il suffit de montrer que l'application canonique
$$
\mathcal{G} \longrightarrow i_* i^*\mathcal{G}
$$
est un isomorphisme\footnote{Cela a été prouvé dans une situation
plus générale dans la démonstration de Modules, Lemme \ref{modules-lemma-i-star-equivalence}.}. Dans
le cas de schémas et modules quasi-cohérents, en travaillant localement sur
les affines de $X$ et en utilisant Lemme \ref{lemma-closed-immersion}
et Schémas, Lemme \ref{schemes-lemma-widetilde-pullback} il suffit de prouver l'énoncé
algébrique suivant : Étant donné un anneau $R$, un idéal $I$
et un $R$-module $N$ tel que $IN = 0$ l'application canonique
$$
N \longrightarrow N \otimes_R R/I,\quad
n \longmapsto n \otimes 1
$$
est un isomorphisme de $R$-modules. La démonstration de ce fait algébrique simple est omise.
\end{proof}

\noindent
Soit $i : Z \to X$ une immersion fermée. En raison du lemme ci-dessus, nous désignons
souvent, par abus de notation, $\mathcal{F}$ le faisceau $i_*\mathcal{F}$ sur $X$.

\begin{lemma}
\label{lemma-largest-quasi-coherent-subsheaf}
Soit $X$ un schéma. Soit $\mathcal{F}$ un module
$\mathcal{O}_X$ quasi-cohérent. Soit $\mathcal{G} \subset \mathcal{F}$ un $\mathcal{O}_X$-sous-module.
Il existe un unique $\mathcal{O}_X$-sous-module
$\mathcal{G}' \subset \mathcal{G}$ quasi-cohérent avec la propriété
suivante : Pour tout $\mathcal{O}_X$-module $\mathcal{H}$ quasi-cohérent
l'application
$$
\Hom_{\mathcal{O}_X}(\mathcal{H}, \mathcal{G}')
\longrightarrow
\Hom_{\mathcal{O}_X}(\mathcal{H}, \mathcal{G})
$$
est bijective. En particulier, $\mathcal{G}'$ est le plus grand $\mathcal{O}_X$-sous-module
quasi-cohérent de $\mathcal{F}$ contenu dans $\mathcal{G}$.
\end{lemma}

\begin{proof}
Soit $\mathcal{G}_a$, $a \in A$ l'ensemble des $\mathcal{O}_X$-sous-modules
quasi-cohérents contenus dans $\mathcal{G}$.
Alors l'image $\mathcal{G}'$ de
$$
\bigoplus\nolimits_{a \in A} \mathcal{G}_a \longrightarrow \mathcal{F}
$$
est quasi-cohérente comme l'image d'une application de faisceaux quasi-cohérents
sur $X$ est quasi-cohérente et puisqu'une somme directe
de faisceaux quasi-cohérents
est quasi-cohérente, voir Schémas, section \ref{schemes-section-quasi-coherent}.
Le module $\mathcal{G}'$ est contenu dans $\mathcal{G}$. Il s'agit donc
du plus grand module $\mathcal{O}_X$ quasi-cohérent contenu dans $\mathcal{G}$.

\medskip\noindent
Pour prouver la formule, soit $\mathcal{H}$ un module
$\mathcal{O}_X$ quasi-cohérent et soit $\alpha : \mathcal{H} \to \mathcal{G}$ une application
de module $\mathcal{O}_X$. L'image de la composition $\mathcal{H} \to \mathcal{G} \to \mathcal{F}$
est quasi-cohérente comme l'image d'une application
de faisceaux quasi-cohérents. Il est donc contenu dans $\mathcal{G}'$.
Ainsi, $\alpha$ se factorise par $\mathcal{G}'$ comme
souhaité.
\end{proof}

\begin{lemma}
\label{lemma-i-upper-shriek}
Soit $i : Z \to X$ une immersion fermée de schémas.
Il existe un foncteur\footnote{Il s'agit probablement
d'une notation non standard.}
$i^! : \QCoh(\mathcal{O}_X) \to \QCoh(\mathcal{O}_Z)$ qui est un adjoint droit de
$i_*$. (Comparer Modules, Lemme \ref{modules-lemma-i-star-right-adjoint}.)
\end{lemma}

\begin{proof}
Étant donné le $\mathcal{O}_X$-module $\mathcal{G}$ quasi-cohérent,
nous considérons le sous-faisceau $\mathcal{H}_Z(\mathcal{G})$ de $\mathcal{G}$ de sections
locales annihilées par
$\mathcal{I}$. Par le Lemme \ref{lemma-largest-quasi-coherent-subsheaf} il existe
un $\mathcal{O}_X$-sous-module quasi-cohérent canonique le plus grand
$\mathcal{H}_Z(\mathcal{G})'$. Par construction, nous avons
$$
\Hom_{\mathcal{O}_X}(i_*\mathcal{F}, \mathcal{H}_Z(\mathcal{G})')
=
\Hom_{\mathcal{O}_X}(i_*\mathcal{F}, \mathcal{G})
$$
pour tout $\mathcal{O}_Z$-module quasi-cohérent $\mathcal{F}$.
Nous pouvons donc définir $i^!\mathcal{G} = i^*(\mathcal{H}_Z(\mathcal{G})')$. Détails
omis.
\end{proof}

\noindent
À l'aide des $1$-to-$1$ correspondant entre faisceaux
quasi-cohérents d'idéaux et sous-schémas
fermés (voir Lemme \ref{lemma-closed-immersion-bijection-ideals}) on peut définir
des intersections schématiques et des unions de sous-schémas
fermés.

\begin{definition}
\label{definition-scheme-theoretic-intersection-union}
Soit $X$ un schéma. Soient $Z, Y \subset X$ des sous-schémas fermés
correspondant à des faisceaux d'idéaux quasi-cohérents
$\mathcal{I}, \mathcal{J} \subset \mathcal{O}_X$. L'{\it intersection
schématique} de $Z$ et $Y$ est le
sous-schéma fermé de $X$ découpé par $\mathcal{I} + \mathcal{J}$. La {\it
réunion schématique} de $Z$
et $Y$ est le sous-schéma fermé de $X$
découpé par $\mathcal{I} \cap \mathcal{J}$.
\end{definition}

\begin{lemma}
\label{lemma-scheme-theoretic-intersection}
Soit $X$ un schéma. Soient $Z, Y \subset X$ des sous-schémas fermés.
Soit $Z \cap Y$ l'intersection schématique de $Z$ et $Y$.
Alors $Z \cap Y \to Z$ et $Z \cap Y \to Y$ sont des immersions fermées
et
$$
\xymatrix{
Z \cap Y \ar[r] \ar[d] & Z \ar[d] \\
Y \ar[r] & X
}
$$
est un diagramme cartésien de schémas, c'est-à-dire $Z \cap Y = Z \times_X Y$.
\end{lemma}

\begin{proof}
Les morphismes $Z \cap Y \to Z$ et $Z \cap Y \to Y$ sont des immersions fermées
par le lemme \ref{lemma-closed-immersion-ideals}. Soit $U = \Spec(A)$
un ouvert affine de $X$ et soit $Z \cap U$ et $Y \cap U$
correspondre aux idéaux $I \subset A$ et $J \subset A$.
Alors $Z \cap Y \cap U$ correspond à $I + J \subset A$. Puisque
$A/I \otimes_A A/J = A/(I + J)$ nous voyons que le schéma est cartésien
par notre description des produits fibrés
de schémas dans Schémas, section \ref{schemes-section-fibre-products}.
\end{proof}

\begin{lemma}
\label{lemma-scheme-theoretic-union}
Soit $S$ un schéma. Soient $X, Y \subset S$ des sous-schémas fermés.
Soit $X \cup Y$ la réunion schématique de $X$ et $Y$.
Soit $X \cap Y$ l'intersection schématique de $X$ et $Y$.
Alors $X \to X \cup Y$ et $Y \to X \cup Y$ sont des immersions fermées, il existe une
suite exacte courte
$$
0 \to \mathcal{O}_{X \cup Y} \to \mathcal{O}_X \times \mathcal{O}_Y
\to \mathcal{O}_{X \cap Y} \to 0
$$
de $\mathcal{O}_S$-modules, et le diagramme
$$
\xymatrix{
X \cap Y \ar[r] \ar[d] & X \ar[d] \\
Y \ar[r] & X \cup Y
}
$$
est cocartésien dans la catégorie des schémas,
c'est-à-dire $X \cup Y = X \amalg_{X \cap Y} Y$.
\end{lemma}

\begin{proof}
Les morphismes $X \to X \cup Y$ et $Y \to X \cup Y$ sont des immersions fermées par
le lemme \ref{lemma-closed-immersion-ideals}. Dans la suite exacte courte nous utilisons l'équivalence
du lemme \ref{lemma-i-star-equivalence} pour penser les modules quasi-cohérents
sur sous-schémas fermés de $S$ comme des modules quasi-cohérents
sur $S$. Pour la première application de la séquence
on utilise les applications canoniques $\mathcal{O}_{X \cup Y} \to \mathcal{O}_X$
et $\mathcal{O}_{X \cup Y} \to \mathcal{O}_Y$ et pour la deuxième
application on utilise l'application canonique
$\mathcal{O}_X \to \mathcal{O}_{X \cap Y}$ et le négatif de l'application
canonique $\mathcal{O}_Y \to \mathcal{O}_{X \cap Y}$. Ensuite,
pour vérifier l'exactitude, nous pouvons travailler
localement sur les affines.
Soit $U = \Spec(A)$ un ouvert affine de $S$ et soit $X \cap U$ et
$Y \cap U$ correspondre aux idéaux $I \subset A$ et $J \subset A$.
Alors $(X \cup Y) \cap U$ correspond à $I \cap J \subset A$
et $X \cap Y \cap U$ correspond à $I + J \subset A$. Ainsi l'exactitude
découle de l'exactitude de
$$
0 \to A/I \cap J \to A/I \times A/J \to A/(I + J) \to 0
$$
Pour montrer que le diagramme est cocartésien, supposons qu'on nous donne
un schéma $T$ et des morphismes de schémas $f : X \to T$, $g : Y \to T$
qui s'accordent comme morphismes $X \cap Y \to T$. Il s'agit de montrer
qu'il existe un morphisme unique $h : X \cup Y \to T$ en accord
avec $f$ et $g$. Pour construire $h$ on peut
travailler localement sur les affines de $X \cup Y$, voir
Schémas, section \ref{schemes-section-glueing-schemes}. Si $s \in X$, $s \not \in Y$, alors
$X \to X \cup Y$ est un isomorphisme dans un voisinage de $s$
et il est clair comment construire $h$. De même pour
$s \in Y$, $s \not \in X$. Pour $s \in X \cap Y$ nous pouvons
choisir un $V = \Spec(B) \subset T$ ouvert affine contenant $f(s) = g(s)$.
On peut alors choisir un $U = \Spec(A) \subset S$ ouvert affine
contenant $s$ de telle sorte que $f(X \cap U)$ et $g(Y \cap U)$ soient
contenus dans $V$. Les morphismes $f|_{X \cap U}$ et $g|_{Y \cap V}$
dans $V$ correspondent à des homomorphismes d'anneaux
$$
B \to A/I
\quad\text{et}\quad
B \to A/J
$$
dont les composés vers $A/(I + J)$ coïncident. Par la suite exacte courte affichée
ci-dessus, il y a un relèvement unique de ces homomorphismes d'anneaux vers
un homomorphisme d'anneaux $B \to A/I \cap J$ comme souhaité.
\end{proof}









\section{Supports des modules}
\label{section-support}

\noindent
Dans cette section nous collectons quelques résultats élémentaires
sur les supports de modules quasi-cohérents
sur des schémas. Rappelons que le support d'un faisceau
de modules a été défini dans Modules, section \ref{modules-section-support}.
En revanche, le support d'un module a été défini
dans Algèbre, section \ref{algebra-section-support}. Ceux-ci
correspondent.

\begin{lemma}
\label{lemma-support-affine-open}
Soit $X$ un schéma. Soit $\mathcal{F}$ un faisceau quasi-cohérent sur
$X$. Soit $\Spec(A) = U \subset X$ un ouvert affine et définissons
$M = \Gamma(U, \mathcal{F})$. Soit $x \in U$,
et soit $\mathfrak p \subset A$ l'idéal premier correspondant. Les conditions suivantes
sont équivalentes
\begin{enumerate}
\item $\mathfrak p$ est dans le support de $M$, et
\item $x$ est dans le support de $\mathcal{F}$.
\end{enumerate}
\end{lemma}

\begin{proof}
Cela résulte de l'égalité $\mathcal{F}_x = M_{\mathfrak p}$, voir Schémas,
Lemme \ref{schemes-lemma-spec-sheaves} et les
définitions.
\end{proof}

\begin{lemma}
\label{lemma-support-closed-specialization}
Soit $X$ un schéma.
Soit $\mathcal{F}$ un faisceau quasi-cohérent sur $X$.
Le support de $\mathcal{F}$ est clôturé sous spécialisation.
\end{lemma}

\begin{proof}
Si $x' \leadsto x$ est une spécialisation et $\mathcal{F}_x = 0$
alors $\mathcal{F}_{x'}$ vaut zéro, car $\mathcal{F}_{x'}$ est une localisation
du module $\mathcal{F}_x$. Le complément
de $\text{Supp}(\mathcal{F})$ est donc fermé sous génération.
\end{proof}

\noindent
Pour les modules de type fini quasi-cohérents le support est fermé,
vérifiable sur fibres, et commute avec changement de base.

\begin{lemma}
\label{lemma-support-finite-type}
Soit $\mathcal{F}$ un module de type fini quasi-cohérent sur un schéma
$X$. Ensuite,
\begin{enumerate}
\item Le support de $\mathcal{F}$ est fermé.
\item Pour $x \in X$ nous avons
$$
x \in \text{Supp}(\mathcal{F})
\Leftrightarrow
\mathcal{F}_x \not = 0
\Leftrightarrow
\mathcal{F}_x \otimes_{\mathcal{O}_{X, x}} \kappa(x) \not = 0.
$$
\item Pour tout morphisme de schémas $f : Y \to X$
le pullback $f^*\mathcal{F}$ est également
de type fini et nous avons $\text{Supp}(f^*\mathcal{F}) = f^{-1}(\text{Supp}(\mathcal{F}))$.
\end{enumerate}
\end{lemma}

\begin{proof}
La partie (1) est une reformulation
des modules, Lemme \ref{modules-lemma-support-finite-type-closed}. Vous pouvez
également combiner
Lemme \ref{lemma-support-affine-open}, Propriétés,
Lemme \ref{properties-lemma-finite-type-module} et Algèbre, Lemme \ref{algebra-lemma-support-closed}
pour
voir cela. La première équivalence
dans (2) est la définition du support, et la deuxième
équivalence découle du lemme de Nakayama,
voir Algèbre,
Lemme \ref{algebra-lemma-NAK}. Soit $f : Y \to X$
un morphisme de schémas. Notez que
$f^*\mathcal{F}$ est de type fini par
Modules, Lemme \ref{modules-lemma-pullback-finite-type}. Pour l'assertion finale,
soit $y \in Y$ avec l'image $x \in X$. Rappelons
que
$$
(f^*\mathcal{F})_y =
\mathcal{F}_x \otimes_{\mathcal{O}_{X, x}} \mathcal{O}_{Y, y},
$$
voir
Faisceaux, Lemme \ref{sheaves-lemma-stalk-pullback-modules}. Donc $(f^*\mathcal{F})_y \otimes \kappa(y)$
est différent de zéro si et seulement
si $\mathcal{F}_x \otimes \kappa(x)$ est différent de zéro. Par
(2) cela implique $x \in \text{Supp}(\mathcal{F})$ si et seulement
si $y \in \text{Supp}(f^*\mathcal{F})$, qui est le contenu de l'assertion
(3).
\end{proof}

\begin{lemma}
\label{lemma-scheme-theoretic-support}
Soit $\mathcal{F}$ un module de type fini quasi-cohérent
sur un schéma $X$. Il existe un plus petit
sous-schéma fermé $i : Z \to X$ tel qu'il existe
un $\mathcal{O}_Z$-module $\mathcal{G}$ quasi-cohérent
avec $i_*\mathcal{G} \cong \mathcal{F}$. De plus :
\begin{enumerate}
\item Si $\Spec(A) \subset X$ est un ouvert
affine, et $\mathcal{F}|_{\Spec(A)} = \widetilde{M}$
alors $Z \cap \Spec(A) = \Spec(A/I)$ où $I = \text{Ann}_A(M)$.
\item Le faisceau quasi-cohérent $\mathcal{G}$ est unique à isomorphisme près
unique.
\item Le faisceau quasi-cohérent $\mathcal{G}$ est de type fini.
\item Le support de $\mathcal{G}$ et de $\mathcal{F}$ est $Z$.
\end{enumerate}
\end{lemma}

\begin{proof}
Supposons que $i' : Z' \to X$ soit un sous-schéma fermé qui satisfait
à la description sur ouverts affines
du lemme. Puis par le lemme \ref{lemma-i-star-equivalence}
on voit que $\mathcal{F} \cong i'_*\mathcal{G}'$ pour un faisceau unique
quasi-cohérent $\mathcal{G}'$ sur $Z'$. De plus, il
est clair que $Z'$ est le plus petit sous-schéma fermé
avec cette propriété
(par le même lemme). Enfin, en utilisant Propriétés,
Lemme
\ref{properties-lemma-finite-type-module} et Algèbre, Lemme \ref{algebra-lemma-finite-over-subring}
Il s'ensuit que $\mathcal{G}'$ est de type fini.
Nous avons $\text{Supp}(\mathcal{G}') = Z$ par
Algèbre, Lemme \ref{algebra-lemma-support-closed}. Par conséquent, pour
prouver le lemme, il suffit de montrer que la caractérisation
en (1) définit effectivement un sous-schéma fermé. Et,
pour ce faire, il suffit de prouver que la règle donnée
produit un faisceau quasi-cohérent d'idéaux,
voir Lemme \ref{lemma-closed-immersion-bijection-ideals}. Cela se résume au fait
algébrique suivant : si $A$ est un anneau, $f \in A$ et $M$
est un module $A$
fini, alors $\text{Ann}_A(M)_f = \text{Ann}_{A_f}(M_f)$. Nous omettons la
démonstration.
\end{proof}

\begin{definition}
\label{definition-scheme-theoretic-support}
Soit $X$ un schéma. Soit $\mathcal{F}$ un $\mathcal{O}_X$-module
quasi-cohérent de type fini. Le {\it support
schématique de $\mathcal{F}$} est le sous-schéma fermé $Z \subset X$
construit dans le lemme \ref{lemma-scheme-theoretic-support}.
\end{definition}

\noindent
Dans cette situation on pense souvent à $\mathcal{F}$ comme un faisceau
quasi-cohérent de type fini sur $Z$ (via l'équivalence
des catégories du lemme \ref{lemma-i-star-equivalence}).













\section{Image schématique}
\label{section-scheme-theoretic-image}

\noindent
Attention : Certains éléments de cette section sont ultra-généraux et se comportent
différemment de ce à quoi vous pourriez vous attendre.

\begin{lemma}
\label{lemma-scheme-theoretic-image}
Soit $f : X \to Y$ un morphisme de schémas. Il existe un sous-schéma fermé
$Z \subset Y$ tel que $f$ se factorise par $Z$ et tel que
pour tout autre sous-schéma fermé $Z' \subset Y$ tel que $f$
se factorise par $Z'$ on a $Z \subset Z'$.
\end{lemma}

\begin{proof}
Soit $\mathcal{I} = \Ker(\mathcal{O}_Y \to f_*\mathcal{O}_X)$. Si $\mathcal{I}$ est quasi-cohérent
alors on prend simplement $Z$ comme étant le
sous-schéma fermé déterminé par $\mathcal{I}$,
voir Lemme \ref{lemma-closed-immersion-bijection-ideals}. Cela fonctionne par le lemme
\ref{lemma-closed-immersion-ideals}. De manière générale,
le même lemme nous oblige à montrer qu'il existe
un plus grand faisceau quasi-cohérent d'idéaux $\mathcal{I}'$
contenu
dans $\mathcal{I}$. Cela résulte du lemme \ref{lemma-largest-quasi-coherent-subsheaf}.
\end{proof}

\begin{definition}
\label{definition-scheme-theoretic-image}
Soit $f : X \to Y$ un morphisme de schémas. L'{\it image schématique}
de $f$ est le plus petit sous-schéma fermé $Z \subset Y$ à travers lequel
$f$ se factorise par, voir Lemme \ref{lemma-scheme-theoretic-image} ci-dessus.
\end{definition}

\noindent
Pour un morphisme $f : X \to Y$ de schémas à image schématique $Z$ on note
souvent $f : X \to Z$ la factorisation de $f$ à travers
son image schématique. Si le morphisme $f$ n'est pas
quasi-compact, alors (en général)
\begin{enumerate}
\item l'inclusion théorique des ensembles $\overline{f(X)} \subset Z$ n'est pas une
égalité, c'est-à-dire que $f(X) \subset Z$ n'est pas un sous-ensemble dense, et
\item la construction de l'image schématique ne commute pas avec restriction
aux sous-schémas ouverts à $Y$.
\end{enumerate}
Dans Exemples, section \ref{examples-section-scheme-theoretic-image}, le lecteur trouve un exemple
pour les deux phénomènes. Ces phénomènes
peuvent survenir même en immersion, voir Exemples,
section \ref{examples-section-strange-immersion}. Cependant, le lemme suivant
montre que les deux désastres sont évités lorsque le
morphisme est quasi-compact.

\begin{lemma}
\label{lemma-quasi-compact-scheme-theoretic-image}
Soit $f : X \to Y$ un morphisme de schémas. Soit
$Z \subset Y$ l'image schématique de $f$. Si $f$
est quasi-compact alors
\begin{enumerate}
\item le faisceau
d'idéaux $\mathcal{I} = \Ker(\mathcal{O}_Y \to f_*\mathcal{O}_X)$ est
quasi-cohérent,
\item l'image schématique $Z$ est le sous-schéma fermé
déterminé par $\mathcal{I}$,
\item pour tout ouvert $U \subset Y$ l'image schématique
de $f|_{f^{-1}(U)} : f^{-1}(U) \to U$ est égale à $Z \cap U$, et
\item l'image $f(X) \subset Z$ est un sous-ensemble dense de $Z$,
autrement dit le morphisme $X \to Z$ est
dominant (voir Définition \ref{definition-dominant}).
\end{enumerate}
\end{lemma}

\begin{proof}
La partie (4) découle de la partie (3). Pour montrer (3) il suffit
de prouver (1) puisque la formation de $\mathcal{I}$ commute avec restriction aux
sous-schémas ouverts de $Y$. Et si (1) est vérifié alors dans
la démonstration du lemme \ref{lemma-scheme-theoretic-image} nous
avons montré (2). Il suffit donc de prouver que $\mathcal{I}$ est quasi-cohérent. Puisque
la propriété d'être quasi-cohérent est locale, nous
pouvons supposer que $Y$ est affine. Comme $f$
est quasi-compact, on peut trouver un recouvrement
ouvert fini affine $X = \bigcup_{i = 1, \ldots, n} U_i$. Notons $f'$ la composition
$$
X' = \coprod U_i \longrightarrow X \longrightarrow Y.
$$
Alors $f_*\mathcal{O}_X$ est un sous-faisceau de $f'_*\mathcal{O}_{X'}$,
et donc $\mathcal{I} = \Ker(\mathcal{O}_Y \to f'_*\mathcal{O}_{X'})$. Par Schémas, Lemme
\ref{schemes-lemma-push-forward-quasi-coherent} le faisceau $f'_*\mathcal{O}_{X'}$ est quasi-cohérent
sur $Y$. On obtient donc le résultat.
\end{proof}

\begin{example}
\label{example-scheme-theoretic-image}
Si $A \to B$ est un homomorphisme d'anneaux à noyau $I$,
alors l'image schématique de $\Spec(B) \to \Spec(A)$ est
le sous-schéma fermé $\Spec(A/I)$ de
$\Spec(A)$. Cela résulte du lemme \ref{lemma-quasi-compact-scheme-theoretic-image}.
\end{example}

\noindent
Si le morphisme est quasi-compact, alors l'image schématique ajoute uniquement
des points qui sont des spécialisations de points dans l'image.

\begin{lemma}
\label{lemma-reach-points-scheme-theoretic-image}
Soit $f : X \to Y$ un morphisme quasi-compact. Soit $Z$
l'image schématique de $f$. Soit $z \in Z$\footnote{Par
Lemme
\ref{lemma-quasi-compact-scheme-theoretic-image}, au sens ensembliste $Z$ coïncide avec l'adhérence de
$f(X)$ dans $Y$.}. Il existe un anneau
de valorisation $A$ avec corps des
fractions $K$ et un diagramme commutatif
$$
\xymatrix{
\Spec(K) \ar[rr] \ar[d] & & X \ar[d] \ar[ld] \\
\Spec(A) \ar[r] & Z \ar[r] & Y
}
$$
tel que le point fermé de $\Spec(A)$ s'envoie sur $z$. En particulier tout
point de $Z$ est la spécialisation d'un point de $f(X)$.
\end{lemma}

\begin{proof}
Soit $z \in \Spec(R) = V \subset Y$ un voisinage ouvert affine de
$z$. Par le Lemme
\ref{lemma-quasi-compact-scheme-theoretic-image} l'intersection $Z \cap V$ est l'image
schématique de $f^{-1}(V) \to V$. Nous pouvons donc
remplacer $Y$ par $V$ et supposer
que $Y = \Spec(R)$ est affine.
Dans ce cas, $X$ est quasi-compact comme
$f$ est quasi-compact.
Disons que $X = U_1 \cup \ldots \cup U_n$ est un recouvrement ouvert
affine fini. Écrivez $U_i = \Spec(A_i)$. Soit
$I = \Ker(R \to A_1 \times \ldots \times A_n)$. Par le Lemme \ref{lemma-quasi-compact-scheme-theoretic-image} encore on voit
que $Z$ correspond au sous-schéma fermé
$\Spec(R/I)$ sur $Y$. Si $\mathfrak p \subset R$
est l'idéal premier correspondant à
$z$, alors on voit que $I_{\mathfrak p} \subset R_{\mathfrak p}$
n'est pas une égalité. Par conséquent
(la localisation étant exacte, voir Algèbre, Proposition
\ref{algebra-proposition-localization-exact}),
le morphisme $R_{\mathfrak p} \to
(A_1)_{\mathfrak p} \times \ldots \times (A_n)_{\mathfrak p}$ n'est pas nul. Donc
l'un des anneaux $(A_i)_{\mathfrak p}$ n'est pas nul. Il existe
donc un $i$ et un premier $\mathfrak q_i \subset A_i$ superposés
à un premier $\mathfrak p_i \subset \mathfrak p$. Par Algèbre,
Lemme \ref{algebra-lemma-dominate} on peut choisir un anneau de valorisation
$A \subset K = \kappa(\mathfrak q_i)$ dominant l'anneau
local
$R_{\mathfrak p}/\mathfrak p_iR_{\mathfrak p} \subset \kappa(\mathfrak q_i)$. Cela donne le diagramme souhaité.
Certains détails omis.
\end{proof}

\begin{lemma}
\label{lemma-factor-factor}
Soit
$$
\xymatrix{
X_1 \ar[d] \ar[r]_{f_1} & Y_1 \ar[d] \\
X_2 \ar[r]^{f_2} & Y_2
}
$$
un diagramme commutatif de schémas. Soit $Z_i \subset Y_i$, $i = 1, 2$ l'image
schématique de $f_i$. Alors le morphisme $Y_1 \to Y_2$
induit un morphisme $Z_1 \to Z_2$ et un diagramme
commutatif
$$
\xymatrix{
X_1 \ar[r] \ar[d] & Z_1 \ar[d] \ar[r] & Y_1 \ar[d] \\
X_2 \ar[r] & Z_2 \ar[r] & Y_2
}
$$
\end{lemma}

\begin{proof}
L'image réciproque schématique de $Z_2$ dans $Y_1$ est
un sous-schéma fermé de $Y_1$ par lequel
$f_1$ se factorise. Par conséquent, $Z_1$ est contenu dans
ceci. Cela prouve le lemme.
\end{proof}

\begin{lemma}
\label{lemma-scheme-theoretic-image-reduced}
Soit $f : X \to Y$ un morphisme de schémas. Si $X$
est réduit, alors l'image schématique de $f$ est
la structure schématique induite réduite sur $\overline{f(X)}$.
\end{lemma}

\begin{proof}
Ceci est vrai car la structure schématique induite réduite sur $\overline{f(X)}$ est clairement
le plus petit sous-schéma fermé de $Y$ à travers lequel $f$
se
factorise par, voir Schémas, Lemme \ref{schemes-lemma-map-into-reduction}.
\end{proof}

\begin{lemma}
\label{lemma-scheme-theoretic-image-of-partial-section}
Soit $f : X \to Y$ un morphisme de schémas séparé. Soit
$V \subset Y$ un ouvert rétrocompact. Soit $s : V \to X$ un morphisme
tel que $f \circ s = \text{id}_V$. Soit $Y'$ l'image
schématique de $s$. Alors $Y' \to Y$
est un isomorphisme sur $V$.
\end{lemma}

\begin{proof}
L'hypothèse que $V$ est rétrocompact
dans $Y$ (Topologie, Définition \ref{topology-definition-quasi-compact})
signifie que $V \to Y$ est un morphisme
quasi-compact. Par Schémas, Lemme \ref{schemes-lemma-quasi-compact-permanence} le morphisme
$s : V \to X$ est quasi-compact. D'où
la construction de l'image schématique $Y'$
de $s$ navette avec restriction
aux ouverts par le lemme \ref{lemma-quasi-compact-scheme-theoretic-image}. En particulier,
on voit que $Y' \cap f^{-1}(V)$ est l'image schématique
d'une section du morphisme séparé $f^{-1}(V) \to V$.
Puisqu'une section d'un morphisme séparé
est une immersion fermée
(Schémas, Lemme \ref{schemes-lemma-section-immersion}), Nous concluons
que $Y' \cap f^{-1}(V) \to V$
est un isomorphisme souhaité.
\end{proof}







\section{Adhérence schématique et densité}
\label{section-scheme-theoretic-closure}

\noindent
On reprend la définition suivante de \cite[IV, Définition 11.10.2]{EGA}.

\begin{definition}
\label{definition-scheme-theoretically-dense}
Soit $X$ un schéma. Soit $U \subset X$ un sous-schéma ouvert.
\begin{enumerate}
\item L'image schématique du morphisme $U \to X$ est appelée {\it
adhérence schématique de $U$ dans $X$}.
\item On dit que $U$ est {\it schématiquement dense dans
$X$} si pour tout $V \subset X$ ouvert l'adhérence schématique de $U \cap V$
dans $V$ est égale à $V$.
\end{enumerate}
\end{definition}

\noindent
Avec cette définition,
il n'est {\bf pas} que $U$ soit
schématiquement dense dans $X$ si et
seulement si l'adhérence schématique de $U$
est $X$, voir Exemple
\ref{example-scheme-theretically-dense-not-dense}. C'est quelque peu inélégant ; mais voir Lemmes
\ref{lemma-scheme-theoretically-dense-quasi-compact} et \ref{lemma-reduced-scheme-theoretically-dense} ci-dessous.
En revanche, avec cette définition $U$ est schématiquement
dense en $X$ si et seulement si pour chaque $V \subset X$
ouvert l'homomorphisme d'anneaux $\mathcal{O}_X(V) \to \mathcal{O}_X(U \cap V)$ est
injectif, voir Lemme \ref{lemma-characterize-scheme-theoretically-dense} ci-dessous. En
particulier on voit que schématiquement dense implique dense
ce qui est agréable.

\begin{example}
\label{example-scheme-theretically-dense-not-dense}
Voici un exemple où le fait que l'adhérence schématique soit $X$ n'implique
pas la densité des espaces topologiques sous-jacents.
Soit $k$
un corps. Posons $A = k[x, z_1, z_2, \ldots]/(x^n z_n)$.
Posons $I = (z_1, z_2, \ldots) \subset A$. Considérons le schéma
affine $X = \Spec(A)$ et le sous-schéma ouvert $U = X \setminus V(I)$.
Comme $A \to \prod_n A_{z_n}$ est injectif,
l'adhérence schématique de $U$ est $X$. Considérons
le morphisme $X \to \Spec(k[x])$. Ce morphisme
est surjectif (poser tous les $z_n = 0$ permet
de le voir). Mais sa restriction à $U$
n'est pas surjective, car son image est réduite
au point $x = 0$. Par conséquent, $U$ ne peut pas être topologiquement dense dans
$X$.
\end{example}

\begin{lemma}
\label{lemma-scheme-theoretically-dense-quasi-compact}
Soit $X$ un schéma.
Soit $U \subset X$ un sous-schéma ouvert. Si le morphisme
d'inclusion $U \to X$ est quasi-compact, alors $U$ est schématiquement
dense dans $X$ si et seulement si l'adhérence schématique de $U$
dans $X$ est $X$.
\end{lemma}

\begin{proof}
Cela résulte de la partie (3) du Lemme \ref{lemma-quasi-compact-scheme-theoretic-image}.
\end{proof}

\begin{example}
\label{example-scheme-theoretic-closure}
Soit $A$ un anneau et $X = \Spec(A)$.
Soient $f_1, \ldots, f_n \in A$ et posons $U = D(f_1) \cup \ldots \cup D(f_n)$.
Posons $I = \Ker(A \to \prod A_{f_i})$. Alors l'adhérence schématique de $U$
dans $X$ est le sous-schéma fermé
$\Spec(A/I)$ de $X$. Notons que
$U \to X$ est quasi-compact. Ainsi par
le lemme \ref{lemma-scheme-theoretically-dense-quasi-compact} on voit que $U$ est schématiquement
dense en $X$ si et seulement si $I = 0$.
\end{example}

\begin{lemma}
\label{lemma-characterize-scheme-theoretically-dense}
Soit $j : U \to X$ une immersion ouverte de schémas.
Alors $U$ est schématiquement dense dans $X$
si et seulement si $\mathcal{O}_X \to j_*\mathcal{O}_U$ est injectif.
\end{lemma}

\begin{proof}
Si $\mathcal{O}_X \to j_*\mathcal{O}_U$ est injectif, alors il en va de même lorsqu'il
est restreint à n'importe quel $V$ ouvert de $X$.
D'où l'adhérence schématique de $U \cap V$ dans $V$
est égale à $V$, voir démonstration du lemme \ref{lemma-scheme-theoretic-image}.
A l'inverse, supposons que l'adhérence schématique
de $U \cap V$ soit égale à $V$ pour toutes les
ouvertures $V$. Supposons que $\mathcal{O}_X \to j_*\mathcal{O}_U$ ne soit pas injectif.
Ensuite, nous pouvons trouver un ouvert affine, disons $\Spec(A) = V \subset X$
et un élément non nul $f \in A$ tel que l'image de $f$ soit
nulle dans $\Gamma(V \cap U, \mathcal{O}_X)$. Dans ce cas l'adhérence schématique de $V \cap U$
dans $V$ est clairement contenue dans $\Spec(A/(f))$ une
contradiction.
\end{proof}

\begin{lemma}
\label{lemma-intersection-scheme-theoretically-dense}
Soit $X$ un schéma. Si $U$, $V$ sont schématiquement denses
sous-schémas ouverts de $X$, alors $U \cap V$.
\end{lemma}

\begin{proof}
Soit $W \subset X$ un ouvert
quelconque. Considérons le morphisme
$\mathcal{O}_X(W) \to \mathcal{O}_X(W \cap V)
\to \mathcal{O}_X(W \cap V \cap U)$. D'après le Lemme \ref{lemma-characterize-scheme-theoretically-dense},
les deux applications sont injectives.
Le composite est donc injectif. Donc
par le lemme \ref{lemma-characterize-scheme-theoretically-dense} $U \cap V$ est schématiquement
dense en $X$.
\end{proof}

\begin{lemma}
\label{lemma-quasi-compact-immersion}
Soit $h : Z \to X$ une immersion. Supposons que $h$ soit quasi-compact
ou que $Z$ soit réduit. Soit $\overline{Z} \subset X$ l'image schématique de
$h$. Alors le morphisme $Z \to \overline{Z}$ est une immersion ouverte qui identifie
$Z$ à un sous-schéma ouvert schématiquement
dense de $\overline{Z}$. De plus, $Z$ est topologiquement
dense dans $\overline{Z}$.
\end{lemma}

\begin{proof}
Par le Lemme \ref{lemma-factor-quasi-compact-immersion} ou Lemme \ref{lemma-factor-reduced-immersion}, nous
pouvons factoriser $Z \to X$ comme $Z \to \overline{Z}_1 \to X$
avec $Z \to \overline{Z}_1$ ouvert et $\overline{Z}_1 \to X$ fermé. Par contre, soit
$Z \to \overline{Z} \subset X$ l'image schématique de $Z \to X$. Nous
concluons que $\overline{Z} \subset \overline{Z}_1$. Puisque $Z$ est un sous-schéma
ouvert de $\overline{Z}_1$ Il s'ensuit que $Z$
est également un sous-schéma ouvert de $\overline{Z}$.
Dans le cas où $Z$ est réduit nous savons
que $Z \subset \overline{Z}_1$ est topologiquement dense par la construction
de $\overline{Z}_1$ dans la démonstration du lemme \ref{lemma-factor-reduced-immersion}.
Par conséquent, $\overline{Z}_1$ et $\overline{Z}$ ont
les mêmes espaces topologiques sous-jacents
Ainsi $\overline{Z} \subset \overline{Z}_1$ est une immersion fermée dans un schéma réduit
qui induit une bijection sur les espaces topologiques sous-jacents,
et donc un isomorphisme. Dans le cas où $Z \to X$
est quasi-compact nous argumentons comme suit
: L'affirmation selon laquelle $Z$ est
schématiquement dense
dans $\overline{Z}$ découle de la partie (3) du lemme \ref{lemma-quasi-compact-scheme-theoretic-image}.
La dernière assertion
découle de la partie (4) du lemme \ref{lemma-quasi-compact-scheme-theoretic-image}.
\end{proof}

\begin{lemma}
\label{lemma-reduced-scheme-theoretically-dense}
Soit $X$ un schéma réduit et soit $U \subset X$ un sous-schéma ouvert. Alors Les
conditions suivantes sont équivalentes
\begin{enumerate}
\item $U$ est topologiquement dense dans $X$,
\item l'adhérence schématique de $U$ dans $X$ est $X$, et
\item $U$ est schématiquement dense en $X$.
\end{enumerate}
\end{lemma}

\begin{proof}
Cela résulte du lemme
\ref{lemma-quasi-compact-immersion} et le fait qu'un sous-schéma
fermé $Z$ de $X$ dont l'espace topologique
sous-jacent est égal à $X$ doit être égal à $X$ comme
schéma.
\end{proof}

\begin{lemma}
\label{lemma-reduced-subscheme-closure}
Soit $X$ un schéma et soit $U \subset X$ un sous-schéma ouvert réduit. Alors Les
conditions suivantes sont équivalentes
\begin{enumerate}
\item l'adhérence schématique de $U$ dans $X$ est $X$, et
\item $U$ est schématiquement dense dans $X$.
\end{enumerate}
Si cela est vrai alors $X$ est un schéma réduit.
\end{lemma}

\begin{proof}
Cela résulte du
lemme \ref{lemma-quasi-compact-immersion} et du fait que
l'adhérence schématique de $U$ dans $X$
est
réduite par le lemme \ref{lemma-scheme-theoretic-image-reduced}.
\end{proof}


\begin{lemma}
\label{lemma-equality-of-morphisms}
Soit $S$ un schéma. Soient $X$, $Y$ des schémas
sur $S$. Soit $f, g : X \to Y$ des morphismes de schémas sur
$S$. Soit $U \subset X$ un sous-schéma ouvert tel que $f|_U = g|_U$.
Si l'adhérence schématique de $U$ dans $X$
est $X$ et que $Y \to S$ est séparé, alors $f = g$.
\end{lemma}

\begin{proof}
Dérivé des définitions
et Schémas, Lemme \ref{schemes-lemma-where-are-they-equal}.
\end{proof}




\section{Morphismes dominants}
\label{section-dominant}

\noindent
La définition d'un morphisme de schémas dominant est un peu différente
de ce à quoi on pourrait s'attendre si l'on est habitué à la notion
de morphisme dominant des variétés.

\begin{definition}
\label{definition-dominant}
Un morphisme $f : X \to S$ de schémas est dit {\it dominant} si l'image de
$f$ est un sous-ensemble dense de $S$.
\end{definition}

\noindent
Ainsi, par exemple, si $k$ est un corps infini et $\lambda_1, \lambda_2, \ldots$ est une collection
dénombrable d'éléments distincts de $k$, alors le morphisme
$$
\coprod\nolimits_{i = 1, 2, \ldots } \Spec(k)
\longrightarrow
\Spec(k[x])
$$
avec le $i$ème facteur envoyé au point $x = \lambda_i$ est dominant.

\begin{lemma}
\label{lemma-generic-points-in-image-dominant}
Soit $f : X \to S$ un morphisme de schémas. Si chaque
point générique de chaque composante irréductible de $S$ est à l'image
de $f$, alors $f$ est dominant.
\end{lemma}

\begin{proof}
Il s'agit d'un fait topologique qui découle directement du fait
que l'espace topologique sous-jacent à un schéma est
sobre, voir Schémas, Lemme \ref{schemes-lemma-scheme-sober}, et que tout point
de $S$ est contenu dans une composante irréductible
de $S$, voir Topologie, Lemme \ref{topology-lemma-irreducible}.
\end{proof}

\noindent
L'attente selon laquelle les morphismes sont dominants uniquement si des points génériques
de la cible sont dans l'image est valable si le morphisme est quasi-compact.

\begin{lemma}
\label{lemma-quasi-compact-dominant}
\begin{slogan}
Les morphismes dont l'image contient les points génériques sont dominants
\end{slogan}
Soit $f : X \to S$ un morphisme de schémas quasi-compact. Alors $f$
est dominant si et seulement si pour chaque composante irréductible
$Z \subset S$ le point générique de $Z$ est à l'image de
$f$.
\end{lemma}

\begin{proof}
Soit $V \subset S$ un ouvert affine. Parce
que $f$ est quasi-compact, nous pouvons choisir un nombre
fini d'ouverts affines $U_i \subset f^{-1}(V)$, $i = 1, \ldots, n$ couvrant $f^{-1}(V)$.
Considérons le morphisme des affines
$$
f' :
\coprod\nolimits_{i = 1, \ldots, n} U_i
\longrightarrow
V.
$$
Une union disjointe d'affines est affine, voir
Schémas, Lemme \ref{schemes-lemma-disjoint-union-affines}. Les points génériques des composantes
irréductibles de $V$ sont exactement
les points génériques des composantes irréductibles de $S$ qui
répondent à $V$. De plus, $f$ est dominant si et seulement
si $f'$ est dominant quels que soient les choix de $V, n, U_i$
que nous faisons ci-dessus. Nous avons ainsi réduit le lemme au cas
d'un morphisme de
schémas affines. Le cas affine est l'Algèbre, Lemme \ref{algebra-lemma-image-dense-generic-points}.
\end{proof}

\begin{lemma}
\label{lemma-base-change-dominant}
Soit $f : X \to S$ un morphisme de schémas dominant quasi-compact. Soit $g : S' \to S$
un morphisme de schémas, et soit $f' : X' \to S'$ le changement de base de $f$
par $g$. Si les générisations s'élèvent le long de $g$, alors
$f'$ est dominant.
\end{lemma}

\noindent
Nous verrons que les générisations s'élèvent le long de $g$
si $g$ est ouvert
(Lemme \ref{lemma-open-generizing}) ou plat (Lemme \ref{lemma-generalizations-lift-flat}).

\begin{proof}
Observez que $f'$ est quasi-compact par Schémas,
Lemme \ref{schemes-lemma-quasi-compact-preserved-base-change}. Soit $\eta' \in S'$ le point générique
d'une composante irréductible de $S'$.
Si les générisations s'élèvent le long de $g$,
alors $\eta = g(\eta')$ est le point générique d'une composante
irréductible de $S$. Par le Lemme \ref{lemma-quasi-compact-dominant}
on voit que $\eta$ est à l'image de
$f$. Ainsi $\eta'$ est à
l'image de $f'$ d'après Schémas, Lemme \ref{schemes-lemma-points-fibre-product}.
Il s'ensuit que $f'$ est dominant par le lemme \ref{lemma-quasi-compact-dominant}.
\end{proof}

\noindent
Cela vaut en partie pour les morphismes dominants arbitraires (pas nécessairement
quasi-compacts).

\begin{lemma}
\label{lemma-open-base-change-dominant}
Soit $f : X \to S$ un morphisme de schémas dominant. Soit $g : S' \to S$
un morphisme de schémas ouvert, et soit $f' : X' \to S'$ le changement de base
de $f$ par $g$. Alors $f'$ est dominant.
\end{lemma}

\begin{proof}
Si $f'$ n'est pas dominant, alors $f'(X')$ est contenu dans un sous-ensemble
fermé $B$ différent de $S'$. Ensuite, $S'-B$ est ouvert dans $S'$
et non vide. Puisque $g$ est ouvert, $g(S'-B)$ n'est pas vide et ouvert
dans $S$. Puisque $f(X)$ est dense dans $S$, il existe un
point $x\in X$ avec $f(x) \in g(S'-B)$. Soit
$W = \Spec(\kappa (f(x)))$, et écrivez $X_W$, $S_W$, $S'_W$
pour les retraits de ces schémas de $S$ vers $W$. Alors $X_W$ et
$(S'-B)_W$ ne sont pas vides, et donc $X_W \times_W (S'-B)_W$ n'est pas vide, car $W$
est le spectre d'un corps. Cela contredit le fait que $f'(X \times_S S')$ soit contenu
dans $B \subset S'$. Donc en fait $f'$ est dominant.
\end{proof}

\begin{lemma}
\label{lemma-quasi-compact-generic-point-not-in-image}
Soit $f : X \to S$ un morphisme de schémas quasi-compact.
Soit $\eta \in S$ un point générique d'une composante
irréductible de $S$. Si $\eta \not \in f(X)$ alors
il existe un voisinage ouvert $V \subset S$ de
$\eta$ tel que $f^{-1}(V) = \emptyset$.
\end{lemma}

\begin{proof}
Soit $Z \subset S$ l'image schématique de $f$.
Nous devons montrer que $\eta \not \in Z$.
Cela résulte du
lemme \ref{lemma-reach-points-scheme-theoretic-image} mais peut aussi être vu comme suit.
Par le Lemme \ref{lemma-quasi-compact-scheme-theoretic-image} le morphisme
$X \to Z$ est dominant, ce qui par le lemme
\ref{lemma-quasi-compact-dominant} signifie que tous les points
génériques de toutes les composantes
irréductibles de $Z$ sont à l'image de $X \to Z$.
Par hypothèse on voit que $\eta \not \in Z$ puisque $\eta$
serait le point générique d'une composante
irréductible de $Z$ s'il était dans $Z$.
\end{proof}

\noindent
Il existe un autre cas où la dominante revient à avoir tous les points
génériques des composantes irréductibles dans l'image.

\begin{lemma}
\label{lemma-dominant-finite-number-irreducible-components}
Soit $f : X \to S$ un morphisme de schémas. Supposons que
$X$ ait un nombre fini de composantes irréductibles. Alors $f$
est dominant (si et) seulement si pour toute composante irréductible
$Z \subset S$ le point générique de $Z$ est à l'image de $f$.
Si tel est le cas, alors $S$ a également un nombre fini de composantes
irréductibles.
\end{lemma}

\begin{proof}
Supposons que $f$
soit dominant. Disons que $X = Z_1 \cup Z_2 \cup \ldots \cup Z_n$ est la décomposition de $X$
en composantes irréductibles. Soit $\xi_i \in Z_i$ son
point générique, donc $Z_i = \overline{\{\xi_i\}}$. Notez que $f(Z_i)$
est un sous-ensemble irréductible de
$S$.
$$
S = \overline{f(X)} = \bigcup \overline{f(Z_i)} =
\bigcup \overline{\{f(\xi_i)\}}
$$
est donc une union finie de sous-ensembles irréductibles dont les points
génériques sont à l'image de $f$. Le lemme suit.
\end{proof}

\begin{lemma}
\label{lemma-dominant-between-integral}
Soit $f : X \to Y$ un morphisme de schémas intègres. Les conditions suivantes sont
équivalentes
\begin{enumerate}
\item $f$ est dominant,
\item $f$ envoie le point générique de $X$ sur le point générique de $Y$,
\item il existe des ouverts affines non vides $U \subset X$ et
$V \subset Y$ tels que $f(U) \subset V$ et l'homomorphisme d'anneaux $\mathcal{O}_Y(V) \to \mathcal{O}_X(U)$ est
injectif,
\item pour tous les ouverts affines non vides $U \subset X$ et
$V \subset Y$ tels que $f(U) \subset V$, l'homomorphisme d'anneaux $\mathcal{O}_Y(V) \to \mathcal{O}_X(U)$ est
injectif,
\item il existe $x \in X$, d'image $y = f(x) \in Y$, tel que l'homomorphisme
d'anneaux locaux $\mathcal{O}_{Y, y} \to \mathcal{O}_{X, x}$ soit injectif, et
\item pour tout $x \in X$, d'image $y = f(x) \in Y$, l'homomorphisme
d'anneaux locaux $\mathcal{O}_{Y, y} \to \mathcal{O}_{X, x}$ est injectif.
\end{enumerate}
\end{lemma}

\begin{proof}
L'équivalence de (1) et (2) découle du lemme
\ref{lemma-dominant-finite-number-irreducible-components}. Soient $U \subset X$ et $V \subset Y$ des ouverts affines
non vides avec $f(U) \subset V$. Rappelons que les anneaux $A = \mathcal{O}_X(U)$
et $B = \mathcal{O}_Y(V)$ sont des domaines intègres. Le morphisme
$f|_U : U \to V$ correspond à un homomorphisme
d'anneaux $\varphi : B \to A$. Les points génériques de $X$
et $Y$ correspondent aux idéaux premiers $(0) \subset A$
et $(0) \subset B$. Ainsi (2) équivaut à la condition $(0) = \varphi^{-1}((0))$,
c'est-à-dire à la condition que $\varphi$ soit injectif.
De cette façon, nous voyons que (2), (3) et
(4) sont équivalents. De même, étant donné $x$
et $y$ comme dans (5), les anneaux locaux
$\mathcal{O}_{X, x}$ et $\mathcal{O}_{Y, y}$ sont des domaines et les idéaux
premiers $(0) \subset \mathcal{O}_{X, x}$ et $(0) \subset \mathcal{O}_{Y, y}$ correspondent
aux points génériques de $X$ et $Y$
(via l'identification du spectre de l'anneau local à
$x$ avec l'ensemble des points
spécialisés à $x$, voir Schémas,
Lemme \ref{schemes-lemma-specialize-points}). Nous pouvons donc argumenter exactement
de la même manière que ci-dessus pour voir
que (2), (5) et (6) sont équivalents.
\end{proof}







\section{Morphismes surjectifs}
\label{section-surjective}

\begin{definition}
\label{definition-surjective}
Un morphisme de schémas est dit {\it surjectif}
s'il est surjectif sur les espaces topologiques
sous-jacents.
\end{definition}

\begin{lemma}
\label{lemma-composition-surjective}
La composition des morphismes surjectifs est surjectif.
\end{lemma}

\begin{proof}
Omis.
\end{proof}

\begin{lemma}
\label{lemma-when-point-maps-to-pair}
Soient $X$ et $Y$ des schémas sur un schéma de base $S$. Étant donné les points $x \in X$
et $y \in Y$, il y a un point de $X \times_S Y$ envoyé sur $x$ et $y$ sous les projections
si et seulement si $x$ et $y$ se trouvent au-dessus du même point de $S$.
\end{lemma}

\begin{proof}
La condition est évidemment nécessaire, et la réciproque découle de
la démonstration d'après Schémas, Lemme \ref{schemes-lemma-points-fibre-product}.
\end{proof}

\begin{lemma}
\label{lemma-base-change-surjective}
Le changement de base d'un morphisme surjectif est surjectif.
\end{lemma}

\begin{proof}
Soit $f: X \to Y$ un morphisme de schémas sur un schéma de base $S$.
Si $S' \to S$ est un morphisme de schémas, soit $p: X_{S'} \to X$
et $q: Y_{S'} \to Y$ les projections canoniques. Le carré
commutatif
$$
\xymatrix{
X_{S'} \ar[d]_{f_{S'}} \ar[r]_p & X \ar[d]^{f} \\
Y_{S'} \ar[r]^{q} & Y.
}
$$
identifie $X_{S'}$ comme un produit fibré de $X \to Y$
et $Y_{S'} \to Y$. Soit $Z$ un sous-ensemble de l'espace
topologique sous-jacent de $X$. Puis $q^{-1}(f(Z)) = f_{S'}(p^{-1}(Z))$,
car $y' \in q^{-1}(f(Z))$ si et seulement si $q(y') = f(x)$ pour certains $x \in Z$,
si et seulement si, par le lemme \ref{lemma-when-point-maps-to-pair}, il existe des $x' \in X_{S'}$
tels que $f_{S'}(x') = y'$ et $p(x') = x$. En particulier en prenant $Z = X$
on voit que si $f$ est surjectif le changement de base
$f_{S'}: X_{S'} \to Y_{S'}$ l'est aussi.
\end{proof}

\begin{example}
\label{example-injective-not-preserved-base-change}
La bijectivité n'est pas stable par changement
de base, et l'injectivité ne l'est
pas non plus. Par exemple, considérons
la bijection
$\Spec(\mathbf{C}) \to \Spec(\mathbf{R})$. Le changement de base
$\Spec(\mathbf{C} \otimes_{\mathbf{R}} \mathbf{C}) \to
\Spec(\mathbf{C})$ n'est pas injectif, car il existe un isomorphisme
$\mathbf{C} \otimes_{\mathbf{R}} \mathbf{C} \cong \mathbf{C} \times \mathbf{C}$
(la décomposition
provient de l'idempotent $\frac{1 \otimes 1 + i \otimes i}{2}$) ;
ainsi, $\Spec(\mathbf{C} \otimes_{\mathbf{R}} \mathbf{C})$ a deux points.
\end{example}

\begin{lemma}
\label{lemma-surjection-from-quasi-compact}
Soit
$$
\xymatrix{
X \ar[rr]_f \ar[rd]_p & &
Y \ar[dl]^q \\
& Z
}
$$
être un diagramme commutatif de morphismes de schémas. Si $f$
est surjectif et $p$ est quasi-compact, alors $q$ est quasi-compact.
\end{lemma}

\begin{proof}
Soit $W \subset Z$ un ouvert quasi-compact. Par hypothèse $p^{-1}(W)$
est quasi-compact.
Donc par Topologie, Lemme \ref{topology-lemma-image-quasi-compact} l'image
inverse $q^{-1}(W) = f(p^{-1}(W))$ est également quasi-compacte. Cela
prouve le lemme.
\end{proof}




\section{Morphismes radiciels et universellement injectifs}
\label{section-radicial}

\noindent
Dans cette section, nous définissons ce que signifie pour un morphisme de schémas d'être {\it
radical} et ce que signifie pour un morphisme de schémas d'être {\it
universellement injectif}. Nous montrons ensuite que ces notions
concordent. La raison pour laquelle nous introduisons les deux est que dans le cas des espaces algébriques,
il existe des notions correspondantes qui ne concordent pas toujours.

\begin{definition}
\label{definition-universally-injective}
Soit $f : X \to S$ un morphisme.
\begin{enumerate}
\item On dit que $f$ est {\it universellement injectif}
si et seulement si pour tout morphisme de schémas $S' \to S$ le changement
de base $f' : X_{S'} \to S'$ est injectif (sur les espaces topologiques sous-jacents).
\item On dit que $f$ est {\it radical} si $f$
est injectif comme application d'espaces topologiques, et pour tout
$x \in X$ l'extension du corps $\kappa(x)/\kappa(f(x))$ est purement inséparable.
\end{enumerate}
\end{definition}

\begin{lemma}
\label{lemma-universally-injective}
\begin{reference}
\cite[Proposition 3.7.1]{EGA1-second}
\end{reference}
Soit $f : X \to S$ un morphisme de schémas. Les conditions
suivantes sont équivalentes :
\begin{enumerate}
\item Pour chaque corps $K$ l'application
induite $\Mor(\Spec(K), X) \to \Mor(\Spec(K), S)$ est
injectif.
\item Le morphisme $f$ est universellement injectif.
\item Le morphisme $f$ est radical.
\item Le morphisme diagonal $\Delta_{X/S} : X \longrightarrow X \times_S X$ est
surjectif.
\end{enumerate}
\end{lemma}

\begin{proof}
Soit $K$ un corps, et $s : \Spec(K) \to S$ un morphisme.
Donner un morphisme $x : \Spec(K) \to X$ tel que $f \circ x = s$ revient
à donner une section de la projection
$X_K = \Spec(K) \times_S X \to \Spec(K)$, ce qui revient à donner un point
$x \in X_K$ dont le corps résiduel est $K$. On voit
donc que (2) implique (1).

\medskip\noindent Inversement,
 supposons que (1) soit vérifié. Supposons que $x, x' \in X_{S'}$ s'envoient
sur le même point $s' \in S'$. Choisissons un diagramme commutatif
$$
\xymatrix{
K & \kappa(x) \ar[l] \\
\kappa(x') \ar[u] & \kappa(s') \ar[l] \ar[u]
}
$$
de corps. Par Schémas, Lemme \ref{schemes-lemma-characterize-points} on obtient deux morphismes
$a, a' : \Spec(K) \to X_{S'}$. L'un correspondant au point $x$ et à l'intégration
$\kappa(x) \subset K$ et l'autre correspondant au point $x'$
et à l'intégration $\kappa(x') \subset K$. Nous avons également $f' \circ a = f' \circ a'$.
La condition (1) implique désormais que les compositions
de $a$ et $a'$ avec $X_{S'} \to X$ sont égales. Puisque
$X_{S'}$ est le produit fibré de $S'$ et $X$
sur $S$ on voit que $a = a'$. D'où $x = x'$. Ainsi
(1) implique (2).

\medskip\noindent
S'il y a deux points différents $x, x' \in X$ envoyés au même point de $s$
alors (2) est violé.
Si pour certains $s = f(x)$, $x \in X$ l'extension
de corps $\kappa(x)/\kappa(s)$ n'est pas purement inséparable,
alors on peut trouver une extension de corps $K/\kappa(s)$
telle que $\kappa(x)$ possède deux $\kappa(s)$-homomorphismes
en $K$. Par Schémas, Lemme \ref{schemes-lemma-characterize-points} cela
implique que l'application
$\Mor(\Spec(K), X) \to \Mor(\Spec(K), S)$ n'est pas injectif,
et donc (1) est violé. Nous voyons donc
que les conditions équivalentes (1) et (2) impliquent que
$f$ est radical, c'est-à-dire qu'elles impliquent (3).

\medskip\noindent
Supposons
(3). Par Schémas, Lemme \ref{schemes-lemma-characterize-points} un morphisme
$\Spec(K) \to X$ est donné par un couple $(x, \kappa(x) \to K)$. La propriété
(3) dit exactement qu'en associant au couple $(x, \kappa(x) \to K)$
le couple $(s, \kappa(s) \to \kappa(x) \to K)$ est injectif. En d’autres termes,
(1) est valable. À ce stade, nous savons que (1), (2)
et (3) sont tous équivalents.

\medskip\noindent
Enfin, nous prouvons l'équivalence de (4) avec (1),
(2) et (3). Un point de $X \times_S X$ est
donné par un quadruple $(x_1, x_2, s, \mathfrak p)$, où $x_1, x_2 \in X$,
$f(x_1) = f(x_2) = s$ et $\mathfrak p \subset \kappa(x_1) \otimes_{\kappa(s)} \kappa(x_2)$
est un idéal premier, voir Schémas, Lemme \ref{schemes-lemma-points-fibre-product}.
Si $f$
est universellement injectif, alors en prenant
$S'=X$ dans la définition
de universellement injectif, $\Delta_{X/S}$ doit être
surjectif puisqu'il s'agit d'une section du morphisme
injectif
$X \times_S X  \longrightarrow X$. Inversement, si $\Delta_{X/S}$
est surjectif,
alors toujours $x_1 = x_2 = x$ et il existe exactement un tel idéal
premier $\mathfrak p$, ce qui signifie que $\kappa(s) \subset \kappa(x)$ est purement
inséparable. Donc $f$ est radical.
Alternativement, si $\Delta_{X/S}$
est surjectif, alors pour tout $S' \to S$
le changement de base $\Delta_{X_{S'}/S'}$
est surjectif ce qui implique que $f$ est universellement
injectif. Ceci termine la démonstration du lemme.
\end{proof}

\begin{lemma}
\label{lemma-universally-injective-separated}
Un morphisme universellement injectif est séparé.
\end{lemma}

\begin{proof}
Combiner
Lemme \ref{lemma-universally-injective} avec la remarque
que $X \to S$ est séparé si et seulement si l'image
de $\Delta_{X/S}$ est fermée dans $X \times_S X$,
voir Schémas, Définition \ref{schemes-definition-separated} et la
discussion qui en suit.
\end{proof}

\begin{lemma}
\label{lemma-base-change-universally-injective}
Un changement de base d'un morphisme universellement injectif est universellement injectif.
\end{lemma}

\begin{proof}
Ceci est formel.
\end{proof}

\begin{lemma}
\label{lemma-composition-universally-injective}
Une composition de morphismes radiciels est radicale, et il en va donc de même
pour la condition équivalente d'être universellement injectif.
\end{lemma}

\begin{proof}
Omis.
\end{proof}









\section{Morphismes affines}
\label{section-affine}

\begin{definition}
\label{definition-affine}
Un morphisme de schémas $f : X \to S$ est appelé {\it affine}
si l'image réciproque de chaque ouvert affine de $S$ est un ouvert affine
de $X$.
\end{definition}

\begin{lemma}
\label{lemma-affine-separated}
Un morphisme affine est séparé et quasi-compact.
\end{lemma}

\begin{proof}
Soit $f : X \to S$ affine. La quasi-compacité est immédiate chez Schémas,
Lemme \ref{schemes-lemma-quasi-compact-affine}. Nous allons montrer que $f$
est séparé à l'aide d'après
Schémas, Lemme \ref{schemes-lemma-characterize-separated}. Soient $x_1, x_2 \in X$ des points de $X$
qui ont la même image $s \in S$. Choisissons un
ouvert affine quelconque $W \subset S$ contenant $s$. Par hypothèse,
$f^{-1}(W)$ est affine. Appliquer le lemme cité avec $U = V = f^{-1}(W)$.
\end{proof}

\begin{lemma}
\label{lemma-characterize-affine}
\begin{reference}
\cite[II, Corollaire 1.3.2]{EGA}
\end{reference}
Soit $f : X \to S$ un morphisme de schémas. Les conditions
suivantes sont équivalentes
\begin{enumerate}
\item Le morphisme $f$ est affine.
\item Il existe un recouvrement ouvert affine $S = \bigcup W_j$ tel
que chaque $f^{-1}(W_j)$ soit affine.
\item Il existe un faisceau quasi-cohérent de $\mathcal{O}_S$-algèbres
$\mathcal{A}$ et un
isomorphisme $X \cong \underline{\Spec}_S(\mathcal{A})$ de schémas
sur $S$. Voir
Constructions, section \ref{constructions-section-spec} pour la notation.
\end{enumerate}
De plus, dans ce cas $X = \underline{\Spec}_S(f_*\mathcal{O}_X)$.
\end{lemma}

\begin{proof}
Il est évident que (1) implique (2).

\medskip\noindent
Supposons que $S = \bigcup_{j \in J} W_j$ soit un recouvrement ouvert
affine tel que chaque $f^{-1}(W_j)$
soit affine. Par Schémas, Lemme \ref{schemes-lemma-quasi-compact-affine}
on voit que $f$
est quasi-compact. Par Schémas, Lemme \ref{schemes-lemma-characterize-quasi-separated}
on voit que le morphisme $f$ est
quasi-séparé. Ainsi par Schémas, Lemme \ref{schemes-lemma-push-forward-quasi-coherent} le
faisceau $\mathcal{A} = f_*\mathcal{O}_X$ est un faisceau quasi-cohérent
de $\mathcal{O}_S$-algèbres. On a donc le schéma
$g : Y = \underline{\Spec}_S(\mathcal{A}) \to S$ sur $S$. L'application
d'identité
$\text{id} : \mathcal{A} = f_*\mathcal{O}_X \to f_*\mathcal{O}_X$ fournit, via la définition du spectre
relatif, un morphisme $can : X \to Y$ sur $S$,
voir Constructions, Lemme
\ref{constructions-lemma-canonical-morphism}. Par hypothèse et le lemme qui vient
d'être cité la restriction
$can|_{f^{-1}(W_j)} : f^{-1}(W_j) \to g^{-1}(W_j)$ est un isomorphisme. Ainsi $can$
est un isomorphisme. Nous avons montré
que (2) implique (3).

\medskip\noindent
Supposons (3). Par Constructions, Lemme \ref{constructions-lemma-spec-properties} nous voyons que l'image
inverse de tout ouvert affine est affine, et donc le morphisme
est affine par définition.
\end{proof}

\begin{remark}
\label{remark-direct-argument}
Nous pouvons également affirmer directement
que (2) implique (1) dans le lemme \ref{lemma-characterize-affine}
ci-dessus comme suit. Supposons que $S = \bigcup W_j$
soit un recouvrement ouvert
affine tel que chaque $f^{-1}(W_j)$ soit affine. Démontrons
d’abord que $\mathcal{A} = f_*\mathcal{O}_X$
est quasi-cohérent comme dans
la démonstration ci-dessus. Soit $\Spec(R) = V \subset S$
ouvert affine. Il faut montrer que $f^{-1}(V)$ est affine.
Posons $A = \mathcal{A}(V) = f_*\mathcal{O}_X(V) = \mathcal{O}_X(f^{-1}(V))$. Par Schémas, Lemme \ref{schemes-lemma-morphism-into-affine} il
existe un morphisme canonique $\psi : f^{-1}(V) \to \Spec(A)$ sur
$\Spec(R) = V$.
Par Schémas, Lemme \ref{schemes-lemma-good-subcover} il existe un entier $n \geq 0$,
un recouvrement standard ouvert
$V = \bigcup_{i = 1, \ldots, n} D(h_i)$, $h_i \in R$, et une application $a : \{1, \ldots, n\} \to J$
telle que chaque $D(h_i)$ est aussi un ouvert
principal du schéma affine $W_{a(i)}$. L'image réciproque
d'un principal ouvert sous un morphisme de schémas affines
est principal ouvert, voir Algèbre, Lemme \ref{algebra-lemma-spec-functorial}.
On voit donc que $f^{-1}(D(h_i))$ est un principe ouvert
de $f^{-1}(W_{a(i)})$, en particulier que $f^{-1}(D(h_i))$ est affine.
Parce que $\mathcal{A}$
est quasi-cohérent, nous avons $A_{h_i} = \mathcal{A}(D(h_i)) = \mathcal{O}_X(f^{-1}(D(h_i)))$, donc $f^{-1}(D(h_i))$
est le spectre de $A_{h_i}$.
Il s'ensuit que le morphisme $\psi$ induit un isomorphisme
de l'ouvert $f^{-1}(D(h_i))$ avec l'ouvert
$\Spec(A_{h_i})$ de $\Spec(A)$. Puisque $f^{-1}(V) = \bigcup f^{-1}(D(h_i))$
et $\Spec(A) = \bigcup \Spec(A_{h_i})$ on obtient le résultat.
\end{remark}

\begin{lemma}
\label{lemma-affine-equivalence-algebras}
Soit $S$ un schéma. Il existe une anti-équivalence de catégories
$$
\begin{matrix}
\text{Schémas affines} \\
\text{sur }S
\end{matrix}
\longleftrightarrow
\begin{matrix}
\text{faisceaux quasi-cohérents} \\
\text{de }\mathcal{O}_S\text{-algèbres}
\end{matrix}
$$
qui associe à $f : X \to S$ le faisceau $f_*\mathcal{O}_X$. De plus, cette
équivalence est compatible avec un changement de base arbitraire.
\end{lemma}

\begin{proof}
Le foncteur de droite à gauche est donné par $\underline{\Spec}_S$.
Les deux foncteurs sont mutuellement
inverses par le lemme
\ref{lemma-characterize-affine} et Constructions, Lemme \ref{constructions-lemma-spec-properties} partie (3).
L'instruction
finale est Constructions, Lemme \ref{constructions-lemma-spec-properties} partie (2).
\end{proof}

\begin{lemma}
\label{lemma-quasi-coherence-scalar-restriction}
\begin{reference}
\cite[I, Proposition 9.6.1]{EGA}
\end{reference}
Soit $S$ un schéma et $\mathcal{A}$ une $\mathcal{O}_S$-algèbre
quasi-cohérente. Un $\mathcal{A}$-module est quasi-cohérent
comme un $\mathcal{O}_S$-module si et seulement si il
est quasi-cohérent comme un $\mathcal{A}$-module.
\end{lemma}

\begin{proof}
Soit $\mathcal{F}$ un module $\mathcal{A}$. Si $\mathcal{F}$
est quasi-cohérent en tant que module $\mathcal{A}$, alors pour chaque
$s \in S$ il existe un voisinage ouvert $U$ de $s$
et une suite exacte
$$
\bigoplus\nolimits_J\mathcal{A}|_U\to
\bigoplus\nolimits_I\mathcal{A}|_U\to
\mathcal{F}|_U
$$
de $\mathcal{A}|_U$-modules. Alors
c'est aussi une suite exacte de $\mathcal{O}_U$-modules. Ainsi
$\mathcal{F}|_U$ est quasi-cohérent comme le conoyau d'un morphisme de $\mathcal{O}_U$-modules
quasi-cohérents sur un schéma. Il s'ensuit que
$\mathcal{F}$ est quasi-cohérent en tant que $\mathcal{O}_X$-module.

\medskip\noindent
Inversement, supposons que $\mathcal{F}$ est quasi-cohérent en tant que module $\mathcal{O}_X$.
Choisissez un $\Spec(R) = U \subset S$ ouvert affine. Nous avons des isomorphismes
de $\mathcal{O}_U$-modules $\mathcal{A}|_U \cong \widetilde{A}$ et $\mathcal{F}|_U \cong \widetilde{M}$, pour certaines $R$-algèbre
$A$ et certains $R$-module $M$.
La structure de module $\mathcal{A}$ sur $\mathcal{F}$ se traduit par une structure
de module $A$ sur $M$ compatible avec la structure
de module $R$ donnée (détails omis). Choisissez une suite exacte
$$
\bigoplus\nolimits_J A \to
\bigoplus\nolimits_I A \to
M \to 0
$$
de $A$-modules. Puisque le foncteur $\widetilde{\ }$ est exact, cela
produit une suite exacte
$$
\bigoplus\nolimits_J \widetilde{A}\to
\bigoplus\nolimits_I \widetilde{A}\to
\widetilde{M} \to 0
$$
de $\widetilde{A}$-modules. Cela signifie que $\mathcal{F}$
est quasi-cohérent en tant que module $\mathcal{A}$.
\end{proof}

\begin{lemma}
\label{lemma-affine-equivalence-modules}
Soit $f : X \to S$ un morphisme affine de
schémas. Soit $\mathcal{A} = f_*\mathcal{O}_X$.
Le foncteur $\mathcal{F} \mapsto f_*\mathcal{F}$ induit une équivalence
de catégories
$$
\left\{
\begin{matrix}
\text{catégorie des}\\
\mathcal{O}_X\text{-modules quasi-cohérents}
\end{matrix}
\right\}
\longrightarrow
\left\{
\begin{matrix}
\text{catégorie des}\\
\mathcal{A}\text{-modules quasi-cohérents}
\end{matrix}
\right\}
$$
De plus, un $\mathcal{A}$-module est quasi-cohérent
comme un $\mathcal{O}_S$-module si et seulement si il
est quasi-cohérent comme un $\mathcal{A}$-module.
\end{lemma}

\begin{proof}
L'instruction finale est Lemme \ref{lemma-quasi-coherence-scalar-restriction}. Par le Lemme \ref{lemma-affine-separated}
et Schémas, Lemme \ref{schemes-lemma-push-forward-quasi-coherent}
l'image directe $f_*\mathcal{F}$ d'un $\mathcal{O}_X$-module
quasi-cohérent $\mathcal{F}$ est un $\mathcal{O}_S$-module
quasi-cohérent.
D'où un foncteur comme dans
l'énoncé du lemme.

\medskip\noindent Nous
allons construire un $h$ quasi-inverse à ce foncteur. Soit $\mathcal{G}$ un
module $\mathcal{A}$ quasi-cohérent. Ensuite, nous définissons
$$
h(\mathcal{G}) = f^*\mathcal{G} \otimes_{f^*\mathcal{A}} \mathcal{O}_X =
f^*\mathcal{G} \otimes_{f^*f_*\mathcal{O}_X} \mathcal{O}_X
$$
Élucidation : le retrait $f^*\mathcal{A} = f^*f_*\mathcal{O}_X$ est une $\mathcal{O}_X$-algèbre,
l'application d'adjonction $f^*f_*\mathcal{O}_X \to \mathcal{O}_X$ est
un homomorphisme d'algèbre et le retrait $f^*\mathcal{G}$ est un $f^*\mathcal{A}$-module.
Observez que $h(\mathcal{G})$ est quasi-cohérent car la quasi-cohérence
est préservée par les retraits et le changement
d'anneaux. Observons qu'il existe un morphisme fonctoriel
$$
h(f_*\mathcal{F}) =
f^*f_*\mathcal{F} \otimes_{f^*f_*\mathcal{O}_X} \mathcal{O}_X \to \mathcal{F}
$$
provenant du morphisme d'adjonction $f^*f_*\mathcal{F} \to \mathcal{F}$ et un morphisme
fonctoriel
$$
\mathcal{G} \to f_*h(\mathcal{G}) \quad\text{adjoint du morphisme}\quad
f^*\mathcal{G} \to f^*\mathcal{G} \otimes_{f^*f_*\mathcal{O}_X} \mathcal{O}_X
$$
qui envoie une section locale $s$ de $f^*\mathcal{F}$ à $s \otimes 1$. Pour terminer la
démonstration il suffit de montrer que ces applications sont des isomorphismes pour
$\mathcal{F}$ et $\mathcal{G}$ comme ci-dessus. Ceci peut être vérifié sur les membres d'un revêtement
affine, c'est-à-dire lorsque $X$ et $S$ sont affines.

\medskip\noindent L'observation
algébrique clé qui fait ce travail est la suivante : Soit
$R \to A$ un homomorphisme d'anneaux. Soit $N$ un module $A$. Alors
$$
(N \otimes_R A) \otimes_{(A \otimes_R A)} A = N
$$
À savoir, le côté gauche de cette égalité est l'effet de l'application de
$h$ au module $\widetilde{A}$ quasi-cohérent $\widetilde{N}$ sur $\Spec(R)$. Nous omettons
les détails.
\end{proof}

\begin{lemma}
\label{lemma-composition-affine}
La composition des morphismes affines est affine.
\end{lemma}

\begin{proof}
Soient $f : X \to Y$ et $g : Y \to Z$ des morphismes affines.
Soit $U \subset Z$ ouvert affine. Alors $g^{-1}(U)$ est affine
par hypothèse sur $g$. Sur quoi $f^{-1}(g^{-1}(U))$ est affine
par hypothèse sur $f$. Donc $(g \circ f)^{-1}(U)$ est affine.
\end{proof}

\begin{lemma}
\label{lemma-base-change-affine}
Le changement de base d'un morphisme affine est affine.
\end{lemma}

\begin{proof}
Soit $f : X \to S$ un morphisme affine. Soit $S' \to S$ n'importe quel
morphisme. Notons $f' : X_{S'} = S' \times_S X \to S'$ le changement de base de $f$.
Pour chaque $s' \in S'$ il existe un voisinage ouvert affine
$s' \in V \subset S'$ qui s'envoie dans un certain ouvert affine $U \subset S$.
Par hypothèse, $f^{-1}(U)$ est affine. D'après le
matériel de Schémas, section \ref{schemes-section-fibre-products}, nous voyons que
$f^{-1}(U)_V = V \times_U f^{-1}(U)$ est affine et égal à $(f')^{-1}(V)$. Cela prouve
que $S'$ a un recouvrement ouvert par affines dont
l'image réciproque par $f'$ est affine. Nous concluons
par le lemme \ref{lemma-characterize-affine} ci-dessus.
\end{proof}

\begin{lemma}
\label{lemma-closed-immersion-affine}
Une immersion fermée est affine.
\end{lemma}

\begin{proof}
La première indication
en est Schémas, Lemme \ref{schemes-lemma-closed-immersion-affine-case}. Voir Schémas,
Lemme \ref{schemes-lemma-closed-subspace-scheme} pour un énoncé
complet.
\end{proof}

\begin{lemma}
\label{lemma-affine-s-open}
Soit $X$ un
schéma. Soit $\mathcal{L}$ un module $\mathcal{O}_X$ inversible.
Soit $s \in \Gamma(X, \mathcal{L})$. Le morphisme
d'inclusion $j : X_s \to X$ est affine.
\end{lemma}

\begin{proof}
Cela résulte de Propriétés, Lemme \ref{properties-lemma-affine-cap-s-open} et la
définition.
\end{proof}

\begin{lemma}
\label{lemma-affine-permanence}
Supposons que $g : X \to Y$ soit un morphisme de schémas sur $S$.
\begin{enumerate}
\item Si $X$ est affine sur $S$ et que $\Delta : Y \to Y \times_S Y$ est affine, alors $g$
est affine.
\item Si $X$ est affine sur $S$ et que $Y$ est séparé sur $S$, alors
$g$ est affine.
\item Un morphisme d'un schéma affine vers un schéma à diagonale
affine est affine.
\item Un morphisme d'un schéma affine vers un schéma séparé est affine.
\end{enumerate}
\end{lemma}

\begin{proof}
Démonstration de (1). Le changement de base $X \times_S Y \to Y$ est affine
par le lemme \ref{lemma-base-change-affine}. Le morphisme
$(1, g) : X \to X \times_S Y$ est le changement de base de $Y \to Y \times_S Y$
par le morphisme $X \times_S Y \to Y \times_S Y$. Il est donc affine par le lemme
\ref{lemma-base-change-affine}. La
composition des morphismes affines
est affine (voir Lemme \ref{lemma-composition-affine})
et (1) s'ensuit. La partie (2) découle de (1) car une
immersion fermée est affine (voir Lemme \ref{lemma-closed-immersion-affine})
et $Y/S$ séparé signifie $\Delta$ est une immersion
fermée. Les parties (3) et (4) sont des cas particuliers
de (1) et (2).
\end{proof}

\begin{lemma}
\label{lemma-morphism-affines-affine}
Un morphisme entre schémas affines est affine.
\end{lemma}

\begin{proof}
Immédiat du lemme \ref{lemma-affine-permanence} avec $S = \Spec(\mathbf{Z})$.
Cela découle également directement de l'équivalence
de (1) et (2) dans le lemme \ref{lemma-characterize-affine}.
\end{proof}

\begin{lemma}
\label{lemma-Artinian-affine}
Soit $S$ un schéma.
Soit $A$ un anneau artinien.
Tout morphisme $\Spec(A) \to S$ est affine.
\end{lemma}

\begin{proof}
Omis.
\end{proof}

\begin{lemma}
\label{lemma-get-affine}
Soit $j : Y \to X$ une immersion de schémas. Supposons
qu'il existe un $U \subset X$ ouvert avec un complément
$Z = X \setminus U$ tel que
\begin{enumerate}
\item $U \to X$ est affine,
\item $j^{-1}(U) \to U$ est affine et
\item $j(Y) \cap Z$ est fermé.
\end{enumerate}
Alors $j$ est affine. En particulier, si $X$ est affine, $Y$ l'est également.
\end{lemma}

\begin{proof}
Par Schémas, Définition \ref{schemes-definition-immersion} il existe un sous-schéma ouvert $W \subset X$
tel que $j$ se factorise en une immersion fermée $i : Y \to W$
suivi du morphisme d'inclusion $W \to X$. Puisqu'une
immersion fermée est affine (Lemme
\ref{lemma-closed-immersion-affine}), on voit que pour tout $x \in W$
il existe un voisinage ouvert affine de $x$
dans $X$ dont l'image réciproque par $j$ est affine.
Si $x \in U$, alors la même chose est vraie par l'hypothèse
(2). Enfin, supposons $x \in Z$ et $x \not \in W$. Puis $x \not \in j(Y) \cap Z$.
Par l'hypothèse (3) on peut trouver un voisinage affine
ouvert $V \subset X$ de $x$ qui ne satisfait pas
à $j(Y) \cap Z$. Puis $j^{-1}(V) = j^{-1}(V \cap U)$ qui est affine
par les hypothèses (1) et (2). Il s'ensuit que $j$
est affine par le lemme \ref{lemma-characterize-affine}.
\end{proof}







\section{Familles de modules inversibles amples}
\label{section-families-ample-invertible-modules}

\noindent
Une courte section sur la notion de famille de modules amples inversibles.

\begin{definition}
\label{definition-family-ample-invertible-modules}
\begin{reference}
\cite[II Définition 2.2.4]{SGA6}
\end{reference}
Soit $X$ un schéma. Soit $\{\mathcal{L}_i\}_{i \in I}$ une famille de $\mathcal{O}_X$-modules
inversibles. Nous disons que $\{\mathcal{L}_i\}_{i \in I}$ est
une {\it grande famille
de modules inversibles sur $X$} si
\begin{enumerate}
\item $X$ est quasi-compact, et
\item pour chaque $x \in X$ il existe un $i \in I$, un $n \geq 1$
et $s \in \Gamma(X, \mathcal{L}_i^{\otimes n})$ de telle sorte que $x \in X_s$
et $X_s$ soient affines.
\end{enumerate}
\end{definition}

\noindent
Si $\{\mathcal{L}_i\}_{i \in I}$ est une ample famille de modules inversibles sur un
schéma $X$, alors il existe un sous-ensemble fini $I' \subset I$ tel que
$\{\mathcal{L}_i\}_{i \in I'}$ est une ample famille de modules inversibles sur $X$
(découle immédiatement de la quasi-compacité). Un schéma
ayant une ample famille de modules inversibles a une diagonale
affine par le lemme suivant et est donc a fortiori quasi-séparé.

\begin{lemma}
\label{lemma-affine-diagonal}
Soit $X$ un schéma tel que pour tout point $x \in X$ il existe
un $\mathcal{O}_X$-module $\mathcal{L}$ inversible et une section
globale $s \in \Gamma(X, \mathcal{L})$ telle que $x \in X_s$ et $X_s$ sont affines.
Alors la diagonale de $X$ est un morphisme affine.
\end{lemma}

\begin{proof}
Étant donné les $\mathcal{O}_X$-modules $\mathcal{L}$, $\mathcal{M}$
et les sections globales $s \in \Gamma(X, \mathcal{L})$,
$t \in \Gamma(X, \mathcal{M})$ tels que $X_s$ et $X_t$ sont
affines, nous devons prouver que $X_s \cap X_t$
est affine. À savoir, alors
Lemme \ref{lemma-characterize-affine} s'applique à $\Delta : X \to X \times X$
et le fait que $\Delta^{-1}(X_s \times X_t) = X_s \cap X_t$ montre que $\Delta$
est affine. Le fait que $X_s \cap X_t$ soit affine
découle de Propriétés, Lemme \ref{properties-lemma-affine-cap-s-open}.
\end{proof}

\begin{remark}
\label{remark-affine-s-opens-cover-family}
Dans Propriétés, Lemme \ref{properties-lemma-affine-s-opens-cover-quasi-separated}, on voit qu'un schéma qui possède un module inversible ample
est séparé. Ceci est faux pour les schémas possédant
une famille ample de modules inversibles. En effet,
soit $X$ comme dans Schémas, Exemple \ref{schemes-example-affine-space-zero-doubled},
avec $n = 1$, c'est-à-dire la droite affine à zéro doublée.
Nous utilisons la notation de cet exemple sauf que nous écrivons $x$
pour $x_1$ et $y$ pour $y_1$. Il existe,
pour tout entier $n$, un faisceau inversible $\mathcal{L}_n$
sur $X$ qui est trivial sur $X_1$ et $X_2$
et dont la fonction de transition $U_{12} \to U_{21}$ est $f(x) \mapsto y^n f(y)$.
Les sections globales de $\mathcal{L}_n$ sont
des couples $(f(x), g(y)) \in k[x] \oplus
k[y]$ tels que $y^n f(y) = g(y)$. Les sections $s = (1,
y)$ de $\mathcal{L}_1$ et $t = (x, 1)$ de
$\mathcal{L}_{-1}$ déterminent un recouvrement par des ouverts affines, car
$X_s = X_1$ et $X_t = X_2$. Ainsi, $X$ possède
une famille ample de modules inversibles, mais il n'est pas séparé.
\end{remark}









\section{Morphismes quasi-affines}
\label{section-quasi-affine}

\noindent
Rappelons qu'un schéma $X$ est dit {\it quasi-affine}
s'il est quasi-compact et isomorphe à un sous-schéma ouvert
d'un schéma affine, voir Propriétés, Définition \ref{properties-definition-quasi-affine}.

\begin{definition}
\label{definition-quasi-affine}
Un morphisme de schémas $f : X \to S$ est appelé {\it quasi-affine} si
l'image réciproque de chaque ouvert affine de $S$ est un schéma quasi-affine.
\end{definition}

\begin{lemma}
\label{lemma-quasi-affine-separated}
Un morphisme quasi-affine est séparé et quasi-compact.
\end{lemma}

\begin{proof}
Soit $f : X \to S$ quasi-affine. La
quasi-compacité est immédiate
chez Schémas, Lemme \ref{schemes-lemma-quasi-compact-affine}. Soit $U \subset S$ un
ouvert affine. Si l'on peut montrer que $f^{-1}(U)$
est un schéma séparé, alors $f$ est séparé
(Schémas, Lemme \ref{schemes-lemma-characterize-separated} montre que la séparation
est locale sur la base). Par hypothèse $f^{-1}(U)$
est isomorphe à un sous-schéma ouvert d'un schéma
affine. Un schéma affine est séparé et donc chaque sous-schéma
ouvert d'un schéma affine est séparé à volonté.
\end{proof}

\begin{lemma}
\label{lemma-characterize-quasi-affine}
Soit $f : X \to S$ un morphisme de schémas. Les conditions
suivantes sont équivalentes
\begin{enumerate}
\item Le morphisme $f$ est quasi-affine.
\item Il existe un recouvrement ouvert affine $S = \bigcup W_j$ tel
que chaque $f^{-1}(W_j)$ soit quasi-affine.
\item Il existe un faisceau quasi-cohérent de $\mathcal{O}_S$-algèbres $\mathcal{A}$
et une immersion ouverte quasi-compacte
$$
\xymatrix{
X \ar[rr] \ar[rd] & & \underline{\Spec}_S(\mathcal{A}) \ar[dl] \\
& S &
}
$$
sur $S$.
\item Idem qu'en (3) mais avec $\mathcal{A} = f_*\mathcal{O}_X$ et
la flèche horizontale le morphisme
canonique des Constructions, Lemme \ref{constructions-lemma-canonical-morphism}.
\end{enumerate}
\end{lemma}

\begin{proof}
Il est évident que (1) implique (2) et que (4) implique (3).

\medskip\noindent
Supposons que $S = \bigcup_{j \in J} W_j$ soit un recouvrement ouvert
affine tel que chaque $f^{-1}(W_j)$
soit quasi-affine. Par Schémas, Lemme \ref{schemes-lemma-quasi-compact-affine}
on voit que $f$
est quasi-compact. Par Schémas, Lemme \ref{schemes-lemma-characterize-quasi-separated}
on voit que le morphisme $f$ est
quasi-séparé. Ainsi par Schémas, Lemme \ref{schemes-lemma-push-forward-quasi-coherent}
le faisceau $\mathcal{A} = f_*\mathcal{O}_X$ est un faisceau quasi-cohérent
de $\mathcal{O}_X$-algèbres. On a donc le schéma
$g : Y = \underline{\Spec}_S(\mathcal{A}) \to S$ sur $S$. L'application
d'identité
$\text{id} : \mathcal{A} = f_*\mathcal{O}_X \to f_*\mathcal{O}_X$ fournit, via la définition du spectre relatif,
un morphisme $can : X \to Y$ sur $S$,
voir Constructions, Lemme
\ref{constructions-lemma-canonical-morphism}. Par hypothèse, le lemme qui vient d'être
cité, et Propriétés, Lemme
\ref{properties-lemma-quasi-affine} la restriction $can|_{f^{-1}(W_j)} : f^{-1}(W_j) \to g^{-1}(W_j)$ est
une immersion ouverte quasi-compacte. Ainsi $can$
est une immersion ouverte quasi-compacte. Nous
avons montré que (2) implique (4).

\medskip\noindent
Supposons (3). Choisissez n’importe quel $U \subset S$
ouvert affine. Par Constructions, Lemme \ref{constructions-lemma-spec-properties} on voit que l'image
inverse de $U$ dans le spectre relatif est affine.
Nous concluons donc que $f^{-1}(U)$ est quasi-affine (notez
que la quasi-compacité est également codée dans (3)). Ainsi
(3) implique (1).
\end{proof}

\begin{lemma}
\label{lemma-composition-quasi-affine}
La composition des morphismes quasi-affines est quasi-affine.
\end{lemma}

\begin{proof}
Soient $f : X \to Y$ et $g : Y \to Z$ des morphismes quasi-affines.
Soit $U \subset Z$ ouvert affine. Alors $g^{-1}(U)$
est quasi-affine par hypothèse sur $g$.
Soit $j : g^{-1}(U) \to V$ une immersion ouverte
quasi-compacte dans un schéma affine
$V$. Par le Lemme \ref{lemma-characterize-quasi-affine}
ci-dessus on voit que $f^{-1}(g^{-1}(U))$ est un
sous-schéma ouvert quasi-compact du
spectre relatif $\underline{\Spec}_{g^{-1}(U)}(\mathcal{A})$ pour un faisceau quasi-cohérent
de $\mathcal{O}_{g^{-1}(U)}$-algèbres
$\mathcal{A}$. Par Schémas, Lemme \ref{schemes-lemma-push-forward-quasi-coherent}
le faisceau $\mathcal{A}' = j_*\mathcal{A}$
est un faisceau quasi-cohérent de $\mathcal{O}_V$-algèbres
avec la propriété que $j^*\mathcal{A}' = \mathcal{A}$. On obtient
donc un diagramme commutatif
$$
\xymatrix{
f^{-1}(g^{-1}(U)) \ar[r] &
\underline{\Spec}_{g^{-1}(U)}(\mathcal{A})
\ar[r] \ar[d] &
\underline{\Spec}_V(\mathcal{A}') \ar[d] \\
& g^{-1}(U) \ar[r]^j & V
}
$$
avec le carré étant un carré de
fibre, voir Constructions, Lemme \ref{constructions-lemma-spec-properties}. Notez que
le coin supérieur droit est un schéma affine.
Donc $(g \circ f)^{-1}(U)$ est quasi-affine.
\end{proof}

\begin{lemma}
\label{lemma-base-change-quasi-affine}
Le changement de base d'un morphisme quasi-affine est quasi-affine.
\end{lemma}

\begin{proof}
Soit $f : X \to S$ un morphisme quasi-affine.
Par le Lemme \ref{lemma-characterize-quasi-affine} ci-dessus on trouve un
faisceau quasi-cohérent de $\mathcal{O}_S$-algèbres
$\mathcal{A}$ et une immersion ouverte quasi-compacte
$X \to \underline{\Spec}_S(\mathcal{A})$ sur $S$. Soit $g : S' \to S$
n'importe
quel morphisme. Notons
$f' : X_{S'} = S' \times_S X \to S'$ le changement de base de $f$. Puisque le
changement de base d'une immersion ouverte quasi-compacte
est une immersion ouverte quasi-compacte
on voit que $X_{S'} \to \underline{\Spec}_{S'}(g^*\mathcal{A})$ est une immersion
ouverte quasi-compacte
(nous avons utilisé
Schémas, Lemmes \ref{schemes-lemma-quasi-compact-preserved-base-change} et \ref{schemes-lemma-base-change-immersion}
et Constructions,
Lemme \ref{constructions-lemma-spec-properties}). Par le Lemme \ref{lemma-characterize-quasi-affine}
encore Nous concluons que $X_{S'} \to S'$
est quasi-affine.
\end{proof}

\begin{lemma}
\label{lemma-quasi-compact-immersion-quasi-affine}
Une immersion quasi-compacte est quasi-affine.
\end{lemma}

\begin{proof}
Soit $X \to S$ une immersion quasi-compacte. Il faut montrer que l'image
l'image réciproque de tout ouvert affine est quasi-affine. Ainsi, en
supposant que $S$ est un schéma affine, il faut
montrer que $X$ est quasi-affine. Par le Lemme \ref{lemma-quasi-compact-immersion} le morphisme
$X \to S$ se factorise en $X \to Z \to S$ où $Z$ est un sous-schéma fermé
de $S$ et $X \subset Z$ est un ouvert quasi-compact. Puisque
$S$ est affine, le Lemme \ref{lemma-closed-immersion} implique que $Z$ est
affine, ce qui conclut.
\end{proof}

\begin{lemma}
\label{lemma-affine-quasi-affine}
Soit $S$ un schéma. Soit $X$ un schéma affine.
Un morphisme $f : X \to S$ est quasi-affine si et seulement si il est quasi-compact.
En particulier, tout morphisme d'un schéma affine à un schéma quasi-séparé
est quasi-affine.
\end{lemma}

\begin{proof}
Soit $V \subset S$ un ouvert affine. Alors $f^{-1}(V)$ est un sous-schéma
ouvert du schéma affine $X$, donc quasi-affine si et
seulement si il est quasi-compact. Cela prouve la première affirmation.
La quasi-compacité de tout $f : X \to S$ où $X$ est affine
et $S$ quasi-séparé découle d'après Schémas, Lemme
\ref{schemes-lemma-quasi-compact-permanence} appliqué à $X \to S \to \Spec(\mathbf{Z})$.
\end{proof}

\begin{lemma}
\label{lemma-quasi-affine-permanence}
Supposons que $g : X \to Y$ soit un morphisme de schémas sur $S$.
Si $X$ est quasi-affine sur $S$ et $Y$ est quasi-séparé sur
$S$, alors $g$ est quasi-affine. En particulier, tout morphisme
d'un schéma quasi-affine à un schéma quasi-séparé est quasi-affine.
\end{lemma}

\begin{proof}
Le changement de base $X \times_S Y \to Y$ est quasi-affine
par le lemme \ref{lemma-base-change-quasi-affine}. Le morphisme
$X \to X \times_S Y$ est une immersion
quasi-compacte comme $Y \to S$ est quasi-séparé,
voir Schémas, Lemme \ref{schemes-lemma-section-immersion}. Une
immersion quasi-compacte est quasi-affine
par le lemme \ref{lemma-quasi-compact-immersion-quasi-affine} et la composition
des morphismes quasi-affines est quasi-affine
(voir Lemme \ref{lemma-composition-quasi-affine}). Ainsi on obtient le résultat.
\end{proof}











\section{Types de morphismes définis par des propriétés d'homomorphismes d'anneaux}
\label{section-properties-ring-maps}

\noindent
Dans cette section nous étudions quelles propriétés des homomorphismes d'anneaux
permettent de définir des propriétés locales des morphismes de schémas.

\begin{definition}
\label{definition-property-local}
Soit $P$ une propriété d'homomorphismes d'anneaux.
\begin{enumerate}
\item On dit que $P$ est {\it local} si ce qui suit est valable :
\begin{enumerate}
\item Pour tout homomorphisme d'anneaux $R \to A$, et
tout $f \in R$ on a $P(R \to A) \Rightarrow P(R_f \to A_f)$.
\item Pour tous les anneaux $R$, $A$, tout $f \in R$, $a\in A$
et tout homomorphisme d'anneaux $R_f \to A$ nous avons $P(R_f \to A) \Rightarrow P(R \to A_a)$.
\item Pour tout homomorphisme d'anneaux $R \to A$,
et $a_i \in A$
tel que $(a_1, \ldots, a_n) = A$ puis $\forall i, P(R \to A_{a_i}) \Rightarrow P(R \to A)$.
\end{enumerate}
\item On dit que $P$ est {\it stable sous changement
de base} si pour tout homomorphismes
d'anneaux $R \to A$, $R \to R'$ on a $P(R \to A) \Rightarrow P(R' \to R' \otimes_R A)$.
\item On dit que $P$ est {\it stable sous composition}
si pour tout homomorphismes
d'anneaux $A \to B$, $B \to C$ on a $P(A \to B) \wedge P(B \to C) \Rightarrow P(A \to C)$.
\end{enumerate}
\end{definition}

\begin{definition}
\label{definition-locally-P}
Soit $P$ une propriété d'homomorphismes
d'anneaux. Soit $f : X \to S$ un morphisme de schémas. On dit
que $f$ est {\it localement de type $P$}
si pour tout $x \in X$ il existe un voisinage ouvert affine $U$
de $x$ dans $X$ qui s'envoie dans un voisinage ouvert affine
$V \subset S$ tel que l'homomorphisme d'anneaux induit $\mathcal{O}_S(V) \to \mathcal{O}_X(U)$ possède la
propriété $P$.
\end{definition}

\noindent
Il ne s'agit pas d'une ``bonne définition'' sauf si la propriété $P$
est une propriété locale. Même si $P$ est une propriété locale, nous n'utiliserons
pas automatiquement cette définition pour dire qu'un morphisme est ``localement
de type $P$'' à moins que nous n'indiquions également explicitement
la définition ailleurs.

\begin{lemma}
\label{lemma-locally-P}
Soit $f : X \to S$ un morphisme de schémas. Soit $P$
une propriété des homomorphismes d'anneaux.
Soit $U$ un ouvert affine
de $X$, et $V$ un ouvert affine de
$S$ telle que
$f(U) \subset V$. Si $f$ est localement de type $P$
et que $P$ est local, alors $P(\mathcal{O}_S(V) \to \mathcal{O}_X(U))$ est valable.
\end{lemma}

\begin{proof}
Comme $f$ est localement de type $P$ pour
chaque $u \in U$, il existe un mappage $U_u \subset X$ ouvert affine
dans un $V_u \subset S$ ouvert affine tel que $P(\mathcal{O}_S(V_u) \to \mathcal{O}_X(U_u))$
est valable. Choisissez un voisinage ouvert $U'_u \subset U \cap U_u$
de $u$ qui est ouvert affine standard
dans $U$ et $U_u$, voir Schémas, Lemme
\ref{schemes-lemma-standard-open-two-affines}. Par Définition \ref{definition-property-local} (1)(b)
nous voyons que $P(\mathcal{O}_S(V_u) \to \mathcal{O}_X(U'_u))$ est vrai. Nous pouvons donc
supposer que $U_u \subset U$ est un ouvert affine standard.
Choisissez un voisinage ouvert $V'_u \subset V \cap V_u$ de
$f(u)$ qui est ouvert affine standard dans $V$
et $V_u$, voir Schémas, Lemme \ref{schemes-lemma-standard-open-two-affines}.
Alors $U'_u = f^{-1}(V'_u) \cap U_u$ est un ouvert affine standard de $U_u$
(donc de $U$) et
on a $P(\mathcal{O}_S(V'_u) \to \mathcal{O}_X(U'_u))$ par définition \ref{definition-property-local} (1)(a).
Par conséquent, nous pouvons supposer
que $U_u \subset U$ et $V_u \subset V$ sont tous deux affines
ouverts standard. En appliquant
la Définition \ref{definition-property-local} (1)(b) une fois
de plus Nous concluons que $P(\mathcal{O}_S(V) \to \mathcal{O}_X(U_u))$ est valable. Parce que $U$
est quasi-compact, nous pouvons choisir un nombre
fini de points $u_1, \ldots, u_n \in U$ tel que
$$
U = U_{u_1} \cup \ldots \cup U_{u_n}.
$$
Par Définition \ref{definition-property-local} (1)(c)
Nous concluons que $P(\mathcal{O}_S(V) \to \mathcal{O}_X(U))$ est vrai.
\end{proof}

\begin{lemma}
\label{lemma-locally-P-characterize}
Soit $P$ une propriété locale des homomorphismes
d'anneaux. Soit $f : X \to S$ un morphisme de schémas. Les conditions
suivantes sont équivalentes
\begin{enumerate}
\item Le morphisme $f$ est localement de type $P$.
\item Pour chaque ouverture affine $U \subset X$,
$V \subset S$ avec $f(U) \subset V$ nous avons $P(\mathcal{O}_S(V) \to \mathcal{O}_X(U))$.
\item Il existe un recouvrement ouvert $S = \bigcup_{j \in J} V_j$ et des
recouvrements ouverts $f^{-1}(V_j) = \bigcup_{i \in I_j} U_i$ tels que chacun des
morphismes $U_i \to V_j$, $j\in J, i\in I_j$ est localement de
type $P$.
\item Il existe un recouvrement ouvert affine $S = \bigcup_{j \in J} V_j$
et un recouvrement ouvert affine $f^{-1}(V_j) = \bigcup_{i \in I_j} U_i$
tel que $P(\mathcal{O}_S(V_j) \to \mathcal{O}_X(U_i))$ est valable, pour tout
$j\in J, i\in I_j$.
\end{enumerate}
De plus, si $f$ est localement de type $P$
alors pour tout sous-schémas ouverts $U \subset X$, $V \subset S$ avec $f(U) \subset V$
la restriction $f|_U : U \to V$ est localement de type $P$.
\end{lemma}

\begin{proof}
Cela résulte du lemme \ref{lemma-locally-P} ci-dessus.
\end{proof}

\begin{lemma}
\label{lemma-composition-type-P}
Soit $P$ une propriété d'homomorphismes d'anneaux.
Supposons que $P$ soit local et stable en termes de composition.
La composition des morphismes localement de type $P$ est localement
de type $P$.
\end{lemma}

\begin{proof}
Soit $f : X \to Y$ et $g : Y \to Z$ des morphismes localement de type
$P$. Soit $x \in X$. Choisir un voisinage affine
ouvert $W \subset Z$ de $g(f(x))$. Choisir un voisinage affine
ouvert $V \subset g^{-1}(W)$ de $f(x)$. Choisir un voisinage affine
ouvert $U \subset f^{-1}(V)$ de $x$. Par le Lemme \ref{lemma-locally-P-characterize} les
homomorphismes d'anneaux $\mathcal{O}_Z(W) \to \mathcal{O}_Y(V)$ et
$\mathcal{O}_Y(V) \to \mathcal{O}_X(U)$ satisfont $P$. Par conséquent,
$\mathcal{O}_Z(W) \to \mathcal{O}_X(U)$ satisfait à $P$ car $P$
est supposé stable sous composition.
\end{proof}

\begin{lemma}
\label{lemma-base-change-type-P}
Soit $P$ une propriété d'homomorphismes d'anneaux.
Supposons que $P$ soit local et stable sous changement de base.
Le changement de base d'un morphisme localement de type $P$ est
localement de type $P$.
\end{lemma}

\begin{proof}
Soit $f : X \to S$ un morphisme localement de type $P$.
Soit $S' \to S$ n'importe quel
morphisme. Notons $f' : X_{S'} = S' \times_S X \to S'$ le changement de base de
$f$. Pour chaque $s' \in S'$ il existe un voisinage affine
ouvert $s' \in V' \subset S'$ qui s'envoie dans un ouvert affine $V \subset S$.
D'après le lemme \ref{lemma-locally-P-characterize}, l'ouvert $f^{-1}(V)$ est une
réunion d'ouverts affines $U_i$ telle que les homomorphismes d'anneaux
$\mathcal{O}_S(V) \to \mathcal{O}_X(U_i)$ vérifient tous $P$.
D'après Schémas, section \ref{schemes-section-fibre-products},
$f^{-1}(U)_{V'} = V' \times_V f^{-1}(V)$ est la réunion
des ouverts affines $V' \times_V U_i$.
Comme $\mathcal{O}_{X_{S'}}(V' \times_V U_i) =
\mathcal{O}_{S'}(V') \otimes_{\mathcal{O}_S(V)} \mathcal{O}_X(U_i)$,
nous voyons que les homomorphismes d'anneaux
$\mathcal{O}_{S'}(V') \to \mathcal{O}_{X_{S'}}(V' \times_V U_i)$
vérifient $P$, puisque $P$ est supposée stable par changement de base.
\end{proof}

\begin{lemma}
\label{lemma-properties-local}
Les propriétés suivantes d'un homomorphisme d'anneaux $R \to A$ sont locales.
\begin{enumerate}
\item (Isomorphisme sur anneaux locaux.)
Pour tout premier $\mathfrak q$ de $A$ couché sur $\mathfrak p \subset R$
l'homomorphisme d'anneaux $R \to A$
induit un isomorphisme $R_{\mathfrak p} \to A_{\mathfrak q}$.
\item (Immersion ouverte.)
Pour tout $\mathfrak q$ premier de $A$ il existe un
$f \in R$, $\varphi(f) \not \in \mathfrak q$ tel que l'homomorphisme d'anneaux $\varphi : R \to A$ induit
un isomorphisme $R_f \to A_f$.
\item (Fibres réduites.)
Pour chaque nombre premier $\mathfrak p$ de $R$,
l'anneau de fibres $A \otimes_R \kappa(\mathfrak p)$ est réduit.
\item (Fibres de dimension au plus $n$.)
Pour chaque nombre premier $\mathfrak p$ de $R$ l'anneau
de fibres $A \otimes_R \kappa(\mathfrak p)$ a une dimension de Krull au plus $n$.
\item (Localement noéthérien sur la cible.) L'homomorphisme
d'anneaux $R \to A$ a la propriété que $A$ est noéthérien.
\item Ajoutez-en davantage ici si nécessaire\footnote{Mais uniquement les propriétés
qui ne sont pas déjà traitées séparément ailleurs.}.
\end{enumerate}
\end{lemma}

\begin{proof}
Omis.
\end{proof}

\begin{lemma}
\label{lemma-properties-base-change}
Les propriétés suivantes des homomorphismes d'anneaux sont stables sous changement de base.
\begin{enumerate}
\item (Isomorphisme sur anneaux locaux.)
Pour tout premier $\mathfrak q$ de $A$ couché sur $\mathfrak p \subset R$
l'homomorphisme d'anneaux $R \to A$
induit un isomorphisme $R_{\mathfrak p} \to A_{\mathfrak q}$.
\item (Immersion ouverte.)
Pour tout $\mathfrak q$ premier de $A$ il existe un
$f \in R$, $\varphi(f) \not \in \mathfrak q$ tel que l'homomorphisme d'anneaux $\varphi : R \to A$ induit
un isomorphisme $R_f \to A_f$.
\item Ajoutez-en davantage ici si nécessaire\footnote{Mais uniquement les propriétés
qui ne sont pas déjà traitées séparément ailleurs.}.
\end{enumerate}
\end{lemma}

\begin{proof}
Omis.
\end{proof}

\begin{lemma}
\label{lemma-properties-composition}
Les propriétés suivantes des homomorphismes d'anneaux sont stables sous composition.
\begin{enumerate}
\item (Isomorphisme sur anneaux locaux.)
Pour tout premier $\mathfrak q$ de $A$ couché sur $\mathfrak p \subset R$
l'homomorphisme d'anneaux $R \to A$
induit un isomorphisme $R_{\mathfrak p} \to A_{\mathfrak q}$.
\item (Immersion ouverte.)
Pour tout $\mathfrak q$ premier de $A$ il existe un
$f \in R$, $\varphi(f) \not \in \mathfrak q$ tel que l'homomorphisme d'anneaux $\varphi : R \to A$ induit
un isomorphisme $R_f \to A_f$.
\item (Localement noethérien sur la cible.) L'homomorphisme
d'anneaux $R \to A$ a la propriété que $A$ est noéthérien.
\item Ajoutez-en davantage ici si nécessaire\footnote{Mais uniquement les propriétés
qui ne sont pas déjà traitées séparément ailleurs.}.
\end{enumerate}
\end{lemma}

\begin{proof}
Omis.
\end{proof}








\section{Morphismes de type fini}
\label{section-finite-type}

\noindent
Rappelons qu'un homomorphisme d'anneaux $R \to A$ est dit de type
fini si $A$ est isomorphe à un quotient de $R[x_1, \ldots, x_n]$ comme une $R$-algèbre,
voir Algèbre, Définition \ref{algebra-definition-finite-type}.

\begin{definition}
\label{definition-finite-type}
Soit $f : X \to S$ un morphisme de schémas.
\begin{enumerate}
\item On dit que $f$ est de type {\it fini
à $x \in X$} s'il existe un voisinage ouvert affine $\Spec(A) = U \subset X$
de $x$ et un ouvert affine $\Spec(R) = V \subset S$
à $f(U) \subset V$ tel que l'homomorphisme d'anneaux induit
$R \to A$ est de type fini.
\item On dit que $f$ est {\it localement de type fini}
s'il est de type fini en tout point de $X$.
\item On dit que $f$ est de type {\it fini} s'il est localement
de type fini et quasi-compact.
\end{enumerate}
\end{definition}

\begin{lemma}
\label{lemma-locally-finite-type-characterize}
Soit $f : X \to S$ un morphisme de schémas. Les conditions
suivantes sont équivalentes
\begin{enumerate}
\item Le morphisme $f$ est localement de type fini.
\item Pour tous les ouverts affines $U \subset X$, $V \subset S$
avec $f(U) \subset V$ l'homomorphisme
d'anneaux $\mathcal{O}_S(V) \to \mathcal{O}_X(U)$ est de type fini.
\item Il existe un recouvrement ouvert $S = \bigcup_{j \in J} V_j$ et
des recouvrements ouverts $f^{-1}(V_j) = \bigcup_{i \in I_j} U_i$ tels que chacun
des morphismes $U_i \to V_j$, $j\in J, i\in I_j$ est localement
de type fini.
\item Il existe un recouvrement ouvert affine $S = \bigcup_{j \in J} V_j$
et des recouvrements ouverts affines $f^{-1}(V_j) = \bigcup_{i \in I_j} U_i$ tels
que l'homomorphisme d'anneaux $\mathcal{O}_S(V_j) \to \mathcal{O}_X(U_i)$ soit de
type fini, pour tout $j\in J, i\in I_j$.
\end{enumerate}
De plus, si $f$ est localement de type fini alors
pour tout sous-schémas ouverts $U \subset X$, $V \subset S$ avec $f(U) \subset V$
la restriction $f|_U : U \to V$ est localement de type fini.
\end{lemma}

\begin{proof}
Cela résulte du lemme \ref{lemma-locally-P} si l'on montre que la propriété
``$R \to A$ est de type fini''
est locale. On vérifie les conditions (a), (b)
et (c) de la Définition \ref{definition-property-local}.
Par Algèbre, Lemme \ref{algebra-lemma-base-change-finiteness}, la propriété d'être de type fini est
stable par changement de base ; on en déduit
la condition (a). Par Algèbre, Lemme
\ref{algebra-lemma-compose-finite-type}, elle est stable par composition,
et pour tout anneau $R$ l'homomorphisme
d'anneaux $R \to R_f$ est de type fini.
On en déduit la condition (b). Enfin, la
condition (c) résulte d'Algèbre, Lemme \ref{algebra-lemma-cover-upstairs}.
\end{proof}

\begin{lemma}
\label{lemma-composition-finite-type}
La composition de deux morphismes qui sont localement de type fini est localement
de type fini. Il en va de même pour les morphismes de type fini.
\end{lemma}

\begin{proof}
Dans la démonstration du lemme \ref{lemma-locally-finite-type-characterize} nous avons vu qu'être
de type fini est une propriété locale des homomorphismes
d'anneaux. D'où la première affirmation du lemme
découle du lemme \ref{lemma-composition-type-P} combinée
au fait qu'être de type fini est une propriété des homomorphismes
d'anneaux qui est stable sous
composition, voir Algèbre, Lemme \ref{algebra-lemma-compose-finite-type}. Par
ce qui précède et le fait que les compositions
de morphismes quasi-compacts sont
quasi-compacts, voir Schémas, Lemme \ref{schemes-lemma-composition-quasi-compact} on voit
que la composition de morphismes de type fini est de
type fini.
\end{proof}

\begin{lemma}
\label{lemma-base-change-finite-type}
Le changement de base d'un morphisme qui est localement de type fini
est localement de type fini. Il en va de même pour les morphismes de
type fini.
\end{lemma}

\begin{proof}
Dans la démonstration du lemme \ref{lemma-locally-finite-type-characterize} nous avons vu qu'être
de type fini est une propriété locale des homomorphismes
d'anneaux. D'où le premier énoncé du lemme découle
du lemme \ref{lemma-base-change-type-P} combiné au fait
qu'être de type fini est une propriété des homomorphismes d'anneaux
qui est stable par changement
de base, voir Algèbre, Lemme \ref{algebra-lemma-base-change-finiteness}. Par
ce qui précède et le fait qu'un changement de
base d'un morphisme quasi-compact est
quasi-compact, voir Schémas, Lemme \ref{schemes-lemma-quasi-compact-preserved-base-change} on voit que le changement
de base d'un morphisme de type fini est un morphisme
de type fini.
\end{proof}

\begin{lemma}
\label{lemma-immersion-locally-finite-type}
Une immersion fermée est de type fini. Une
immersion est localement de type fini.
\end{lemma}

\begin{proof}
Ceci est vrai car une immersion ouverte est un isomorphisme local,
et une immersion fermée est évidemment de type fini.
\end{proof}

\begin{lemma}
\label{lemma-finite-type-noetherian}
Soit $f : X \to S$ un morphisme. Si $S$
est (localement) noéthérien et $f$ (localement) de type fini alors
$X$ est (localement) noéthérien.
\end{lemma}

\begin{proof}
Cela découle immédiatement du fait qu'un anneau
de type fini sur un anneau noéthérien est noéthérien,
voir Algèbre, Lemme \ref{algebra-lemma-Noetherian-permanence}. (Aussi : utiliser
le fait que la source d'un morphisme quasi-compact à cible
quasi-compacte est quasi-compacte.)
\end{proof}

\begin{lemma}
\label{lemma-finite-type-Noetherian-quasi-separated}
Soit $f : X \to S$ localement de type fini avec $S$ localement noethérien. Alors
$f$ est quasi-séparé.
\end{lemma}

\begin{proof}
En fait, il est vrai que $X$ est
quasi-séparé, voir Propriétés, Lemme \ref{properties-lemma-locally-Noetherian-quasi-separated} et Lemme
\ref{lemma-finite-type-noetherian} ci-dessus. Appliquez ensuite
Schémas, Lemme \ref{schemes-lemma-compose-after-separated} pour conclure que
$f$ est quasi-séparé.
\end{proof}

\begin{lemma}
\label{lemma-permanence-finite-type}
Soit $X \to Y$ un morphisme de schémas sur un schéma de base $S$. Si
$X$ est localement de type fini sur $S$, alors $X \to Y$
est localement de type fini.
\end{lemma}

\begin{proof}
Via Lemme \ref{lemma-locally-finite-type-characterize} cela se traduit par le fait algébrique suivant
: Étant donné les
homomorphismes d'anneaux $A \to B \to C$ tels que $A \to C$
est de type fini, alors $B \to C$ est
de
type fini. (Voir Algèbre, Lemme \ref{algebra-lemma-compose-finite-type}).
\end{proof}







\section{Points de type fini et schémas de Jacobson}
\label{section-points-finite-type}

\noindent
Soit $S$ un schéma. Un point de type fini $s$ de $S$
est un point tel que le morphisme $\Spec(\kappa(s)) \to S$ est de type fini. La raison
de cette étude est que les points de type fini peuvent remplacer les
points fermés dans un certain sens et dans certaines situations.
Il y en a toujours assez par exemple. De plus, un schéma est Jacobson
si et seulement si tous les points de type fini sont des points fermés.

\begin{lemma}
\label{lemma-point-finite-type}
Soit $S$ un schéma. Soit $k$ un corps.
Soit $f : \Spec(k) \to S$ un morphisme. Les conditions
suivantes sont équivalentes :
\begin{enumerate}
\item Le morphisme $f$ est de type fini.
\item Le morphisme $f$ est localement de type fini.
\item Il existe un $U = \Spec(R)$ ouvert affine de $S$ tel que
$f$ correspond à un homomorphisme d'anneaux fini $R \to k$.
\item Il existe un $U = \Spec(R)$ ouvert affine de $S$ tel que
l'image de $f$ est constitué d'un point fermé $u$ dans $U$
et l'extension du corps $k/\kappa(u)$ est finie.
\end{enumerate}
\end{lemma}

\begin{proof}
L'équivalence de (1) et (2) est évidente car $\Spec(k)$ est un
singleton et donc tout morphisme qui en découle est quasi-compact.

\medskip\noindent
Supposons que $f$ soit localement de type fini.
Choisissez n'importe quel $\Spec(R) = U \subset S$ affine
ouvert tel que l'image de $f$ soit
contenue dans $U$, et l'homomorphisme d'anneaux
$R \to k$ soit de type fini. Soit $\mathfrak p \subset R$
le noyau. Alors $R/\mathfrak p \subset k$ est de type fini. En algèbre,
Lemme \ref{algebra-lemma-field-finite-type-over-domain} il existe un $\overline{f} \in R/\mathfrak p$ tel que
$(R/\mathfrak p)_{\overline{f}}$ est un corps et $(R/\mathfrak p)_{\overline{f}} \to k$ est
une extension de corps fini. Si $f \in R$ est un
lift de $\overline{f}$, alors nous voyons que $k$
est un $R_f$-module fini. Ainsi (2) $\Rightarrow$ (3).

\medskip\noindent
Supposons que $\Spec(R) = U \subset S$ soit un ouvert affine tel que
$f$ correspond à un homomorphisme d'anneaux fini
$R \to k$. Alors $f$ est
localement de type fini par le lemme \ref{lemma-locally-finite-type-characterize}.
Ainsi (3) $\Rightarrow$ (2).

\medskip\noindent
Supposons que $R \to k$ soit fini. L'image de $R \to k$ est
un corps sur lequel $k$
est fini par l'Algèbre, Lemme \ref{algebra-lemma-integral-under-field}. Le
noyau de $R \to k$ est donc un idéal maximal. Ainsi
(3) $\Rightarrow$ (4).

\medskip\noindent
L'implication (4) $\Rightarrow$ (3) est immédiate.
\end{proof}

\begin{lemma}
\label{lemma-artinian-finite-type}
Soit $S$ un schéma.
Soit $A$ un anneau local artinien à corps résiduel
$\kappa$. Soit $f : \Spec(A) \to S$ un morphisme de
schémas. Alors $f$ est de type fini si et seulement
si la composition $\Spec(\kappa) \to \Spec(A) \to S$ est de type fini.
\end{lemma}

\begin{proof}
Puisque le morphisme $\Spec(\kappa) \to \Spec(A)$ est de type fini, il
est clair que si $f$ est de type fini la composition
$\Spec(\kappa) \to S$ l'est également (voir Lemme \ref{lemma-composition-finite-type}). A l'inverse,
notons que les applications $\Spec(A) \to S$ dans certains affines
ouverts $U = \Spec(B)$ de $S$ comme $\Spec(A)$ n'ont
qu'un seul point.
Pour finir, appliquez l'algèbre, Lemme \ref{algebra-lemma-essentially-of-finite-type-into-artinian-local} à
$B \to A$.
\end{proof}

\noindent
Rappelons que étant donné un point $s$ d'un schéma $S$
il existe un morphisme canonique $\Spec(\kappa(s)) \to S$,
voir Schémas, section \ref{schemes-section-points}.

\begin{definition}
\label{definition-finite-type-point}
Soit $S$ un schéma.
Disons qu'un point $s$ de $S$ est un {\it
point de type fini} si le morphisme canonique $\Spec(\kappa(s)) \to S$ est de type
fini. On note $S_{\text{ft-pts}}$ l'ensemble des points de type fini $S$.
\end{definition}

\noindent
Nous pouvons décrire l'ensemble de points de type fini comme suit.

\begin{lemma}
\label{lemma-identify-finite-type-points}
Soit $S$ un schéma. On a
$$
S_{\text{ft-pts}} = \bigcup\nolimits_{U \subset S\text{ ouvert}} U_0
$$
où $U_0$ est l'ensemble des points fermés de $U$. Ici, nous pouvons
faire parcourir à $U$ tous les ouverts, ou seulement tous les ouverts affines de $S$.
\end{lemma}

\begin{proof}
Immédiat du lemme \ref{lemma-point-finite-type}.
\end{proof}

\begin{lemma}
\label{lemma-finite-type-points-morphism}
Soit $f : T \to S$ un morphisme de schémas.
Si $f$ est localement
de type fini, alors $f(T_{\text{ft-pts}}) \subset S_{\text{ft-pts}}$.
\end{lemma}

\begin{proof}
Si $T$ est le spectre d'un corps, c'est Lemme \ref{lemma-point-finite-type}. En général,
cela s'ensuit puisque la composition des morphismes localement
de type fini est localement de type fini (Lemme \ref{lemma-composition-finite-type}).
\end{proof}

\begin{lemma}
\label{lemma-finite-type-points-surjective-morphism}
Soit $f : T \to S$ un morphisme de schémas.
Si $f$ est localement de type fini
et surjectif, alors $f(T_{\text{ft-pts}}) = S_{\text{ft-pts}}$.
\end{lemma}

\begin{proof}
Nous avons $f(T_{\text{ft-pts}}) \subset S_{\text{ft-pts}}$ par le lemme \ref{lemma-finite-type-points-morphism}. Soit
$s \in S$ un point de type fini.
Comme $f$ est surjectif le schéma $T_s = \Spec(\kappa(s)) \times_S T$ est
non vide, donc a un point de type fini $t \in T_s$
par le lemme \ref{lemma-identify-finite-type-points}. Or
$T_s \to T$ est un morphisme de type fini
comme changement de base de $s \to S$ (Lemme
\ref{lemma-base-change-finite-type}).
Ainsi l'image de $t$ dans
$T$ est un point de type fini par le
lemme \ref{lemma-finite-type-points-morphism} qui s'applique à $s$
par construction.
\end{proof}

\begin{lemma}
\label{lemma-enough-finite-type-points}
Soit $S$ un schéma.
Pour tout sous-ensemble localement fermé $T \subset S$ on a
$$
T \not = \emptyset
\Rightarrow
T \cap S_{\text{ft-pts}} \not = \emptyset.
$$
En particulier, pour tout sous-ensemble fermé $T \subset S$
on voit que $T \cap S_{\text{ft-pts}}$ est dense dans $T$.
\end{lemma}

\begin{proof}
Notons que $T$ porte une structure
schématique (voir Schémas, Lemme \ref{schemes-lemma-reduced-closed-subscheme}) telle
que $T \to S$ est une immersion localement
fermée. Toute immersion localement fermée
est localement de type fini, voir Lemme
\ref{lemma-immersion-locally-finite-type}. Ainsi par le lemme \ref{lemma-finite-type-points-morphism} nous
voyons $T_{\text{ft-pts}} \subset S_{\text{ft-pts}}$. Enfin, tout ouvert affine
non vide de $T$ possède au moins un point fermé
qui est un point de type fini de
$T$ par le lemme \ref{lemma-identify-finite-type-points}.
\end{proof}

\noindent
Il s'ensuit que la plupart du matériel de
Topologie, section \ref{topology-section-space-jacobson} passe en revue avec l'ensemble des points
fermés remplacé par l'ensemble des points de type fini. En
fait, si $S$
est Jacobson alors on récupère les points fermés comme
points de type fini.

\begin{lemma}
\label{lemma-jacobson-finite-type-points}
Soit $S$ un schéma. Les conditions suivantes sont équivalentes :
\begin{enumerate}
\item le schéma $S$ est Jacobson,
\item $S_{\text{ft-pts}}$ est l'ensemble des points fermés de $S$,
\item pour tout morphisme $T \to S$ localement de type fini, les points
fermés s'envoient sur des points fermés, et
\item pour tout morphisme $T \to S$ localement de type
fini, les points fermés $t \in T$ s'envoient sur des points fermés
$s \in S$ tels que l'extension $\kappa(s) \subset \kappa(t)$ soit finie.
\end{enumerate}
\end{lemma}

\begin{proof}
On a trivialement (4) $\Rightarrow$ (3) $\Rightarrow$ (2).
Lemme \ref{lemma-enough-finite-type-points} montre que (2) implique (1). Il suffit donc de
montrer que (1) implique (4). Supposons
que $T \to S$ soit localement de type fini.
Choisissons un point fermé $t \in T$ et notons $s \in S$
son image. Choisissons des voisinages ouverts affines $\Spec(R) = U \subset S$
de $s$ et $\Spec(A) = V \subset T$ de $t$, tels que
$V$ s'envoie dans $U$. L'homomorphisme
d'anneaux induit $R \to A$ est de type fini (voir Lemme
\ref{lemma-locally-finite-type-characterize}) et $R$ est Jacobson par Propriétés,
Lemme \ref{properties-lemma-locally-jacobson}. Ainsi
le résultat découle de l'Algèbre, Proposition \ref{algebra-proposition-Jacobson-permanence}.
\end{proof}

\begin{lemma}
\label{lemma-Jacobson-universally-Jacobson}
Soit $S$ un Schéma de Jacobson.
Tout schéma localisé de type fini sur $S$ est Jacobson.
\end{lemma}

\begin{proof}
Cela ressort
clairement de l'algèbre, proposition
\ref{algebra-proposition-Jacobson-permanence} (et des propriétés, Lemme \ref{properties-lemma-locally-jacobson}
et Lemme \ref{lemma-locally-finite-type-characterize}).
\end{proof}

\begin{lemma}
\label{lemma-ubiquity-Jacobson-schemes}
Les types de schémas suivants sont de Jacobson.
\begin{enumerate}
\item Tout schéma local de type fini sur un corps.
\item Tout schéma localisé de type fini sur $\mathbf{Z}$.
\item Tout schéma local de type fini sur un domaine noéthérien $1$-dimensionnel
avec une infinité de nombres premiers.
\item Un schéma de la forme $\Spec(R) \setminus \{\mathfrak m\}$ où $(R, \mathfrak m)$
est un anneau noéthérien local. Également tout
schéma local de type fini dessus.
\end{enumerate}
\end{lemma}

\begin{proof}
Nous utiliserons Lemme \ref{lemma-Jacobson-universally-Jacobson} sans mention. Le spectre
d’un corps est clairement
Jacobson. Le spectre de $\mathbf{Z}$
est Jacobson, voir Algèbre,
Lemme
\ref{algebra-lemma-pid-jacobson}. Pour (3) voir Algèbre, Lemme
\ref{algebra-lemma-noetherian-dim-1-Jacobson}.
Pour (4) voir Propriétés, Lemme \ref{properties-lemma-complement-closed-point-Jacobson}.
\end{proof}







\section{Schémas universellement caténaires}
\label{section-universally-catenary}

\noindent
Rappelons qu'un espace topologique $X$ est appelé {\it caténaire}
si pour chaque paire de sous-ensembles fermés irréductibles $T \subset T'$
il existe une chaîne maximale de sous-ensembles fermés irréductibles
$$
T = T_0 \subset T_1 \subset \ldots \subset T_e = T'
$$
et que chaque chaîne a la même longueur.
Voir Topologie, Définition \ref{topology-definition-catenary}. Rappelons
qu'un schéma est caténaire si son espace topologique sous-jacent
est caténaire. Voir Propriétés, Définition \ref{properties-definition-catenary}.

\begin{definition}
\label{definition-universally-catenary}
Soit $S$ un schéma. Supposons que $S$ soit localement noethérien.
On dit que $S$ est {\it universellement caténaire}
si pour tout morphisme $X \to S$ localement de type fini le schéma $X$ est caténaire.
\end{definition}

\noindent
C'est une ``meilleure'' notion que la caténaire car il existe
des schémas noéthériens qui sont caténaires mais pas
universellement caténaires. Voir Exemples, section \ref{examples-section-non-catenary-Noetherian-local}. De nombreux
schémas sont universellement caténaire,
voir Lemme \ref{lemma-ubiquity-uc} ci-dessous.

\medskip\noindent Rappelons
qu'un anneau $A$ est appelé {\it
caténaire} si pour toute paire d'idéaux premiers $\mathfrak p \subset \mathfrak q$ il existe
une chaîne maximale de nombres premiers
$$
\mathfrak p =
\mathfrak p_0 \subset \ldots \subset \mathfrak p_e
= \mathfrak q
$$
et tous ont la même longueur. Voir Algèbre,
Définition \ref{algebra-definition-catenary}. Nous avons vu la relation entre les
schémas caténaires et les anneaux caténaires dans Propriétés,
section \ref{properties-section-catenary}. Rappelons qu'un anneau $A$ est appelé {\it
universellement caténaire} si $A$
est noéthérien et pour tout type fini homomorphisme d'anneaux
$A \to B$ l'anneau $B$
est caténaire. Voir Algèbre, Définition \ref{algebra-definition-universally-catenary}. De nombreux anneaux
intéressants rencontrés en géométrie
algébrique satisfont à cette propriété.

\begin{lemma}
\label{lemma-universally-catenary-local}
Soit $S$ un schéma localement noethérien. Les conditions suivantes sont équivalentes
\begin{enumerate}
\item $S$ est universellement caténaire,
\item il existe un recouvrement ouvert de $S$ dont tous les membres sont des
schémas universellement caténaires,
\item pour chaque ouvert affine $\Spec(R) = U \subset S$ l'anneau $R$ est
universellement caténaire, et
\item il existe un recouvrement ouvert affine $S = \bigcup U_i$ tel que chaque
$U_i$ soit le spectre d'un anneau universellement caténaire.
\end{enumerate}
De plus, dans ce cas tout schéma localisé de type fini sur $S$ est également
universellement caténaire.
\end{lemma}

\begin{proof}
Par le Lemme \ref{lemma-immersion-locally-finite-type} une immersion ouverte est localement de type
fini. Une composition de morphismes localement de type fini est
localement de type fini (Lemme \ref{lemma-composition-finite-type}).
Il est donc clair que si $S$ est universellement caténaire
alors tout ouvert et tout schéma localisé de type fini sur $S$
est également universellement caténaire. Cela prouve l’énoncé
final du lemme et que (1) implique (2).

\medskip\noindent Si
$\Spec(R)$ est un schéma universellement caténaire, alors tout schéma
$\Spec(A)$ avec $A$ un type fini $R$-algèbre
est caténaire. Donc tous ces anneaux $A$ sont caténaires
par Algèbre, Lemme \ref{algebra-lemma-catenary}. Ainsi $R$
est universellement caténaire. Combiné avec les remarques ci-dessus Nous concluons
que (1) implique (3), et (2) implique (4). Bien sûr, (3) implique
(4) de manière triviale.

\medskip\noindent
Pour terminer la démonstration nous montrons
que (4) implique (1). Supposons (4) et soit $X \to S$ un morphisme
localement de type fini. On peut trouver un recouvrement
recouvrement par des ouverts affines $X = \bigcup V_j$ tel que chaque $V_j \to S$
ait son image dans l'un des $U_i$. Par le Lemme
\ref{lemma-locally-finite-type-characterize} l'homomorphisme d'anneaux induit $\mathcal{O}(U_i) \to \mathcal{O}(V_j)$ est
de type fini. Donc $\mathcal{O}(V_j)$ est une caténaire. Par
conséquent, $X$ est caténaire par Propriétés, Lemme \ref{properties-lemma-catenary-local}.
\end{proof}

\begin{lemma}
\label{lemma-universally-catenary-local-rings-universally-catenary}
Soit $S$ un schéma localement noethérien. Les conditions
suivantes sont équivalentes :
\begin{enumerate}
\item $S$ est universellement caténaire, et
\item tous les anneaux locaux $\mathcal{O}_{S, s}$ de $S$ sont universellement caténaire.
\end{enumerate}
\end{lemma}

\begin{proof}
Supposons que tous les anneaux locaux de $S$ soient universellement
caténaires. Soit $f : X \to S$ localement
de type fini. On sait que $X$ est caténaire si et
seulement si $\mathcal{O}_{X, x}$ est caténaire pour tous $x \in X$. Si $f(x) = s$,
alors $\mathcal{O}_{X, x}$ est essentiellement de type fini sur $\mathcal{O}_{S, s}$.
Donc $\mathcal{O}_{X, x}$ est caténaire en supposant que
$\mathcal{O}_{S, s}$ est universellement caténaire.

\medskip\noindent A
l’inverse, supposons que $S$ soit universellement caténaire. Soit
$s \in S$. On peut remplacer $S$ par un voisinage affine
ouvert de $s$ par le lemme \ref{lemma-universally-catenary-local}. Disons que $S = \Spec(R)$
et $s$ correspond à l'idéal premier $\mathfrak p$. Tout type
fini $R_{\mathfrak p}$-algèbre $A'$ est de la forme
$A_{\mathfrak p}$ pour certains types fini $R$-algèbre $A$.
Par hypothèse (et Lemme \ref{lemma-universally-catenary-local} si vous préférez) l'anneau $A$
est caténaire, et donc $A'$ (une localisation de $A$) est
caténaire. Ainsi $R_{\mathfrak p}$ est universellement caténaire.
\end{proof}

\begin{lemma}
\label{lemma-catenary-check-irreducible}
Soit $S$ un schéma localement noethérien. Alors $S$ est universellement caténaire
si et seulement si les composantes irréductibles de $S$ sont universellement caténaire.
\end{lemma}

\begin{proof}
Omis. Pour le cas affine, veuillez
consulter Algèbre, Lemme \ref{algebra-lemma-catenary-check-irreducible}.
\end{proof}

\begin{lemma}
\label{lemma-ubiquity-uc}
Les types de schémas suivants sont universellement caténaire.
\begin{enumerate}
\item Tout schéma local de type fini sur un corps.
\item Tout schéma localisé de type fini sur un schéma De Cohen--Macaulay.
\item Tout schéma localisé de type fini sur $\mathbf{Z}$.
\item Tout schéma localisé de type fini sur un domaine noéthérien
$1$-dimensionnel.
\item Et ainsi de suite.
\end{enumerate}
\end{lemma}

\begin{proof}
Tout cela découle du fait qu'un anneau
de De Cohen--Macaulay est universellement
caténaire, voir Algèbre, Lemme \ref{algebra-lemma-CM-ring-catenary}.
Utilisez également la dernière
assertion du lemme \ref{lemma-universally-catenary-local}. Certains
détails omis.
\end{proof}























\section{Schémas de Nagata, reprise}
\label{section-nagata}

\noindent
Voir Propriétés, section \ref{properties-section-nagata} pour les définitions et propriétés
de base des schémas Nagata et universellement japonais.

\begin{lemma}
\label{lemma-finite-type-nagata}
Soit $f : X \to S$ un morphisme. Si $S$
est Nagata et $f$ localement de type fini alors $X$ est Nagata. Si $S$
est universellement japonais et
$f$ localement de type fini alors $X$ est universellement japonais.
\end{lemma}

\begin{proof}
Pour ``universellement japonais''
Cela résulte de Algèbre, Lemme \ref{algebra-lemma-universally-japanese}.
Pour ``Nagata''
Cela résulte de Algèbre, Proposition \ref{algebra-proposition-nagata-universally-japanese}.
\end{proof}

\begin{lemma}
\label{lemma-ubiquity-nagata}
Les types de schémas suivants sont des Nagata.
\begin{enumerate}
\item Tout schéma local de type fini sur un corps.
\item Tout schéma local de type fini sur un anneau local complet noéthérien.
\item Tout schéma localisé de type fini sur $\mathbf{Z}$.
\item Tout schéma localement de type fini sur un anneau de Dedekind de
caractéristique zéro.
\item Et ainsi de suite.
\end{enumerate}
\end{lemma}

\begin{proof}
Par le Lemme \ref{lemma-finite-type-nagata} il suffit de montrer que les anneaux
mentionnés ci-dessus sont des anneaux Nagata.
Pour cela, voir Algèbre, Proposition \ref{algebra-proposition-ubiquity-nagata}.
\end{proof}


\section{Le lieu singulier, reprise}
\label{section-singular-locus}

\noindent
Nous recherchons un critère qui implique l'ouverture du lieu régulier
pour tout schéma local de type fini sur la base. Voici la définition.

\begin{definition}
\label{definition-J}
Soit $X$ un schéma localement noethérien. On dit que $X$ est {\it
J-2} si pour tout morphisme $Y \to X$ qui est localement de type
fini le locus régulier $\text{Reg}(Y)$ est ouvert dans $Y$.
\end{definition}

\noindent
C'est l'analogue de la notion correspondante pour les anneaux noéthériens,
voir Compléments d'algèbre, Définition \ref{more-algebra-definition-J}.

\begin{lemma}
\label{lemma-J}
Soit $X$ un schéma localement noethérien. Les conditions suivantes sont équivalentes
\begin{enumerate}
\item $X$ est J-2,
\item il existe un recouvrement ouvert de $X$ dont tous les membres sont des
schémas J-2,
\item pour chaque ouvert affine $\Spec(R) = U \subset X$ l'anneau $R$
est J-2, et
\item il existe un recouvrement ouvert affine $S = \bigcup U_i$ tel que chaque
$\mathcal{O}(U_i)$ soit J-2 pour tous les $i$.
\end{enumerate}
De plus, dans ce cas, tout schéma localisé de type fini sur $X$ est
également J-2.
\end{lemma}

\begin{proof}
Par le Lemme \ref{lemma-immersion-locally-finite-type} une immersion ouverte est localement de type
fini. Une composition de morphismes localement de type fini
est localement de type fini (Lemme
\ref{lemma-composition-finite-type}). Il est donc clair que si $X$ est J-2 alors tout
ouvert et tout schéma localisé de type fini sur
$X$ est également J-2. Cela prouve l’énoncé
final du lemme.

\medskip\noindent Si
$\Spec(R)$ est J-2, alors pour tout type fini $R$-algèbre $A$
le lieu régulier du schéma $\Spec(A)$ est ouvert. Donc $R$
est J-2, par définition
(voir Compléments d'algèbre, Définition \ref{more-algebra-definition-J}).
Combiné avec les remarques ci-dessus Nous concluons que (1) implique (3),
et (2) implique (4). Bien sûr (1) $\Rightarrow$ (2) et
(3) $\Rightarrow$ (4) trivialement.

\medskip\noindent
Pour terminer la démonstration nous montrons
que (4) implique (1). Supposons (4) et soit $Y \to X$
un morphisme localement de type fini. On peut trouver un
recouvrement par des ouverts affines $Y = \bigcup V_j$ tel
que chaque $V_j \to X$ ait son image dans l'un des $U_i$.
Par le Lemme \ref{lemma-locally-finite-type-characterize}, l'homomorphisme d'anneaux
induit $\mathcal{O}(U_i) \to \mathcal{O}(V_j)$ est de type fini.
Le lieu régulier de $V_j = \Spec(\mathcal{O}(V_j))$ est donc
ouvert. Puisque $\text{Reg}(Y) \cap V_j = \text{Reg}(V_j)$ Nous concluons que $\text{Reg}(Y)$
est ouvert comme souhaité.
\end{proof}

\begin{lemma}
\label{lemma-ubiquity-J-2}
Les types de schémas suivants sont J-2.
\begin{enumerate}
\item Tout schéma local de type fini sur un corps.
\item Tout schéma local de type fini sur un anneau local complet noéthérien.
\item Tout schéma localisé de type fini sur $\mathbf{Z}$.
\item Tout schéma local de type fini sur un anneau noéthérien local de dimension
$1$.
\item Tout schéma local de type fini sur un anneau Nagata de dimension $1$.
\item Tout schéma localisé de type fini sur un anneau de Dedekind de caractéristique
zéro.
\item Et ainsi de suite.
\end{enumerate}
\end{lemma}

\begin{proof}
Par le Lemme \ref{lemma-J} il suffit de montrer que
les anneaux mentionnés ci-dessus sont J-2. Pour
cela, voir Compléments d'algèbre, Proposition \ref{more-algebra-proposition-ubiquity-J-2}.
\end{proof}




\section{Schémas excellents}
\label{section-excellent}

\noindent
On rappelle qu'un anneau est quasi-excellent s'il est un G-anneau et J-2.
Un anneau est excellent s'il est quasi-excellent et universellement caténaire.

\begin{definition}
\label{definition-excellent}
Soit $X$ un schéma.
\begin{enumerate}
\item On dit que $X$ est {\it quasi-excellent}
si pour tout $x \in X$ il existe un voisinage ouvert affine $x \in U \subset X$
tel que l'anneau $\mathcal{O}_X(U)$ est quasi-excellent
(voir Compléments d'algèbre, Définition \ref{more-algebra-definition-excellent}).
\item On dit que $X$ est {\it excellent} si
pour tout $x \in X$ il existe un voisinage ouvert affine $x \in U \subset X$
tel que l'anneau $\mathcal{O}_X(U)$ est
excellent (voir Compléments d'algèbre, Définition \ref{more-algebra-definition-excellent}).
\end{enumerate}
\end{definition}

\noindent
Les schémas quasi-excellents sont localement noethériens (les anneaux G sont noethériens).

\begin{lemma}
\label{lemma-characterize-quasi-excellent}
Soit $X$ un schéma. Les conditions suivantes sont équivalentes
\begin{enumerate}
\item $X$ est quasi-excellent, et
\item $X$ est un G-schéma et J-2.
\end{enumerate}
\end{lemma}

\begin{proof}
L'implication (1) $\Rightarrow$ (2) résulte de la combinaison
de la définition des quasi-schémas
excellents (Définition \ref{definition-excellent}),
de la définition des anneaux quasi-excellents
(Compléments d'algèbre, Définition \ref{more-algebra-definition-excellent}), de
la définition des G-schémas
(Propriétés, Définition \ref{properties-definition-G-scheme}) et du lemme \ref{lemma-J}.
Inversement, si
$X$ est un G-schéma et J-2, alors pour tout $x \in X$
on peut trouver un voisinage ouvert affine $x \in U \subset X$
tel que $\mathcal{O}_X(U)$ soit un G-anneau. Par le lemme \ref{lemma-J},
l'anneau $\mathcal{O}_X(U)$ est également J-2, donc quasi-excellent.
\end{proof}

\begin{lemma}
\label{lemma-quasi-excellent-nagata}
Un quasi-schéma excellent est Nagata.
\end{lemma}

\begin{proof}
Compléments d'algèbre, Lemme \ref{more-algebra-lemma-quasi-excellent-nagata}.
\end{proof}

\begin{lemma}
\label{lemma-characterize-excellent}
Soit $X$ un schéma. Les conditions suivantes sont équivalentes
\begin{enumerate}
\item $X$ est excellent, et
\item $X$ est quasi-excellent et universellement caténaire.
\end{enumerate}
\end{lemma}

\begin{proof}
Supposons (1). Puisque les anneaux excellents sont quasi-excellents,
il est clair que $X$ est quasi-excellent. Puisque les anneaux
excellents sont universellement caténaire, il s'ensuit également
que $X$ est universellement caténaire par le lemme
\ref{lemma-universally-catenary-local}. Supposons (2). Soit $x \in X$. Puisque $X$
est quasi-excellent il existe un voisinage
ouvert affine $x \in U \subset X$ tel que $\mathcal{O}_X(U)$ est quasi-excellent.
Puisque $X$ est universellement caténaire, l'anneau
$\mathcal{O}_X(U)$ est aussi universellement caténaire d'après le lemme \ref{lemma-universally-catenary-local},
donc excellent.
\end{proof}

\begin{lemma}
\label{lemma-locally-excellent}
Soit $X$ un schéma. Les conditions suivantes sont équivalentes :
\begin{enumerate}
\item Le schéma $X$ est (quasi-)excellent.
\item Pour chaque $U \subset X$ ouvert affine, l'anneau $\mathcal{O}_X(U)$ est
(quasi-)excellent.
\item Il existe un recouvrement ouvert affine $X = \bigcup U_i$ tel que
chaque $\mathcal{O}_X(U_i)$ soit (quasi-)excellent.
\item Il existe un recouvrement ouvert $X = \bigcup X_j$ tel que
chaque sous-schéma ouvert $X_j$ est (quasi-)excellent.
\end{enumerate}
De plus, si $X$ est (quasi-)excellent alors tout sous-schéma ouvert
est (quasi-)excellent.
\end{lemma}

\begin{proof}
Par le Lemme \ref{lemma-characterize-quasi-excellent} être quasi-excellent
équivaut à être un G-schéma et J-2. Ainsi
pour le cas quasi-excellent, le lemme
découle du lemme
\ref{lemma-J} et Propriétés, Lemme \ref{properties-lemma-locally-G-scheme}.
Pour le cas excellent, utilisez Lemme \ref{lemma-characterize-excellent},
le cas quasi-excellent, et Lemme \ref{lemma-universally-catenary-local}.
\end{proof}

\begin{lemma}
\label{lemma-finite-type-excellent}
Soit $f : X \to S$ un morphisme. Si $S$
est (quasi-)excellent et $f$ localement de type fini alors $X$
est (quasi-)excellent.
\end{lemma}

\begin{proof}
Compléments d'algèbre, Lemme \ref{more-algebra-lemma-finite-type-over-excellent}.
\end{proof}

\begin{lemma}
\label{lemma-ubiquity-excellent}
Les types de schémas suivants sont excellents.
\begin{enumerate}
\item Tout schéma local de type fini sur un corps.
\item Tout schéma local de type fini sur un anneau local complet noéthérien.
\item Tout schéma localisé de type fini sur $\mathbf{Z}$.
\item Tout schéma localement de type fini sur un anneau de Dedekind de
caractéristique zéro.
\end{enumerate}
\end{lemma}

\begin{proof}
Par les Lemmes \ref{lemma-locally-excellent} et \ref{lemma-finite-type-excellent} il suffit de montrer que
les anneaux mentionnés ci-dessus sont excellents.
Pour cela,
voir Compléments d'algèbre, Proposition \ref{more-algebra-proposition-ubiquity-excellent}.
\end{proof}






\section{Morphismes quasi-finis}
\label{section-quasi-finite}

\noindent
Un traitement solide des morphismes quasi-finis est à la base de nombreux développements
ultérieurs. Cela conduira à différentes versions du Théorème principal
de Zariski, comportement des dimensions des fibres, descente pour morphismes
étales, etc, etc. Avant de lire cette section, il peut être judicieux
de jeter un œil aux résultats d'algèbre dans Algèbre, section \ref{algebra-section-quasi-finite}.

\medskip\noindent
Rappelons qu'un type fini homomorphisme d'anneaux $R \to A$
est quasi-fini en un nombre premier $\mathfrak q$ si $\mathfrak q$
définit un point isolé de sa fibre, voir Algèbre, Définition \ref{algebra-definition-quasi-finite}.

\begin{definition}
\label{definition-quasi-finite}
\begin{reference}
\cite[II Définition 6.2.3]{EGA}
\end{reference}
Soit $f : X \to S$ un morphisme de schémas.
\begin{enumerate}
\item On dit que $f$ est {\it quasi-fini en un point
$x \in X$} s'il existe un voisinage affine $\Spec(A) = U \subset X$
de $x$ et un ouvert affine $\Spec(R) = V \subset S$ tel que $f(U) \subset V$,
l'homomorphisme d'anneaux $R \to A$ est de type fini, et $R \to A$
est quasi-fini au premier de $A$ correspondant à $x$ (voir
ci-dessus).
\item On dit que $f$ est {\it localement quasi-fini} si $f$
est quasi-fini en tout point $x$ de $X$.
\item On dit que $f$ est {\it quasi-fini} si $f$ est
de type fini et chaque point $x$ est un point isolé de sa fibre.
\end{enumerate}
\end{definition}

\noindent
Trivialement, un morphisme quasi-fini localement est localement de type fini. Nous
verrons plus loin qu'un morphisme $f$ qui est localement de type fini
est quasi-fini à $x$ si et seulement si $x$ est isolé dans sa
fibre. De plus, l'ensemble des points auxquels un morphisme est quasi-fini est ouvert
; nous le verrons dans la section \ref{section-Zariski} sur le Théorème principal de Zariski.

\begin{lemma}
\label{lemma-algebraic-residue-field-extension-closed-point-fibre}
Soit $f : X \to S$ un morphisme de schémas. Soit
$x \in X$ un point. Posons $s = f(x)$.
Si $\kappa(x)/\kappa(s)$ est une
extension de corps algébrique, alors
\begin{enumerate}
\item $x$ est un point fermé de sa fibre, et
\item si en plus $s$ est un point fermé de $S$, alors $x$
est un point fermé de $X$.
\end{enumerate}
\end{lemma}

\begin{proof}
La deuxième instruction découle de la première par topologie élémentaire.
D'après Schémas, Lemme \ref{schemes-lemma-fibre-topological} pour prouver la première
affirmation on peut
remplacer $X$ par $X_s$ et $S$ par
$\Spec(\kappa(s))$. On peut donc supposer que $S = \Spec(k)$ est le spectre
d'un corps. Dans ce cas, soit $\Spec(A) = U \subset X$ un ouvert affine quelconque
contenant $x$. Le point $x$ correspond
à un idéal premier $\mathfrak q \subset A$ tel que $\kappa(\mathfrak q)/k$ est
une extension algébrique de corps.
D'après Algèbre, Lemme \ref{algebra-lemma-finite-residue-extension-closed}, $\mathfrak q$ est
un idéal maximal, c'est-à-dire que $x \in U$ est un point fermé.
Puisque les ouverts affines forment
une base de la topologie de $X$, nous concluons que $\{x\}$ est fermé.
\end{proof}

\noindent
Le lemme suivant est une version du Hilbert Nullstellensatz.

\begin{lemma}
\label{lemma-closed-point-fibre-locally-finite-type}
Soit $f : X \to S$ un morphisme de schémas. Soit $x \in X$
un point. Posons $s = f(x)$. Supposons
que $f$ soit localement de type
fini. Alors $x$ est un point fermé de
sa fibre si et seulement si $\kappa(x)/\kappa(s)$ est
une extension finie de corps.
\end{lemma}

\begin{proof}
Si l'extension est finie, alors $x$ est
un point
fermé de la fibre par le lemme \ref{lemma-algebraic-residue-field-extension-closed-point-fibre} ci-dessus. À l'inverse,
supposons que $x$ soit un point
fermé de sa fibre. Choisissons des ouverts affines
$\Spec(A) = U \subset X$ et $\Spec(R) = V \subset S$ tels que $f(U) \subset V$.
Par le Lemme \ref{lemma-locally-finite-type-characterize} l'homomorphisme d'anneaux $R \to A$
est de type fini. Soit $\mathfrak q \subset A$, resp.\ $\mathfrak p \subset R$
l'idéal premier correspondant à
$x$, resp.\ $s$.
Considérez l'anneau de fibre $\overline{A} = A \otimes_R \kappa(\mathfrak p)$.
Soit $\overline{\mathfrak q}$ le premier de $\overline{A}$
correspondant à $\mathfrak q$. L'hypothèse selon
laquelle $x$ est un point fermé de sa fibre
implique que $\overline{\mathfrak q}$ est un idéal maximal
de $\overline{A}$. Puisque $\overline{A}$ est une algèbre
de type fini sur le corps
$\kappa(\mathfrak p)$ nous voyons par le Hilbert Nullstellensatz,
voir Algèbre, Théorème \ref{algebra-theorem-nullstellensatz},
que $\kappa(\overline{\mathfrak q})$
est une extension finie de $\kappa(\mathfrak p)$.
Puisque $\kappa(s) = \kappa(\mathfrak p)$ et $\kappa(x) = \kappa(\mathfrak q) = \kappa(\overline{\mathfrak q})$ nous
gagnons.
\end{proof}

\begin{lemma}
\label{lemma-base-change-closed-point-fibre-locally-finite-type}
Soit $f : X \to S$ un morphisme de schémas qui est localement de type fini. Soit $g : S' \to S$
un morphisme quelconque. Notons $f' : X' \to S'$ le changement de base. Si $x' \in X'$
s'envoie sur un point $x \in X$ qui est fermé dans $X_{f(x)}$, alors
$x'$ est fermé dans $X'_{f'(x')}$.
\end{lemma}

\begin{proof}
Le corps résiduel $\kappa(x')$ est un quotient
de $\kappa(f'(x')) \otimes_{\kappa(f(x))} \kappa(x)$, voir Schémas, Lemme
\ref{schemes-lemma-points-fibre-product}. Il s'agit donc d'une extension finie
de $\kappa(f'(x'))$ comme $\kappa(x)$ est une
extension finie de $\kappa(f(x))$ par le lemme
\ref{lemma-closed-point-fibre-locally-finite-type}. On voit donc que $x'$ est
fermé dans sa fibre en appliquant une nouvelle fois
ce lemme.
\end{proof}

\begin{lemma}
\label{lemma-residue-field-quasi-finite}
Soit $f : X \to S$ un morphisme de schémas. Soit
$x \in X$ un point. Posons $s = f(x)$. Si
$f$ est quasi-fini en $x$, alors l'extension
du corps résiduel $\kappa(x)/\kappa(s)$ est finie.
\end{lemma}

\begin{proof}
Cela ressort clairement de l'algèbre, Définition \ref{algebra-definition-quasi-finite}.
\end{proof}

\begin{lemma}
\label{lemma-quasi-finite-at-point-characterize}
Soit $f : X \to S$ un morphisme de schémas. Soit $x \in X$
un point. Posons $s = f(x)$. Soit $X_s$ la fibre
de $f$ en $s$. Supposons que $f$
soit localement de type fini. Les conditions
suivantes sont équivalentes :
\begin{enumerate}
\item Le morphisme $f$ est quasi-fini en $x$.
\item Le point $x$ est isolé dans $X_s$.
\item Le point $x$ est fermé dans $X_s$
et il n'y a aucun point $x' \in X_s$, $x' \not = x$ qui se
spécialise dans $x$.
\item Pour toute paire d'affines
ouvertes $\Spec(A) = U \subset X$, $\Spec(R) = V \subset S$ avec $f(U) \subset V$
et $x \in U$ correspondant à $\mathfrak q \subset A$ l'homomorphisme d'anneaux
$R \to A$ est quasi-fini en $\mathfrak q$.
\end{enumerate}
\end{lemma}

\begin{proof}
Supposons que $f$ soit quasi-fini à $x$. Par hypothèse il existe
des ouvertures $U \subset X$, $V \subset S$ telles que $f(U) \subset V$, $x \in U$
et $x$ un point isolé de $U_s$. Par conséquent, $\{x\} \subset U_s$ est un sous-ensemble
ouvert. Puisque $U_s = U \cap X_s \subset X_s$ est également ouvert, nous concluons que $\{x\} \subset X_s$
est également un sous-ensemble ouvert. Ainsi Nous concluons que $x$
est un point isolé de $X_s$.

\medskip\noindent Notez
que $X_s$ est un Schéma de Jacobson
du lemme \ref{lemma-ubiquity-Jacobson-schemes} (et Lemme \ref{lemma-base-change-finite-type}).
Si
$x$ est isolé dans $X_s$, c'est-à-dire
que $\{x\} \subset X_s$ est ouvert, alors $\{x\}$ contient un point
fermé (par la propriété Jacobson), donc $x$ est fermé dans $X_s$.
Il est clair qu'il n'y a aucun intérêt pour $x' \in X_s$, distinct de $x$,
de se spécialiser dans $x$.

\medskip\noindent
Supposons que $x$ est fermé dans $X_s$ et qu'il n'y
a aucun point $x' \in X_s$, distinct de $x$, spécialisé dans
$x$. Considérons une paire d'ouverts affines
$\Spec(A) = U \subset X$, $\Spec(R) = V \subset S$ avec $f(U) \subset V$ et $x \in U$. Soit
$\mathfrak q \subset A$ correspond à $x$ et $\mathfrak p \subset R$
correspond à $s$. Par le Lemme \ref{lemma-locally-finite-type-characterize} l'homomorphisme
d'anneaux $R \to A$ est de type fini.
Considérez l'anneau de fibre $\overline{A} = A \otimes_R \kappa(\mathfrak p)$. Soit
$\overline{\mathfrak q}$ le premier de $\overline{A}$ correspondant à $\mathfrak q$.
Puisque $\Spec(\overline{A})$ est un sous-schéma ouvert de la fibre $X_s$
on voit que $\overline{q}$ est un idéal maximal
de $\overline{A}$ et qu'il ne sert à rien de spécialiser $\Spec(\overline{A})$
sur $\overline{\mathfrak q}$. Cela implique
que $\dim(\overline{A}_{\overline{q}}) = 0$. Ainsi par Algèbre, Définition
\ref{algebra-definition-quasi-finite}
nous voyons que $R \to A$ est quasi-fini
à $\mathfrak q$, c'est-à-dire que $X \to S$ est quasi-fini
à $x$ par définition.

\medskip\noindent À
ce stade, nous avons montré que les conditions (1) à (3) sont toutes équivalentes.
Il est clair que (4) implique (1). Et il est clair aussi que (2) implique
(4) puisque si $x$ est un point isolé de $X_s$ alors
c'est aussi un point isolé de $U_s$ pour tout $U$ ouvert
qui le contient.
\end{proof}

\begin{lemma}
\label{lemma-finite-fibre}
Soit $f : X \to S$ un morphisme de schémas. Soit
$s \in S$. Supposons que
\begin{enumerate}
\item $f$ est localement de type fini et que
\item $f^{-1}(\{s\})$ est un ensemble fini.
\end{enumerate}
Alors $X_s$ est un espace topologique fini discret, et $f$
est quasi-fini en chaque point de $X$ recouvrant $s$.
\end{lemma}

\begin{proof}
Supposons que $T$ soit un schéma qui (a) est localement
de type fini sur un corps $k$, et (b) comporte
un nombre fini de points. Alors Lemme \ref{lemma-ubiquity-Jacobson-schemes} montre
que $T$ est un Schéma de Jacobson. Un espace
de Jacobson fini est discret, voir Topologie, Lemme
\ref{topology-lemma-finite-jacobson}. Appliquez cette remarque à la fibre $X_s$ qui est localement
de type fini sur $\Spec(\kappa(s))$ pour voir la première instruction.
Enfin, appliquez Lemme \ref{lemma-quasi-finite-at-point-characterize} pour voir la seconde.
\end{proof}

\begin{lemma}
\label{lemma-locally-quasi-finite-fibres}
\begin{slogan}
Les morphismes de type fini avec des fibres discrètes sont quasi-finis.
\end{slogan}
Soit $f : X \to S$ un morphisme de schémas. Supposons que $f$
soit localement de type fini. Alors Les
conditions suivantes sont équivalentes
\begin{enumerate}
\item $f$ est localement quasi-fini,
\item pour tout $s \in S$ la fibre $X_s$ est un espace topologique discret, et
\item pour tout morphisme $\Spec(k) \to S$ où $k$ est un corps le changement
de base $X_k$ a un espace topologique discret sous-jacent.
\end{enumerate}
\end{lemma}

\begin{proof}
Il est immédiat que (3) implique
(2). Lemme \ref{lemma-quasi-finite-at-point-characterize} montre que (2) est équivalent
à (1). Supposons (2)
et soit $\Spec(k) \to S$ comme dans (3).
Notons $s \in S$ l'image de $\Spec(k) \to S$.
Alors $X_k$ est le changement
de base de $X_s$ via $\Spec(k) \to \Spec(\kappa(s))$.
Ainsi chaque point
de $X_k$ est fermé par le lemme \ref{lemma-base-change-closed-point-fibre-locally-finite-type}. Comme $X_k \to \Spec(k)$
est localement de type fini (par
le lemme \ref{lemma-base-change-finite-type}), nous pouvons
appliquer
Lemme \ref{lemma-quasi-finite-at-point-characterize} pour conclure que chaque point de
$X_k$ est isolé, c'est-à-dire que $X_k$ a un espace
topologique discret sous-jacent.
\end{proof}

\begin{lemma}
\label{lemma-quasi-finite-locally-quasi-compact}
Soit $f : X \to S$ un morphisme de schémas. Alors $f$
est quasi-fini si et seulement si $f$
est localement quasi-fini et quasi-compact.
\end{lemma}

\begin{proof}
Supposons que $f$ soit quasi-fini. Il est quasi-compact
par Définition \ref{definition-finite-type}. Soit $x \in X$.
On voit que $f$ est quasi-fini
en $x$ par le lemme \ref{lemma-quasi-finite-at-point-characterize}. Donc $f$
est quasi-compact et localement quasi-fini.

\medskip\noindent
Supposons que $f$ soit quasi-compact et localement
quasi-fini. Alors $f$ est de type fini.
Soit $x \in X$ un point. Par le Lemme \ref{lemma-quasi-finite-at-point-characterize} on voit
que $x$ est un point isolé de sa fibre. Le
lemme est prouvé.
\end{proof}

\begin{lemma}
\label{lemma-quasi-finite}
Soit $f : X \to S$ un morphisme de schémas. Les conditions
suivantes sont équivalentes :
\begin{enumerate}
\item $f$ est quasi-fini, et
\item $f$ est localement de type fini, quasi-compact, et possède des fibres finies.
\end{enumerate}
\end{lemma}

\begin{proof}
Supposons que $f$ soit quasi-fini. En particulier $f$ est localement
de type fini et quasi-compact (puisqu'il est de type fini). Soit $s \in S$.
Puisque chaque $x \in X_s$ est isolé dans $X_s$ on
voit que $X_s = \bigcup_{x \in X_s} \{x\}$ est un recouvrement ouvert. Comme $f$ est
quasi-compact, la fibre $X_s$ est quasi-compacte. On voit donc que
$X_s$ est fini.

\medskip\noindent
À l'inverse, supposons que $f$ est localement de type
fini, quasi-compact et possède des fibres finies. Alors
il est localement quasi-fini par le lemme \ref{lemma-finite-fibre}.
Il est donc quasi-fini par le lemme \ref{lemma-quasi-finite-locally-quasi-compact}.
\end{proof}

\noindent
Rappelons qu'un homomorphisme d'anneaux $R \to A$ est quasi-fini
s'il est de type fini et quasi-fini à {\it tous les}
premiers de $A$, voir Algèbre, Définition \ref{algebra-definition-quasi-finite}.

\begin{lemma}
\label{lemma-locally-quasi-finite-characterize}
Soit $f : X \to S$ un morphisme de schémas. Les conditions
suivantes sont équivalentes
\begin{enumerate}
\item Le morphisme $f$ est localement quasi-fini.
\item Pour toute paire d'ouverts affines $U \subset X$, $V \subset S$
avec $f(U) \subset V$ l'homomorphisme
d'anneaux $\mathcal{O}_S(V) \to \mathcal{O}_X(U)$ est quasi-fini.
\item Il existe un recouvrement ouvert $S = \bigcup_{j \in J} V_j$ et
des recouvrements ouverts $f^{-1}(V_j) = \bigcup_{i \in I_j} U_i$ tels que chacun
des morphismes $U_i \to V_j$, $j\in J, i\in I_j$ est localement
quasi-fini.
\item Il existe un recouvrement ouvert affine $S = \bigcup_{j \in J} V_j$
et des recouvrements ouverts affines $f^{-1}(V_j) = \bigcup_{i \in I_j} U_i$ tels
que l'homomorphisme d'anneaux $\mathcal{O}_S(V_j) \to \mathcal{O}_X(U_i)$ est quasi-fini,
pour tout $j\in J, i\in I_j$.
\end{enumerate}
De plus, si $f$ est localement quasi-fini alors
pour tout sous-schémas ouverts $U \subset X$, $V \subset S$ avec $f(U) \subset V$
la restriction $f|_U : U \to V$ est localement quasi-finie.
\end{lemma}

\begin{proof}
Pour un homomorphisme d'anneaux $R \to A$
définissons $P(R \to A)$ comme signifiant que ``$R \to A$
est quasi-fini''
(voir remarque ci-dessus lemme). Nous affirmons
que $P$ est une propriété locale des homomorphismes
d'anneaux. On vérifie les conditions (a), (b) et
(c) de la Définition \ref{definition-property-local}. Dans la démonstration
du lemme \ref{lemma-locally-finite-type-characterize} nous avons vu que (a), (b) et (c) sont
valables pour la propriété d'être ``de type fini''. A noter
que, pour un homomorphisme d'anneaux de type fini $R \to A$, la propriété
$R \to A$ est quasi-fini en $\mathfrak q$ ne dépend que
de l'anneau local $A_{\mathfrak q}$ comme une algèbre sur $R_{\mathfrak p}$ où $\mathfrak p = R \cap \mathfrak q$
(abus de notation habituel). À l'aide de ces remarques (a), (b) et
(c) de la Définition \ref{definition-property-local} suivent immédiatement. Par exemple,
supposons que $R \to A$ soit un homomorphisme
d'anneaux tel que tous les homomorphismes d'anneaux $R \to A_{a_i}$
soient quasi-finis pour $a_1, \ldots, a_n \in A$ générant l'idéal
unitaire. Nous concluons que $R \to A$ est de type
fini. De plus, pour tout $\mathfrak q \subset A$ premier, l'anneau local $A_{\mathfrak q}$
est isomorphe en tant qu'algèbre $R$
à l'anneau local $(A_{a_i})_{\mathfrak q_i}$ pour certains $i$
et certains $\mathfrak q_i \subset A_{a_i}$. D'où Nous concluons que
$R \to A$ est quasi-fini en $\mathfrak q$.

\medskip\noindent Nous
concluons que Lemme \ref{lemma-locally-P} s'applique à $P$ comme dans
le paragraphe précédent. Il
suffit donc de prouver que $f$ est localement quasi-fini équivaut
à $f$ est localement de type $P$. Puisque $P(R \to A)$
est ``$R \to A$ est quasi-fini'' ce qui
signifie que $R \to A$ est quasi-fini à chaque nombre
premier de $A$, cela résulte du lemme \ref{lemma-quasi-finite-at-point-characterize}.
\end{proof}

\begin{lemma}
\label{lemma-composition-quasi-finite}
La composition de deux morphismes localement quasi-finis est localement
quasi-finie. Il en va de même pour les morphismes quasi-finis.
\end{lemma}

\begin{proof}
Dans la démonstration du lemme \ref{lemma-locally-quasi-finite-characterize} nous avons vu que $P = $``quasi-fini''
est une propriété locale des
homomorphismes d'anneaux, et qu'un morphisme de schémas est localement
quasi-fini si et seulement si il est localement de type
$P$ comme dans Définition \ref{definition-locally-P}. D'où
le premier énoncé du lemme découle du lemme
\ref{lemma-composition-type-P} combiné au fait qu'être quasi-fini est une propriété
des homomorphismes d'anneaux
qui est stable sous composition, voir Algèbre, Lemme \ref{algebra-lemma-quasi-finite-composition}.
Par ce qui précède, Lemme \ref{lemma-quasi-finite-locally-quasi-compact} et le fait que
les compositions de morphismes
quasi-compacts sont quasi-compacts, voir
Schémas, Lemme \ref{schemes-lemma-composition-quasi-compact} on voit que la composition
de morphismes quasi-finis est
quasi-fini.
\end{proof}

\noindent
Nous verrons plus loin (Lemme \ref{lemma-quasi-finite-points-open}) que l'ensemble
$U$ du lemme suivant est ouvert.

\begin{lemma}
\label{lemma-base-change-quasi-finite}
\begin{slogan}
(Localement) les morphismes quasi-finis sont stables sous changement de base.
\end{slogan}
Soit $f : X \to S$ un morphisme de schémas. Soit $g : S' \to S$
un morphisme de schémas. Notons $f' : X' \to S'$
le changement de base de $f$ par $g$ et notons
$g' : X' \to X$ la projection. Supposons que $X$
soit localement de type fini sur $S$.
\begin{enumerate}
\item Soit $U \subset X$ (resp.\ $U' \subset X'$)
l'ensemble des points où $f$ (resp.\ $f'$)
est quasi-fini. Puis $U' = U \times_S S' = (g')^{-1}(U)$.
\item Le changement de base d'un morphisme localement quasi-fini est
localement quasi-fini.
\item Le changement de base d'un morphisme quasi-fini est
quasi-fini.
\end{enumerate}
\end{lemma}

\begin{proof}
La première et la deuxième assertion découlent du
résultat algébrique
correspondant, voir Algèbre, Lemme \ref{algebra-lemma-quasi-finite-base-change}
(combiné au fait que $f'$ est également localement
de type fini par le lemme \ref{lemma-base-change-finite-type}).
Par ce qui précède, Lemme \ref{lemma-quasi-finite-locally-quasi-compact} et le fait
qu'un changement de base d'un
morphisme quasi-compact est quasi-compact,
voir Schémas, Lemme \ref{schemes-lemma-quasi-compact-preserved-base-change} on voit que le changement
de base d'un morphisme quasi-fini est
quasi-fini.
\end{proof}

\begin{lemma}
\label{lemma-quasi-finite-at-a-finite-number-of-points}
Soit $f : X \to S$ un morphisme de schémas de type fini. Soit $s \in S$.
Il y a au plus un nombre fini de points de $X$ situés
sur $s$ auxquels $f$ est quasi-fini.
\end{lemma}

\begin{proof}
La fibre $X_s$ est un schéma de type fini
sur un corps, donc noéthérien (Lemme \ref{lemma-finite-type-noetherian}).
Ainsi la topologie sur $X_s$ est noéthérienne
(Propriétés, Lemme \ref{properties-lemma-Noetherian-topology})
et peut avoir au plus un nombre fini de points isolés (par
topologie élémentaire). Ainsi
notre lemme découle du lemme \ref{lemma-quasi-finite-at-point-characterize}.
\end{proof}

\begin{lemma}
\label{lemma-monomorphism-loc-finite-type-loc-quasi-finite}
Soit $f : X \to Y$ un morphisme de schémas. Si $f$
est localement de type fini et un monomorphisme, alors $f$
est séparé et localement quasi-fini.
\end{lemma}

\begin{proof}
Un monomorphisme est séparé par Schémas,
Lemme \ref{schemes-lemma-monomorphism-separated}. Un monomorphisme
est injectif, on obtient donc que
$f$ est quasi-fini à chaque $x \in X$
par exemple par le lemme \ref{lemma-quasi-finite-at-point-characterize}.
\end{proof}

\begin{lemma}
\label{lemma-immersion-locally-quasi-finite}
Toute immersion est localement quasi-finie.
\end{lemma}

\begin{proof}
Cela est vrai car une immersion ouverte est un isomorphisme local et
une immersion fermée est clairement quasi-finie.
\end{proof}

\begin{lemma}
\label{lemma-permanence-quasi-finite}
Soit $X \to Y$ un morphisme de schémas sur un schéma de base $S$. Soit
$x \in X$. Si $X \to S$ est quasi-fini à $x$, alors $X \to Y$
est quasi-fini à $x$. Si $X$
est localement quasi-fini sur $S$, alors $X \to Y$
est localement quasi-fini.
\end{lemma}

\begin{proof}
Via Lemme \ref{lemma-locally-quasi-finite-characterize} cela se traduit par le fait algébrique suivant
: Étant donné des homomorphismes
d'anneaux $A \to B \to C$ tels que $A \to C$ est
quasi-fini, alors $B \to C$ est quasi-fini.
Cela résulte
de Algèbre, Lemme \ref{algebra-lemma-four-rings} avec $R = A$,
$S = S' = C$ et $R' = B$.
\end{proof}

\begin{lemma}
\label{lemma-quasi-finite-local-source}
Soient $f : X \to Y$ et $g : Y \to S$ des morphismes de schémas. Si $f$
est surjectif, $g \circ f$ localement quasi-fini, et $g$
localement de type fini, alors $g : Y \to S$ est localement quasi-fini.
\end{lemma}

\begin{proof}
Soit $x \in X$ avec les images $y \in Y$ et
$s \in S$. Puisque $g \circ f$ est localement
quasi-fini par le lemme \ref{lemma-residue-field-quasi-finite}
l'extension $\kappa(x)/\kappa(s)$ est finie.
Donc $\kappa(y)/\kappa(s)$ est fini. Donc $y$
est un point fermé de $Y_s$
par le lemme \ref{lemma-algebraic-residue-field-extension-closed-point-fibre}. Puisque $f$ est surjectif,
on voit que chaque point de $Y$ est fermé
dans sa fibre sur $S$.
Ainsi par le lemme \ref{lemma-quasi-finite-at-point-characterize} Nous concluons
que $g$ est quasi-fini en tout point.
\end{proof}














\section{Morphismes de présentation finie}
\label{section-finite-presentation}

\noindent
Rappelons qu'un homomorphisme d'anneaux $R \to A$ est de présentation
finie si $A$ est isomorphe à $R[x_1, \ldots, x_n]/(f_1, \ldots, f_m)$ comme une
$R$-algèbre pour certains $n, m$ et certains polynômes
$f_j$, voir Algèbre, Définition \ref{algebra-definition-finite-type}.

\begin{definition}
\label{definition-finite-presentation}
Soit $f : X \to S$ un morphisme de schémas.
\begin{enumerate}
\item On dit que $f$ est de {\it présentation finie
à $x \in X$} s'il existe un voisinage ouvert affine $\Spec(A) = U \subset X$
de $x$ et ouvert affine $\Spec(R) = V \subset S$ avec
$f(U) \subset V$ tel que l'homomorphisme d'anneaux induit
$R \to A$ est de présentation finie.
\item On dit que $f$ est {\it localement de présentation finie}
s'il est de présentation finie en tout point de $X$.
\item On dit que $f$ est de {\it présentation finie} s'il est
localement de présentation finie, quasi-compact et quasi-séparé.
\end{enumerate}
\end{definition}

\noindent
Notez qu'un morphisme de présentation finie n'est {\bf pas} simplement
un morphisme quasi-compact qui est localement de présentation
finie. Plus tard, nous caractériserons les morphismes
qui sont localement de présentation finie comme les morphismes tels que
$$
\colim \Mor_S(T_i, X) = \Mor_S(\lim T_i, X)
$$
pour tout système orienté de schémas affines $T_i$ sur
$S$.
Voir Limites, Proposition \ref{limits-proposition-characterize-locally-finite-presentation}. Dans Limites, section \ref{limits-section-descending-relative}
nous montrons que, si $S = \lim_i S_i$ est une limite de schémas
affines, tout schéma $X$ de présentation finie
sur $S$ descend à un schéma $X_i$ sur $S_i$
pour certains $i$.

\begin{lemma}
\label{lemma-locally-finite-presentation-characterize}
Soit $f : X \to S$ un morphisme de schémas. Les conditions
suivantes sont équivalentes
\begin{enumerate}
\item Le morphisme $f$ est localement de présentation finie.
\item Pour chaque ouvert affines $U \subset X$, $V \subset S$
avec $f(U) \subset V$ l'homomorphisme
d'anneaux $\mathcal{O}_S(V) \to \mathcal{O}_X(U)$ est de présentation finie.
\item Il existe un recouvrement ouvert $S = \bigcup_{j \in J} V_j$ et
des recouvrements ouverts $f^{-1}(V_j) = \bigcup_{i \in I_j} U_i$ tels que chacun
des morphismes $U_i \to V_j$, $j\in J, i\in I_j$ est localement
de présentation finie.
\item Il existe un recouvrement ouvert affine $S = \bigcup_{j \in J} V_j$
et des recouvrements ouverts affines $f^{-1}(V_j) = \bigcup_{i \in I_j} U_i$ tels
que l'homomorphisme d'anneaux $\mathcal{O}_S(V_j) \to \mathcal{O}_X(U_i)$ soit de présentation
finie, pour tout $j\in J, i\in I_j$.
\end{enumerate}
De plus, si $f$ est localement de présentation finie alors
pour tout sous-schémas ouverts $U \subset X$, $V \subset S$ avec $f(U) \subset V$
la restriction $f|_U : U \to V$ est localement de présentation finie.
\end{lemma}

\begin{proof}
Cela résulte du lemme \ref{lemma-locally-P-characterize} si l'on montre que la propriété ``$R \to A$
est de présentation finie'' est locale.
On vérifie les conditions (a), (b) et (c) de
la Définition \ref{definition-property-local}. Par
Algèbre, Lemme \ref{algebra-lemma-base-change-finiteness}, la propriété d'être de présentation finie est
stable par changement de base ; on en déduit la condition (a).
Par Algèbre, Lemme \ref{algebra-lemma-compose-finite-type},
elle est stable par composition, et
pour tout anneau $R$ l'homomorphisme
d'anneaux $R \to R_f$ est de présentation finie.
On en déduit la condition (b). Enfin, la
condition (c) résulte d'Algèbre, Lemme \ref{algebra-lemma-cover-upstairs}.
\end{proof}

\begin{lemma}
\label{lemma-composition-finite-presentation}
La composition de deux morphismes qui sont localement de présentation finie est localement
de présentation finie. Il
en va de même pour les morphismes de présentation finie.
\end{lemma}

\begin{proof}
Dans la démonstration du lemme \ref{lemma-locally-finite-presentation-characterize} nous avons vu qu'être de
présentation finie est une propriété locale des homomorphismes d'anneaux.
D'où la première affirmation du lemme
découle du lemme \ref{lemma-composition-type-P} combinée au
fait qu'être de présentation finie est une
propriété des homomorphismes
d'anneaux qui est stable
sous composition, voir Algèbre, Lemme \ref{algebra-lemma-compose-finite-type}.
Par ce qui précède et le fait que les compositions
de quasi-compacts, quasi-morphismes
séparés sont quasi-compacts
et quasi-séparés, voir Schémas, Lemmes \ref{schemes-lemma-composition-quasi-compact}
et \ref{schemes-lemma-separated-permanence} on voit que la composition
de morphismes de présentation finie est de présentation
finie.
\end{proof}

\begin{lemma}
\label{lemma-base-change-finite-presentation}
Le changement de base d'un morphisme qui est localement de présentation finie
est localement de présentation finie. Il en va de même pour les morphismes de
présentation finie.
\end{lemma}

\begin{proof}
Dans la démonstration du lemme \ref{lemma-locally-finite-presentation-characterize} nous avons vu qu'être de
présentation finie est une propriété locale des homomorphismes
d'anneaux. D'où le premier énoncé du lemme
découle du lemme \ref{lemma-composition-type-P} combiné au
fait qu'être de présentation finie est une
propriété des homomorphismes
d'anneaux qui est stable
sous changement de base, voir Algèbre, Lemme \ref{algebra-lemma-base-change-finiteness}.
Par ce qui précède et le fait qu'un
changement de base d'un quasi-compact, quasi-morphisme
séparé est quasi-compact
et quasi-séparé, voir Schémas, Lemmes \ref{schemes-lemma-quasi-compact-preserved-base-change}
et \ref{schemes-lemma-separated-permanence} on voit que le changement
de base d'un morphisme de présentation finie est un morphisme
de présentation finie.
\end{proof}

\begin{lemma}
\label{lemma-open-immersion-locally-finite-presentation}
Toute immersion ouverte est localement de présentation finie.
\end{lemma}

\begin{proof}
Ceci est vrai car une immersion ouverte est un isomorphisme local.
\end{proof}

\begin{lemma}
\label{lemma-quasi-compact-open-immersion-finite-presentation}
Toute immersion ouverte est de présentation finie si et seulement si elle
est quasi-compacte.
\end{lemma}

\begin{proof}
Nous avons vu (Lemme \ref{lemma-open-immersion-locally-finite-presentation}) qu'une immersion ouverte
est localement de présentation finie.
Nous avons vu (Schémas, Lemme \ref{schemes-lemma-immersions-monomorphisms}) qu'une immersion
est séparée et donc quasi-séparée. De ceci
et de la Définition \ref{definition-finite-presentation} découle le lemme.
\end{proof}

\begin{lemma}
\label{lemma-closed-immersion-finite-presentation}
\begin{slogan}
Les immersions fermées de présentation finie correspondent
à des faisceaux quasi-cohérents d'idéaux de type fini.
\end{slogan}
Une immersion fermée $i : Z \to X$ est de présentation finie si
et seulement si le faisceau quasi-cohérent
d'idéaux associé $\mathcal{I} = \Ker(\mathcal{O}_X \to i_*\mathcal{O}_Z)$ est de type
fini (en tant que $\mathcal{O}_X$-module).
\end{lemma}

\begin{proof}
Sur tout $\Spec(R) \subset X$ ouvert affine, nous avons
$i^{-1}(\Spec(R)) = \Spec(R/I)$ et $\mathcal{I} = \widetilde{I}$. De plus,
$\mathcal{I}$ est de type fini si et seulement si $I$
est un $R$-module fini pour chaque ouvert affine
de ce type (voir Propriétés,
Lemme \ref{properties-lemma-finite-type-module}). Et $R/I$ est de présentation finie
sur $R$ si et seulement si
$I$ est un $R$-module fini. On obtient donc le résultat.
\end{proof}

\begin{lemma}
\label{lemma-finite-presentation-finite-type}
Un morphisme qui est localement de présentation finie est localement de type fini.
Un morphisme de présentation finie est de type fini.
\end{lemma}

\begin{proof}
Omis.
\end{proof}

\begin{lemma}
\label{lemma-noetherian-finite-type-finite-presentation}
\begin{slogan}
Sur une base localement noethérien, le type fini est présentation finie.
\end{slogan}
Soit $f : X \to S$ un morphisme.
\begin{enumerate}
\item Si $S$ est localement noethérien et $f$ localement de type fini
alors $f$ est localement de présentation finie.
\item Si $S$ est localement noethérien et $f$ de type fini
alors $f$ est de présentation finie.
\end{enumerate}
\end{lemma}

\begin{proof}
La première affirmation découle du fait qu'un
anneau de type fini sur un anneau noéthérien est de présentation
finie, voir Algèbre, Lemme \ref{algebra-lemma-Noetherian-finite-type-is-finite-presentation}. Supposons que $f$
soit de type fini et que $S$ soit localement
noethérien. Alors $f$ est quasi-compact et
localement de fini présenté par (1). Il suffit
donc de prouver que $f$ est quasi-séparé. Cela résulte
du lemme \ref{lemma-finite-type-Noetherian-quasi-separated} (et Lemme \ref{lemma-finite-presentation-finite-type}).
\end{proof}

\begin{lemma}
\label{lemma-finite-presentation-quasi-compact-quasi-separated}
Soit $S$ un schéma quasi-compact et quasi-séparé. Si $X$
est de présentation finie sur $S$, alors $X$ est quasi-compact
et quasi-séparé.
\end{lemma}

\begin{proof}
Omis.
\end{proof}

\begin{lemma}
\label{lemma-finite-presentation-permanence}
Soit $f : X \to Y$ un morphisme de schémas sur $S$.
\begin{enumerate}
\item Si $X$ est localement de présentation finie sur $S$ et $Y$
est localement de type fini sur $S$, alors $f$ est localement
de présentation finie.
\item Si $X$ est de présentation finie sur $S$ et $Y$ est quasi-séparé
et localement de type fini sur $S$, alors $f$ est de présentation finie.
\end{enumerate}
\end{lemma}

\begin{proof}
Démonstration de (1). Via Lemme \ref{lemma-locally-finite-presentation-characterize} cela se traduit par le fait algébrique
suivant : Étant donné les homomorphismes
d'anneaux $A \to B \to C$ tels que $A \to C$
est de présentation finie et $A \to B$ est de
type fini, alors $B \to C$ est de présentation
finie. Voir Algèbre, Lemme \ref{algebra-lemma-compose-finite-type}.

\medskip\noindent
La partie (2)
découle de (1) et des Schémas, Lemmes
\ref{schemes-lemma-compose-after-separated} et \ref{schemes-lemma-quasi-compact-permanence}.
\end{proof}

\begin{lemma}
\label{lemma-diagonal-morphism-finite-type}
Soit $f : X \to Y$ un morphisme de schémas de diagonale $\Delta : X \to X \times_Y X$.
Si $f$ est localement de type fini alors $\Delta$
est localement de présentation finie. Si $f$ est quasi-séparé et
localement de type fini, alors $\Delta$ est de présentation finie.
\end{lemma}

\begin{proof}
Notez que $\Delta$ est un morphisme de schémas sur $X$
(via la deuxième projection $X \times_Y X \to X$). Supposons que $f$
soit localement de type fini. Notez que $X$ est de présentation
finie sur $X$ et $X \times_Y X$ est localement de type fini sur $X$
(par le lemme \ref{lemma-base-change-finite-type}).
Ainsi, la première affirmation du lemme
\ref{lemma-finite-presentation-permanence} est valable. Le deuxième énoncé découle du premier,
des définitions, et du fait qu'un morphisme diagonal est un monomorphisme,
donc séparé (Schémas, Lemme \ref{schemes-lemma-monomorphism-separated}).
\end{proof}








\section{Ensembles constructibles}
\label{section-constructible}

\noindent
Les ensembles constructibles et localement constructibles des schémas ont été discutés
dans Propriétés, section \ref{properties-section-constructible}. Dans cette section nous prouvons
quelques résultats concernant les images et les images inverses d'ensembles
(localement) constructibles. Le résultat principal est le théorème de Chevalley
qui affirme que l'image d'un ensemble localement constructible sous un morphisme
de présentation finie est localement constructible.

\begin{lemma}
\label{lemma-inverse-image-constructible}
Soit $f : X \to Y$ un morphisme de schémas. Soit $E \subset Y$
un sous-ensemble. Si $E$ est
(localement) constructible dans $Y$, alors $f^{-1}(E)$ est (localement) constructible
dans $X$.
\end{lemma}

\begin{proof}
Pour montrer que l'image réciproque de chaque sous-ensemble constructible est constructible,
il suffit de montrer que l'image réciproque de chaque $V$ ouvert rétrocompact
de $Y$ est rétrocompacte
dans $X$, voir Topologie, Lemme \ref{topology-lemma-inverse-images-constructibles}. L'importance
du fait que $V$ soit rétrocompact dans
$Y$ est simplement que l'immersion ouverte $V \to Y$ est quasi-compacte.
Ainsi le changement de base $f^{-1}(V) = X \times_Y V \to X$ est également quasi-compact,
voir
Schémas, Lemme \ref{schemes-lemma-quasi-compact-preserved-base-change}. Nous voyons donc que $f^{-1}(V)$ est rétrocompact
dans $X$. Supposons que $E$ soit localement
constructible dans $Y$. Choisissez
$x \in X$. Choisir un voisinage affine $V$ de $f(x)$ et
un voisinage affine $U \subset X$ de $x$ tel que $f(U) \subset V$. Ainsi
nous pensons à $f|_U : U \to V$ comme un morphisme en $V$.
Par Propriétés, Lemme \ref{properties-lemma-locally-constructible} nous voyons
que $E \cap V$ est constructible dans $V$. Par le cas constructible
nous voyons que $(f|_U)^{-1}(E \cap V)$ est constructible
dans $U$. Puisque $(f|_U)^{-1}(E \cap V) = f^{-1}(E) \cap U$ on obtient le résultat.
\end{proof}

\begin{lemma}
\label{lemma-chevalley}
Soit $f : X \to Y$ un morphisme de schémas. Supposons
que
\begin{enumerate}
\item $f$ est quasi-compact et localement de présentation finie, et que
\item $Y$ est quasi-compact et quasi-séparé.
\end{enumerate}
Alors l'image de chaque sous-ensemble constructible de $X$ est constructible dans $Y$.
\end{lemma}

\begin{proof}
Par
Propriétés,
Lemme \ref{properties-lemma-constructible-quasi-compact-quasi-separated} il suffit de prouver ce lemme dans le cas
où $Y$ est affine. Dans ce cas, $X$
est quasi-compact. On peut donc écrire
$X = U_1 \cup \ldots \cup U_n$ avec chaque $U_i$ ouvert affine dans $X$.
Si $E \subset X$ est constructible, alors chaque $E \cap U_i$
est
également
constructible, voir Topologie, Lemme \ref{topology-lemma-open-immersion-constructible-inverse-image}. Donc,
puisque $f(E) = \bigcup f(E \cap U_i)$ et que les unions finies d'ensembles constructibles
sont constructibles, cela nous ramène au cas où
$X$ est affine. Dans ce cas, le
résultat est Algèbre, Théorème \ref{algebra-theorem-chevalley}.
\end{proof}

\begin{theorem}[Théorème de Chevalley]
\label{theorem-chevalley}
\begin{reference}
\cite[IV, Théorème 1.8.4]{EGA}
\end{reference}
Soit $f : X \to Y$ un morphisme de schémas. Supposons
que $f$ soit quasi-compact et localement de présentation finie. Alors
l’image de chaque sous-ensemble localement constructible est localement constructible.
\end{theorem}

\begin{proof}
Soit $E \subset X$ localement constructible. Il
faut montrer que $f(E)$ est également constructible
localement. Nous allons montrer que $f(E) \cap V$ est constructible
pour tout $V \subset Y$ ouvert affine. On se réduit donc au cas où
$Y$ est affine. Dans ce cas, $X$ est quasi-compact.
On peut donc écrire $X = U_1 \cup \ldots \cup U_n$ avec chaque $U_i$
ouvert affine dans $X$. Si $E \subset X$ est localement constructible,
alors chaque
$E \cap U_i$ est constructible, voir Propriétés, Lemme \ref{properties-lemma-locally-constructible}.
Donc, puisque $f(E) = \bigcup f(E \cap U_i)$ et que les unions finies d'ensembles constructibles
sont constructibles, cela nous ramène au cas où $X$ est
affine. Dans ce cas, le résultat est
Algèbre, Théorème \ref{algebra-theorem-chevalley}.
\end{proof}

\begin{lemma}
\label{lemma-constructible-containing-open}
Soit $X$ un schéma. Soit $x \in X$. Soit $E \subset X$ un sous-ensemble
localement constructible. Si $\{x' \mid x' \leadsto x\} \subset E$, alors $E$
contient un voisinage ouvert de $x$.
\end{lemma}

\begin{proof}
Supposons $\{x' \mid x' \leadsto x\} \subset E$. Nous pouvons supposer
que $X$ est
affine. Dans ce cas, $E$
est constructible, voir Propriétés, Lemme \ref{properties-lemma-locally-constructible}.
En particulier, le complément $E^c$ est
également constructible. D'après Algèbre, Lemme \ref{algebra-lemma-constructible-is-image},
on peut trouver un morphisme de schémas affines
$f : Y \to X$ tel que $E^c = f(Y)$. Soit $Z \subset X$ l'image schématique
de $f$. Par le Lemme \ref{lemma-reach-points-scheme-theoretic-image}
et l'hypothèse $\{x' \mid x' \leadsto x\} \subset E$ on voit que $x \not \in Z$.
Ainsi $X \setminus Z \subset E$ est un voisinage ouvert de $x$
contenu dans $E$.
\end{proof}




\section{Morphismes ouverts}
\label{section-open}

\begin{definition}
\label{definition-open}
Soit $f : X \to S$ un morphisme.
\begin{enumerate}
\item On dit que $f$ est {\it ouvert} si l'application sur les espaces
topologiques sous-jacents est ouverte.
\item On dit que $f$ est {\it universellement ouvert} si
pour tout morphisme de schémas $S' \to S$ le changement de base $f' : X_{S'} \to S'$ est ouvert.
\end{enumerate}
\end{definition}

\noindent
Selon la topologie,
les générisations du lemme \ref{topology-lemma-closed-open-map-specialization} s'élèvent le long
de certains types d'applications ouvertes d'espaces topologiques.
En fait, les générisations suivent tout morphisme de schémas ouvert
(voir
Lemme \ref{lemma-open-generizing}). Aussi,
nous verrons que les générisations s'élèvent le long de morphismes
de schémas plats (Lemme \ref{lemma-generalizations-lift-flat}). Cela implique
parfois à son tour que le morphisme est ouvert.

\begin{lemma}
\label{lemma-locally-finite-presentation-universally-open}
Soit $f : X \to S$ un morphisme.
\begin{enumerate}
\item Si $f$ est localement de présentation finie et que les générisations s'élèvent le long de $f$,
alors $f$ est ouvert.
\item Si $f$ est localement de présentation finie et que les générisations s'élèvent à chaque
changement de base de $f$, alors $f$ est universellement ouvert.
\end{enumerate}
\end{lemma}

\begin{proof}
Il suffit de prouver la première assertion.
Cela se réduit au cas où $X$ et
$S$ sont affines. Dans
ce cas, le résultat découle de l'algèbre, du Lemme \ref{algebra-lemma-going-up-down-specialization}
et de la proposition \ref{algebra-proposition-fppf-open}.
\end{proof}

\noindent
Voir aussi Lemme \ref{lemma-fppf-open} pour le cas d'un morphisme plat
de présentation finie.

\begin{lemma}
\label{lemma-composition-open}
Une composition de morphismes (universellement) ouverts est (universellement) ouverte.
\end{lemma}

\begin{proof}
Omis.
\end{proof}

\begin{lemma}
\label{lemma-scheme-over-field-universally-open}
Soit $k$ un corps. Soit $X$ un schéma sur $k$.
Le morphisme de structure $X \to \Spec(k)$ est universellement ouvert.
\end{lemma}

\begin{proof}
Soit $S \to \Spec(k)$ un morphisme. Il faut
montrer que le changement de base $X_S \to S$ est ouvert.
La question est locale sur $S$ et $X$, on peut
donc supposer que $S$ et $X$ sont affines.
Dans ce cas, le résultat est Algèbre, Lemme \ref{algebra-lemma-map-into-tensor-algebra-open}.
\end{proof}

\begin{lemma}
\label{lemma-open-generizing}
\begin{reference}
découle de l'implication (a) $\Rightarrow$ (b) dans \cite[IV, Corollaire
1.10.4]{EGA}
\end{reference}
Soit $\varphi : X \to Y$ un morphisme de schémas. Si $\varphi$ est ouvert,
alors $\varphi$ est en train de générer (c'est-à-dire que
les générisations s'élèvent le long de $\varphi$). Si
$\varphi$ est universellement ouvert, alors $\varphi$ est universellement
génératrice.
\end{lemma}

\begin{proof}
Supposons que $\varphi$ soit
ouvert. Soit $y' \leadsto y$ une spécialisation de points de
$Y$. Soit $x \in X$ avec $\varphi(x) = y$.
Choisissons des ouverts affines $U \subset X$ et $V \subset Y$
tels que $\varphi(U) \subset V$ et $x \in U$. Alors $y' \in V$ aussi. Nous pouvons
donc remplacer $X$ par $U$ et $Y$ par $V$
et supposer
que $X$, $Y$ sont affines. Le cas affine
est l'Algèbre,
Lemme \ref{algebra-lemma-open-going-down} (combiné avec l'Algèbre, Lemme \ref{algebra-lemma-going-up-down-specialization}).
\end{proof}

\begin{lemma}
\label{lemma-descent-quasi-compact}
Soit $f : X \to Y$ un morphisme de schémas. Soit
$g : Y' \to Y$ ouvert et surjectif tel que le changement de base $f' : X' \to Y'$
soit quasi-compact. Alors $f$ est quasi-compact.
\end{lemma}

\begin{proof}
Soit $V \subset Y$ un ouvert quasi-compact. Comme $g$ est ouvert
et surjectif on peut trouver un $W' \subset W$ ouvert quasi-compact tel
que $g(W') = V$. Par hypothèse $(f')^{-1}(W')$ est quasi-compact.
L'image de $(f')^{-1}(W')$ dans $X$ est
égale à $f^{-1}(V)$, voir Lemme \ref{lemma-when-point-maps-to-pair}.
Ainsi $f^{-1}(V)$ est quasi-compact comme l'image d'un espace quasi-compact,
voir Topologie, Lemme \ref{topology-lemma-image-quasi-compact}. Ainsi $f$
est quasi-compact.
\end{proof}





\section{Morphismes submersifs}
\label{section-submersive}

\begin{definition}
\label{definition-submersive}
Soit $f : X \to Y$ un morphisme de schémas.
\begin{enumerate}
\item Nous disons que $f$ est {\it submersif}\footnote{Ceci
est très différent de la notion de submersion de
variétés différentielles.} si l'application continue des espaces topologiques
sous-jacents est submersive, voir Topologie, Définition \ref{topology-definition-submersive}.
\item On dit que $f$ est {\it universellement submersif}
si pour tout morphisme de schémas $Y' \to Y$ le
changement de base $Y' \times_Y X \to Y'$ est submersif.
\end{enumerate}
\end{definition}

\noindent
On remarque qu'un morphisme submersif est notamment surjectif.

\begin{lemma}
\label{lemma-base-change-universally-submersive}
Le changement de base d'un morphisme de schémas universellement submersif
par tout morphisme de schémas est universellement submersif.
\end{lemma}

\begin{proof}
Cela découle directement de la définition.
\end{proof}

\begin{lemma}
\label{lemma-composition-universally-submersive}
La composition d’une paire de morphismes de schémas (universellement) submersifs
est (universellement) submersive.
\end{lemma}

\begin{proof}
Omis.
\end{proof}










\section{Morphismes plats}
\label{section-flat}

\noindent
La planéité est l'un des outils techniques les plus importants en
géométrie algébrique. Dans cette section, nous introduisons cette notion.
Nous limitons intentionnellement la discussion à des observations
simples, en dehors du lemme \ref{lemma-fppf-open}. Une classe
de résultats
très importante, à savoir les critères de
planéité, est discutée dans Algèbre,
sections \ref{algebra-section-criteria-flatness}, \ref{algebra-section-flatness-artinian}, \ref{algebra-section-more-flatness-criteria},
et plus sur les
morphismes, section \ref{more-morphisms-section-criterion-flat-fibres}. Il existe un
chapitre consacré au matériel avancé sur les morphismes de
schémas plats, à savoir Compléments sur la platitude, section \ref{flat-section-introduction}.

\medskip\noindent Rappelons
qu'un module $M$ sur un anneau $R$ est {\it plat}
si le foncteur $-\otimes_R M : \text{Mod}_R \to \text{Mod}_R$ est exact. Un homomorphisme d'anneaux $R \to A$
est dit {\it plat} si $A$ est plat comme $R$-module.
Voir
Algèbre, Définition \ref{algebra-definition-flat}.

\begin{definition}
\label{definition-flat}
Soit $f : X \to S$ un morphisme de schémas.
Soit $\mathcal{F}$ un faisceau quasi-cohérent de $\mathcal{O}_X$-modules.
\begin{enumerate}
\item On dit que $f$ est {\it plat en un point
$x \in X$} si l'anneau local $\mathcal{O}_{X, x}$ est plat sur
l'anneau local $\mathcal{O}_{S, f(x)}$.
\item On dit que $\mathcal{F}$ est {\it plat sur $S$ en un point
$x \in X$} si la tige $\mathcal{F}_x$ est un module $\mathcal{O}_{S, f(x)}$ plat.
\item Nous disons que $f$ est {\it plat} si $f$ est plat en chaque point de $X$.
\item On dit que $\mathcal{F}$ est {\it plat sur $S$}
si $\mathcal{F}$ est plat sur $S$ en chaque point $x$ de $X$.
\end{enumerate}
\end{definition}

\noindent
On voit donc que $f$ est plat si et seulement
si le faisceau structurel $\mathcal{O}_X$ est plat sur $S$.

\begin{lemma}
\label{lemma-flat-module-characterize}
Soit $f : X \to S$ un morphisme de schémas. Soit
$\mathcal{F}$ un faisceau quasi-cohérent de $\mathcal{O}_X$-modules. Les conditions
suivantes sont équivalentes
\begin{enumerate}
\item Le faisceau $\mathcal{F}$ est plat sur $S$.
\item Pour chaque ouverture affine $U \subset X$, $V \subset S$
avec $f(U) \subset V$ le $\mathcal{O}_S(V)$-module $\mathcal{F}(U)$ est plat.
\item Il existe un recouvrement ouvert $S = \bigcup_{j \in J} V_j$ et
des revêtements ouverts $f^{-1}(V_j) = \bigcup_{i \in I_j} U_i$ tels que chacun des
modules $\mathcal{F}|_{U_i}$ soit à plat sur $V_j$,
pour tous les $j\in J, i\in I_j$.
\item Il existe un recouvrement ouvert affine $S = \bigcup_{j \in J} V_j$
et des recouvrements ouverts affines $f^{-1}(V_j) = \bigcup_{i \in I_j} U_i$ tels
que $\mathcal{F}(U_i)$ soit un $\mathcal{O}_S(V_j)$-module plat, pour tout
$j\in J, i\in I_j$.
\end{enumerate}
De plus, si $\mathcal{F}$ est plat sur $S$ alors
pour tout sous-schémas ouverts $U \subset X$, $V \subset S$ avec $f(U) \subset V$ la
restriction $\mathcal{F}|_U$ est plate sur $V$.
\end{lemma}

\begin{proof}
Soit $R \to A$ un homomorphisme d'anneaux. Soit $M$
un module $A$. Si $M$ est $R$-plat,
alors pour tous les nombres premiers $\mathfrak q$ le module $M_{\mathfrak q}$
est plat sur $R_{\mathfrak p}$ avec $\mathfrak p$ le premier de $R$ se trouvant sous $\mathfrak q$.
Inversement, si $M_{\mathfrak q}$ est plat sur $R_{\mathfrak p}$ pour tous les nombres premiers
$\mathfrak q$ de $A$, alors $M$
est plat sur $R$. Voir Algèbre, Lemme \ref{algebra-lemma-flat-localization}.
Cette équivalence implique facilement les énoncés du lemme.
\end{proof}

\begin{lemma}
\label{lemma-flat-characterize}
Soit $f : X \to S$ un morphisme de schémas. Les conditions
suivantes sont équivalentes
\begin{enumerate}
\item Le morphisme $f$ est plat.
\item Pour chaque ouvert affines $U \subset X$, $V \subset S$
avec $f(U) \subset V$ l'homomorphisme
d'anneaux $\mathcal{O}_S(V) \to \mathcal{O}_X(U)$ est plat.
\item Il existe un recouvrement ouvert $S = \bigcup_{j \in J} V_j$ et
des recouvrements ouverts $f^{-1}(V_j) = \bigcup_{i \in I_j} U_i$ tels que chacun
des morphismes $U_i \to V_j$, $j\in J, i\in I_j$ soit
plat.
\item Il existe un recouvrement ouvert affine $S = \bigcup_{j \in J} V_j$
et un recouvrement ouvert affine $f^{-1}(V_j) = \bigcup_{i \in I_j} U_i$
tel que $\mathcal{O}_S(V_j) \to \mathcal{O}_X(U_i)$ soit plat, pour tout
$j\in J, i\in I_j$.
\end{enumerate}
De plus, si $f$ est plat alors
pour tout sous-schémas ouverts $U \subset X$, $V \subset S$ avec $f(U) \subset V$
la restriction $f|_U : U \to V$ est plate.
\end{lemma}

\begin{proof}
Ceci est un cas particulier du lemme \ref{lemma-flat-module-characterize}
ci-dessus.
\end{proof}

\begin{lemma}
\label{lemma-pushforward-flat-affine}
Soit $f : X \to Y$ un morphisme affine de schémas sur un schéma de base $S$.
Soit $\mathcal{F}$ un module $\mathcal{O}_X$ quasi-cohérent. Alors $\mathcal{F}$
est plat sur $S$ si et seulement si $f_*\mathcal{F}$
est plat sur $S$.
\end{lemma}

\begin{proof}
Par le Lemme \ref{lemma-flat-module-characterize} et le fait que $f$ soit un morphisme
affine, cela nous ramène au cas affine. Supposons que $X \to Y \to S$
corresponde aux homomorphismes d'anneaux $C \leftarrow B \leftarrow A$. Soit $N$
le module $C$ correspondant à $\mathcal{F}$.
Rappelons que $f_*\mathcal{F}$ correspond à $N$ vu comme un $B$-module,
voir Schémas, Lemme \ref{schemes-lemma-widetilde-pullback}. Le résultat
est donc clair.
\end{proof}

\begin{lemma}
\label{lemma-composition-module-flat}
Soit $X \to Y \to Z$ des morphismes de schémas. Soit $\mathcal{F}$ un module $\mathcal{O}_X$ quasi-cohérent.
Soit $x \in X$ avec l'image $y$ dans $Y$. Si $\mathcal{F}$ est plat
sur $Y$ à $x$ et que $Y$ est plat sur $Z$ à $y$,
alors $\mathcal{F}$ est plat sur $Z$ à $x$.
\end{lemma}

\begin{proof}
Voir Algèbre, Lemme \ref{algebra-lemma-composition-flat}.
\end{proof}

\begin{lemma}
\label{lemma-composition-flat}
La composition des morphismes plats est plate.
\end{lemma}

\begin{proof}
Ceci est un cas particulier du lemme \ref{lemma-composition-module-flat}.
\end{proof}

\begin{lemma}
\label{lemma-base-change-module-flat}
Soit $f : X \to S$ un morphisme de schémas. Soit
$\mathcal{F}$ un faisceau quasi-cohérent de $\mathcal{O}_X$-modules. Soit
$g : S' \to S$ un morphisme de schémas. Notons
$g' : X' = X_{S'} \to X$ la projection. Soit $x' \in X'$ un
point d'image $x = g'(x') \in X$. Si $\mathcal{F}$ est plat sur
$S$ à $x$, alors $(g')^*\mathcal{F}$ est
plat sur $S'$ à $x'$. En particulier,
si $\mathcal{F}$ est plat sur $S$, alors
$(g')^*\mathcal{F}$ est plat sur $S'$.
\end{lemma}

\begin{proof}
Voir Algèbre, Lemme \ref{algebra-lemma-flat-base-change}.
\end{proof}

\begin{lemma}
\label{lemma-base-change-flat}
Le changement de base d'un morphisme plat est plat.
\end{lemma}

\begin{proof}
Ceci est un cas particulier du lemme \ref{lemma-base-change-module-flat}.
\end{proof}

\begin{lemma}
\label{lemma-generalizations-lift-flat}
Soit $f : X \to S$ un morphisme de schémas plat.
Puis les générisations s'élèvent le
long de $f$, voir Topologie, Définition \ref{topology-definition-lift-specializations}.
\end{lemma}

\begin{proof}
Voir Algèbre, section \ref{algebra-section-going-up}.
\end{proof}

\begin{lemma}
\label{lemma-fppf-open}
A morphisme plat localement de fini présentation est universellement ouvert.
\end{lemma}

\begin{proof}
Cela résulte de Lemmes \ref{lemma-generalizations-lift-flat} et Lemme \ref{lemma-locally-finite-presentation-universally-open}
ci-dessus. Nous pouvons également argumenter
directement comme suit.

\medskip\noindent
Soit $f : X \to S$ soit plat et localement de présentation finie.
Par les Lemmes \ref{lemma-base-change-flat} et \ref{lemma-base-change-finite-presentation}
tout changement de base de $f$ est plat et localement de présentation
finie. Il suffit donc de montrer que $f$ est ouvert. Pour
montrer que $f$ est ouvert, il suffit de montrer
que l'on peut couvrir $X$ par des ouverts affines
$X = \bigcup U_i$ telles que $U_i \to S$ soit ouvert. On peut recouvrir $X$
par des ouverts affines $U_i \subset X$ tels que chaque
$U_i$ ait son image dans un ouvert affine $V_i \subset S$ et tels
que l'homomorphisme d'anneaux induit $\mathcal{O}_S(V_i) \to \mathcal{O}_X(U_i)$ soit plat et de présentation
finie (Lemmes \ref{lemma-flat-characterize} et \ref{lemma-locally-finite-presentation-characterize}).
Alors $U_i \to V_i$ est ouvert
par Algèbre, Proposition \ref{algebra-proposition-fppf-open} et la démonstration
est complète.
\end{proof}

\begin{lemma}
\label{lemma-pf-flat-module-open}
Soit $f : X \to Y$ un morphisme de schémas. Soit
$\mathcal{F}$ un module $\mathcal{O}_X$ quasi-cohérent. Supposons
que $f$ localement présentation finie, $\mathcal{F}$
de type fini, $X = \text{Supp}(\mathcal{F})$ et $\mathcal{F}$ plat sur
$Y$. Alors $f$ est universellement ouvert.
\end{lemma}

\begin{proof}
Par les Lemmes \ref{lemma-base-change-module-flat}, \ref{lemma-base-change-finite-presentation},
et \ref{lemma-support-finite-type} les hypothèses
sont conservées sous
changement de base. Par le Lemme
\ref{lemma-locally-finite-presentation-universally-open} il suffit de montrer que les générisations
s'élèvent le long de $f$.
Cela résulte de Algèbre, Lemme \ref{algebra-lemma-going-down-flat-module}.
\end{proof}

\begin{lemma}
\label{lemma-fpqc-quotient-topology}
\begin{reference}
\cite[Expose VIII, Corollaire 4.3]{SGA1} et \cite[IV,
Corollaire 2.3.12]{EGA}
\end{reference}
Soit $f : X \to Y$ être un quasi-compact, surjectif, morphisme plat. Un sous-ensemble
$T \subset Y$ est ouvert (resp.\ fermé) si et seulement $f^{-1}(T)$
est ouvert (resp.\ fermé). Autrement dit, $f$ est
un morphisme submersif.
\end{lemma}

\begin{proof}
La question est locale sur $Y$, on peut donc supposer que $Y$
est affine. Dans ce cas, $X$ est quasi-compact
tout comme $f$ est quasi-compact. Écrivez $X = X_1 \cup \ldots \cup X_n$ comme
une union finie d'ouverts affines. Alors $f' : X' = X_1 \amalg \ldots \amalg X_n \to Y$ est un
morphisme surjectif plat de schémas affines. Notez que pour $T \subset Y$
nous avons $(f')^{-1}(T) = f^{-1}(T) \cap X_1 \amalg \ldots \amalg f^{-1}(T) \cap X_n$. Par conséquent, $f^{-1}(T)$ est ouvert si et seulement
si $(f')^{-1}(T)$ est ouvert. Ainsi, nous pouvons supposer
que $X$ et $Y$ sont affines.

\medskip\noindent
Soit $f : \Spec(B) \to \Spec(A)$ un morphisme surjectif de schémas affines
correspondant à un homomorphisme plat d'anneaux $A \to B$.
Supposons que $f^{-1}(T)$ soit fermé, disons $f^{-1}(T) = V(J)$
pour $J \subset B$ un idéal. Alors $T = f(f^{-1}(T)) = f(V(J))$ est l'image
de $\Spec(B/J) \to \Spec(A)$ (ici on utilise que $f$ est surjectif).
En revanche, les générisations s'élèvent le long
de $f$ (Lemme \ref{lemma-generalizations-lift-flat}). Ainsi par Topologie, Lemme
\ref{topology-lemma-lift-specializations-images} on voit que $Y \setminus T = f(X \setminus f^{-1}(T))$ est stable par générisation.
Ainsi $T$ est stable par spécialisation
(Topologie, Lemme \ref{topology-lemma-open-closed-specialization}). Ainsi $T$
est fermé par
Algèbre, Lemme \ref{algebra-lemma-image-stable-specialization-closed}.
\end{proof}

\begin{lemma}
\label{lemma-flat-permanence}
Soit $h : X \to Y$ un morphisme de schémas sur $S$. Soit
$\mathcal{G}$ un faisceau quasi-cohérent sur $Y$. Soit
$x \in X$ avec $y = h(x) \in Y$. Si $h$ est plat en $x$, alors
$$
\mathcal{G}\text{ est plat sur }S\text{ en }y
\Leftrightarrow
h^*\mathcal{G}\text{ est plat sur }S\text{ en }x.
$$
En particulier : Si $h$ est surjectif et plat, alors $\mathcal{G}$
est plat sur $S$, si et seulement si $h^*\mathcal{G}$
est plat sur $S$. Si $h$ est surjectif et plat, et que $X$
est plat sur $S$, alors $Y$ est plat sur $S$.
\end{lemma}

\begin{proof}
Vous pouvez le prouver
en appliquant l'algèbre, Lemme \ref{algebra-lemma-flatness-descends-more-general}.
Voici une démonstration directe. Soit $s \in S$
l'image de
$y$. Considérez les applications locales
anneau $\mathcal{O}_{S, s} \to \mathcal{O}_{Y, y} \to \mathcal{O}_{X, x}$. Par hypothèse l'homomorphisme
d'anneaux $\mathcal{O}_{Y, y} \to \mathcal{O}_{X, x}$
est fidèlement plat, voir
Algèbre, Lemme \ref{algebra-lemma-local-flat-ff}.
Soit $N = \mathcal{G}_y$. A noter que $h^*\mathcal{G}_x = N \otimes_{\mathcal{O}_{Y, y}} \mathcal{O}_{X, x}$, voir
Faisceaux, Lemme \ref{sheaves-lemma-stalk-pullback-modules}. Soit $M' \to M$
une injection de $\mathcal{O}_{S, s}$-modules.
Par la planéité fidèle mentionnée ci-dessus, nous avons
\begin{align*}
\Ker(
M' \otimes_{\mathcal{O}_{S, s}} N \to M \otimes_{\mathcal{O}_{S, s}} N)
\otimes_{\mathcal{O}_{Y, y}} \mathcal{O}_{X, x} \\
=
\Ker(
M' \otimes_{\mathcal{O}_{S, s}} N
\otimes_{\mathcal{O}_{Y, y}} \mathcal{O}_{X, x}
\to
M \otimes_{\mathcal{O}_{S, s}} N
\otimes_{\mathcal{O}_{Y, y}} \mathcal{O}_{X, x})
\end{align*}
Par conséquent, l'équivalence du lemme découle de la deuxième caractérisation
de la platitude donnée dans
Algèbre, Lemme \ref{algebra-lemma-flat}.
\end{proof}

\begin{lemma}
\label{lemma-flat-pullback-support}
Soit $f : Y \to X$ un morphisme de schémas. Soit $\mathcal{F}$ un
$\mathcal{O}_X$-module quasi-cohérent de type fini avec support
schématique $Z \subset X$. Si $f$ est plat,
alors $f^{-1}(Z)$ est le support schématique de $f^*\mathcal{F}$.
\end{lemma}

\begin{proof}
En utilisant la caractérisation schématique des supports
sur affines telle que donnée dans le lemme \ref{lemma-scheme-theoretic-support}
nous réduisons à l'Algèbre, Lemme \ref{algebra-lemma-annihilator-flat-base-change}.
\end{proof}

\begin{lemma}
\label{lemma-flat-morphism-scheme-theoretically-dense-open}
Soit $f : X \to Y$ un morphisme de schémas plat. Soit $V \subset Y$ un ouvert
rétrocompact qui est schématiquement dense. Alors $f^{-1}V$ est schématiquement
dense en $X$.
\end{lemma}

\begin{proof}
Nous utiliserons la caractérisation
du lemme \ref{lemma-characterize-scheme-theoretically-dense}. Il faut montrer
que pour tout $U \subset X$ ouvert l'application
$\mathcal{O}_X(U) \to \mathcal{O}_X(U \cap f^{-1}V)$ est injectif. Il suffit
de le prouver lorsque $U$ est un ouvert affine
qui s'envoie dans un ouvert affine $W \subset Y$.
Posons $W = \Spec(A)$ et $U = \Spec(B)$. Puis
$V \cap W = D(f_1) \cup \ldots \cup D(f_n)$
pour certains $f_i \in A$, voir
Algèbre, Lemme \ref{algebra-lemma-qc-open}.
Il faut donc montrer que $B \to B_{f_1} \times \ldots \times B_{f_n}$
est injectif. On sait que $A \to A_{f_1} \times \ldots \times A_{f_n}$ est injectif
et que $A \to B$ est plat. Puisque $B_{f_i} = A_{f_i} \otimes_A B$ on obtient le résultat.
\end{proof}

\begin{lemma}
\label{lemma-flat-base-change-scheme-theoretic-image}
\begin{slogan}
La formation des images schématiques commute aux changements de base plats
dans le cas quasi-compact
\end{slogan}
Soit $f : X \to Y$ un morphisme plat de schémas. Soit $g : V \to Y$ un
morphisme de schémas quasi-compact. Soit $Z \subset Y$ l'image schématique
de $g$ et soit $Z' \subset X$ l'image schématique du changement
de base $V \times_Y X \to X$. Puis $Z' = f^{-1}Z$.
\end{lemma}

\begin{proof}
Rappelons que $Z$ est défini par
$\mathcal{I} = \Ker(\mathcal{O}_Y \to g_*\mathcal{O}_V)$, tandis que $Z'$ est défini par
$\mathcal{I}' = \Ker(\mathcal{O}_X \to
(V \times_Y X \to X)_*\mathcal{O}_{V \times_Y X})$,
voir le Lemme \ref{lemma-quasi-compact-scheme-theoretic-image}. La question est locale
sur $X$ et $Y$, et l'on peut donc supposer
que $X$ et $Y$ sont affines. Nous pouvons remplacer $V$
par $\coprod V_i$ où $V = V_1 \cup \ldots \cup V_n$ est un recouvrement
ouvert affine fini. Nous pouvons donc supposer
que $g$ est affine. Dans ce cas,
$(V \times_Y X \to X)_*\mathcal{O}_{V \times_Y X}$ est le retrait de $g_*\mathcal{O}_V$ par $f$.
Puisque $f$ est plat, nous concluons que
$f^*\mathcal{I} = \mathcal{I}'$ et le lemme est vrai.
\end{proof}




\section{Immersions fermées plates}
\label{section-flat-closed-immersions}

\noindent
Les composants connectés des schémas ne sont pas toujours ouverts. Mais ils ont toujours
une structure schématique canonique. Nous expliquons cela dans cette section.

\begin{lemma}
\label{lemma-characterize-flat-closed-immersions}
Soit $X$ un schéma. La règle qui associe à un sous-schéma fermé de $X$
son sous-ensemble fermé sous-jacent définit une bijection
$$
\left\{
\begin{matrix}
\text{sous-schémas fermés }Z \subset X \\
\text{tels que }Z \to X\text{ soit plat}
\end{matrix}
\right\}
\leftrightarrow
\left\{
\begin{matrix}
\text{parties fermées }Z \subset X \\
\text{stables par générisation}
\end{matrix}
\right\}
$$
Si $Z \subset X$ est un tel sous-schéma fermé, tout morphisme de schémas $g : Y \to X$
tel que $g(Y) \subset Z$ ensemblistement se factorise
(schématiquement) par $Z$.
\end{lemma}

\begin{proof}
Le cas affine de la bijection est
l'Algèbre, Lemme \ref{algebra-lemma-pure-open-closed-specializations}. Pour les schémas généraux $X$
la bijection s'ensuit en recouvrant $X$ par affines et collage.
Détails omis. Pour l'assertion finale, notons que la projection
$Z \times_{X, g} Y \to Y$ est une immersion fermée plate
(Lemme \ref{lemma-base-change-flat}) qui est bijective sur les espaces topologiques
sous-jacents et doit donc être un isomorphisme
par la bijection établie dans la première partie de la démonstration.
\end{proof}

\begin{lemma}
\label{lemma-flat-closed-immersions-finite-presentation}
Une immersion fermée plate de présentation finie est
l'immersion ouverte d'un ouvert et sous-schéma fermé.
\end{lemma}

\begin{proof}
Le cas affine est
l'Algèbre, Lemme \ref{algebra-lemma-finitely-generated-pure-ideal}. En général, le lemme suit
en couvrant $X$ par des affines. Détails
omis.
\end{proof}

\noindent
Notons qu'un composant connexe $T$ d'un schéma $X$ est un
sous-ensemble fermé stable par génération. La définition suivante est donc
logique.

\begin{definition}
\label{definition-scheme-structure-connected-component}
Soit $X$ un schéma. Soit $T \subset X$ un composant connexe. La structure
de schéma canonique {\it sur $T$}
est la structure de schéma unique sur $T$ telle
que l'immersion
fermée $T \to X$ est plate, voir Lemme \ref{lemma-characterize-flat-closed-immersions}.
\end{definition}

\noindent
Il s'avère que nous pouvons déterminer quand chaque module $\mathcal{O}_X$
plat fini est fini localement libre en utilisant le lemme précédent.

\begin{lemma}
\label{lemma-finite-flat-is-finite-locally-free}
Soit $X$ un schéma. Les conditions suivantes sont équivalentes
\begin{enumerate}
\item tout module $\mathcal{O}_X$ plat fini quasi-cohérent est fini
localement libre, et
\item tout sous-ensemble fermé $Z \subset X$ qui est fermé sous générisations est
ouvert.
\end{enumerate}
\end{lemma}

\begin{proof}
Dans le cas affine, il
s'agit de l'algèbre, Lemme \ref{algebra-lemma-finite-flat-module-finitely-presented}. Le cas schéma ne découle
pas directement du cas affine, nous répétons donc simplement
les arguments.

\medskip\noindent
Supposons (1). Considérons une immersion fermée $i : Z \to X$ telle que
$i$ soit plate. Alors $i_*\mathcal{O}_Z$ est quasi-cohérent et plat,
donc fini localement libre par (1). Ainsi $Z = \text{Supp}(i_*\mathcal{O}_Z)$ est également
ouvert et on voit que (2) est vrai. D'où l'implication (1) $\Rightarrow$
(2) découle de la caractérisation des immersions
plates fermées dans le lemme \ref{lemma-characterize-flat-closed-immersions}.

\medskip\noindent
À l’inverse, supposons que $X$
satisfait (2). Soit $\mathcal{F}$ un module $\mathcal{O}_X$
plat fini, quasi-cohérent. Le support $Z = \text{Supp}(\mathcal{F})$ de $\mathcal{F}$
est fermé, voir Modules, Lemme \ref{modules-lemma-support-finite-type-closed}.
En revanche, si $x \leadsto x'$ est une spécialisation,
alors par Algèbre, Lemme \ref{algebra-lemma-finite-flat-local}
le module $\mathcal{F}_{x'}$ est libre sur $\mathcal{O}_{X, x'}$, et
$$
\mathcal{F}_x =
\mathcal{F}_{x'} \otimes_{\mathcal{O}_{X, x'}} \mathcal{O}_{X, x}.
$$
Donc
$x' \in \text{Supp}(\mathcal{F}) \Rightarrow x \in \text{Supp}(\mathcal{F})$, autrement dit, le support est fermé sous
générisation. Comme $X$ satisfait
(2), nous voyons que le support de $\mathcal{F}$
est ouvert et fermé. Les modules $\wedge^i(\mathcal{F})$, $i = 1, 2, 3, \ldots$
sont également des modules $\mathcal{O}_X$
plats
finis et quasi-cohérents, voir
Modules,
section \ref{modules-section-symmetric-exterior}. Notons que
$\text{Supp}(\wedge^{i + 1}(\mathcal{F})) \subset
\text{Supp}(\wedge^i(\mathcal{F}))$. On voit donc qu'il
existe une décomposition
$$
X = U_0 \amalg U_1 \amalg U_2 \amalg \ldots
$$
par sous-ensembles ouverts et fermés telle
que le support de $\wedge^i(\mathcal{F})$ est $U_i \cup U_{i + 1} \cup \ldots$ pour tout $i$.
Soit $x$ un point de $X$
et supposons $x \in U_r$. Notons que
$\wedge^i(\mathcal{F})_x \otimes \kappa(x) =
\wedge^i(\mathcal{F}_x \otimes \kappa(x))$. Par conséquent, $x \in U_r$
implique que $\mathcal{F}_x \otimes \kappa(x)$ est un espace vectoriel de dimension
$r$. Par le lemme de Nakayama, voir
Algèbre, Lemme \ref{algebra-lemma-NAK}, on peut choisir un voisinage ouvert affine $U \subset U_r \subset X$ de $x$
et des sections $s_1, \ldots, s_r \in \mathcal{F}(U)$ tels que l'application
induite
$$
\mathcal{O}_U^{\oplus r} \longrightarrow \mathcal{F}|_U, \quad
(f_1, \ldots, f_r) \longmapsto \sum f_i s_i
$$
soit surjectif. Cela signifie
que $\wedge^r(\mathcal{F}|_U)$ est un module $\mathcal{O}_U$
plat fini et quasi-cohérent dont le
support est entièrement $U$. Par ce
qui précède, il est généré
par un seul élément, à savoir $s_1 \wedge \ldots \wedge s_r$. D'où
$\wedge^r(\mathcal{F}|_U) \cong \mathcal{O}_U/\mathcal{I}$ pour un faisceau quasi-cohérent
d'idéaux $\mathcal{I}$ tel que $\mathcal{O}_U/\mathcal{I}$ est plat
sur $\mathcal{O}_U$ et tel
que $V(\mathcal{I}) = U$. Il s'ensuit que
$\mathcal{I} = 0$ en appliquant le Lemme \ref{lemma-characterize-flat-closed-immersions}.
Ainsi $s_1 \wedge \ldots \wedge s_r$ est une base pour
$\wedge^r(\mathcal{F}|_U)$ et Il s'ensuit que l'application affichée est aussi bien injective
que surjective. Cela prouve que $\mathcal{F}$ est fini localement libre comme
souhaité.
\end{proof}






\section{Platitude générique}
\label{section-generic-flatness}

\noindent
Un schéma de type fini sur une base intègre est plat sur un ouvert dense
ouverte de la base.
En Algèbre, section \ref{algebra-section-generic-flatness} nous avons prouvé une version
noéthérienne, une version pour les morphismes de présentation finie et une
version générale. Nous ne faisons ici qu’énoncer et prouver la version générale.
Cependant, il s'avère que cela sera remplacé
par la proposition \ref{proposition-generic-flatness-reduced} qui montre que le résultat
est valable si nous supposons uniquement que la base est réduite.

\begin{proposition}[Platitude générique]
\label{proposition-generic-flatness}
Soit $f : X \to S$ un morphisme de schémas. Soit
$\mathcal{F}$ un faisceau quasi-cohérent de $\mathcal{O}_X$-modules. Supposons
que
\begin{enumerate}
\item $S$ est intègre,
\item $f$ est de type fini et
\item $\mathcal{F}$ est un module de type fini $\mathcal{O}_X$.
\end{enumerate}
Alors il existe un sous-schéma dense ouvert $U \subset S$ tel que
$X_U \to U$ est plat et de présentation finie et tel
que $\mathcal{F}|_{X_U}$ est plat sur $U$ et de présentation
finie sur $\mathcal{O}_{X_U}$.
\end{proposition}

\begin{proof}
Comme $S$ est intègre, il est irréductible
(voir Propriétés, Lemme \ref{properties-lemma-characterize-integral}) et tout ouvert non vide
est dense. On peut donc remplacer $S$ par un
ouvert affine de $S$ et supposer que $S = \Spec(A)$ est affine.
Comme $S$ est intègre, nous voyons que $A$
est un domaine. Comme $f$ est de type fini, il est quasi-compact,
donc $X$ est quasi-compact. On peut
donc trouver un recouvrement ouvert affine fini $X = \bigcup_{i = 1, \ldots, n} X_i$.
Écrivons $X_i = \Spec(B_i)$. Alors $B_i$ est une algèbre
de type fini $A$, voir Lemme \ref{lemma-locally-finite-type-characterize}. De
plus, il existe des $B_i$-modules
$M_i$ de type fini tels que $\mathcal{F}|_{X_i}$ est le
faisceau quasi-cohérent associé au $B_i$-module $M_i$,
voir Propriétés, Lemme \ref{properties-lemma-finite-type-module}. Ensuite,
pour chaque paire d'indices $i, j$ choisissez un $I_{ij} \subset B_i$
idéal tel que $X_i \setminus X_i \cap X_j = V(I_{ij})$ à l'intérieur de $X_i = \Spec(B_i)$.
Définissez $M_{ij} = B_i/I_{ij}$ et considérez-le comme
un module $B_i$. Alors $V(I_{ij}) = \text{Supp}(M_{ij})$ et $M_{ij}$
sont un module $B_i$ fini.

\medskip\noindent
À ce stade, nous
appliquons l'algèbre, Lemme \ref{algebra-lemma-generic-flatness} aux
paires $(A \to B_i, M_{ij})$ et aux paires $(A \to B_i, M_i)$.
Ainsi on obtient un $f_{ij}, f_i \in A$
non nul tel que (a) $A_{f_{ij}} \to B_{i, f_{ij}}$ est plat et de présentation
finie et $M_{ij, f_{ij}}$ est plat sur $A_{f_{ij}}$ et de
présentation finie sur $B_{i, f_{ij}}$, et (b) $B_{i, f_i}$ est plat
et de présentation finie sur $A_f$ et $M_{i, f_i}$
est plat et de présentation finie sur $B_{i, f_i}$. Posons $f = (\prod f_i) (\prod f_{ij})$.
Nous affirmons que
prendre $U = D(f)$ fonctionne.

\medskip\noindent
Pour prouver notre affirmation, nous pouvons remplacer
$A$ par $A_f$, c'est-à-dire effectuer
le changement de base par $U = \Spec(A_f) \to S$. Après ce changement
de base on voit que chacun des $A \to B_i$ est plat et de présentation
finie et que $M_i$, $M_{ij}$
sont plats sur $A$ et de présentation finie
sur $B_i$. Cela prouve déjà que $X \to S$ est quasi-compact,
localement de présentation finie, plate, et que $\mathcal{F}$
est plat sur $S$ et de présentation finie sur
$\mathcal{O}_X$,
voir Lemme \ref{lemma-locally-finite-presentation-characterize} et Propriétés, Lemme \ref{properties-lemma-finite-presentation-module}.
Puisque $M_{ij}$ est de présentation finie
sur $B_i$ on voit que $X_i \cap X_j = X_i \setminus \text{Supp}(M_{ij})$ est un ouvert
quasi-compact
de $X_i$, voir Algèbre, Lemme \ref{algebra-lemma-support-finite-presentation-constructible}. On voit
donc que $X \to S$ est quasi-séparé
par Schémas, Lemme \ref{schemes-lemma-characterize-quasi-separated}. Cela prouve
la proposition.
\end{proof}

\noindent
Il s'avère effectivement qu'il existe également une version de platitude
générique sur une base réduite arbitraire. C'est ici.

\begin{proposition}[Platitude générique, cas réduit]
\label{proposition-generic-flatness-reduced}
Soit $f : X \to S$ un morphisme de schémas. Soit
$\mathcal{F}$ un faisceau quasi-cohérent de $\mathcal{O}_X$-modules. Supposons
que
\begin{enumerate}
\item $S$ est réduit,
\item $f$ est de type fini et
\item $\mathcal{F}$ est un module de type fini $\mathcal{O}_X$-module.
\end{enumerate}
Il existe alors un sous-schéma dense ouvert $U \subset S$ tel que
$X_U \to U$ est plat et de présentation finie et tel
que $\mathcal{F}|_{X_U}$ est plat sur $U$ et de présentation
finie sur $\mathcal{O}_{X_U}$.
\end{proposition}

\begin{proof}
Pour le lecteur impatient : cette démonstration est une
répétition de la démonstration de la
proposition
\ref{proposition-generic-flatness} utilisant l'algèbre, Lemme \ref{algebra-lemma-generic-flatness-reduced}
au
lieu de l'algèbre, Lemme \ref{algebra-lemma-generic-flatness}.

\medskip\noindent
Puisque le fait d'être plat et d'être de présentation finie est
local
sur la base, voir Lemmes \ref{lemma-flat-module-characterize} et \ref{lemma-locally-finite-presentation-characterize},
on peut travailler localement sur les affines
de $S$. Ainsi on peut supposer que $S = \Spec(A)$,
où $A$ est un anneau réduit (voir Propriétés,
Lemme \ref{properties-lemma-characterize-reduced}). Comme $f$ est de type fini,
il est quasi-compact, donc $X$ est quasi-compact. On peut
donc trouver un recouvrement ouvert affine
fini $X = \bigcup_{i = 1, \ldots, n} X_i$. Écrivons $X_i = \Spec(B_i)$. Alors $B_i$ est une
algèbre de type fini $A$, voir
Laisse-moi \ref{lemma-locally-finite-type-characterize}. De plus il existe des
$B_i$-modules $M_i$
de type fini tels que $\mathcal{F}|_{X_i}$ est le faisceau
quasi-cohérent associé au $B_i$-module $M_i$,
voir Propriétés, Lemme \ref{properties-lemma-finite-type-module}. Ensuite, pour
chaque paire d'indices $i, j$ choisissez un idéal $I_{ij} \subset B_i$
tel que $X_i \setminus X_i \cap X_j = V(I_{ij})$ à l'intérieur de $X_i = \Spec(B_i)$.
Définissez $M_{ij} = B_i/I_{ij}$ et considérez-le comme
un module $B_i$. Alors $V(I_{ij}) = \text{Supp}(M_{ij})$ et $M_{ij}$
sont un module $B_i$ fini.

\medskip\noindent
À ce stade, nous
appliquons l'algèbre, Lemme \ref{algebra-lemma-generic-flatness-reduced} aux paires
$(A \to B_i, M_{ij})$ et aux paires $(A \to B_i, M_i)$. On obtient ainsi
des ouverts denses $U(A \to B_i, M_{ij}) \subset S$
et des ouverts denses $U(A \to B_i, M_i) \subset S$
avec une notation comme en Algèbre, Equation
(\ref{algebra-equation-good-locus}). Puisqu’une intersection
finie d’ouverts denses est ouverte dense, on voit que
$$
U =
\bigcap\nolimits_{i, j} U(A \to B_i, M_{ij})
\quad\cap\quad
\bigcap\nolimits_i U(A \to B_i, M_i)
$$
est ouverte et dense dans $S$. Nous affirmons que $U$ est l'ouverture souhaitée.

\medskip\noindent
Choisissons $u \in U$. Par définition des lieux $U(A \to B_i, M_{ij})$
et $U(A \to B, M_i)$, il existe $f_{ij}, f_i \in A$ tels que (a)
$u \in D(f_i)$ et $u \in D(f_{ij})$, (b) $A_{f_{ij}} \to B_{i, f_{ij}}$
est plat et de présentation finie et $M_{ij, f_{ij}}$ est plat
sur $A_{f_{ij}}$ et de présentation finie sur $B_{i, f_{ij}}$, et (c)
$B_{i, f_i}$ est plat
et de présentation finie sur $A_{f_i}$ et $M_{i, f_i}$ est de
présentation finie sur $B_{i, f_i}$ et plat sur $A_{f_i}$.
Définissez $f = (\prod f_i) (\prod f_{ij})$.
Il suffit maintenant
de prouver que $X \to S$ est plat et de présentation finie
sur $D(f)$ et que $\mathcal{F}$ restreint à $X_{D(f)}$ est plat
sur $D(f)$ et de présentation finie sur le faisceau structural de
$X_{D(f)}$.

\medskip\noindent
On peut donc remplacer $A$ par
$A_f$, c'est-à-dire effectuer le changement
de base par $\Spec(A_f) \to S$. Après ce changement de
base on voit que chacun des $A \to B_i$ est plat et de présentation finie
et que $M_i$, $M_{ij}$ sont
plats sur $A$ et de présentation finie sur
$B_i$. Cela prouve déjà que $X \to S$ est quasi-compact,
localement de présentation finie, plate, et que $\mathcal{F}$
est plat sur $S$ et de présentation finie sur
$\mathcal{O}_X$,
voir Lemme \ref{lemma-locally-finite-presentation-characterize} et Propriétés, Lemme \ref{properties-lemma-finite-presentation-module}.
Puisque $M_{ij}$ est de présentation finie
sur $B_i$ on voit que $X_i \cap X_j = X_i \setminus \text{Supp}(M_{ij})$ est un ouvert
quasi-compact
de $X_i$, voir Algèbre, Lemme \ref{algebra-lemma-support-finite-presentation-constructible}. On voit
donc que $X \to S$ est quasi-séparé
par Schémas, Lemme \ref{schemes-lemma-characterize-quasi-separated}. Cela prouve
la proposition.
\end{proof}









\begin{remark}
\label{remark-flattening}
Les résultats ci-dessus sont un premier pas vers des techniques d'aplatissement plus raffinées
des morphismes de schémas. L'article \cite{GruRay} de Raynaud et Gruson contient
de nombreux résultats merveilleux dans ce sens.
\end{remark}







\section{Morphismes et dimensions des fibres}
\label{section-dimension-fibres}

\noindent
Soit $X$ un espace topologique, et $x \in X$.
Rappelons que nous avons défini $\dim_x(X)$ comme le minimum
des dimensions des voisinages ouverts de $x$
dans $X$. Voir Topologie, Définition \ref{topology-definition-Krull}.

\begin{lemma}
\label{lemma-dimension-fibre-at-a-point}
Soit $f : X \to S$ un morphisme de schémas. Soit $x \in X$
et définissez $s = f(x)$. Supposons
que $f$ soit localement de type fini.
Puis
$$
\dim_x(X_s) =
\dim(\mathcal{O}_{X_s, x}) + \text{trdeg}_{\kappa(s)}(\kappa(x)).
$$
\end{lemma}

\begin{proof}
Cela se réduit immédiatement au cas $S = s$, et
$X$ affine. Dans ce
cas, le résultat découle de l'algèbre, Lemme \ref{algebra-lemma-dimension-at-a-point-finite-type-field}.
\end{proof}

\begin{lemma}
\label{lemma-dimension-fibre-at-a-point-additive}
Soient $f : X \to Y$ et $g : Y \to S$ des morphismes de schémas. Soit $x \in X$
et définissez $y = f(x)$, $s = g(y)$. Supposons que $f$
et $g$ localement de type fini.
Alors
$$
\dim_x(X_s) \leq \dim_x(X_y) + \dim_y(Y_s).
$$
De plus, l'égalité est vraie si $\mathcal{O}_{X_s, x}$ est
plat sur $\mathcal{O}_{Y_s, y}$, ce qui est vrai par exemple si $\mathcal{O}_{X, x}$
est plat sur $\mathcal{O}_{Y, y}$.
\end{lemma}

\begin{proof}
Notons que $\text{trdeg}_{\kappa(s)}(\kappa(x)) =
\text{trdeg}_{\kappa(y)}(\kappa(x)) + \text{trdeg}_{\kappa(s)}(\kappa(y))$. Ainsi, par le Lemme \ref{lemma-dimension-fibre-at-a-point},
l'instruction est équivalente
à
$$
\dim(\mathcal{O}_{X_s, x})
\leq
\dim(\mathcal{O}_{X_y, x}) + \dim(\mathcal{O}_{Y_s, y}).
$$
Pour cela voir Algèbre, Lemme \ref{algebra-lemma-dimension-base-fibre-total}.
Pour le cas
plat, voir Algèbre, Lemme \ref{algebra-lemma-dimension-base-fibre-equals-total}.
\end{proof}

\begin{lemma}
\label{lemma-dimension-fibre-after-base-change}
Soit
$$
\xymatrix{
X' \ar[r]_{g'} \ar[d]_{f'} & X \ar[d]^f \\
S' \ar[r]^g & S
}
$$
un diagramme produit fibré de schémas. Supposons que $f$ localement de
type fini. Supposons que $x' \in X'$, $x = g'(x')$, $s' = f'(x')$
et $s = g(s') = f(x)$. Alors
\begin{enumerate}
\item $\dim_x(X_s) = \dim_{x'}(X'_{s'})$,
\item si $F$ est la fibre du morphisme $X'_{s'} \to X_s$ sur $x$,
alors
$$
\dim(\mathcal{O}_{F, x'}) =
\dim(\mathcal{O}_{X'_{s'}, x'}) - \dim(\mathcal{O}_{X_s, x}) =
\text{trdeg}_{\kappa(s)}(\kappa(x)) -
\text{trdeg}_{\kappa(s')}(\kappa(x'))
$$
notamment $\dim(\mathcal{O}_{X'_{s'}, x'}) \geq \dim(\mathcal{O}_{X_s, x})$
et $\text{trdeg}_{\kappa(s)}(\kappa(x)) \geq
\text{trdeg}_{\kappa(s')}(\kappa(x'))$.
\item étant donné $s', s, x$,
il existe un choix de $x'$
tel que $\dim(\mathcal{O}_{X'_{s'}, x'}) = \dim(\mathcal{O}_{X_s, x})$ et $\text{trdeg}_{\kappa(s)}(\kappa(x)) = \text{trdeg}_{\kappa(s')}(\kappa(x'))$.
\end{enumerate}
\end{lemma}

\begin{proof}
La partie (1) découle
immédiatement
de l'algèbre, Lemme \ref{algebra-lemma-dimension-at-a-point-preserved-field-extension}.
Parties (2)
et (3) d'Algèbre, Lemme \ref{algebra-lemma-inequalities-under-field-extension}.
\end{proof}

\noindent
Le lemme suivant découle d’un résultat algébrique non trivial. À savoir
la version algébrique du Théorème principal de Zariski.

\begin{lemma}
\label{lemma-openness-bounded-dimension-fibres}
\begin{reference}
\cite[IV Théorème 13.1.3]{EGA}
\end{reference}
Soit $f : X \to S$ un morphisme de schémas. Soit $n \geq 0$.
Supposons que $f$ soit localement de type fini.
L'ensemble
$$
U_n = \{x \in X \mid \dim_x X_{f(x)} \leq n\}
$$
est ouvert dans $X$.
\end{lemma}

\begin{proof}
Ceci est immédiat
de
l'Algèbre, Lemme \ref{algebra-lemma-dimension-fibres-bounded-open-upstairs}
\end{proof}

\begin{lemma}
\label{lemma-morphism-finite-type-bounded-dimension}
Soit $f : X \to Y$ un morphisme de type fini avec $Y$ quasi-compact. Alors
la dimension des fibres de $f$ est bornée.
\end{lemma}

\begin{proof}
Par le Lemme \ref{lemma-openness-bounded-dimension-fibres} l'ensemble $U_n \subset X$ des points où
la dimension de la fibre est $\leq n$ est ouvert. Puisque $f$
est de type fini, chaque point est contenu dans un $U_n$ (car
la dimension d'une algèbre de type fini sur un corps est finie).
Puisque $Y$ est quasi-compact et $f$ est de type fini, on voit
que $X$ est quasi-compact. D'où $X = U_n$ pour certains
$n$.
\end{proof}

\begin{lemma}
\label{lemma-openness-bounded-dimension-fibres-finite-presentation}
Soit $f : X \to S$ un morphisme de schémas. Soit $n \geq 0$.
Supposons que $f$ soit localement de présentation finie.
Le
$$
U_n = \{x \in X \mid \dim_x X_{f(x)} \leq n\}
$$
ouvert du lemme \ref{lemma-openness-bounded-dimension-fibres} est rétrocompact dans $X$.
(Voir Topologie, Définition \ref{topology-definition-quasi-compact}.)
\end{lemma}

\begin{proof}
L'espace topologique $X$ a une base pour sa topologie
constituée d'ouverts affines $U \subset X$
telles que le morphisme induit $f|_U : U \to S$ se factorise
par un ouvert affine $V \subset S$. Il suffit donc de montrer
que $U \cap U_n$ est quasi-compact pour
un tel $U$. Notez que $U_n \cap U$ est identique
au $\{x \in U \mid \dim_x U_{f(x)} \leq n\}$ ouvert. Cela nous ramène au cas où $X$
et
$S$ sont affines. Dans ce cas, le lemme découle de
l'algèbre, Lemme \ref{algebra-lemma-dimension-fibres-bounded-quasi-compact-open-upstairs} (et Lemme \ref{lemma-locally-finite-presentation-characterize}).
\end{proof}

\begin{lemma}
\label{lemma-dimension-fibre-specialization}
Soit $f : X \to S$ un morphisme de schémas. Soit $x \leadsto x'$
une spécialisation non triviale de points dans $X$ situés sur
le même point $s \in S$. Supposons que $f$ soit localement de type fini.
Puis
\begin{enumerate}
\item $\dim_x(X_s) \leq \dim_{x'}(X_s)$,
\item $\dim(\mathcal{O}_{X_s, x}) < \dim(\mathcal{O}_{X_s, x'})$ et
\item $\text{trdeg}_{\kappa(s)}(\kappa(x)) >
\text{trdeg}_{\kappa(s)}(\kappa(x'))$.
\end{enumerate}
\end{lemma}

\begin{proof}
La partie (1) découle du fait que toute ouverture de $X_s$ contenant
$x'$ contient également $x$. La partie (2) suit puisque
$\mathcal{O}_{X_s, x}$ est une localisation de $\mathcal{O}_{X_s, x'}$ à un idéal premier,
donc toute chaîne d'idéaux premiers dans $\mathcal{O}_{X_s, x}$ fait partie d'une
chaîne strictement plus longue de nombres premiers dans $\mathcal{O}_{X_s, x'}$. La dernière
inégalité découle de l'algèbre, Lemme \ref{algebra-lemma-tr-deg-specialization}.
\end{proof}






\section{Morphismes de dimension relative donnée}
\label{section-relative-dimension}

\noindent
Afin de pouvoir parler confortablement de morphismes d'une dimension
relative donnée nous introduisons la notion suivante.

\begin{definition}
\label{definition-relative-dimension-d}
Soit $f : X \to S$ un morphisme de schémas. Supposons que
$f$ soit localement de type fini.
\begin{enumerate}
\item Nous disons que $f$ est de dimension {\it relative $\leq d$ à
$x$} si $\dim_x(X_{f(x)}) \leq d$.
\item Nous disons que $f$ est de dimension {\it relative
$\leq d$} si $\dim_x(X_{f(x)}) \leq d$ pour tout $x \in X$.
\item On dit que $f$ est de dimension {\it relative $d$}
si toutes les fibres non vides $X_s$ sont équidimensionnelles de dimension $d$.
\end{enumerate}
\end{definition}

\noindent
Ce n'est pas une notion particulièrement efficace, mais elle fonctionne bien dans un
certain nombre de situations.

\begin{lemma}
\label{lemma-base-change-relative-dimension-d}
Soit $f : X \to S$ un morphisme de schémas qui est localement de type fini. Si $f$
a une dimension relative $d$, tout changement de base de $f$ aussi.
Idem pour la cote relative $\leq d$.
\end{lemma}

\begin{proof}
Ceci est immédiat
à partir du lemme \ref{lemma-dimension-fibre-after-base-change}.
\end{proof}

\begin{lemma}
\label{lemma-composition-relative-dimension-d}
Soit $f : X \to Y$, $g : Y \to Z$ être localement de type fini. Si $f$
a une dimension relative $\leq d$ et $g$ a une dimension relative $\leq e$
alors $g \circ f$ a une dimension relative $\leq d + e$.
Si
\begin{enumerate}
\item $f$ a une dimension relative $d$,
\item $g$ a une dimension relative $e$ et que
\item $f$ est plat,
\end{enumerate}
alors $g \circ f$ a une dimension relative $d + e$.
\end{lemma}

\begin{proof}
Ceci est immédiat à partir du lemme \ref{lemma-dimension-fibre-at-a-point-additive}.
\end{proof}

\noindent
En général, il n'est pas possible de décomposer un morphisme
en morceaux où la dimension relative est donnée. Cependant,
c'est possible si le morphisme comporte des fibres de De Cohen--Macaulay
et est plat de présentation finie.

\begin{lemma}
\label{lemma-flat-finite-presentation-CM-fibres-relative-dimension}
\begin{slogan}
Les morphismes de Cohen--Macaulay se décomposent en ouverts et fermés de dimension relative pure
\end{slogan}
Soit $f : X \to S$ un morphisme de schémas. Supposons
que
\begin{enumerate}
\item $f$ est plat,
\item $f$ est localement de présentation finie, et
\item pour tout $s \in S$ la fibre $X_s$ est
De Cohen--Macaulay (Propriétés, Définition \ref{properties-definition-Cohen-Macaulay})
\end{enumerate}
Alors il existe des ouverts et sous-schémas fermés $X_d \subset X$
tels que $X = \coprod_{d \geq 0} X_d$ et $f|_{X_d} : X_d \to S$ ont une dimension
relative $d$.
\end{lemma}

\begin{proof}
Ceci est immédiat
de l'algèbre,
Lemme \ref{algebra-lemma-relative-dimension-CM}.
\end{proof}

\begin{lemma}
\label{lemma-locally-quasi-finite-rel-dimension-0}
Soit $f : X \to S$ un morphisme de schémas. Supposons que
$f$ soit localement de type fini. Soit
$x \in X$ avec $s = f(x)$. Alors $f$
est quasi-fini à $x$ si et seulement si $\dim_x(X_s) = 0$. En particulier,
$f$ est localement quasi-fini si et seulement si $f$ a une dimension relative
$0$.
\end{lemma}

\begin{proof}
Première démonstration.
Si $f$ est quasi-fini en $x$ alors
$\kappa(x)$ est
une extension finie de $\kappa(s)$ (par
le lemme \ref{lemma-residue-field-quasi-finite}) et $x$
est isolé dans $X_s$ (par le lemme
\ref{lemma-quasi-finite-at-point-characterize}), donc $\dim_x(X_s) = 0$
par le lemme \ref{lemma-dimension-fibre-at-a-point}. A l'inverse,
si $\dim_x(X_s) = 0$ alors par le
lemme \ref{lemma-dimension-fibre-at-a-point} on voit que $\kappa(s) \subset \kappa(x)$
est algébrique et il n'y a pas d'autres points
de $X_s$ spécialisés dans $x$.
Ainsi $x$ est fermé
dans sa fibre par le lemme \ref{lemma-algebraic-residue-field-extension-closed-point-fibre} et par le lemme
\ref{lemma-quasi-finite-at-point-characterize}
(3) Nous concluons que $f$
est quasi-fini en $x$.

\medskip\noindent
Deuxième démonstration. La fibre $X_s$ est un schéma
localement de type fini sur un corps, donc localement
noethérien (Lemme
\ref{lemma-finite-type-noetherian}). Le résultat découle maintenant
du lemme \ref{lemma-quasi-finite-at-point-characterize} et de Propriétés, Lemme \ref{properties-lemma-isolated-dim-0-locally-noetherian}.
\end{proof}

\begin{lemma}
\label{lemma-rel-dimension-dimension}
Soit $f : X \to Y$ un morphisme de schémas localement noethériens qui est plat,
localement de type fini et de dimension relative $d$. Pour chaque
point $x$ dans $X$ avec l'image
$y$ dans $Y$ nous avons $\dim_x(X) = \dim_y(Y) + d$.
\end{lemma}

\begin{proof}
Après avoir réduit $X$ et $Y$ pour ouvrir les
voisinages de $x$ et $y$, on peut supposer que
$\dim(X) = \dim_x(X)$ et $\dim(Y) = \dim_y(Y)$, par définition de la dimension
d'un schéma en un point (Propriétés, Définition
\ref{properties-definition-dimension}). Le morphisme $f$
est ouvert par les Lemmes \ref{lemma-noetherian-finite-type-finite-presentation} et \ref{lemma-fppf-open}.
On peut donc réduire
$Y$ pour faire en sorte que $f$ soit
surjectif. Reste à montrer que $\dim(X) = \dim(Y) + d$.

\medskip\noindent
Soit $a$ un point dans $X$ avec l'image
$b$ dans $Y$. Par Algèbre, Lemme \ref{algebra-lemma-dimension-base-fibre-equals-total},
$$
\dim(\mathcal{O}_{X,a}) = \dim(\mathcal{O}_{Y,b}) + \dim(\mathcal{O}_{X_b, a}).
$$
Prenant la suprématie sur tous les points $a$ dans
$X$, Il s'ensuit que $\dim(X) = \dim(Y) + d$, comme
on veut, voir Propriétés, Lemme \ref{properties-lemma-dimension}.
\end{proof}











\section{Morphismes syntomiques}
\label{section-syntomic}

\noindent
Une algèbre $A$ sur un corps $k$ est appelée
une {\it intersection globale
complète sur $k$} si $A \cong k[x_1, \ldots, x_n]/(f_1, \ldots, f_c)$ et $\dim(A) = n - c$.
Une algèbre $A$ sur un corps $k$ est appelée une {\it
intersection complète locale} si $\Spec(A)$
peut être couverte par des ouvertures standard dont chacune est
une intersection complète globale sur $k$. Voir Algèbre,
section \ref{algebra-section-lci}. Rappelons qu'un homomorphisme d'anneaux $R \to A$
est {\it syntomique} s'il est
de présentation finie, plat avec des anneaux d'intersection complets
locaux comme fibres, voir Algèbre, Définition \ref{algebra-definition-lci}.

\begin{definition}
\label{definition-syntomic}
Soit $f : X \to S$ un morphisme de schémas.
\begin{enumerate}
\item On dit que $f$ est {\it syntomique
à $x \in X$} s'il existe un voisinage ouvert affine $\Spec(A) = U \subset X$
de $x$ et ouvert affine $\Spec(R) = V \subset S$ avec
$f(U) \subset V$ tel que l'homomorphisme d'anneaux induit $R \to A$
soit syntomique.
\item On dit que $f$ est {\it syntomique} s'il est syntomique
en tout point de $X$.
\item Si $S = \Spec(k)$ et $f$ sont syntomiques, alors on dit que $X$ est une intersection
locale complète {\it sur $k$}.
\item Un morphisme de schémas affines $f : X \to S$
est appelé {\it syntomique
standard} s'il
existe une intersection complète relative
globale
$R \to R[x_1, \ldots, x_n]/(f_1, \ldots, f_c)$ (voir Algèbre, Définition \ref{algebra-definition-relative-global-complete-intersection}) telle que
$X \to S$ est isomorphe à
$$
\Spec(R[x_1, \ldots, x_n]/(f_1, \ldots, f_c)) \to \Spec(R).
$$
\end{enumerate}
\end{definition}

\noindent
Dans la littérature, un morphisme syntomique est parfois appelé {\it
morphisme d'intersection complet local plat}. Il
s’avère qu’il s’agit d’une classe de morphismes pratique. Par exemple, on peut définir
une topologie syntomique en utilisant celles-ci, qui est plus fine que les topologies
lisses et étales, mais possède bon nombre des mêmes propriétés formelles.

\medskip\noindent Une
intersection complète relative globale (que nous avons utilisée pour définir les homomorphismes
syntomiques d'anneaux standards) est particulièrement
plate. Dans En savoir plus sur les morphismes, section \ref{more-morphisms-section-lci},
nous considérerons les morphismes $X \to S$ qui sont localement de la forme
$$
\Spec(R[x_1, \ldots, x_n]/(f_1, \ldots, f_c)) \to \Spec(R).
$$
pour certaines séquences régulières de Koszul $f_1, \ldots, f_r$ dans $R[x_1, \ldots, x_n]$. Un tel morphisme
sera appelé un {\it morphisme d'intersection complet local}.
Une fois que nous aurons cette définition en place, un morphisme sera syntomique
si et seulement si c'est un morphisme d'intersection complet, plat et local.

\medskip\noindent Notons
qu'il n'y a pas d'hypothèse de séparation ou de quasi-compacité dans la
définition d'un morphisme syntomique. La question du caractère syntomique est
donc de nature locale sur la source. Voici le résultat précis.

\begin{lemma}
\label{lemma-syntomic-characterize}
Soit $f : X \to S$ un morphisme de schémas. Les conditions
suivantes sont équivalentes
\begin{enumerate}
\item Le morphisme $f$ est syntomique.
\item Pour toute ouverture affine $U \subset X$, $V \subset S$
avec $f(U) \subset V$ l'homomorphisme
d'anneaux $\mathcal{O}_S(V) \to \mathcal{O}_X(U)$ est syntomique.
\item Il existe un recouvrement ouvert $S = \bigcup_{j \in J} V_j$ et
des recouvrements ouverts $f^{-1}(V_j) = \bigcup_{i \in I_j} U_i$ tels que chacun
des morphismes $U_i \to V_j$, $j\in J, i\in I_j$ est
syntomique.
\item Il existe un recouvrement ouvert affine $S = \bigcup_{j \in J} V_j$
et ouvert affine revêtements $f^{-1}(V_j) = \bigcup_{i \in I_j} U_i$ tel que l'homomorphisme
d'anneaux $\mathcal{O}_S(V_j) \to \mathcal{O}_X(U_i)$ soit syntomique,
pour tout $j\in J, i\in I_j$.
\end{enumerate}
De plus, si $f$ est syntomique
alors pour tout sous-schémas ouverts $U \subset X$, $V \subset S$ avec $f(U) \subset V$
la restriction $f|_U : U \to V$ est syntomique.
\end{lemma}

\begin{proof}
Cela résulte du lemme \ref{lemma-locally-P} si l'on montre que la propriété
``$R \to A$ est syntomique''
est locale. On vérifie les conditions (a), (b)
et (c) de la Définition \ref{definition-property-local}.
Par Algèbre, Lemme \ref{algebra-lemma-base-change-syntomic}, la propriété d'être syntomique est stable
par changement de base ; on en déduit
la condition (a). Par
Algèbre, Lemme \ref{algebra-lemma-composition-syntomic}, elle est stable
par composition, et pour tout anneau $R$ l'homomorphisme
d'anneaux $R \to R_f$ est syntomique.
On en déduit la condition (b). Enfin, la
condition (c) résulte d'Algèbre, Lemme \ref{algebra-lemma-local-syntomic}.
\end{proof}

\begin{lemma}
\label{lemma-composition-syntomic}
La composition de deux morphismes syntomiques est syntomique.
\end{lemma}

\begin{proof}
Dans la démonstration du lemme \ref{lemma-syntomic-characterize} nous avons
vu qu'être syntomique est une propriété locale des homomorphismes
d'anneaux. D'où la première affirmation
du lemme découle du lemme \ref{lemma-composition-type-P} combinée
au fait qu'être syntomique est une propriété des homomorphismes
d'anneaux qui est stable
sous composition, voir Algèbre, Lemme \ref{algebra-lemma-composition-syntomic}.
\end{proof}

\begin{lemma}
\label{lemma-base-change-syntomic}
Le changement de base d'un morphisme qui est syntomique est syntomique.
\end{lemma}

\begin{proof}
Dans la démonstration du lemme \ref{lemma-syntomic-characterize} nous avons
vu qu'être syntomique est une propriété locale des homomorphismes
d'anneaux. D'où
le lemme découle du lemme \ref{lemma-composition-type-P} combiné
au fait qu'être syntomique est une propriété des homomorphismes
d'anneaux stable par
changement de base, voir Algèbre, Lemme \ref{algebra-lemma-base-change-syntomic}.
\end{proof}

\begin{lemma}
\label{lemma-open-immersion-syntomic}
Toute immersion ouverte est syntomique.
\end{lemma}

\begin{proof}
Ceci est vrai car une immersion ouverte est un isomorphisme local.
\end{proof}

\begin{lemma}
\label{lemma-syntomic-locally-finite-presentation}
Un morphisme syntomique est localement de présentation finie.
\end{lemma}

\begin{proof}
Vrai car un homomorphisme syntomique d'anneaux est de présentation finie par
définition.
\end{proof}

\begin{lemma}
\label{lemma-syntomic-flat}
Un morphisme syntomique est plat.
\end{lemma}

\begin{proof}
Vrai car un homomorphisme syntomique d'anneaux est plat par définition.
\end{proof}

\begin{lemma}
\label{lemma-syntomic-open}
Un morphisme syntomique est universellement ouvert.
\end{lemma}

\begin{proof}
Combine
Lemmes \ref{lemma-syntomic-locally-finite-presentation},
\ref{lemma-syntomic-flat} et
\ref{lemma-fppf-open}.
\end{proof}

\noindent
Soit $k$ un corps. Soit $A$ une algèbre $k$ locale
essentiellement de type fini sur $k$. Rappelons que
$A$ est appelé une {\it intersection
complète sur $k$} si l'on peut écrire $A \cong R/(f_1, \ldots, f_c)$ où $R$
est un anneau local régulier essentiellement de type fini sur $k$,
et $f_1, \ldots, f_c$ est un séquence
régulière dans $R$, voir Algèbre, Définition \ref{algebra-definition-lci-local-ring}.

\begin{lemma}
\label{lemma-local-complete-intersection}
Soit $k$ un corps.
Soit $X$ un schéma local de type fini sur $k$. Les conditions
suivantes sont équivalentes :
\begin{enumerate}
\item $X$ est une intersection complète locale sur $k$,
\item pour tout $x \in X$ il existe un voisinage
ouvert affine $U = \Spec(R) \subset X$ de $x$
tel que $R \cong k[x_1, \ldots, x_n]/(f_1, \ldots, f_c)$ est une intersection
complète globale sur $k$, et
\item pour chaque $x \in X$ l'anneau local $\mathcal{O}_{X, x}$ est une intersection
complète sur $k$.
\end{enumerate}
\end{lemma}

\begin{proof}
Les résultats d'algèbre correspondants peuvent
être trouvés dans Algèbre, Lemmes \ref{algebra-lemma-lci-at-prime}
et \ref{algebra-lemma-lci-global}.
\end{proof}

\noindent
Le lemme suivant dit localement que tout morphisme syntomique est standard syntomique. Nous pouvons
donc utiliser les morphismes syntomiques standards comme {\it modèle
local} pour un morphisme syntomique. De plus, il dit qu'un morphisme plat de présentation
finie est syntomique si et seulement si les fibres sont des schémas d'intersection
complets locaux.

\begin{lemma}
\label{lemma-syntomic-locally-standard-syntomic}
Soit $f : X  \to S$ un morphisme de schémas. Soit $x \in X$ un point d'image
$s = f(x)$. Soit $V \subset S$ un voisinage ouvert affine de $s$. Les conditions
suivantes sont équivalentes
\begin{enumerate}
\item Le morphisme $f$ est syntomique en $x$.
\item Il existe un $U \subset X$ ouvert affine avec $x \in U$ et
$f(U) \subset V$ tel que $f|_U : U \to V$ soit syntomique standard.
\item Le morphisme $f$ est de présentation finie
à $x$, l'homomorphisme d'anneaux locaux
$\mathcal{O}_{S, s} \to \mathcal{O}_{X, x}$ est plat et $\mathcal{O}_{X, x}/\mathfrak m_s \mathcal{O}_{X, x}$ est
une intersection complète sur $\kappa(s)$
(voir Algèbre, Définition \ref{algebra-definition-lci-local-ring}).
\end{enumerate}
\end{lemma}

\begin{proof}
Dérivé des définitions et
de l'algèbre, Lemme \ref{algebra-lemma-syntomic}.
\end{proof}

\begin{lemma}
\label{lemma-syntomic-flat-fibres}
Soit $f : X \to S$ un morphisme de schémas. Si $f$
est plat, localement de présentation finie, et que toutes les fibres $X_s$
sont des intersections locales complètes, alors $f$ est
syntomique.
\end{lemma}

\begin{proof}
Cela résulte des Lemmes
\ref{lemma-local-complete-intersection} et
\ref{lemma-syntomic-locally-standard-syntomic}, ainsi que des
isomorphismes d'anneaux locaux
$
\mathcal{O}_{X, x}/\mathfrak m_s \mathcal{O}_{X, x}
\cong
\mathcal{O}_{X_s, x}
$.
\end{proof}

\begin{lemma}
\label{lemma-set-points-where-fibres-lci}
Soit $f : X \to S$ un morphisme de schémas. Supposons
que $f$ soit localement de type fini. La formation de l'ensemble
$$
T = \{x \in X \mid \mathcal{O}_{X_{f(x)}, x}
\text{ est d'intersection complète sur }\kappa(f(x))\}
$$
commute à tout changement de base :
pour tout morphisme $g : S' \to S$, considérons
le changement de base $f' : X' \to S'$ de $f$ et la
projection $g' : X' \to X$. Alors l'ensemble correspondant $T'$
pour le morphisme $f'$ est égal à $T' = (g')^{-1}(T)$. En particulier,
si $f$ est supposé plat, et localement de présentation finie
alors il en va de même pour l'ensemble ouvert de points où $f$
est syntomique.
\end{lemma}

\begin{proof}
Soit $s' \in S'$ un point et soit $s = g(s')$. On a alors
$$
X'_{s'} =
\Spec(\kappa(s')) \times_{\Spec(\kappa(s))} X_s
$$
Autrement dit les fibres du changement de base sont les changements
de base des fibres. La première partie
est donc équivalente à l'Algèbre, Lemme \ref{algebra-lemma-lci-field-change-local}.
La deuxième partie découle de la première car dans
ce cas $T$ est l'ensemble des points où $f$
est syntomique d'après le Lemme \ref{lemma-syntomic-locally-standard-syntomic}.
\end{proof}

\begin{lemma}
\label{lemma-standard-syntomic-relative-dimension}
Soit $R$ un
anneau. Soit $R \to A = R[x_1, \ldots, x_n]/(f_1, \ldots, f_c)$ une intersection complète relative
globale. Définissez $S = \Spec(R)$ et $X = \Spec(A)$.
Considérons le morphisme $f : X \to S$
associé à l'homomorphisme d'anneaux $R \to A$.
La fonction $x \mapsto \dim_x(X_{f(x)})$ est constante de valeur $n - c$.
\end{lemma}

\begin{proof}
En algèbre,
la Définition \ref{algebra-definition-relative-global-complete-intersection} $R \to A$ étant une intersection complète
globale relative signifie que tous les anneaux de
fibres non nuls ont une dimension $n - c$.
Ainsi, pour un $\mathfrak p$ premier de $R$
l'anneau de fibres $\kappa(\mathfrak p)[x_1, \ldots, x_n]/(\overline{f}_1, \ldots, \overline{f}_c)$ est soit nul, soit un anneau d'intersection
complet global de dimension $n - c$. D'après la discussion
qui suit l'Algèbre,
Définition \ref{algebra-definition-lci-field} cela implique qu'elle
est équidimensionnelle de dimension $n - c$. D'où
le lemme.
\end{proof}

\begin{lemma}
\label{lemma-syntomic-relative-dimension}
Soit $f : X \to S$ un morphisme syntomique. La fonction
$x \mapsto \dim_x(X_{f(x)})$ est localement constante sur $X$.
\end{lemma}

\begin{proof}
Par le Lemme \ref{lemma-syntomic-locally-standard-syntomic} le morphisme $f$
ressemble localement à un morphisme
syntomique d'affines standard. Le
résultat découle donc du lemme \ref{lemma-standard-syntomic-relative-dimension}.
\end{proof}

\noindent
Lemme \ref{lemma-syntomic-relative-dimension} dit que la définition
suivante a du sens.

\begin{definition}
\label{definition-syntomic-relative-dimension}
Soit $d \geq 0$ un nombre entier. On dit qu'un morphisme de schémas $f : X \to S$ est {\it
syntomique de dimension relative $d$} si $f$
est syntomique et la fonction $\dim_x(X_{f(x)}) = d$ pour tout $x \in X$.
\end{definition}

\noindent
Autrement dit, $f$ est syntomique et les fibres non vides sont équidimensionnelles de dimension
$d$.

\begin{lemma}
\label{lemma-syntomic-permanence}
Soit
$$
\xymatrix{
X \ar[rr]_f \ar[rd]_p & &
Y \ar[dl]^q \\
& S
}
$$
un diagramme commutatif de morphismes de schémas. Supposons que
\begin{enumerate}
\item $f$ est surjectif et syntomique,
\item $p$ est syntomique, et
\item $q$ est localement de présentation finie\footnote{En fait,
si $f$ est surjectif, plat et localement de présentation finie et $p$
est syntomique, alors $q$ et $f$
sont syntomiques, voir Descente, Lemme \ref{descent-lemma-syntomic-permanence}.}.
\end{enumerate}
Alors $q$ est syntomique.
\end{lemma}

\begin{proof}
Par le Lemme \ref{lemma-flat-permanence} on voit que $q$ est plat. Il suffit
donc de montrer que les fibres de $Y \to S$ sont
des intersections locales complètes, voir Lemme \ref{lemma-syntomic-flat-fibres}. Soit $s \in S$.
Considérons le morphisme $X_s \to Y_s$. Il
s'agit d'un changement de base du morphisme $X \to Y$
et donc surjectif, et syntomique (Lemme \ref{lemma-base-change-syntomic}). Pour la
même raison, $X_s$ est syntomique par rapport à
$\kappa(s)$. De plus, $Y_s$ est localement de type
fini sur $\kappa(s)$ (Lemme \ref{lemma-base-change-finite-type}). On se réduit
ainsi au cas où $S$ est le spectre d'un corps.

\medskip\noindent
Supposons que $S = \Spec(k)$. Soit $y \in Y$.
Choisissez un voisinage $\Spec(A) \subset Y$ ouvert affine
de $y$. Soit $\Spec(B) \subset X$ un ouvert
affine telle que $f(\Spec(B)) \subset \Spec(A)$, contenant
un point $x \in X$ tel que $f(x) = y$. Choisissez
une surjection $k[x_1, \ldots, x_n] \to A$ avec
noyau $I$. Choisissez une surjection $A[y_1, \ldots, y_m] \to B$,
qui donne lieu à son tour à une surjection $k[x_i, y_j] \to B$
à noyau $J$. Soit $\mathfrak q \subset k[x_i, y_j]$ le premier
correspondant à $y \in \Spec(B)$ et soit $\mathfrak p \subset k[x_i]$ le premier
correspondant à $x \in \Spec(A)$.
Puisque $x$ s'envoie sur $y$, nous avons $\mathfrak p = \mathfrak q \cap k[x_i]$. Considérons
le diagramme commutatif d'anneaux locaux suivant :
$$
\xymatrix{
\mathcal{O}_{X, x} \ar@{=}[r] &
B_{\mathfrak q} &
k[x_1, \ldots, x_n, y_1, \ldots, y_m]_{\mathfrak q} \ar[l] \\
\mathcal{O}_{Y, y} \ar@{=}[r] \ar[u] & A_{\mathfrak p} \ar[u] &
k[x_1, \ldots, x_n]_{\mathfrak p} \ar[l] \ar[u]
}
$$
Nous affirmons que les hypothèses
d'algèbre du lemme \ref{algebra-lemma-lci-permanence-initial} sont satisfaites. Les
conditions (1) et (2) sont triviales. La condition
(4) s'ensuit car $X \to Y$ est plat. La
condition (3) s'ensuit
car les anneaux $\mathcal{O}_{Y, y}$ et $\mathcal{O}_{X_y, x} = \mathcal{O}_{X, x}/\mathfrak m_y\mathcal{O}_{X, x}$ sont des anneaux
d'intersection complets d'après nos hypothèses
selon lesquelles $f$
et $p$ sont syntomiques, voir
Lemme \ref{lemma-syntomic-locally-standard-syntomic}. Le résultat d'Algèbre, Lemme \ref{algebra-lemma-lci-permanence-initial}
est exactement que $\mathcal{O}_{Y, y}$ est un anneau d'intersection
complet ! Par conséquent, en Lemme \ref{lemma-syntomic-locally-standard-syntomic}
, nous voyons à nouveau que $Y$ est syntomique sur $k$ à $y$ comme souhaité.
\end{proof}









\section{Faisceau conormal d'une immersion}
\label{section-conormal-sheaf}

\noindent
Soit $i : Z \to X$ une immersion fermée.
Soit $\mathcal{I} \subset \mathcal{O}_X$ le faisceau quasi-cohérent d'idéaux correspondant.
Considérons la suite exacte courte
$$
0 \to \mathcal{I}^2 \to \mathcal{I} \to \mathcal{I}/\mathcal{I}^2 \to 0
$$
de faisceaux quasi-cohérents sur $X$. Le faisceau $\mathcal{I}/\mathcal{I}^2$ étant
annihilé par $\mathcal{I}$ cela correspond à un faisceau sur $Z$
par le lemme \ref{lemma-i-star-equivalence}. Ce module $\mathcal{O}_Z$ quasi-cohérent
est appelé le {\it faisceau conormal de
$Z$ dans $X$}
et est souvent simplement noté $\mathcal{I}/\mathcal{I}^2$ par l'abus
de notation mentionné dans la section \ref{section-closed-immersions-quasi-coherent}.

\medskip\noindent
Dans le cas où $i : Z \to X$ est une immersion (localement
fermée) nous définissons le faisceau conormal
de $i$ comme le faisceau conormal
de l'immersion fermée $i : Z \to X \setminus \partial Z$, où $\partial Z = \overline{Z} \setminus Z$.
Il est souvent noté $\mathcal{I}/\mathcal{I}^2$ où $\mathcal{I}$ est le faisceau
d'idéaux de l'immersion fermée $i : Z \to X \setminus \partial Z$.

\begin{definition}
\label{definition-conormal-sheaf}
Soit $i : Z \to X$ une immersion. Le {\it faisceau conormal $\mathcal{C}_{Z/X}$
de $Z$ dans $X$} ou le {\it faisceau conormal
de $i$} est le faisceau conormal quasi-cohérent $\mathcal{O}_Z$-module $\mathcal{I}/\mathcal{I}^2$
décrit ci-dessus.
\end{definition}

\noindent
Dans \cite[IV Définition 16.1.2]{EGA} ce faisceau
est noté $\mathcal{N}_{Z/X}$. Nous ne suivrons pas cette convention puisque nous souhaiterions
réserver la notation $\mathcal{N}_{Z/X}$ au {\it
faisceau normal de l'immersion}. Il est défini comme
$$
\mathcal{N}_{Z/X} =
\SheafHom_{\mathcal{O}_Z}(\mathcal{C}_{Z/X}, \mathcal{O}_Z) =
\SheafHom_{\mathcal{O}_Z}(\mathcal{I}/\mathcal{I}^2, \mathcal{O}_Z)
$$
à condition que le faisceau conormal soit de présentation finie (sinon le faisceau
normal pourrait même ne pas être quasi-cohérent). Nous reviendrons plus tard sur
le faisceau normal (insérer la référence future ici).

\begin{lemma}
\label{lemma-affine-conormal}
Soit $i : Z \to X$ une immersion. Le faisceau conormal de $i$
a les propriétés suivantes :
\begin{enumerate}
\item Soit $U \subset X$ tout sous-schéma ouvert tel
que $i$ se factorise en $Z \xrightarrow{i'} U \to X$ où $i'$
est une immersion fermée. Soit $\mathcal{I} = \Ker((i')^\sharp) \subset \mathcal{O}_U$.
Alors
$$
\mathcal{C}_{Z/X} = (i')^*\mathcal{I}\quad\text{et}\quad
i'_*\mathcal{C}_{Z/X} = \mathcal{I}/\mathcal{I}^2
$$
\item
Pour tout $\Spec(R) = U \subset X$ affine
ouvert tel que $Z \cap U = \Spec(R/I)$ il
existe un isomorphisme
canonique $\Gamma(Z \cap U, \mathcal{C}_{Z/X}) = I/I^2$.
\end{enumerate}
\end{lemma}

\begin{proof}
Généralement clair d'après les définitions. Notez que étant donné un anneau $R$
et un $I$ idéal de $R$ nous avons $I/I^2 = I \otimes_R R/I$. Détails omis.
\end{proof}

\begin{lemma}
\label{lemma-conormal-functorial}
Soit
$$
\xymatrix{
Z \ar[r]_i \ar[d]_f & X \ar[d]^g \\
Z' \ar[r]^{i'} & X'
}
$$
un diagramme commutatif dans la catégorie des schémas. Supposons
des immersions $i$, $i'$. Il existe une application
canonique de $\mathcal{O}_Z$-modules
$$
f^*\mathcal{C}_{Z'/X'}
\longrightarrow
\mathcal{C}_{Z/X}
$$
caractérisée par la propriété suivante : Pour toute
paire d'ouverts affines $(\Spec(R) = U \subset X, \Spec(R') = U' \subset X')$ à
$g(U) \subset U'$ telles que
$Z \cap U = \Spec(R/I)$ et $Z' \cap U' = \Spec(R'/I')$ l'application
induite
$$
\Gamma(Z' \cap U', \mathcal{C}_{Z'/X'}) = I'/I'^2
\longrightarrow
I/I^2 = \Gamma(Z \cap U, \mathcal{C}_{Z/X})
$$
est celle induite par l'homomorphisme d'anneaux $g^\sharp : R' \to R$
qui possède la propriété $g^\sharp(I') \subset I$.
\end{lemma}

\begin{proof}
Soient $\partial Z' = \overline{Z'} \setminus Z'$ et $\partial Z = \overline{Z} \setminus Z$. Ce sont
des sous-ensembles fermés de $X'$ et de $X$. En remplaçant
$X'$ par $X' \setminus \partial Z'$ et $X$ par
$X \setminus \big(g^{-1}(\partial Z') \cup \partial Z\big)$ on voit que l'on peut supposer que $i$
et $i'$ sont des immersions fermées.

\medskip\noindent
Le fait que $g \circ i$ se factorise
par $i'$ implique que les
$g^*\mathcal{I}'$ applications en
$\mathcal{I}$
sous l'application canonique $g^*\mathcal{I}' \to \mathcal{O}_X$,
voir Schémas, Lemmes
\ref{schemes-lemma-characterize-closed-subspace} et \ref{schemes-lemma-restrict-map-to-closed}. On
obtient donc une application induite de
faisceaux quasi-cohérents
$g^*(\mathcal{I}'/(\mathcal{I}')^2) \to \mathcal{I}/\mathcal{I}^2$. Retirer de $i$ donne
$i^*g^*(\mathcal{I}'/(\mathcal{I}')^2) \to i^*(\mathcal{I}/\mathcal{I}^2)$. Notez que $i^*(\mathcal{I}/\mathcal{I}^2) = \mathcal{C}_{Z/X}$.
D'autre part,
$i^*g^*(\mathcal{I}'/(\mathcal{I}')^2) =
f^*(i')^*(\mathcal{I}'/(\mathcal{I}')^2) = f^*\mathcal{C}_{Z'/X'}$. Cela donne l'application
souhaitée.

\medskip\noindent
Vérifier que l'application est décrite localement comme l'application
donnée $I'/(I')^2 \to I/I^2$ consiste à dérouler les définitions et
est omise. Une autre observation est que étant donné tout
$x \in i(Z)$ il existe des voisinages ouverts affines $U$,
$U'$ avec $f(U) \subset U'$ et $Z \cap U$ ainsi que $U' \cap Z'$ fermés
tels que $x \in U$. Démonstration omise. L’exigence du lemme
caractérise donc bien l’application (et aurait pu être utilisée pour
la définir).
\end{proof}

\begin{lemma}
\label{lemma-conormal-functorial-flat}
Soit
$$
\xymatrix{
Z \ar[r]_i \ar[d]_f & X \ar[d]^g \\
Z' \ar[r]^{i'} & X'
}
$$
un diagramme produit fibré dans la catégorie des schémas
à immersions $i$, $i'$.
Alors l'application canonique $f^*\mathcal{C}_{Z'/X'} \to \mathcal{C}_{Z/X}$ du
lemme \ref{lemma-conormal-functorial} est surjectif. Si
$g$ est plat, alors c'est un isomorphisme.
\end{lemma}

\begin{proof}
Soit $R' \to R$ un homomorphisme d'anneaux, et $I' \subset R'$
un idéal. Posons $I = I'R$. Alors $I'/(I')^2 \otimes_{R'} R \to I/I^2$ est
surjectif. Si $R' \to R$ est plat, alors $I = I' \otimes_{R'} R$ et
$I^2 = (I')^2 \otimes_{R'} R$ et on voit que l'application est un
isomorphisme.
\end{proof}

\begin{lemma}
\label{lemma-transitivity-conormal}
Soit $Z \to Y \to X$ des immersions de schémas. Il existe alors une suite exacte
canonique
$$
i^*\mathcal{C}_{Y/X} \to
\mathcal{C}_{Z/X} \to
\mathcal{C}_{Z/Y} \to 0
$$
dont les applications proviennent
du lemme \ref{lemma-conormal-functorial} et
$i : Z \to Y$ est le premier morphisme.
\end{lemma}

\begin{proof}
Via
Lemme \ref{lemma-conormal-functorial} cela se traduit
par le fait algébrique suivant. Supposons que $C \to B \to A$
soient des homomorphismes surjectifs d'anneaux. Soit $I = \Ker(B \to A)$,
$J = \Ker(C \to A)$ et $K = \Ker(C \to B)$. Il existe
alors une suite exacte
$$
K/K^2 \otimes_B A \to J/J^2 \to I/I^2 \to 0.
$$
Cela découle immédiatement de l'observation que $I = J/K$.
\end{proof}










\section{Faisceau des différentielles d'un morphisme}
\label{section-sheaf-differentials}

\noindent
Nous suggérons au lecteur de consulter la section correspondante
dans le chapitre sur l'algèbre
commutative (Algèbre, section \ref{algebra-section-differentials})
et la section correspondante dans le chapitre sur les faisceaux
de modules (Modules, section \ref{modules-section-differentials}).

\begin{definition}
\label{definition-sheaf-differentials}
Soit $f : X \to S$ un morphisme de schémas. Le {\it
faisceau des différentielles $\Omega_{X/S}$ de $X$ sur $S$}
est le faisceau des différentielles de $f$ vu comme un morphisme d'espaces
annelés (Modules, Définition \ref{modules-definition-differentials}) équipé de son {\it
dérivation $S$ universelle}
$$
\text{d}_{X/S} : \mathcal{O}_X \longrightarrow \Omega_{X/S}.
$$
\end{definition}

\noindent
Il s'avère que $\Omega_{X/S}$ est un $\mathcal{O}_X$-module quasi-cohérent par exemple
car il est isomorphe au faisceau conormal de la diagonale
du morphisme $\Delta : X \to X \times_S X$ (Lemme \ref{lemma-differentials-diagonal}). Nous avons
défini le module des différentielles de $X$
sur $S$ en utilisant une propriété universelle, à savoir
comme réceptacle de la dérivation universelle. Si vous avez une autre
construction du faisceau de différentielles relatives qui satisfait cette propriété
universelle alors, par le lemme de Yoneda, elle sera canoniquement
isomorphe à celle définie ci-dessus. Pour plus de commodité, nous reformulons
ici la propriété universelle.

\begin{lemma}
\label{lemma-universal-derivation-universal}
Soit $f : X \to S$ un morphisme de schémas. L'application
$$
\Hom_{\mathcal{O}_X}(\Omega_{X/S}, \mathcal{F})
\longrightarrow
\text{Der}_S(\mathcal{O}_X, \mathcal{F}),\quad
\alpha \longmapsto \alpha \circ \text{d}_{X/S}
$$
est un isomorphisme de foncteurs $\textit{Mod}(\mathcal{O}_X) \to \textit{Ensembles}$.
\end{lemma}

\begin{proof}
Il s'agit simplement d'une reformulation de la définition.
\end{proof}

\begin{lemma}
\label{lemma-differentials-restrict-open}
Soit $f : X \to S$ un morphisme de schémas.
Soit $U \subset X$, $V \subset S$ être des sous-schémas
ouverts tels que $f(U) \subset V$. Il existe
alors un unique isomorphisme $\Omega_{X/S}|_U = \Omega_{U/V}$ des $\mathcal{O}_U$-modules
tel que $\text{d}_{X/S}|_U = \text{d}_{U/V}$.
\end{lemma}

\begin{proof}
Ceci est un cas particulier
de Modules, Lemme \ref{modules-lemma-localize-differentials}
si l'on utilise l'identification
canonique $f^{-1}\mathcal{O}_S|_U = (f|_U)^{-1}\mathcal{O}_V$.
\end{proof}

\noindent
Désormais nous utiliserons ces identifications canoniques et écrirons
simplement $\Omega_{U/S}$ ou $\Omega_{U/V}$ pour la restriction de $\Omega_{X/S}$
à $U$.

\begin{lemma}
\label{lemma-affine-case-derivation}
Soit $R \to A$ un homomorphisme d'anneaux. Soit
$\mathcal{F}$ un faisceau de $\mathcal{O}_X$-modules
sur $X = \Spec(A)$. Posons $S = \Spec(R)$.
La règle qui associe à une $S$-dérivation sur $\mathcal{F}$ son
action sur les sections globales définit une bijection entre
l'ensemble des $S$-dérivations de $\mathcal{F}$ et l'ensemble
des $R$-dérivations sur $M = \Gamma(X, \mathcal{F})$.
\end{lemma}

\begin{proof}
Soit $D : A \to M$ une dérivation $R$. Nous devons montrer qu'il existe
une unique dérivation $S$ sur $\mathcal{F}$ qui donne naissance
à $D$ sur les sections globales. Soit $U = D(f) \subset X$ un ouvert affine standard.
Tout élément de $\Gamma(U, \mathcal{O}_X)$ est de la forme $a/f^n$ pour
certains $a \in A$ et $n \geq 0$. Par la règle de Leibniz, nous
avons
$$
D(a)|_U = a/f^n D(f^n)|_U + f^n D(a/f^n)
$$
dans $\Gamma(U, \mathcal{F})$. Puisque $f$ agit de manière
inversible sur $\Gamma(U, \mathcal{F})$, cela détermine complètement
la valeur de $D(a/f^n) \in \Gamma(U, \mathcal{F})$. Cela prouve l’unicité.
L'existence s'ensuit en définissant simplement
$$
D(a/f^n) := (1/f^n) D(a)|_U - a/f^{2n} D(f^n)|_U
$$
et en prouvant qu'il possède toutes les propriétés souhaitées (sur la base des
ouvertures standard de $X$). Détails omis.
\end{proof}

\begin{lemma}
\label{lemma-differentials-affine}
Soit $f : X \to S$ un morphisme de schémas. Pour toute paire d'affines
ouvertes $\Spec(A) = U \subset X$, $\Spec(R) = V \subset S$ avec $f(U) \subset V$ il existe
un isomorphisme unique
$$
\Gamma(U, \Omega_{X/S}) = \Omega_{A/R}.
$$
compatible avec $\text{d}_{X/S}$ et $\text{d} : A \to \Omega_{A/R}$.
\end{lemma}

\begin{proof}
Par le Lemme \ref{lemma-differentials-restrict-open} on peut remplacer $X$ et $S$
par $U$ et $V$. Ainsi nous pouvons supposer
$X = \Spec(A)$ et $S = \Spec(R)$ et nous devons montrer le
lemme avec $U = X$ et $V = S$. Considérons le module
$A$ $M = \Gamma(X, \Omega_{X/S})$ avec la dérivation $R$ $\text{d}_{X/S} : A \to M$.
Soit $N$ un autre $A$-module et désignons $\widetilde{N}$
le module $\mathcal{O}_X$ quasi-cohérent
associé à $N$, voir Schémas, section \ref{schemes-section-quasi-coherent-affine}.
En précomposant par $\text{d}_{X/S} : A \to M$ on obtient une flèche
$$
\alpha : \Hom_A(M, N) \longrightarrow \text{Der}_R(A, N)
$$
En utilisant les lemmes \ref{lemma-universal-derivation-universal} et \ref{lemma-affine-case-derivation}
on obtient les identifications
$$
\Hom_{\mathcal{O}_X}(\Omega_{X/S}, \widetilde{N}) =
\text{Der}_S(\mathcal{O}_X, \widetilde{N}) =
\text{Der}_R(A, N)
$$
La prise de sections globales détermine
une flèche $\Hom_{\mathcal{O}_X}(\Omega_{X/S}, \widetilde{N}) \to \Hom_R(M, N)$. En combinant cette flèche et les
identifications ci-dessus nous obtenons une flèche
$$
\beta : \text{Der}_R(A, N) \longrightarrow \Hom_R(M, N)
$$
En vérifiant ce qui se passe sur les sections globales, nous
constatons que $\alpha$ et $\beta$ sont l'un
l'autre inverses. Nous voyons donc
que $\text{d}_{X/S} : A \to M$ satisfait la même propriété universelle que
$\text{d} : A \to \Omega_{A/R}$, voir Algèbre, Lemme \ref{algebra-lemma-universal-omega}. Ainsi le lemme de Yoneda (Catégories,
Lemme \ref{categories-lemma-yoneda}) implique qu'il existe
un unique isomorphisme des $A$-modules
$M \cong \Omega_{A/R}$ compatible avec les dérivations.
\end{proof}

\begin{remark}
\label{remark-differentials-glue}
Le lemme ci-dessus donne une deuxième façon de construire le module des
différentielles. En effet, soit $f : X \to S$ un morphisme de schémas. Considérons
la famille de tous les ouverts affines $U \subset X$ dont l'image est contenue
dans un ouvert affine de $S$. Ceux-ci forment une base
de la topologie sur $X$. Il suffit donc de définir $\Gamma(U, \Omega_{X/S})$
pour un tel $U$. Nous définissons simplement $\Gamma(U, \Omega_{X/S}) = \Omega_{A/R}$ si $A$,
$R$ sont comme dans le lemme \ref{lemma-differentials-affine} ci-dessus. Cela fonctionne,
mais il faut des préliminaires un peu plus algébriques pour
construire les mappages de restriction et vérifier la condition
du faisceau avec cet ansatz.
\end{remark}

\noindent
Le lemme suivant donne encore une autre façon de définir le faisceau des différentielles
et il montre en particulier que $\Omega_{X/S}$ est quasi-cohérent
si $X$ et $S$ sont des schémas.

\begin{lemma}
\label{lemma-differentials-diagonal}
Soit $f : X \to S$ un morphisme de schémas. Il existe
un isomorphisme canonique entre $\Omega_{X/S}$ et le
faisceau conormal du morphisme diagonal $\Delta_{X/S} : X \longrightarrow X \times_S X$.
\end{lemma}

\begin{proof}
Nous établissons d'abord l'existence de quelques ``globaux'' faisceaux
et applications globales de faisceaux, et plus loin nous décrivons
les constructions sur certaines ouverts affines.

\medskip\noindent
Rappelons que $\Delta = \Delta_{X/S} : X \to X \times_S X$ est une immersion,
voir Schémas, Lemme \ref{schemes-lemma-diagonal-immersion}. Soit $\mathcal{J}$ le faisceau d'idéaux
de l'immersion qui habite un sous-schéma
ouvert $W$ de $X \times_S X$ tel que $\Delta(X) \subset W$
soit fermé. Prenons celui que l'on retrouve dans la
démonstration d'après
Schémas, Lemme \ref{schemes-lemma-diagonal-immersion}. Notez que le faisceau
d'anneaux $\mathcal{O}_W/\mathcal{J}^2$ est supporté sur $\Delta(X)$.
De plus il se situe dans une suite exacte
courte de faisceaux
$$
0 \to \mathcal{J}/\mathcal{J}^2
\to \mathcal{O}_W/\mathcal{J}^2
\to \Delta_*\mathcal{O}_X
\to 0.
$$
En utilisant $\Delta^{-1}$ on peut penser à cela comme
une surjection de faisceaux de $f^{-1}\mathcal{O}_S$-algèbres
avec noyau le faisceau conormal de $\Delta$ (voir Définition
\ref{definition-conormal-sheaf} et Lemme \ref{lemma-affine-conormal}).
$$
0 \to \mathcal{C}_{X/X \times_S X}
\to \Delta^{-1}(\mathcal{O}_W/\mathcal{J}^2)
\to \mathcal{O}_X
\to 0
$$
Cela nous place dans la
situation des Modules, Lemme \ref{modules-lemma-double-structure-gives-derivation}. Les morphismes
de projection $p_i : X \times_S X \to X$, $i = 1, 2$ induisent
des applications
de faisceaux d'anneaux $(p_i)^\sharp : (p_i)^{-1}\mathcal{O}_X \to \mathcal{O}_{X \times_S X}$. Nous pouvons
nous limiter à $W$ et quotient par
$\mathcal{J}^2$ pour obtenir $(p_i)^{-1}\mathcal{O}_X \to \mathcal{O}_W/\mathcal{J}^2$. Puisque
$\Delta^{-1}p_i^{-1}\mathcal{O}_X = \mathcal{O}_X$, nous obtenons des
applications
$$
s_i : \mathcal{O}_X \to \Delta^{-1}(\mathcal{O}_W/\mathcal{J}^2).
$$
$s_1$ et $s_2$ sont des
sections de l'application $\Delta^{-1}(\mathcal{O}_W/\mathcal{J}^2) \to \mathcal{O}_X$, comme dans
Modules, Lemme \ref{modules-lemma-double-structure-gives-derivation}. Nous obtenons ainsi une dérivation
$S$ $\text{d} = s_2 - s_1 : \mathcal{O}_X \to \mathcal{C}_{X/X \times_S X}$.
Par la propriété universelle du module des
différentielles nous trouvons une application unique
$\mathcal{O}_X$-linéaire
$$
\Omega_{X/S} \longrightarrow \mathcal{C}_{X/X \times_S X},\quad
f\text{d}g \longmapsto fs_2(g) - fs_1(g)
$$
Pour voir que l'application est un isomorphisme, calculons cela sur
des ouverts affines appropriées. On peut couvrir $X$
par des ouverts affines $\Spec(A) = U \subset X$ dont l'image est contenue
dans un ouvert affine $\Spec(R) = V \subset S$. D'après la démonstration d'après
Schémas, Lemme \ref{schemes-lemma-diagonal-immersion} $U \times_V U \subset X \times_S X$ est un ouvert
affine contenue dans l'ouvert $W$
mentionné ci-dessus. Également
$U \times_V U = \Spec(A \otimes_R A)$. Le faisceau $\mathcal{J}$ correspond
à l'idéal $J = \Ker(A \otimes_R A \to A)$. La
suite exacte courte à la suite exacte courte
des $A \otimes_R A$-modules
$$
0 \to J/J^2 \to (A \otimes_R A)/J^2 \to A \to 0
$$
Les sections $s_i$ correspondent aux homomorphismes d'anneaux
$$
A \longrightarrow (A \otimes_R A)/J^2,\quad
s_1 : a \mapsto a \otimes 1,\quad
s_2 : a \mapsto 1 \otimes a.
$$
Par le Lemme \ref{lemma-affine-conormal} nous avons
$\Gamma(U, \mathcal{C}_{X/X \times_S X}) = J/J^2$ et par le lemme
\ref{lemma-differentials-affine} nous avons $\Gamma(U, \Omega_{X/S}) = \Omega_{A/R}$.
L'application ci-dessus
est l'application $a \text{d}b \mapsto a \otimes b - ab \otimes 1$, dont
Algèbre, Lemme \ref{algebra-lemma-differentials-diagonal}, montre qu'elle est
un isomorphisme.
\end{proof}

\begin{lemma}
\label{lemma-functoriality-differentials}
Soit
$$
\xymatrix{
X' \ar[d] \ar[r]_f & X \ar[d] \\
S' \ar[r] & S
}
$$
être un diagramme commutatif de schémas. L'application
canonique $\mathcal{O}_X \to f_*\mathcal{O}_{X'}$ composée avec l'application
$f_*\text{d}_{X'/S'} : f_*\mathcal{O}_{X'} \to f_*\Omega_{X'/S'}$ est une $S$-dérivation. On obtient
donc une application canonique des $\mathcal{O}_X$-modules $\Omega_{X/S} \to f_*\Omega_{X'/S'}$,
et par adjointité de $f_*$
et $f^*$ un homomorphisme
canonique des $\mathcal{O}_{X'}$-modules
$$
c_f : f^*\Omega_{X/S} \longrightarrow \Omega_{X'/S'}.
$$
Il est uniquement caractérisé par la propriété que
les applications $f^*\text{d}_{X/S}(h)$ à $\text{d}_{X'/S'}(f^* h)$ pour toute
section locale $h$ de $\mathcal{O}_X$.
\end{lemma}

\begin{proof}
Ceci est un cas particulier
de Modules, Lemme \ref{modules-lemma-functoriality-differentials-ringed-spaces}. Dans le cas des schémas, on peut
aussi utiliser la fonctorialité des faisceaux conormaux
(voir le lemme \ref{lemma-conormal-functorial}) et le lemme \ref{lemma-differentials-diagonal} pour définir
$c_f$. Ou nous pouvons utiliser la caractérisation
de la dernière ligne du lemme pour coller des
applications définies sur des
ouverts affines (voir Algèbre, Equation (\ref{algebra-equation-functorial-omega})).
\end{proof}

\begin{lemma}
\label{lemma-triangle-differentials}
Soient $f : X \to Y$, $g : Y \to S$ des morphismes de schémas. Il
existe alors une suite exacte canonique
$$
f^*\Omega_{Y/S} \to \Omega_{X/S} \to \Omega_{X/Y} \to 0
$$
où les applications proviennent
d'applications du lemme \ref{lemma-functoriality-differentials}.
\end{lemma}

\begin{proof}
Il s'agit de la version simplifiée
de l'algèbre, Lemme \ref{algebra-lemma-exact-sequence-differentials}. Alternativement,
il existe une version générale pour les morphismes
des espaces annelés, voir Modules, Lemme \ref{modules-lemma-triangle-differentials}.
\end{proof}

\begin{lemma}
\label{lemma-base-change-differentials}
Soit $X \to S$ un morphisme de schémas. Soit
$g : S' \to S$ un morphisme de schémas. Soit $X' = X_{S'}$
le changement de base de $X$.
Notons $g' : X' \to X$ la projection. Alors
l'application
$$
(g')^*\Omega_{X/S} \to \Omega_{X'/S'}
$$
du lemme \ref{lemma-functoriality-differentials} est un isomorphisme.
\end{lemma}

\begin{proof}
Il s'agit de la version
simplifiée de l'algèbre, Lemme \ref{algebra-lemma-differentials-base-change}.
\end{proof}

\begin{lemma}
\label{lemma-differential-product}
Soient $f : X \to S$ et $g : Y \to S$ des morphismes de schémas
de même cible. Soit $p : X \times_S Y \to X$ et $q : X \times_S Y \to Y$ les morphismes
de projection. Les
applications du lemme \ref{lemma-functoriality-differentials}
$$
p^*\Omega_{X/S} \oplus q^*\Omega_{Y/S}
\longrightarrow
\Omega_{X \times_S Y/S}
$$
donnent un isomorphisme.
\end{lemma}

\begin{proof}
Par le Lemme \ref{lemma-base-change-differentials} la composition
$p^*\Omega_{X/S} \to \Omega_{X \times_S Y/S} \to \Omega_{X \times_S Y/Y}$ est un isomorphisme, et de même
pour $q$. De plus, le conoyau
de $p^*\Omega_{X/S} \to \Omega_{X \times_S Y/S}$ est $\Omega_{X \times_S Y/X}$
par le lemme
\ref{lemma-triangle-differentials}. Le résultat suit.
\end{proof}

\begin{lemma}
\label{lemma-finite-type-differentials}
Soit $f : X \to S$ un morphisme de schémas. Si $f$
est localement de type fini, alors $\Omega_{X/S}$ est
un $\mathcal{O}_X$-module de type fini.
\end{lemma}

\begin{proof}
Immédiat
de l'algèbre, Lemme \ref{algebra-lemma-differentials-finitely-generated},
Lemme \ref{lemma-differentials-affine},
Lemme \ref{lemma-locally-finite-type-characterize} et Propriétés,
Lemme \ref{properties-lemma-finite-type-module}.
\end{proof}

\begin{lemma}
\label{lemma-finite-presentation-differentials}
Soit $f : X \to S$ un morphisme de schémas. Si
$f$ est localement de présentation finie, alors $\Omega_{X/S}$
est un $\mathcal{O}_X$-module de présentation finie.
\end{lemma}

\begin{proof}
Immédiat
de l'algèbre, Lemme \ref{algebra-lemma-differentials-finitely-presented},
Lemme \ref{lemma-differentials-affine},
Lemme \ref{lemma-locally-finite-presentation-characterize} et Propriétés,
Lemme \ref{properties-lemma-finite-presentation-module}.
\end{proof}

\begin{lemma}
\label{lemma-immersion-differentials}
Si $X \to S$ est une immersion, ou plus généralement un monomorphisme, alors
$\Omega_{X/S}$ est nul.
\end{lemma}

\begin{proof}
Ceci est vrai car $\Delta_{X/S}$ est un isomorphisme
dans ce cas et a donc un faisceau conormal trivial.
D'où $\Omega_{X/S} = 0$ par le lemme \ref{lemma-differentials-diagonal}. La version
algébrique est Algèbre, Lemme \ref{algebra-lemma-trivial-differential-surjective}.
\end{proof}

\begin{lemma}
\label{lemma-differentials-relative-immersion}
Soit $i : Z \to X$ une immersion de schémas sur $S$. Il
existe une suite exacte canonique
$$
\mathcal{C}_{Z/X} \to i^*\Omega_{X/S} \to \Omega_{Z/S} \to 0
$$
où la première flèche est induite par $\text{d}_{X/S}$
et la deuxième flèche vient du lemme \ref{lemma-functoriality-differentials}.
\end{lemma}

\begin{proof}
Il s'agit de la version simplifiée
de l'algèbre, Lemme \ref{algebra-lemma-differential-seq}. Cependant, nous
devons nous assurer que nous pouvons définir la
première flèche de manière globale. Nous expliquons donc ici
la signification de `` induit par $\text{d}_{X/S}$''.
À savoir, on peut supposer
que $i$ est une immersion fermée par rétrécissement
$X$. Soit $\mathcal{I} \subset \mathcal{O}_X$ le
faisceau d'idéaux correspondant à $Z \subset X$. Puis
$\text{d}_{X/S} : \mathcal{I} \to \Omega_{X/S}$ applique le sous-faisceau $\mathcal{I}^2 \subset \mathcal{I}$
à $\mathcal{I}\Omega_{X/S}$. Il induit donc une application $\mathcal{I}/\mathcal{I}^2 \to \Omega_{X/S}/\mathcal{I}\Omega_{X/S}$
qui est $\mathcal{O}_X/\mathcal{I}$-linéaire.
Par le Lemme \ref{lemma-i-star-equivalence} cela correspond à une application
$\mathcal{C}_{Z/X} \to i^*\Omega_{X/S}$ souhaitée.
\end{proof}

\begin{lemma}
\label{lemma-differentials-relative-immersion-section}
Soit $i : Z \to X$ une immersion de schémas sur $S$, et supposons
que $i$ (localement) a un inverse gauche. Alors la séquence
canonique
$$
0 \to \mathcal{C}_{Z/X} \to i^*\Omega_{X/S} \to \Omega_{Z/S} \to 0
$$
du
lemme \ref{lemma-differentials-relative-immersion} est (localement) divisée
exactement. En particulier, si $s : S \to X$ est une section
du morphisme de structure $X \to S$ alors
l'application $\mathcal{C}_{S/X} \to s^*\Omega_{X/S}$ induite par $\text{d}_{X/S}$
est un isomorphisme.
\end{lemma}

\begin{proof}
Il résulte de l'algèbre, Lemme \ref{algebra-lemma-differential-seq-split} que le résultat est valable
lorsque $X, Z, S$ sont affines. En général, nous concluons
que la séquence est exacte localement divisée. Si $g : X \to Z$
est un inverse gauche de $i$, alors
$i^*c_g$ est un inverse droit
de l'application $i^*\Omega_{X/S} \to \Omega_{Z/S}$ par Modules, Lemme \ref{modules-lemma-check-functoriality-differentials}.
Cela implique que la séquence est
divisée par homologie, Lemme \ref{homology-lemma-ses-split}.
Enfin, si $s$ est une section, alors c'est une immersion
$s : Z = S \to X$ sur $S$ (voir Schémas, Lemme \ref{schemes-lemma-section-immersion})
et dans ce cas $\Omega_{Z/S} = 0$.
\end{proof}

\begin{remark}
\label{remark-differentials-diagonal}
Soit $X \to S$ un morphisme de schémas.
D'après le Lemme \ref{lemma-differential-product}
on a
$$
\Omega_{X \times_S X/S} =
\text{pr}_1^*\Omega_{X/S} \oplus \text{pr}_2^*\Omega_{X/S}
$$
Par contre, le morphisme diagonal $\Delta : X \to X \times_S X$ est une immersion,
qui a localement un inverse gauche.
Ainsi par le lemme \ref{lemma-differentials-relative-immersion-section} on obtient une
suite canonique exacte courte
$$
0 \to \mathcal{C}_{X/X \times_S X} \to \Omega_{X/S} \oplus \Omega_{X/S}
\to \Omega_{X/S} \to 0
$$
Notons que la flèche droite est $(1, 1)$ qui
est bien une surjection fractionnée. Par
contre, par le lemme \ref{lemma-differentials-diagonal} nous avons une
identification $\Omega_{X/S} = \mathcal{C}_{X/X \times_S X}$. Comme nous avons choisi $\text{d}_{X/S}(f) = s_2(f) - s_1(f)$ dans cette
identification
il apparaît que la flèche de gauche est l'application
$(-1, 1)$\footnote{Plus précisément,
la section locale $\text{d}_{X/S}(f) = 1 \otimes f - f \otimes 1$ du faisceau d'idéaux de $\Delta$
s'envoie par $\text{d}_{X \times_S X/X}$ sur la section locale
$1 \otimes 1 \otimes 1 \otimes f - 1 \otimes f \otimes 1 \otimes 1
-1 \otimes 1 \otimes f \otimes 1 + f \otimes 1 \otimes 1 \otimes 1 =
\text{pr}_2^*\text{d}_{X/S}(f) - \text{pr}_1^*\text{d}_{X/S}(f)$.}.
\end{remark}

\begin{lemma}
\label{lemma-two-immersions}
Soit
$$
\xymatrix{
Z \ar[r]_i \ar[rd]_j & X \ar[d] \\
& Y
}
$$
un diagramme commutatif de schémas où $i$ et $j$ sont des immersions.
Il existe alors une suite exacte canonique
$$
\mathcal{C}_{Z/Y} \to
\mathcal{C}_{Z/X} \to
i^*\Omega_{X/Y} \to 0
$$
où la première flèche
vient du lemme \ref{lemma-conormal-functorial}
et la
seconde du lemme \ref{lemma-differentials-relative-immersion}.
\end{lemma}

\begin{proof}
La version algébrique de
ceci est Algèbre, Lemme \ref{algebra-lemma-application-NL}.
\end{proof}










\section{Opérateurs différentiels d'ordre fini}
\label{section-differential-operators}

\noindent
Nous suggérons au lecteur de consulter la section correspondante
dans le chapitre sur
l'algèbre commutative (Algèbre, section \ref{algebra-section-differential-operators})
et la section correspondante dans le chapitre sur les
faisceaux de modules (Modules, section \ref{modules-section-differential-operators}).

\begin{lemma}
\label{lemma-affine-case-differential-operators}
Soit $R \to A$ un homomorphisme d'anneaux. Notons $f : X \to S$ le morphisme
correspondant des schémas affines. Soit $\mathcal{F}$ et $\mathcal{G}$ être
des $\mathcal{O}_X$-modules. Si $\mathcal{F}$ est quasi-cohérent alors
l'application
$$
\text{Diff}^k_{X/S}(\mathcal{F}, \mathcal{G}) \to
\text{Diff}^k_{A/R}(\Gamma(X, \mathcal{F}), \Gamma(X, \mathcal{G}))
$$
envoyant un opérateur différentiel à son action sur les sections globales
est bijective.
\end{lemma}

\begin{proof}
Écrivez $\mathcal{F} = \widetilde{M}$ pour certains $A$-module $M$.
Posons $N = \Gamma(X, \mathcal{G})$. Soit $D : M \to N$ un opérateur différentiel
d'ordre $k$. Nous devons montrer
qu'il existe un unique opérateur différentiel $\mathcal{F} \to \mathcal{G}$ d'ordre $k$
qui donne naissance à $D$ sur les sections globales.
Soit $U = D(f) \subset X$ un ouvert affine standard. Alors $\mathcal{F}(U) = M_f$
est la localisation. Par Algèbre, Lemme \ref{algebra-lemma-invert-system-differential-operators} l'opérateur différentiel
$D$ s'étend à un unique opérateur différentiel
$$
D_f : \mathcal{F}(U) = \widetilde{M}(U) = M_f \to N_f = \widetilde{N}(U)
$$
L'unicité montre que ces applications $D_f$ collent pour donner
une application de faisceaux $\widetilde{M} \to \widetilde{N}$ sur la base de toutes
les ouvertures standards de $X$.
Nous obtenons donc une application unique de faisceaux $\widetilde{D} : \widetilde{M} \to \widetilde{N}$
en accord avec ces applications par le matériel de Faisceaux,
section \ref{sheaves-section-bases}. Puisque $\widetilde{D}$ est donné par des
opérateurs différentiels d'ordre $k$ sur le standard ouvre,
on constate que $\widetilde{D}$ est un
opérateur différentiel d'ordre $k$ (petit détail omis). Enfin,
on peut post-composer avec
l'application canonique $\mathcal{O}_X$-module $c : \widetilde{N} \to \mathcal{G}$
(Schémas, Lemme \ref{schemes-lemma-compare-constructions}) pour obtenir $c \circ \widetilde{D} : \mathcal{F} \to \mathcal{G}$
qui est un opérateur différentiel d'ordre $k$
par Modules, Lemme \ref{modules-lemma-composition-differential-operators}. Cela prouve l'existence.
Nous omettons la démonstration de l'unicité.
\end{proof}

\begin{lemma}
\label{lemma-base-change-differential-operators}
Soient $a : X \to S$ et $b : Y \to S$ des morphismes de schémas.
Soit $\mathcal{F}$ et $\mathcal{F}'$ des modules $\mathcal{O}_X$ quasi-cohérents. Soit
$D : \mathcal{F} \to \mathcal{F}'$ un opérateur différentiel d'ordre $k$
sur $X/S$. Soit $\mathcal{G}$ un module $\mathcal{O}_Y$ quasi-cohérent.
Il existe alors un unique opérateur
différentiel
$$
D' :
\text{pr}_1^*\mathcal{F} \otimes_{\mathcal{O}_{X \times_S Y}}
\text{pr}_2^*\mathcal{G}
\longrightarrow
\text{pr}_1^*\mathcal{F}' \otimes_{\mathcal{O}_{X \times_S Y}}
\text{pr}_2^*\mathcal{G}
$$
d'ordre $k$ sur $X \times_S Y / Y$ tel que
$
D'(s \otimes t) = D(s) \otimes t
$
pour les sections locales $s$ de $\mathcal{F}$ et $t$ de $\mathcal{G}$.
\end{lemma}

\begin{proof}
Dans le cas où $X$, $Y$ et $S$
sont affines, cela découle, via Lemme \ref{lemma-affine-case-differential-operators},
du résultat algébrique correspondant,
voir Algèbre, Lemme \ref{algebra-lemma-base-change-differential-operators}. En général, on
utilise des revêtements par affines
(par exemple comme dans Schémas,
Lemme \ref{schemes-lemma-affine-covering-fibre-product}) pour construire $D'$
globalement. Détails omis.
\end{proof}

\begin{remark}
\label{remark-base-change-differential-operators}
Soient $a : X \to S$ et $b : Y \to S$ des morphismes de schémas.
Notons $p : X \times_S Y \to X$ et $q : X \times_S Y \to Y$ les projections.
Dans cette remarque, étant donné un $\mathcal{O}_X$-module $\mathcal{F}$
et un $\mathcal{O}_Y$-module $\mathcal{G}$ fixons
$$
\mathcal{F} \boxtimes \mathcal{G} =
p^*\mathcal{F} \otimes_{\mathcal{O}_{X \times_S Y}} q^*\mathcal{G}
$$
Désignons $\mathcal{A}_{X/S}$ la catégorie additive
dont les objets sont des $\mathcal{O}_X$-modules quasi-cohérents
et dont les morphismes sont des opérateurs différentiels
d'ordre fini sur $X/S$. De même pour $\mathcal{A}_{Y/S}$
et $\mathcal{A}_{X \times_S Y/S}$. La construction du lemme \ref{lemma-base-change-differential-operators} détermine
un foncteur
$$
\boxtimes : 
\mathcal{A}_{X/S} \times \mathcal{A}_{Y/S} \longrightarrow
\mathcal{A}_{X \times_S Y/S},
\quad
(\mathcal{F}, \mathcal{G}) \longmapsto \mathcal{F} \boxtimes \mathcal{G}
$$
qui est bilinéaire sur les morphismes. Si $X = \Spec(A)$, $Y = \Spec(B)$,
et $S = \Spec(R)$, alors via l'identification de
faisceaux quasi-cohérents avec des modules ce foncteur
est donné par $(M, N) \mapsto M \otimes_R N$ sur les objets et envoie
le morphisme $(D, D') : (M, N) \to (M', N')$ à $D \otimes D' : M \otimes_R N \to M' \otimes_R N'$.
\end{remark}








\section{Morphismes lisses}
\label{section-smooth}

\noindent
Soit $f : X \to Y$ une application continue d'espaces topologiques.
Considérons la condition
suivante : Pour tout $x \in X$ il existe des voisinages
ouverts $x \in U \subset X$ et $f(x) \in V \subset Y$, et un entier $d$
tel que $f(U) \subset V$ et tel que l'on obtient un diagramme
commutatif
$$
\xymatrix{
X \ar[d] &
U \ar[l] \ar[d] \ar[r]_-\pi &
V \times \mathbf{R}^d \ar[ld] \\
Y & V \ar[l]
}
$$
où $\pi$ est un homéomorphisme sur un sous-ensemble
ouvert. Les morphismes de schémas lisses
sont l'analogue de ces applications
dans la catégorie des schémas.
Voir
Lemme \ref{lemma-smooth-locally-standard-smooth} et Lemme \ref{lemma-smooth-etale-over-affine-space}.

\medskip\noindent Contrairement
aux attentes (peut-être) la notion d'homomorphisme
d'anneaux lisse n'est pas définie uniquement en termes
de module des différentielles. À savoir, rappelons que $R \to A$
est un {\it homomorphisme d'anneaux lisse} si $A$ est de
présentation finie sur $R$ et si le complexe cotangent naïf de $A$ sur
$R$ est quasi-isomorphe à un module projectif placé
en degré $0$, voir Algèbre, Définition \ref{algebra-definition-smooth}.

\begin{definition}
\label{definition-smooth}
Soit $f : X \to S$ un morphisme de schémas.
\begin{enumerate}
\item On dit que $f$ est {\it lisse
à $x \in X$} s'il existe un voisinage ouvert affine $\Spec(A) = U \subset X$
de $x$ et ouvert affine $\Spec(R) = V \subset S$ avec
$f(U) \subset V$ tel que l'homomorphisme d'anneaux induit $R \to A$
soit lisse.
\item On dit que $f$ est {\it lisse} s'il est lisse en chaque point de $X$.
\item Un morphisme de schémas affines $f : X \to S$
est appelé {\it standard lisse} s'il
existe un homomorphisme d'anneaux lisse standard
$R \to R[x_1, \ldots, x_n]/(f_1, \ldots, f_c)$ (voir Algèbre, Définition \ref{algebra-definition-standard-smooth}) tel
que $X \to S$ est isomorphe à
$$
\Spec(R[x_1, \ldots, x_n]/(f_1, \ldots, f_c)) \to \Spec(R).
$$
\end{enumerate}
\end{definition}

\noindent
Une caractéristique intéressante de cette définition est que l'ensemble des points
où un morphisme est lisse est automatiquement ouvert.

\medskip\noindent
A noter qu'il n'y a pas d'hypothèse de séparation ou de quasi-compacité dans
la définition. La question de la fluidité est donc de nature locale sur
la source. Voici le résultat précis.

\begin{lemma}
\label{lemma-smooth-characterize}
Soit $f : X \to S$ un morphisme de schémas. Les conditions
suivantes sont équivalentes
\begin{enumerate}
\item Le morphisme $f$ est lisse.
\item Pour chaque ouverture affine $U \subset X$, $V \subset S$
avec $f(U) \subset V$ l'homomorphisme
d'anneaux $\mathcal{O}_S(V) \to \mathcal{O}_X(U)$ est lisse.
\item Il existe un recouvrement ouvert $S = \bigcup_{j \in J} V_j$ et
des recouvrements ouverts $f^{-1}(V_j) = \bigcup_{i \in I_j} U_i$ tels que chacun
des morphismes $U_i \to V_j$, $j\in J, i\in I_j$ soit
lisse.
\item Il existe un recouvrement ouvert affine $S = \bigcup_{j \in J} V_j$
et ouvert affine revêtements $f^{-1}(V_j) = \bigcup_{i \in I_j} U_i$ tel que
l'homomorphisme d'anneaux $\mathcal{O}_S(V_j) \to \mathcal{O}_X(U_i)$ soit lisse,
pour tout $j\in J, i\in I_j$.
\end{enumerate}
De plus, si $f$ est lisse alors
pour tout sous-schémas ouverts $U \subset X$, $V \subset S$ avec $f(U) \subset V$
la restriction $f|_U : U \to V$ est lisse.
\end{lemma}

\begin{proof}
Cela résulte du lemme \ref{lemma-locally-P} si l'on montre que la propriété
``$R \to A$ est lisse''
est locale. On vérifie les conditions (a), (b)
et (c) de la Définition \ref{definition-property-local}.
Par Algèbre, Lemme \ref{algebra-lemma-base-change-smooth}, la lissité est stable
par changement de base ; on en déduit
la condition (a). Par
Algèbre, Lemme \ref{algebra-lemma-compose-smooth}, la lissité est stable
par composition, et pour tout anneau $R$ l'homomorphisme
d'anneaux $R \to R_f$ est (standard) lisse.
On en déduit la condition (b). Enfin, la
condition (c) résulte d'Algèbre, Lemme \ref{algebra-lemma-locally-smooth}.
\end{proof}

\noindent
Le lemme suivant caractérise un morphisme lisse comme un
morphisme plat, de présentation finie, avec des fibres
lisses. A noter que les schémas lissant un corps sont abordés plus
en détail dans Variétés, section \ref{varieties-section-smooth}.

\begin{lemma}
\label{lemma-smooth-flat-smooth-fibres}
Soit $f : X \to S$ un morphisme de schémas. Si $f$
est plat, localement de présentation finie, et que toutes les fibres
$X_s$ sont lisses, alors $f$
est lisse.
\end{lemma}

\begin{proof}
Cela résulte d'Algèbre, Lemme \ref{algebra-lemma-flat-fibre-smooth}.
\end{proof}

\begin{lemma}
\label{lemma-composition-smooth}
La composition de deux morphismes lisses est lisse.
\end{lemma}

\begin{proof}
Dans la démonstration du lemme \ref{lemma-smooth-characterize} nous avons
vu qu'être lisse est une propriété locale des homomorphismes
d'anneaux. D'où la première affirmation
du lemme découle du lemme \ref{lemma-composition-type-P} combinée
au fait qu'être lisse est une propriété des homomorphismes d'anneaux
qui est stable sous
composition, voir Algèbre, Lemme \ref{algebra-lemma-compose-smooth}.
\end{proof}

\begin{lemma}
\label{lemma-base-change-smooth}
Le changement de base d'un morphisme qui est lisse est lisse.
\end{lemma}

\begin{proof}
Dans la démonstration du lemme \ref{lemma-smooth-characterize} nous avons
vu qu'être lisse est une propriété locale des homomorphismes
d'anneaux. D'où
le lemme découle du lemme \ref{lemma-composition-type-P} combiné
au fait qu'être lisse est une propriété des homomorphismes
d'anneaux stables par changement
de base, voir Algèbre, Lemme \ref{algebra-lemma-base-change-smooth}.
\end{proof}

\begin{lemma}
\label{lemma-open-immersion-smooth}
Toute immersion ouverte est fluide.
\end{lemma}

\begin{proof}
Ceci est vrai car une immersion ouverte est un isomorphisme local.
\end{proof}

\begin{lemma}
\label{lemma-smooth-syntomic}
Un morphisme lisse est syntomique.
\end{lemma}

\begin{proof}
Voir Algèbre, Lemme \ref{algebra-lemma-smooth-syntomic}.
\end{proof}

\begin{lemma}
\label{lemma-smooth-locally-finite-presentation}
Un morphisme lisse est localement de présentation finie.
\end{lemma}

\begin{proof}
Vrai car un homomorphisme d'anneaux lisse est de présentation finie par
définition.
\end{proof}

\begin{lemma}
\label{lemma-smooth-flat}
Un morphisme lisse est plat.
\end{lemma}

\begin{proof}
Cela résulte des lemmes \ref{lemma-syntomic-flat} et \ref{lemma-smooth-syntomic}.
\end{proof}

\begin{lemma}
\label{lemma-smooth-open}
Un morphisme lisse est universellement ouvert.
\end{lemma}

\begin{proof}
Combine
Lemmes \ref{lemma-smooth-flat},
\ref{lemma-smooth-locally-finite-presentation} et
\ref{lemma-fppf-open}.
Ou bien, combinez
Lemmes \ref{lemma-smooth-syntomic},
\ref{lemma-syntomic-open}.
\end{proof}

\noindent
Le lemme suivant dit localement que tout morphisme lisse est standard lisse. Nous pouvons donc utiliser
les morphismes lisses standards comme {\it modèle local}
pour un morphisme lisse.

\begin{lemma}
\label{lemma-smooth-locally-standard-smooth}
\begin{slogan}
Les morphismes lisses sont des lisses localement standards.
\end{slogan}
Soit $f : X  \to S$ un morphisme de schémas. Soit $x \in X$
un point. Soit $V \subset S$
un voisinage ouvert affine de $f(x)$. Les conditions
suivantes sont équivalentes
\begin{enumerate}
\item Le morphisme $f$ est lisse en $x$.
\item Il existe un $U \subset X$ ouvert affine,
avec $x \in U$ et $f(U) \subset V$ tel que le morphisme
induit $f|_U : U \to V$ est lisse standard.
\end{enumerate}
\end{lemma}

\begin{proof}
Dérivé des définitions
et de l'algèbre, Lemmes \ref{algebra-lemma-standard-smooth}
et \ref{algebra-lemma-smooth-syntomic}.
\end{proof}

\begin{lemma}
\label{lemma-smooth-omega-finite-locally-free}
Soit $f : X \to S$ un morphisme de schémas. Supposons que
$f$ soit fluide.
Alors le module des différentielles $\Omega_{X/S}$ de $X$ sur $S$
est fini localement libre et
$$
\text{rang}_x(\Omega_{X/S}) = \dim_x(X_{f(x)})
$$
pour chaque $x \in X$.
\end{lemma}

\begin{proof}
L'instruction est locale sur $X$
et $S$. D'après le Lemme \ref{lemma-smooth-locally-standard-smooth}
ci-dessus, nous pouvons supposer que $f$ est un morphisme
lisse standard d'affines.
Dans ce cas le résultat découle de l'Algèbre,
Lemme \ref{algebra-lemma-standard-smooth} (et de la définition d'une intersection
complète
relative globale, voir Algèbre, Définition \ref{algebra-definition-relative-global-complete-intersection}).
\end{proof}

\noindent
Lemme \ref{lemma-smooth-omega-finite-locally-free} dit que la définition
suivante a du sens.

\begin{definition}
\label{definition-smooth-relative-dimension}
Soit $d \geq 0$ un nombre entier. On dit qu'un morphisme de schémas $f : X \to S$ est {\it
lisse de dimension relative $d$} si $f$ est
lisse et $\Omega_{X/S}$ est fini localement libre de rang constant $d$.
\end{definition}

\noindent
Autrement dit, $f$ est lisse et les fibres non vides sont équidimensionnelles
de dimension $d$. Par le Lemme \ref{lemma-smooth-at-point} ci-dessous cela revient
également à exiger : (a) $f$ est localement de présentation finie, (b) $f$
est plat, (c) toutes les fibres non vides équidimensionnelles de dimension
$d$, et (d) $\Omega_{X/S}$ fini localement libre de rang $d$.
Il ne suffit pas de simplement supposer que $f$ est plat, de présentation
finie, et que $\Omega_{X/S}$ est fini localement libre de rang
$d$. Un contre-exemple est donné par $\Spec(\mathbf{F}_p[t]) \to \Spec(\mathbf{F}_p[t^p])$.

\medskip\noindent Voici
un critère différentiel de régularité en un point. Il existe de nombreuses
variantes de ce résultat, qui peuvent
toutes être utiles à un moment donné. Nous les ajouterons simplement ici
si nécessaire.

\begin{lemma}
\label{lemma-smooth-at-point}
Soit $f : X \to S$ un morphisme de schémas. Soit $x \in X$.
Posons $s = f(x)$.
Supposons que $f$
soit localement de présentation finie. Les conditions
suivantes sont équivalentes :
\begin{enumerate}
\item Le morphisme $f$ est lisse en $x$.
\item L'homomorphisme d'anneaux locaux $\mathcal{O}_{S, s} \to \mathcal{O}_{X, x}$ est plat
et $X_s \to \Spec(\kappa(s))$ est lisse à $x$.
\item L'homomorphisme local d'anneaux $\mathcal{O}_{S, s} \to \mathcal{O}_{X, x}$ est plat
et le $\mathcal{O}_{X, x}$-module $\Omega_{X/S, x}$ peut être
généré par au plus $\dim_x(X_{f(x)})$ éléments.
\item L'homomorphisme d'anneaux locaux $\mathcal{O}_{S, s} \to \mathcal{O}_{X, x}$ est plat et
l'espace vectoriel $\kappa(x)$-
$$
\Omega_{X_s/s, x} \otimes_{\mathcal{O}_{X_s, x}} \kappa(x) =
\Omega_{X/S, x} \otimes_{\mathcal{O}_{X, x}} \kappa(x)
$$
peut être généré par au plus $\dim_x(X_{f(x)})$ éléments.
\item Il existe des ouverts affines $U \subset X$,
et $V \subset S$ telles que $x \in U$, $f(U) \subset V$ et le
morphisme induit $f|_U : U \to V$ est lisse standard.
\item Il existe des ouverts affines $\Spec(A) = U \subset X$
et $\Spec(R) = V \subset S$ avec $x \in U$ correspondant
à $\mathfrak q \subset A$, et $f(U) \subset V$ telles qu'il
existe une présentation
$$
A = R[x_1, \ldots, x_n]/(f_1, \ldots, f_c)
$$
avec
$$
g =
\det
\left(
\begin{matrix}
\partial f_1/\partial x_1 &
\partial f_2/\partial x_1 &
\ldots &
\partial f_c/\partial x_1 \\
\partial f_1/\partial x_2 &
\partial f_2/\partial x_2 &
\ldots &
\partial f_c/\partial x_2 \\
\ldots & \ldots & \ldots & \ldots \\
\partial f_1/\partial x_c &
\partial f_2/\partial x_c &
\ldots &
\partial f_c/\partial x_c
\end{matrix}
\right)
$$
mappant à un élément de $A$ pas dans $\mathfrak q$.
\end{enumerate}
\end{lemma}

\begin{proof}
Notez que si $f$ est lisse à $x$, alors
nous voyons d'après le Lemme \ref{lemma-smooth-locally-standard-smooth} que (5) est valable et (6)
est une version légèrement affaiblie de (5). De plus, si $f$
est lisse, l'homomorphisme d'anneaux
$\mathcal{O}_{S, s} \to \mathcal{O}_{X, x}$ est plat (voir Lemme \ref{lemma-smooth-flat}) et que
$\Omega_{X/S}$ est fini localement
libre de rang égal à $\dim_x(X_s)$ (voir Lemme \ref{lemma-smooth-omega-finite-locally-free}). Ainsi
(1) implique (3) et (4). Par le Lemme \ref{lemma-base-change-smooth} on voit
aussi que (1) implique (2).

\medskip\noindent
Par le Lemme \ref{lemma-base-change-differentials} le module des différentielles
$\Omega_{X_s/s}$ de la fibre $X_s$ sur $\kappa(s)$
est le retrait du module des différentielles $\Omega_{X/S}$
de $X$ sur $S$. D’où l’égalité affichée
dans la partie (4) du lemme. Par le Lemme \ref{lemma-finite-type-differentials} ces modules
sont de type fini. Par conséquent, le nombre minimal de
générateurs des modules
$\Omega_{X/S, x}$ et $\Omega_{X_s/s, x}$ est le même et égal à la dimension
de cet espace vectoriel $\kappa(x)$ par le Lemme de Nakayama
(Algèbre, Lemme \ref{algebra-lemma-NAK}). Cela montre notamment que (3) et
(4) sont équivalents.

\medskip\noindent
Algèbre, Lemme \ref{algebra-lemma-flat-fibre-smooth} montre que (2)
implique
(1). Algèbre, Lemme \ref{algebra-lemma-characterize-smooth-over-field} montre que (3)
et (4) impliquent (2). Enfin, (6) implique
(5) voir par exemple Algèbre, Exemple \ref{algebra-example-make-standard-smooth}
et (5) implique (1) par Algèbre, Lemme \ref{algebra-lemma-standard-smooth}.
\end{proof}

\begin{lemma}
\label{lemma-set-points-where-fibres-smooth}
Soit
$$
\xymatrix{
X' \ar[r]_{g'} \ar[d]_{f'} & X \ar[d]^f \\
S' \ar[r]^g & S
}
$$
un diagramme cartésien de schémas. Soit $W \subset X$, resp.\ $W' \subset X'$ le sous-schéma
ouvert des points où $f$, resp.\ $f'$ est
lisse. Alors $W' = (g')^{-1}(W)$ si
\begin{enumerate}
\item $f$ est plat et localement de présentation finie, ou
\item $f$ est localement de présentation finie et $g$ est plat.
\end{enumerate}
\end{lemma}

\begin{proof}
Supposons d'abord que $f$ soit localement de type fini. Considérons l'ensemble
$$
T = \{x \in X \mid X_{f(x)}\text{ est lisse sur }\kappa(f(x))\text{ en }x\}
$$
et l'ensemble correspondant $T' \subset X'$ pour $f'$. Nous affirmons alors que
$T' = (g')^{-1}(T)$. En effet, soit $s' \in S'$ un point, et posons
$s = g(s')$. On a alors
$$
X'_{s'} =
\Spec(\kappa(s')) \times_{\Spec(\kappa(s))} X_s
$$
Autrement dit les fibres du changement de base sont les changements
de base des fibres. L'affirmation
est donc équivalente à l'algèbre, Lemme \ref{algebra-lemma-smooth-field-change-local}.

\medskip\noindent
Ainsi le cas (1) en résulte, car $T$ est alors l'ensemble (ouvert)
des points où $f$ est lisse, d'après le lemme \ref{lemma-smooth-at-point}.

\medskip\noindent
Dans le cas (2), soit $x' \in W'$. Alors $g'$
est plat en $x'$ (Lemme \ref{lemma-base-change-module-flat})
et $g \circ f'$ est plat en $x'$ (Lemme \ref{lemma-composition-module-flat}). Il s'ensuit
que $f$ est plat en $x = g'(x')$
par le lemme \ref{lemma-flat-permanence}. D'autre part, puisque $x' \in T'$
(Lemme \ref{lemma-base-change-smooth}) on voit que $x \in T$.
Par conséquent, $f$ est lisse en $x$
par le lemme \ref{lemma-smooth-at-point}.
\end{proof}

\noindent
Voici un lemme qui utilise en fait la disparition de $H^{-1}$ du complexe
cotangent naïf pour un homomorphisme d'anneaux lisse.

\begin{lemma}
\label{lemma-triangle-differentials-smooth}
Soient $f : X \to Y$, $g : Y \to S$ des morphismes de schémas. Supposons que
$f$ soit fluide. Alors
$$
0 \to f^*\Omega_{Y/S} \to \Omega_{X/S} \to \Omega_{X/Y} \to 0
$$
(voir Lemme \ref{lemma-triangle-differentials}) est court et exact.
\end{lemma}

\begin{proof}
La version algébrique de ce lemme est la suivante : Étant donnés
les homomorphismes d'anneaux $A \to B \to C$ avec $B \to C$ lisses, alors la séquence
$$
0 \to C \otimes_B \Omega_{B/A} \to \Omega_{C/A} \to \Omega_{C/B} \to 0
$$
d'Algèbre,
Lemme \ref{algebra-lemma-exact-sequence-differentials} est exacte.
C'est
de l'algèbre, Lemme \ref{algebra-lemma-triangle-differentials-smooth}.
\end{proof}

\begin{lemma}
\label{lemma-differentials-relative-immersion-smooth}
Soit $i : Z \to X$ une immersion de schémas sur $S$. Supposons
que $Z$ soit lisse sur $S$. Alors
la suite exacte canonique
$$
0 \to \mathcal{C}_{Z/X} \to i^*\Omega_{X/S} \to \Omega_{Z/S} \to 0
$$
du
lemme \ref{lemma-differentials-relative-immersion} est exacte
courte.
\end{lemma}

\begin{proof}
La version algébrique de ce lemme est la suivante : Étant donné
les homomorphismes d'anneaux $A \to B \to C$ avec $A \to C$ lisse et $B \to C$
surjectif avec noyau $J$, alors la séquence
$$
0 \to J/J^2 \to C \otimes_B \Omega_{B/A} \to \Omega_{C/A} \to 0
$$
d'Algèbre,
Lemme \ref{algebra-lemma-differential-seq} est exacte.
C'est
de l'algèbre, Lemme \ref{algebra-lemma-differential-seq-smooth}.
\end{proof}

\begin{lemma}
\label{lemma-two-immersions-smooth}
Soit
$$
\xymatrix{
Z \ar[r]_i \ar[rd]_j & X \ar[d] \\
& Y
}
$$
un diagramme commutatif de schémas où $i$ et $j$ sont des immersions
et $X \to Y$ est lisse.
Alors la séquence canonique exacte
$$
0 \to  \mathcal{C}_{Z/Y} \to \mathcal{C}_{Z/X} \to i^*\Omega_{X/Y} \to 0
$$
de
Lemme \ref{lemma-two-immersions}
est exacte.
\end{lemma}

\begin{proof}
La version algébrique de ce lemme est la suivante : Étant donnés
les homomorphismes d'anneaux $A \to B \to C$ avec $A \to C$ surjectif et $A \to B$
lisse, alors la séquence
$$
0 \to I/I^2 \to J/J^2 \to C \otimes_B \Omega_{B/A} \to 0
$$
d'Algèbre,
Lemme \ref{algebra-lemma-application-NL} est
exacte. C'est
de l'algèbre, Lemme \ref{algebra-lemma-application-NL-smooth}.
\end{proof}

\begin{lemma}
\label{lemma-smooth-permanence}
Soit
$$
\xymatrix{
X \ar[rr]_f \ar[rd]_p & &
Y \ar[dl]^q \\
& S
}
$$
un diagramme commutatif de morphismes de schémas. Supposons que
\begin{enumerate}
\item $f$ est surjectif, et lisse,
\item $p$ est lisse, et
\item $q$ est localement de présentation finie\footnote{En
fait, cela est
impliqué par (1) et (2), voir Descente, Lemme \ref{descent-lemma-flat-finitely-presented-permanence}. De plus,
il suffit de supposer que $f$ est surjectif, plat et
localement de fini présentation,
voir Descente, Lemme \ref{descent-lemma-smooth-permanence}.}.
\end{enumerate}
Alors $q$ est lisse.
\end{lemma}

\begin{proof}
Par le Lemme \ref{lemma-flat-permanence} on voit que $q$ est plat. Choisissez
un point $y \in Y$. Choisissez un point $x \in X$ envoyé
sur $y$. Supposons que $f$ ait une dimension relative
$a$ à $x$ et que $p$ ait une dimension relative $b$
à $x$. Par le Lemme \ref{lemma-smooth-omega-finite-locally-free} cela signifie que $\Omega_{X/S, x}$
est libre de rang $b$ et $\Omega_{X/Y, x}$ est
libre de rang $a$. Par la suite exacte courte
du lemme \ref{lemma-triangle-differentials-smooth} cela signifie que $(f^*\Omega_{Y/S})_x$
est libre de rang $b - a$. D'après le Lemme de
Nakayama, cela implique que $\Omega_{Y/S, y}$ peut être engendré
par $b - a$ éléments. De plus, le lemme \ref{lemma-dimension-fibre-at-a-point-additive} montre
que $\dim_y(Y_s) = b - a$. Nous concluons donc que
$Y \to S$ est lisse en $y$ d'après le Lemme \ref{lemma-smooth-at-point}, partie (2).
\end{proof}

\noindent
Dans la situation du lemme suivant l'image de $\sigma$
est localement sur $X$ découpée par
une séquence régulière, voir Diviseurs, Lemme \ref{divisors-lemma-section-smooth-regular-immersion}.

\begin{lemma}
\label{lemma-section-smooth-morphism}
Soit $f : X \to S$ un morphisme de schémas. Soit
$\sigma : S \to X$ une section de $f$.
Soit $s \in S$ un point tel que $f$ soit lisse en $x = \sigma(s)$.
Alors il existe des voisinages affines
ouverts $\Spec(A) = U \subset S$ de $s$ et $\Spec(B) = V \subset X$ de $x$
tels que
\begin{enumerate}
\item $f(V) \subset U$ et $\sigma(U) \subset V$,
\item avec $I = \Ker(\sigma^\# : B \to A)$ le module $I/I^2$ est un module
$A$ libre, et
\item $B^\wedge \cong A[[x_1, \ldots, x_d]]$ comme $A$-algèbres où $B^\wedge$ désigne
l'achèvement de $B$ par rapport à $I$.
\end{enumerate}
\end{lemma}

\begin{proof}
Choisissez un ouvert affine $U \subset S$ contenant
$s$ Choisissez un ouvert affine $V \subset f^{-1}(U)$ contenant
$x$. Choisissez un $U' \subset \sigma^{-1}(V)$ ouvert affine contenant
$s$. Notez que $V' = f^{-1}(U') \cap V$ est affine car il est égal au produit
fibré $V' = U' \times_U V$. Alors $U'$ et $V'$ satisfont (1).
Écrivez $U' = \Spec(A')$ et $V' = \Spec(B')$. Par algèbre,
Lemme \ref{algebra-lemma-section-smooth} le module $I'/(I')^2$ est fini
localement libre en tant que $A'$-module. Ainsi, après
avoir remplacé $U'$ par un ouvert affine plus petit $U'' \subset U'$
et $V'$ par $V'' = V' \cap f^{-1}(U'')$ nous obtenons la situation où $I''/(I'')^2$
est libre, c'est-à-dire que (2) est valable. Dans ce cas (3) est également
valable pour l'algèbre, Lemme \ref{algebra-lemma-section-smooth}.
\end{proof}

\noindent
La dimension d'un schéma $X$ en un
point $x$ (Propriétés, Définition \ref{properties-definition-dimension}) est juste
la dimension de $X$ en $x$ comme espace
topologique, voir Topologie, Définition \ref{topology-definition-Krull}.
Ce n'est pas la dimension de l'anneau local $\mathcal{O}_{X,x}$, en général.

\begin{lemma}
\label{lemma-smoothness-dimension}
Soit $f : X \to Y$ un morphisme lisse de schémas localement noethériens. Pour tout point $x$
de $X$, d'image $y$ dans $Y$,
$$
\dim_x(X) = \dim_y(Y) + \dim_x(X_y),
$$
où $X_y$ désigne la fibre en $y$.
\end{lemma}

\begin{proof}
Après avoir remplacé $X$ par un voisinage
ouvert de $x$, il existe un nombre
naturel $d$ tel que toutes les fibres de
$X \to Y$ ont la dimension $d$ en tout
point, voir Lemme \ref{lemma-smooth-omega-finite-locally-free}. Alors $f$
est plat (Lemme \ref{lemma-smooth-flat}), localement de type fini (Lemme \ref{lemma-smooth-locally-finite-presentation}),
et de dimension relative $d$. Le résultat
découle donc du lemme \ref{lemma-rel-dimension-dimension}.
\end{proof}












\section{Morphismes non ramifiés}
\label{section-unramified}

\noindent
Nous discutons brièvement des morphismes non ramifiés avant la classe (peut-être) plus intéressante
des morphismes étales. Rappelons qu'un homomorphisme d'anneaux $R \to A$
est {\it non ramifié} s'il est de type fini et $\Omega_{A/R} = 0$
(c'est la définition de \cite{Henselian}). Un homomorphisme d'anneaux $R \to A$ est appelé
{\it G-unramified} s'il est de présentation finie et $\Omega_{A/R} = 0$
(c'est la
définition de \cite{EGA}). Voir Algèbre, Définition \ref{algebra-definition-unramified}.

\begin{definition}
\label{definition-unramified}
Soit $f : X \to S$ un morphisme de schémas.
\begin{enumerate}
\item On dit que $f$ est {\it non ramifié
en $x \in X$} s'il existe un voisinage ouvert affine $\Spec(A) = U \subset X$
de $x$ et ouvert affine $\Spec(R) = V \subset S$ avec $f(U) \subset V$
tel que l'homomorphisme d'anneaux induit $R \to A$
ne soit pas ramifié.
\item On dit que $f$ est {\it G-non ramifié
en $x \in X$} s'il existe un voisinage ouvert affine $\Spec(A) = U \subset X$
de $x$ et ouvert affine $\Spec(R) = V \subset S$ avec
$f(U) \subset V$ tel que l'homomorphisme d'anneaux induit $R \to A$
soit G-non ramifié.
\item On dit que $f$ est {\it non ramifié} s'il est non ramifié
en tout point de $X$.
\item On dit que $f$ est {\it G-non ramifié} s'il est G-non ramifié
en tout point de $X$.
\end{enumerate}
\end{definition}

\noindent
Notons qu'un G-morphisme non ramifié est non ramifié. Ainsi, tout résultat pour
les morphismes non ramifiés implique le résultat correspondant pour les G-morphismes
non ramifiés. De plus, si $S$ est localement noethérien alors il n'y a pas
de différence entre G-unramifié et morphismes non ramifiés,
voir Lemme \ref{lemma-noetherian-unramified}. Une caractéristique
intéressante de cette définition est que l'ensemble des points où un morphisme
est non ramifié (resp.\ G-non ramifié) est automatiquement ouvert.

\begin{lemma}
\label{lemma-unramified-omega-zero}
Soit $f : X \to S$ un morphisme de schémas. Ensuite,
\begin{enumerate}
\item $f$ est non ramifié si et seulement si $f$ est localement de
type fini et $\Omega_{X/S} = 0$, et
\item $f$ est G-unramifié si et seulement si $f$ est localement de présentation
finie et $\Omega_{X/S} = 0$.
\end{enumerate}
\end{lemma}

\begin{proof}
Par définition un homomorphisme d'anneaux $R \to A$ est unramifié (resp.\ G-unramifié)
si et seulement si il est de type fini (resp.\ présentation
finie) et $\Omega_{A/R} = 0$.
Le lemme découle donc directement des définitions et du Lemme \ref{lemma-differentials-affine}.
\end{proof}

\noindent
Notez qu'il n'y a pas d'hypothèse de séparation ou de quasi-compacité dans
la définition. La question de la non-ramification est donc de nature locale
sur la source. Voici le résultat précis.

\begin{lemma}
\label{lemma-unramified-characterize}
Soit $f : X \to S$ un morphisme de schémas. Les conditions
suivantes sont équivalentes
\begin{enumerate}
\item Le morphisme $f$ est non ramifié (resp.\ G-non ramifié).
\item Pour tout ouvert affine $U \subset X$, $V \subset S$
avec $f(U) \subset V$ l'homomorphisme
d'anneaux $\mathcal{O}_S(V) \to \mathcal{O}_X(U)$ est non ramifié (resp.\ G-unramifié).
\item Il existe un recouvrement ouvert $S = \bigcup_{j \in J} V_j$ et des
recouvrements ouverts $f^{-1}(V_j) = \bigcup_{i \in I_j} U_i$ tels que chacun des
morphismes $U_i \to V_j$, $j\in J, i\in I_j$ soit non ramifié (resp.\ G-non
ramifié).
\item Il existe un recouvrement ouvert affine $S = \bigcup_{j \in J} V_j$
et ouvert affine revêtements $f^{-1}(V_j) = \bigcup_{i \in I_j} U_i$ tel que l'homomorphisme
d'anneaux $\mathcal{O}_S(V_j) \to \mathcal{O}_X(U_i)$ soit non ramifié (resp.\ G-unramifié),
pour tout $j\in J, i\in I_j$.
\end{enumerate}
De plus, si $f$ est non ramifié (resp.\ G-unramifié)
alors pour tout sous-schémas ouverts $U \subset X$, $V \subset S$ avec $f(U) \subset V$
la restriction $f|_U : U \to V$ est non ramifiée (resp.\ G-non ramifié).
\end{lemma}

\begin{proof}
Cela résulte du lemme \ref{lemma-locally-P} si l'on montre que la propriété
``$R \to A$ est non ramifié''
est locale. On vérifie les conditions (a),
(b) et (c) de la Définition
\ref{definition-property-local}. Ces propriétés sont
prouvées dans Algèbre, Lemme \ref{algebra-lemma-unramified}.
\end{proof}

\begin{lemma}
\label{lemma-composition-unramified}
La composition de deux morphismes non ramifiés est non ramifiée. Il en va
de même pour les G-morphismes non ramifiés.
\end{lemma}

\begin{proof}
La démonstration du lemme \ref{lemma-unramified-characterize} montre qu'être
non ramifié (resp.\ G-non ramifié) est
une propriété locale des homomorphismes
d'anneaux. D'où le premier énoncé du lemme
découle du lemme \ref{lemma-composition-type-P} combiné au fait
qu'être non ramifié (resp.\ G-non ramifié)
est une propriété des homomorphismes d'anneaux stable sous composition,
voir Algèbre, Lemme \ref{algebra-lemma-unramified}.
\end{proof}

\begin{lemma}
\label{lemma-base-change-unramified}
Le changement de base d'un morphisme non ramifié est non ramifié. Il en
va de même pour les G-morphismes non ramifiés.
\end{lemma}

\begin{proof}
La démonstration du lemme \ref{lemma-unramified-characterize} montre qu'être non
ramifié (resp.\ G-non ramifié) est une
propriété locale des homomorphismes d'anneaux. D'où le lemme
découle du lemme \ref{lemma-base-change-type-P} combiné au fait
qu'être non ramifié (resp.\ G-non ramifié) est
une propriété des homomorphismes d'anneaux stable par changement
de base, voir Algèbre, Lemme \ref{algebra-lemma-unramified}.
\end{proof}

\begin{lemma}
\label{lemma-noetherian-unramified}
Soit $f : X \to S$ un morphisme de schémas. Supposons que $S$ soit localement noethérien. Alors $f$
est non ramifié si et seulement si $f$ est G-non ramifié.
\end{lemma}

\begin{proof}
Découle des définitions
et du lemme \ref{lemma-noetherian-finite-type-finite-presentation}.
\end{proof}


\begin{lemma}
\label{lemma-open-immersion-unramified}
Toute immersion ouverte est G-non ramifiée.
\end{lemma}

\begin{proof}
Ceci est vrai car une immersion ouverte est un isomorphisme local.
\end{proof}

\begin{lemma}
\label{lemma-closed-immersion-unramified}
Une immersion fermée $i : Z \to X$ n'est pas
ramifiée. Il est G-non ramifié si et seulement si le faisceau
quasi-cohérent d'idéaux associé $\mathcal{I} = \Ker(\mathcal{O}_X \to i_*\mathcal{O}_Z)$ est de type
fini (en tant que $\mathcal{O}_X$-module).
\end{lemma}

\begin{proof}
Cela résulte du Lemme \ref{lemma-closed-immersion-finite-presentation} et d'Algèbre,
Lemme \ref{algebra-lemma-unramified}.
\end{proof}

\begin{lemma}
\label{lemma-unramified-locally-finite-type}
An morphisme non ramifié est localement de type fini. Un
G-morphisme non ramifié est localement de présentation finie.
\end{lemma}

\begin{proof}
Un homomorphisme d'anneaux non ramifié est de type fini par définition.
Un homomorphisme d'anneaux G-non ramifié est de présentation finie par définition.
\end{proof}

\begin{lemma}
\label{lemma-unramified-quasi-finite}
Soit $f : X \to S$ un morphisme de schémas. Si $f$
est non ramifié en $x$ alors $f$ est quasi-fini en $x$.
En particulier, un morphisme non ramifié est localement quasi-fini.
\end{lemma}

\begin{proof}
Voir Algèbre, Lemme \ref{algebra-lemma-unramified-quasi-finite}.
\end{proof}

\begin{lemma}
\label{lemma-unramified-over-field}
Fibres de morphismes non ramifiés.
\begin{enumerate}
\item Soit $X$ un schéma sur un corps $k$.
Le morphisme de structure $X \to \Spec(k)$ est non ramifié si et seulement
si $X$ est une union disjointe de spectres d'extensions de corps finis
séparables de $k$.
\item Si $f : X \to S$ est un morphisme non ramifié alors pour tout $s \in S$ la fibre
$X_s$ est une union disjointe de spectres d'extensions de corps finis
séparables de $\kappa(s)$.
\end{enumerate}
\end{lemma}

\begin{proof}
La partie (2) découle de la partie
(1) et Lemme \ref{lemma-base-change-unramified}. Démontrons
la partie (1). Nous utilisons
d'abord l'algèbre, Lemme \ref{algebra-lemma-characterize-unramified}. Ce lemme implique que
si $X$ est une union disjointe de spectres d'extensions de corps
finis séparables de $k$ alors $X \to \Spec(k)$ est non ramifié.
À l’inverse, supposons que $X \to \Spec(k)$ ne soit
pas ramifié. Par Algèbre, Lemme \ref{algebra-lemma-unramified-at-prime} pour tout $x \in X$ l'extension
du corps résiduel $\kappa(x)/k$ est fini
séparable. Puisque $X \to \Spec(k)$ est localement quasi-fini
(Lemme \ref{lemma-unramified-quasi-finite}), nous voyons que tous les
points de $X$ sont des points fermés isolés, voir
Lemme \ref{lemma-quasi-finite-at-point-characterize}. Ainsi $X$ est un espace
discret, en particulier l'union disjointe des spectres
de ses anneaux locaux. Par algèbre,
Lemme \ref{algebra-lemma-unramified-at-prime} encore une fois, ces anneaux locaux sont
des corps, et on obtient le résultat.
\end{proof}

\noindent
Le lemme suivant caractérise les morphismes non ramifiés comme des
morphismes localement de type fini à fibres non ramifiées.

\begin{lemma}
\label{lemma-unramified-etale-fibres}
Soit $f : X \to S$ un morphisme de schémas.
\begin{enumerate}
\item Si $f$ est non ramifié, alors pour tout $x \in X$ l'extension
du corps $\kappa(x)/\kappa(f(x))$ est finie séparable.
\item Si $f$ est localement de type fini, et pour chaque
$s \in S$ la fibre $X_s$ est une union disjointe de spectres d'extensions de corps finis
séparables de $\kappa(s)$ alors $f$ est non ramifié.
\item Si $f$ est localement de présentation finie, et pour chaque
$s \in S$ la fibre $X_s$ est une union disjointe de spectres d'extensions de corps
finis séparables de $\kappa(s)$ alors $f$ est G-non ramifié.
\end{enumerate}
\end{lemma}

\begin{proof}
Cela résulte d'Algèbre,
les Lemmes \ref{algebra-lemma-unramified-at-prime}
et \ref{algebra-lemma-characterize-unramified}.
\end{proof}

\noindent
Voici une caractérisation des morphismes non ramifiés en termes de diagonale
du morphisme.

\begin{lemma}
\label{lemma-diagonal-unramified-morphism}
Soit $f : X \to S$ un morphisme.
\begin{enumerate}
\item Si $f$ est non ramifié, alors le morphisme
diagonal $\Delta : X \to X \times_S X$ est une immersion ouverte.
\item Si $f$ est localement de type fini et
$\Delta$ est une immersion ouverte, alors $f$ est non ramifié.
\item Si $f$ est localement de présentation finie et $\Delta$ est une immersion
ouverte, alors $f$ est G-non ramifié.
\end{enumerate}
\end{lemma}

\begin{proof}
La première instruction découle
de l'algèbre, Lemme \ref{algebra-lemma-diagonal-unramified}. La deuxième
affirmation vient du fait que $\Omega_{X/S}$ est le
faisceau conormal du morphisme diagonal
(Lemme \ref{lemma-differentials-diagonal}) et donc clairement
nul si $\Delta$ est une immersion ouverte.
\end{proof}

\begin{lemma}
\label{lemma-unramified-at-point}
Soit $f : X \to S$ un morphisme de schémas. Soit $x \in X$.
Posons $s = f(x)$.
Supposons que $f$
est localement de type fini (resp.\ localement de présentation finie). Les conditions
suivantes sont équivalentes :
\begin{enumerate}
\item Le morphisme $f$ est non ramifié (resp.\ G-non ramifié) en $x$.
\item La fibre $X_s$ est non ramifiée sur $\kappa(s)$ à $x$.
\item Le module $\mathcal{O}_{X, x}$ $\Omega_{X/S, x}$ est nul.
\item Le module $\mathcal{O}_{X_s, x}$ $\Omega_{X_s/s, x}$ est nul.
\item L'espace vectoriel $\kappa(x)$
$$
\Omega_{X_s/s, x} \otimes_{\mathcal{O}_{X_s, x}} \kappa(x) =
\Omega_{X/S, x} \otimes_{\mathcal{O}_{X, x}} \kappa(x)
$$
est nul.
\item Nous avons $\mathfrak m_s\mathcal{O}_{X, x} = \mathfrak m_x$ et l'extension du corps
$\kappa(x)/\kappa(s)$ est séparable de manière
finie.
\end{enumerate}
\end{lemma}

\begin{proof}
Notons que si $f$ est non ramifié en
$x$, alors on voit que $\Omega_{X/S} = 0$ à proximité de
$x$ par les définitions et les résultats sur modules
de différentiels dans la section \ref{section-sheaf-differentials}. Par conséquent (1)
implique (3) et la disparition de l'espace vectoriel de droite
dans (5). Cela implique également
(2) car par le lemme \ref{lemma-base-change-differentials} le module
des différentielles $\Omega_{X_s/s}$ de la fibre $X_s$ sur
$\kappa(s)$ est le retrait du module des différentielles $\Omega_{X/S}$
de $X$ sur $S$. Ce fait sur les modules de différentielles
implique également l'égalité affichée des espaces vectoriels
dans la partie (4). Par le Lemme \ref{lemma-finite-type-differentials}
les modules $\Omega_{X/S, x}$ et $\Omega_{X_s/s, x}$ sont de type fini. D'où les modules
$\Omega_{X/S, x}$ et $\Omega_{X_s/s, x}$ sont nuls si et seulement si l'espace vectoriel
$\kappa(x)$ correspondant dans (4) est nul par le Lemme
de Nakayama (Algèbre,
Lemme \ref{algebra-lemma-NAK}). Cela montre
notamment que (3), (4) et (5) sont équivalents. Le support
de $\Omega_{X/S}$ est terminé dans $X$, voir
Modules, Lemme \ref{modules-lemma-support-finite-type-closed}. L'hypothèse (3) implique que
$x$ n'est pas dans le support. Donc $\Omega_{X/S}$
est nul au voisinage de $x$, ce qui implique
(1). L'équivalence de (1) et (3) appliquée à $X_s \to s$ implique
l'équivalence de (2) et (4).
À ce stade, nous avons vu que (1) -- (5) sont équivalents.

\medskip\noindent
Vous pouvez également utiliser Algèbre, Lemme \ref{algebra-lemma-unramified} pour voir
l'équivalence de (1) -- (5) plus directement.

\medskip\noindent
L'équivalence de (1) et (6)
découle du lemme \ref{lemma-unramified-etale-fibres}.
Cela découle aussi
plus directement de l'Algèbre, de Lemmes
\ref{algebra-lemma-unramified-at-prime} et \ref{algebra-lemma-characterize-unramified}.
\end{proof}

\begin{lemma}
\label{lemma-set-points-where-fibres-unramified}
Soit $f : X \to S$ un morphisme de schémas. Supposons
que $f$ soit localement de type fini. La formation de l'ouvert
\begin{align*}
T
& =
\{x \in X \mid X_{f(x)}\text{ est non ramifié sur }\kappa(f(x))\text{ en }x\} \\
& =
\{x \in X \mid X\text{ est non ramifié sur }S\text{ en }x\}
\end{align*}
commute à tout changement de base :
pour tout morphisme $g : S' \to S$, considérons
le changement de base $f' : X' \to S'$ de $f$
et la projection $g' : X' \to X$. Alors l'ensemble correspondant
$T'$ pour le morphisme $f'$ est égal à $T' = (g')^{-1}(T)$. Si $f$
est supposé localement de présentation finie alors il en va de même pour
l'ensemble ouvert de points où $f$ est G-non ramifié.
\end{lemma}

\begin{proof}
Soit $s' \in S'$ un point et soit $s = g(s')$. On a alors
$$
X'_{s'} =
\Spec(\kappa(s')) \times_{\Spec(\kappa(s))} X_s
$$
Autrement dit les fibres du changement de base sont les changements de base
des fibres. En particulier
$$
\Omega_{X_s/s, x} \otimes_{\mathcal{O}_{X_s, x}} \kappa(x')
=
\Omega_{X'_{s'}/s', x'} \otimes_{\mathcal{O}_{X'_{s'}, x'}} \kappa(x')
$$
voir
Lemme \ref{lemma-base-change-differentials}. D'où $x' \in T'$ si
et seulement si $x \in T$ par le lemme
\ref{lemma-unramified-at-point}. La deuxième partie
découle de la première car dans ce cas $T$
est l'ensemble (ouvert) de points où $f$ est G-non
ramifié d'après le Lemme \ref{lemma-unramified-at-point}.
\end{proof}

\begin{lemma}
\label{lemma-unramified-permanence}
Soit $f : X \to Y$ un morphisme de schémas sur $S$.
\begin{enumerate}
\item Si $X$ est non ramifié sur $S$, alors $f$ est non ramifié.
\item Si $X$ est G-non ramifié sur $S$ et que $Y$ est localement de type fini sur
$S$, alors $f$ est G-non ramifié.
\end{enumerate}
\end{lemma}

\begin{proof}
Supposons que $X$ est non ramifié
sur $S$. Par le Lemme \ref{lemma-permanence-finite-type} on voit que $f$
est localement de type fini.
Par hypothèse, nous avons $\Omega_{X/S} = 0$. D'où $\Omega_{X/Y} = 0$
du lemme \ref{lemma-triangle-differentials}. Ainsi, $f$ est non ramifié. Si
$X$ est G-non ramifié sur $S$ et $Y$ est localement
de type fini sur $S$, alors
par le lemme \ref{lemma-finite-presentation-permanence} on voit que $f$
est localement de présentation finie et Nous concluons que $f$
est G-non ramifié.
\end{proof}

\begin{lemma}
\label{lemma-value-at-one-point}
Soit $S$ un schéma.
Soient $X$, $Y$ des schémas sur
$S$. Soit $f, g : X \to Y$ des morphismes sur $S$. Soit $x \in X$.
Supposons que
\begin{enumerate}
\item le morphisme de structure $Y \to S$ est non ramifié,
\item $f(x) = g(x)$ dans $Y$, disons $y = f(x) = g(x)$, et
\item les applications induites $f^\sharp, g^\sharp : \kappa(y) \to \kappa(x)$ sont
égales.
\end{enumerate}
Alors il existe un voisinage ouvert de $x$ dans $X$ sur lequel $f$
et $g$ sont égaux.
\end{lemma}

\begin{proof}
Considérons le morphisme $(f, g) : X \to Y \times_S Y$. Par l'hypothèse (1) et Lemme \ref{lemma-diagonal-unramified-morphism}
l'image réciproque de $\Delta_{Y/S}(Y)$ est ouverte
dans $X$. Et les hypothèses (2) et (3) impliquent
que $x$ se trouve dans ce sous-ensemble ouvert.
\end{proof}










\section{Morphismes étales}
\label{section-etale}

\noindent
La topologie Zariski d'un schéma est une topologie très grossière.
Cela est particulièrement clair lorsque l'on examine les variétés supérieures
à $\mathbf{C}$. Il s’avère que déclarer un morphisme étale comme l’analogue
d’un isomorphisme local en topologie introduit une topologie beaucoup
plus fine. Sur les variétés supérieures à $\mathbf{C}$, cette topologie donne
lieu aux ``corrects'' nombres de Betti lors du calcul de
cohomologie à coefficients finis. Un autre observable est que si $f : X \to Y$ est un
morphisme étale des variétés sur $\mathbf{C}$, et si
$x$ est un point fermé
de $X$, alors $f$ induit un isomorphisme $\mathcal{O}^{\wedge}_{Y, f(x)} \to \mathcal{O}^{\wedge}_{X, x}$ d'anneaux
locaux complets.

\medskip\noindent Dans
cette section, nous commençons notre étude de ces questions. En fait, nous limitons
délibérément notre discussion au minimum puisque nous discuterons ailleurs de résultats
plus intéressants. Rappelons qu'un homomorphisme d'anneaux $R \to A$ est dit {\it
étale} s'il est lisse
et $\Omega_{A/R} = 0$, voir Algèbre, Définition \ref{algebra-definition-etale}.

\begin{definition}
\label{definition-etale}
Soit $f : X \to S$ un morphisme de schémas.
\begin{enumerate}
\item On dit que $f$ est {\it étale
à $x \in X$} s'il existe un voisinage ouvert affine $\Spec(A) = U \subset X$
de $x$ et ouvert affine $\Spec(R) = V \subset S$ avec $f(U) \subset V$
tel que l'homomorphisme d'anneaux induit $R \to A$
soit étale.
\item On dit que $f$ est {\it étale} s'il est étale en tout point de $X$.
\item Un morphisme de schémas affines $f : X \to S$ est appelé {\it
standard étale} si $X \to S$ est isomorphe à
$$
\Spec(R[x]_h/(g)) \to \Spec(R)
$$
où $R \to R[x]_h/(g)$ est un standard étale homomorphisme d'anneaux,
voir Algèbre, Définition \ref{algebra-definition-standard-etale}, c'est-à-dire que $g$
est monique et $g'$ inversible dans $R[x]_h/(g)$.
\end{enumerate}
\end{definition}

\noindent
Un morphisme est étale si et seulement si il est lisse de dimension relative
$0$ (voir Définition \ref{definition-smooth-relative-dimension}). Une caractéristique intéressante
de la définition est que l'ensemble des points où un morphisme
est étale est automatiquement ouvert.

\medskip\noindent Notez
qu'il n'y a pas d'hypothèse de séparation ou de quasi-compacité dans la définition.
La question de l'existence de étale est donc de nature locale
sur la source. Voici le résultat précis.

\begin{lemma}
\label{lemma-etale-characterize}
Soit $f : X \to S$ un morphisme de schémas. Les conditions
suivantes sont équivalentes
\begin{enumerate}
\item Le morphisme $f$ est étale.
\item Pour chaque ouvert affines $U \subset X$, $V \subset S$
avec $f(U) \subset V$ l'homomorphisme
d'anneaux $\mathcal{O}_S(V) \to \mathcal{O}_X(U)$ est étale.
\item Il existe un recouvrement ouvert $S = \bigcup_{j \in J} V_j$ et
des recouvrements ouverts $f^{-1}(V_j) = \bigcup_{i \in I_j} U_i$ tels que chacun
des morphismes $U_i \to V_j$, $j\in J, i\in I_j$ est
étale.
\item Il existe un recouvrement ouvert affine $S = \bigcup_{j \in J} V_j$
et ouvert affine revêtements $f^{-1}(V_j) = \bigcup_{i \in I_j} U_i$ tel que l'homomorphisme
d'anneaux $\mathcal{O}_S(V_j) \to \mathcal{O}_X(U_i)$ est étale,
pour tout $j\in J, i\in I_j$.
\end{enumerate}
De plus, si $f$ est étale
alors pour tout sous-schémas ouverts $U \subset X$, $V \subset S$ avec $f(U) \subset V$
la restriction $f|_U : U \to V$ est étale.
\end{lemma}

\begin{proof}
Cela résulte du lemme \ref{lemma-locally-P} si l'on montre que la propriété
``$R \to A$ est étale''
est locale. On vérifie les conditions (a),
(b) et (c) de la Définition
\ref{definition-property-local}. Tout cela découle de l'algèbre, Lemme \ref{algebra-lemma-etale}.
\end{proof}

\begin{lemma}
\label{lemma-composition-etale}
La composition de deux morphismes qui sont étale est étale.
\end{lemma}

\begin{proof}
Dans la démonstration du lemme \ref{lemma-etale-characterize} nous avons vu
qu'être étale est une propriété locale des homomorphismes
d'anneaux. D'où la première affirmation du
lemme découle du lemme \ref{lemma-composition-type-P} combinée au fait
qu'être étale est une propriété des homomorphismes d'anneaux
qui est stable sous composition,
voir Algèbre, Lemme \ref{algebra-lemma-etale}.
\end{proof}

\begin{lemma}
\label{lemma-base-change-etale}
Le changement de base d'un morphisme qui est étale est étale.
\end{lemma}

\begin{proof}
Dans la démonstration du lemme \ref{lemma-etale-characterize} nous avons
vu qu'être étale est une propriété locale des homomorphismes
d'anneaux. D'où le
lemme découle du lemme \ref{lemma-composition-type-P} combiné au
fait qu'être étale est une propriété des homomorphismes d'anneaux
stable par changement
de base, voir Algèbre, Lemme \ref{algebra-lemma-etale}.
\end{proof}

\begin{lemma}
\label{lemma-etale-smooth-unramified}
Soit $f : X \to S$ un morphisme de schémas. Soit $x \in X$.
Alors $f$ est étale à $x$ si et seulement si $f$ est
lisse et non ramifié à $x$.
\end{lemma}

\begin{proof}
Cela découle immédiatement des définitions.
\end{proof}

\begin{lemma}
\label{lemma-etale-locally-quasi-finite}
Un morphisme étale est localement quasi-fini.
\end{lemma}

\begin{proof}
Par
le lemme \ref{lemma-etale-smooth-unramified}, un morphisme
étale est non ramifié.
Par le Lemme \ref{lemma-unramified-quasi-finite} un morphisme
non ramifié est localement quasi-fini.
\end{proof}

\begin{lemma}
\label{lemma-etale-over-field}
\begin{slogan}
Description des schémas étales sur corps et fibres de morphismes
étales.
\end{slogan}
Fibres de morphismes étales.
\begin{enumerate}
\item Soit $X$ un schéma sur un corps $k$.
Le morphisme de structure $X \to \Spec(k)$ est étale si et seulement
si $X$ est une union disjointe de spectres d'extensions de corps finis
séparables de $k$.
\item Si $f : X \to S$ est un morphisme étale, alors pour tout $s \in S$ la fibre
$X_s$ est une union disjointe de spectres d'extensions de corps finis
séparables de $\kappa(s)$.
\end{enumerate}
\end{lemma}

\begin{proof}
Vous pouvez le déduire du lemme \ref{lemma-unramified-over-field} via Lemme
\ref{lemma-etale-smooth-unramified} ci-dessus. Voici une démonstration
directe.

\medskip\noindent
Nous utiliserons l'algèbre, Lemme \ref{algebra-lemma-etale-over-field}. Il est donc clair que
si $X$ est une union disjointe de spectres d'extensions de corps finis
séparables de $k$ alors $X \to \Spec(k)$ est étale. À
l’inverse, supposons que $X \to \Spec(k)$ soit étale. Alors pour tout $U \subset X$
ouvert affine, nous voyons que $U$ est une union disjointe
finie de spectres d'extensions de corps finis séparables de $k$.
Ainsi tous les points de
$X$ sont des points fermés (voir Lemme \ref{lemma-algebraic-residue-field-extension-closed-point-fibre} par exemple).
Ainsi $X$ est un espace discret et on obtient le résultat.
\end{proof}

\noindent
Le lemme suivant caractérise un morphisme étale comme un morphisme plat
de présentation finie avec des ``fibres étales''.

\begin{lemma}
\label{lemma-etale-flat-etale-fibres}
Soit $f : X \to S$ un morphisme de schémas. Si $f$
est plat, localement de présentation finie, et pour tout $s \in S$ la fibre $X_s$
est une union disjointe de spectres d'extensions finies de corps séparables
de $\kappa(s)$, alors $f$ est étale.
\end{lemma}

\begin{proof}
Vous pouvez le déduire
de l'algèbre, Lemme \ref{algebra-lemma-characterize-etale}. Voici une autre
démonstration.

\medskip\noindent
Par le Lemme \ref{lemma-etale-over-field} une fibre $X_s$ est étale
et donc lisse sur $s$. Par le Lemme \ref{lemma-smooth-flat-smooth-fibres} on
voit que $X \to S$ est lisse.
Par le Lemme \ref{lemma-unramified-etale-fibres} nous voyons
que $f$ est non ramifié. Nous
concluons par le lemme \ref{lemma-etale-smooth-unramified}.
\end{proof}

\begin{lemma}
\label{lemma-open-immersion-etale}
Toute immersion ouverte est étale.
\end{lemma}

\begin{proof}
Ceci est vrai car une immersion ouverte est un isomorphisme local.
\end{proof}

\begin{lemma}
\label{lemma-etale-syntomic}
Un morphisme étale est syntomique.
\end{lemma}

\begin{proof}
Voir Algèbre, Lemme \ref{algebra-lemma-smooth-syntomic} et utiliser qu'un morphisme
étale est le même qu'un morphisme lisse de dimension relative $0$.
\end{proof}

\begin{lemma}
\label{lemma-etale-locally-finite-presentation}
Un morphisme étale est localement de présentation finie.
\end{lemma}

\begin{proof}
Vrai car un étale homomorphisme d'anneaux est de présentation finie par
définition.
\end{proof}

\begin{lemma}
\label{lemma-etale-flat}
Un morphisme étale est plat.
\end{lemma}

\begin{proof}
Cela résulte des lemmes \ref{lemma-syntomic-flat} et \ref{lemma-etale-syntomic}.
\end{proof}

\begin{lemma}
\label{lemma-etale-open}
Un morphisme étale est ouvert.
\end{lemma}

\begin{proof}
Cela résulte des lemmes \ref{lemma-etale-flat},
\ref{lemma-etale-locally-finite-presentation} et
\ref{lemma-fppf-open}.
\end{proof}

\noindent
Le lemme suivant dit localement que tout morphisme étale est un étale standard.
Il s’agit en fait d’un résultat difficile à prouver en toute généralité. Les parties délicates
sont cachées dans le chapitre sur l'algèbre commutative. Ainsi, un morphisme
étale standard est un {\it modèle local} pour un morphisme
étale général.

\begin{lemma}
\label{lemma-etale-locally-standard-etale}
Soit $f : X  \to S$ un morphisme de schémas. Soit $x \in X$
un point. Soit $V \subset S$
un voisinage ouvert affine de $f(x)$. Les conditions
suivantes sont équivalentes
\begin{enumerate}
\item Le morphisme $f$ est étale à $x$.
\item Il existe un $U \subset X$ ouvert affine avec
$x \in U$ et $f(U) \subset V$ tel que le morphisme
induit $f|_U : U \to V$ est standard étale
(voir Définition \ref{definition-etale}).
\end{enumerate}
\end{lemma}

\begin{proof}
Dérivé des définitions
et de l'algèbre, proposition \ref{algebra-proposition-etale-locally-standard}.
\end{proof}

\noindent
Voici un critère différentiel de étalité en un point. Il existe
de nombreuses variantes de ce résultat, qui peuvent
toutes être utiles à un moment donné. Nous les ajouterons simplement ici
si nécessaire.

\begin{lemma}
\label{lemma-etale-at-point}
Soit $f : X \to S$ un morphisme de schémas. Soit $x \in X$.
Posons $s = f(x)$.
Supposons que $f$
soit localement de présentation finie. Les conditions
suivantes sont équivalentes :
\begin{enumerate}
\item Le morphisme $f$ est étale à $x$.
\item L'application anneau local $\mathcal{O}_{S, s} \to \mathcal{O}_{X, x}$ est plate et
$X_s \to \Spec(\kappa(s))$ est étale à $x$.
\item L'homomorphisme local d'anneaux $\mathcal{O}_{S, s} \to \mathcal{O}_{X, x}$ est plate
et $X_s \to \Spec(\kappa(s))$ est non ramifiée en $x$.
\item L'homomorphisme d'anneaux locaux $\mathcal{O}_{S, s} \to \mathcal{O}_{X, x}$ est
plat et le $\mathcal{O}_{X, x}$-module $\Omega_{X/S, x}$ est
nul.
\item L'homomorphisme d'anneaux locaux $\mathcal{O}_{S, s} \to \mathcal{O}_{X, x}$ est plat et
l'espace vectoriel $\kappa(x)$-
$$
\Omega_{X_s/s, x} \otimes_{\mathcal{O}_{X_s, x}} \kappa(x) =
\Omega_{X/S, x} \otimes_{\mathcal{O}_{X, x}} \kappa(x)
$$
est nul.
\item L'homomorphisme local d'anneaux $\mathcal{O}_{S, s} \to \mathcal{O}_{X, x}$
est plat, nous avons $\mathfrak m_s\mathcal{O}_{X, x} = \mathfrak m_x$ et l'extension
corps $\kappa(x)/\kappa(s)$ est fini
séparable.
\item Il existe des ouverts affines $U \subset X$,
et $V \subset S$ telles que $x \in U$, $f(U) \subset V$ et le morphisme
induit $f|_U : U \to V$ est standard lisse de dimension
relative $0$.
\item Il existe des ouverts affines $\Spec(A) = U \subset X$
et $\Spec(R) = V \subset S$ avec $x \in U$ correspondant
à $\mathfrak q \subset A$, et $f(U) \subset V$ telles qu'il
existe une présentation
$$
A = R[x_1, \ldots, x_n]/(f_1, \ldots, f_n)
$$
avec
$$
g =
\det
\left(
\begin{matrix}
\partial f_1/\partial x_1 &
\partial f_2/\partial x_1 &
\ldots &
\partial f_n/\partial x_1 \\
\partial f_1/\partial x_2 &
\partial f_2/\partial x_2 &
\ldots &
\partial f_n/\partial x_2 \\
\ldots & \ldots & \ldots & \ldots \\
\partial f_1/\partial x_n &
\partial f_2/\partial x_n &
\ldots &
\partial f_n/\partial x_n
\end{matrix}
\right)
$$
mappant à un élément de $A$ pas dans $\mathfrak q$.
\item Il existe des ouverts affines $U \subset X$,
et $V \subset S$ telles que $x \in U$, $f(U) \subset V$ et le morphisme
induit $f|_U : U \to V$ est standard étale.
\item Il existe des ouverts affines $\Spec(A) = U \subset X$
et $\Spec(R) = V \subset S$ avec $x \in U$ correspondant
à $\mathfrak q \subset A$, et $f(U) \subset V$ telles qu'il
existe une présentation
$$
A = R[x]_Q/(P) = R[x, 1/Q]/(P)
$$
avec $P, Q \in R[x]$, $P$ monique et $P' = \text{d}P/\text{d}x$ mappage à un élément
de $A$ pas dans $\mathfrak q$.
\end{enumerate}
\end{lemma}

\begin{proof}
Utilisez Lemme \ref{lemma-etale-locally-standard-etale} et les définitions
pour voir que (1) implique toutes les autres
conditions. Pour chacune des conditions (2)
-- (10), combinons les Lemmes \ref{lemma-smooth-at-point} et \ref{lemma-unramified-at-point}
pour voir que (1) est valable en montrant
que $f$ est à la fois lisse et non ramifié
à $x$ et en appliquant le Lemme \ref{lemma-etale-smooth-unramified}. Certains
détails omis.
\end{proof}

\begin{lemma}
\label{lemma-flat-unramified-etale}
Un morphisme est étale en un point si et seulement s'il est plat et G-non ramifié
en ce point. Un morphisme
est étale si et seulement si il est plat et G-non ramifié.
\end{lemma}

\begin{proof}
Cela ressort clairement des Lemmes \ref{lemma-etale-at-point}
et \ref{lemma-unramified-at-point}.
\end{proof}

\begin{lemma}
\label{lemma-set-points-where-fibres-etale}
Soit
$$
\xymatrix{
X' \ar[r]_{g'} \ar[d]_{f'} & X \ar[d]^f \\
S' \ar[r]^g & S
}
$$
un diagramme cartésien de schémas. Soit $W \subset X$, resp.\ $W' \subset X'$ le sous-schéma
ouvert des points où $f$, resp.\ $f'$ est étale.
Alors $W' = (g')^{-1}(W)$ si
\begin{enumerate}
\item $f$ est plat et localement de présentation finie, ou
\item $f$ est localement de présentation finie et $g$ est plat.
\end{enumerate}
\end{lemma}

\begin{proof}
Supposons d'abord que $f$ localement de type fini. Considérons l'ensemble
$$
T = \{x \in X \mid f\text{ est non ramifié en }x\}
$$
et l'ensemble correspondant $T' \subset X'$ pour
$f'$. Puis
$T' = (g')^{-1}(T)$ par le lemme \ref{lemma-set-points-where-fibres-unramified}.

\medskip\noindent Ainsi
le cas (1) suit car dans le cas (1) $T$ est l'ensemble (ouvert) de
points où $f$ est étale par le lemme \ref{lemma-flat-unramified-etale}.

\medskip\noindent
Dans le cas (2), soit $x' \in W'$. Alors $g'$
est plat en $x'$ (Lemme \ref{lemma-base-change-module-flat})
et $g \circ f'$ est plat en $x'$ (Lemme \ref{lemma-composition-module-flat}). Il s'ensuit
que $f$ est plat à $x = g'(x')$
par le lemme \ref{lemma-flat-permanence}. D'autre part, puisque $x' \in T'$
(Lemme \ref{lemma-base-change-smooth}) on voit que $x \in T$.
Ainsi $f$ est étale à $x$
par le lemme \ref{lemma-etale-at-point}.
\end{proof}

\begin{lemma}
\label{lemma-etale-permanence-better}
\begin{reference}
Voir par exemple \cite[Remarque 18.30 (2).]{GWII}.
\end{reference}
Soit $f : X \to Y$ un morphisme de schémas sur $S$. Si $X$ est étale
sur $S$ et que $Y$ est non ramifié sur $S$, alors $f$
est étale.
\end{lemma}

\begin{proof}
Considérons la factorisation $X \to X \times_S Y \to Y$, où la première flèche est donnée par
$\text{id}_X$ et $f$ et la deuxième flèche est la projection. Nous affirmons
que les deux flèches sont étale et donc $f$
est étale du lemme \ref{lemma-composition-etale}. À savoir, la projection est
étale car c'est le changement de base de $X \to S$, voir Lemme \ref{lemma-base-change-etale}.
La première flèche est le changement de base du morphisme
diagonal $Y \to Y \times_S Y$ car le carré
$$
\xymatrix{
X \ar[d] \ar[r] & X \times_S Y \ar[d] \\
Y \ar[r] & Y \times_S Y
}
$$
est cartésien. La diagonale $Y \to Y \times_S Y$ est une immersion ouverte
du lemme \ref{lemma-diagonal-unramified-morphism}. Le changement de base d'une
immersion ouverte est une immersion
ouverte (Schémas, Lemme \ref{schemes-lemma-base-change-immersion}) et une
immersion ouverte est étale
(Lemme \ref{lemma-open-immersion-etale}).
\end{proof}

\begin{lemma}
\label{lemma-etale-permanence}
\begin{slogan}
Loi d'annulation pour les morphismes étales
\end{slogan}
Soit $f : X \to Y$ un morphisme de schémas sur $S$. Si $X$ et $Y$
sont étales sur $S$, alors $f$
est étale.
\end{lemma}

\begin{proof}
Puisque $Y \to S$ est non ramifié d'après le lemme \ref{lemma-etale-smooth-unramified},
nous pouvons appliquer le Lemme \ref{lemma-etale-permanence-better}.
Voir aussi Algèbre, Lemme \ref{algebra-lemma-map-between-etale}.
\end{proof}

\begin{lemma}
\label{lemma-etale-permanence-two}
Soit
$$
\xymatrix{
X \ar[rr]_f \ar[rd]_p & &
Y \ar[dl]^q \\
& S
}
$$
un diagramme commutatif de morphismes de schémas. Supposons que
\begin{enumerate}
\item $f$ est surjectif et étale,
\item $p$ est étale, et
\item $q$ est localement de présentation finie\footnote{En
fait cela
est impliqué par (1) et (2), voir Descente, Lemme \ref{descent-lemma-flat-finitely-presented-permanence}. De plus,
il suffit de supposer que $f$ est surjectif, plat
et localement de présentation
finie, voir Descente, Lemme \ref{descent-lemma-smooth-permanence}.}.
\end{enumerate}
Alors $q$ est étale.
\end{lemma}

\begin{proof}
Par le Lemme \ref{lemma-smooth-permanence} nous voyons que $q$ est lisse. Il suffit
donc de voir que $q$ a une dimension relative
$0$. Cela résulte du lemme \ref{lemma-dimension-fibre-at-a-point-additive} et du fait que $f$
et $p$ ont une dimension relative $0$.
\end{proof}

\noindent
Une dernière caractérisation des morphismes lisses est qu'un morphisme lisse
$f : X \to S$ est localement la composition d'un morphisme étale par une
projection $\mathbf{A}_S^d \to S$.

\begin{lemma}
\label{lemma-smooth-etale-over-affine-space}
\begin{slogan}
Les schémas lisses sont des espaces affines de type étalement-local.
\end{slogan}
Soit $\varphi : X \to Y$ un morphisme de schémas. Soit $x \in X$. Soit
$V \subset Y$ un voisinage ouvert affine de $\varphi(x)$. Si $\varphi$
est lisse à $x$, alors il existe un entier $d \geq 0$
et un ouvert affine $U \subset X$ avec $x \in U$
et $\varphi(U) \subset V$ tel qu'il existe un diagramme commutatif
$$
\xymatrix{
X \ar[d] & U \ar[l] \ar[d] \ar[r]_-\pi & \mathbf{A}^d_V \ar[ld] \\
Y & V \ar[l]
}
$$
où $\pi$ est étale.
\end{lemma}

\begin{proof}
Par
Lemme \ref{lemma-smooth-locally-standard-smooth} on peut trouver un affine
ouvert $U$ comme dans le lemme tel que $\varphi|_U : U \to V$
est lisse standard. Écrivez $U = \Spec(A)$
et $V = \Spec(R)$ afin que nous puissions écrire
$$
A = R[x_1, \ldots, x_n]/(f_1, \ldots, f_c)
$$
avec le mappage
$$
g =
\det
\left(
\begin{matrix}
\partial f_1/\partial x_1 &
\partial f_2/\partial x_1 &
\ldots &
\partial f_c/\partial x_1 \\
\partial f_1/\partial x_2 &
\partial f_2/\partial x_2 &
\ldots &
\partial f_c/\partial x_2 \\
\ldots & \ldots & \ldots & \ldots \\
\partial f_1/\partial x_c &
\partial f_2/\partial x_c &
\ldots &
\partial f_c/\partial x_c
\end{matrix}
\right)
$$
sur un élément inversible de $A$. Il est alors
clair que $R[x_{c + 1}, \ldots, x_n] \to A$ est un lissage standard de dimension
relative $0$. Il est donc lisse de dimension
relative $0$. Autrement dit l'homomorphisme
d'anneaux $R[x_{c + 1}, \ldots, x_n] \to A$ est étale. Comme $\mathbf{A}^{n - c}_V = \Spec(R[x_{c + 1}, \ldots, x_n])$ le
lemme avec $d = n - c$.
\end{proof}
















\section{Faisceaux relativement amples}
\label{section-relatively-ample}

\noindent
Soit $X$ un schéma et $\mathcal{L}$ un faisceau inversible sur
$X$. Alors $\mathcal{L}$ est suffisant sur $X$ si $X$
est quasi-compact et que chaque point de $X$ est
contenu dans un ouvert affine de la forme $X_s$, où $s \in \Gamma(X, \mathcal{L}^{\otimes n})$ et $n \geq 1$,
voir Propriétés, Définition \ref{properties-definition-ample}. Nous transformons
cela en une notion relative comme suit.

\begin{definition}
\label{definition-relatively-ample}
\begin{reference}
\cite[II Définition 4.6.1]{EGA}
\end{reference}
Soit $f : X \to S$ un morphisme de schémas. Soit $\mathcal{L}$ un
module $\mathcal{O}_X$ inversible. On dit que $\mathcal{L}$ est {\it
relativement ample}, ou {\it $f$-relativement ample}, ou {\it
amplement sur $X/S$}, ou {\it $f$-ample}
si $f : X \to S$ est quasi-compact, et si pour tout ouvert
affine $V \subset S$ la restriction de $\mathcal{L}$ au sous-schéma ouvert $f^{-1}(V)$
de $X$ est suffisant.
\end{definition}

\noindent
On remarque que l'existence d'un faisceau relativement ample sur $X$ ne
force pas le morphisme $X \to S$ à être de type fini.

\begin{lemma}
\label{lemma-ample-power-ample}
Soit $X \to S$ un morphisme de schémas.
Soit $\mathcal{L}$ un module $\mathcal{O}_X$ inversible. Soit $n \geq 1$.
Alors $\mathcal{L}$ est $f$-ample si et seulement
si $\mathcal{L}^{\otimes n}$ est $f$-ample.
\end{lemma}

\begin{proof}
Cela résulte de Propriétés, Lemme \ref{properties-lemma-ample-power-ample}.
\end{proof}

\begin{lemma}
\label{lemma-relatively-ample-separated}
Soit $f : X \to S$ un morphisme de schémas. S'il existe
un $f$-faisceau ample inversible, alors $f$
est séparé.
\end{lemma}

\begin{proof}
Être séparé est local sur la base (voir
Schémas, Lemme \ref{schemes-lemma-characterize-separated} par exemple ; cela découle aussi facilement
de la définition). Nous pouvons donc
supposer que $S$ est affine et
que $X$ a un faisceau ample inversible.
Dans ce cas,
le résultat découle de Propriétés, Lemme \ref{properties-lemma-ample-separated}.
\end{proof}

\noindent
Il existe de nombreuses façons de caractériser des faisceaux inversibles
relativement amples, analogues aux conditions
équivalentes dans Propriétés, Proposition \ref{properties-proposition-characterize-ample}. Nous les
ajouterons ici si nécessaire.

\begin{lemma}
\label{lemma-characterize-relatively-ample}
\begin{reference}
\cite[II, Proposition 4.6.3]{EGA}
\end{reference}
Soit $f : X \to S$ un morphisme de schémas quasi-compact. Soit $\mathcal{L}$
un faisceau inversible sur $X$. Les conditions
suivantes sont équivalentes :
\begin{enumerate}
\item Le faisceau inversible $\mathcal{L}$ est $f$-ample.
\item Il existe un recouvrement ouvert $S = \bigcup V_i$
tel que chaque $\mathcal{L}|_{f^{-1}(V_i)}$ soit ample
par rapport à $f^{-1}(V_i) \to V_i$.
\item Il existe un recouvrement ouvert affine $S = \bigcup V_i$
tel que chaque $\mathcal{L}|_{f^{-1}(V_i)}$ soit ample.
\item Il existe une $\mathcal{O}_S$-algèbre graduée
$\mathcal{A}$ et une application de $\mathcal{O}_X$-algèbres
graduées $\psi : f^*\mathcal{A} \to \bigoplus_{d \geq 0} \mathcal{L}^{\otimes d}$ telle que
$U(\psi) = X$ et
$$
r_{\mathcal{L}, \psi} :
X
\longrightarrow
\underline{\text{Proj}}_S(\mathcal{A})
$$
est une immersion ouverte (voir
Constructions, Lemme \ref{constructions-lemma-invertible-map-into-relative-proj} pour la notation).
\item Le morphisme $f$ est quasi-séparé
et la partie
(4) ci-dessus tient avec $\mathcal{A} = f_*(\bigoplus_{d \geq 0} \mathcal{L}^{\otimes d})$ et $\psi$
le mappage d'adjonction.
\item Identique à (4) mais exigeant simplement que $r_{\mathcal{L}, \psi}$ soit
une immersion.
\end{enumerate}
\end{lemma}

\begin{proof}
Il ressort immédiatement de la définition que (1) implique (2) et
(2) implique (3). Il est clair que (5) implique (4).

\medskip\noindent
Supposons que (3) soit valable pour le recouvrement
ouvert affine $S = \bigcup V_i$.
Nous allons montrer (5) prises. Puisque
chaque $f^{-1}(V_i)$ possède
un faisceau ample inversible on voit que $f^{-1}(V_i)$
est séparé
(Propriétés, Lemme \ref{properties-lemma-ample-separated}). Par conséquent,
$f$ est séparé. Par Schémas, Lemme \ref{schemes-lemma-push-forward-quasi-coherent} nous voyons
que $\mathcal{A} = f_*(\bigoplus_{d \geq 0} \mathcal{L}^{\otimes d})$ est une $\mathcal{O}_S$-algèbre
graduée quasi-cohérente. Désignons $\psi : f^*\mathcal{A} \to \bigoplus_{d \geq 0} \mathcal{L}^{\otimes d}$ le mappage
d'adjonction.
La description du $U(\psi)$
ouvert dans Constructions,
section \ref{constructions-section-invertible-relative-proj} et
la définition de l'ampleur
de $\mathcal{L}|_{f^{-1}(V_i)}$ montrent que $U(\psi) = X$.
De plus,
Constructions, Lemme \ref{constructions-lemma-invertible-map-into-relative-proj} partie (3) montre
que la restriction de $r_{\mathcal{L}, \psi}$ à
$f^{-1}(V_i)$ est la même que le morphisme
de Propriétés, Lemme \ref{properties-lemma-map-into-proj}
qui est une immersion
ouverte d'après Propriétés, Lemme \ref{properties-lemma-ample-immersion-into-proj}.
Donc (5) est valable.

\medskip\noindent
Montrons que (4) implique (1).
Supposons (4). Notons $\pi : \underline{\text{Proj}}_S(\mathcal{A}) \to S$ le morphisme
de structure. Choisissez $V \subset S$ ouvert
affine. Par Constructions, Définition \ref{constructions-definition-relative-proj}
on voit que $\pi^{-1}(V) \subset \underline{\text{Proj}}_S(\mathcal{A})$ est égal à $\text{Proj}(A)$ où
$A = \mathcal{A}(V)$ comme anneau gradué. Ainsi $r_{\mathcal{L}, \psi}$
applique $f^{-1}(V)$ de manière
isomorphe sur un ouvert
quasi-compact de $\text{Proj}(A)$.
De plus, $\mathcal{L}^{\otimes d}$ est isomorphe au retrait de
$\mathcal{O}_{\text{Proj}(A)}(d)$ pour certains $d \geq 1$.
(Voir la partie (3)
de Constructions, Lemme \ref{constructions-lemma-invertible-map-into-relative-proj} et
l'énoncé final de Constructions,
Lemme \ref{constructions-lemma-invertible-map-into-proj}.) Cela implique
que $\mathcal{L}|_{f^{-1}(V)}$ est amplement
rempli par les Propriétés, Lemmes
\ref{properties-lemma-open-in-proj-ample} et \ref{properties-lemma-ample-power-ample}.

\medskip\noindent
Supposons (6). Par l'équivalence de (1) - (5) ci-dessus on
voit que la propriété d'être relativement ample sur $X/S$
est locale sur $S$. On peut donc supposer que
$S$ est affine, et il faut montrer que $\mathcal{L}$
est ample sur $X$. Dans ce cas le morphisme $r_{\mathcal{L}, \psi}$
est identifié au morphisme, également noté $r_{\mathcal{L}, \psi} : X \to \text{Proj}(A)$ associé
à l'application $\psi : A = \mathcal{A}(V) \to \Gamma_*(X, \mathcal{L})$. (Voir les références
ci-dessus.) Comme ci-dessus, nous voyons également
que $\mathcal{L}^{\otimes d}$ est le retrait du faisceau $\mathcal{O}_{\text{Proj}(A)}(d)$
pour certains $d \geq 1$. De plus, puisque
$X$ est quasi-compact nous
voyons que $X$ s'identifie à un sous-schéma fermé
d'un sous-schéma ouvert quasi-compact $Y \subset \text{Proj}(A)$.
Par
Constructions,
Lemme \ref{constructions-lemma-ample-on-proj} (voir aussi Propriétés,
Lemme
\ref{properties-lemma-open-in-proj-ample})
on voit que $\mathcal{O}_Y(d')$ est un faisceau
ample inversible sur $Y$ pour certains $d' \geq 1$.
Puisque la restriction d'un faisceau suffisant à un
sous-schéma fermé est ample, voir Propriétés,
Lemme
\ref{properties-lemma-ample-on-closed} Nous concluons que le
retrait de $\mathcal{O}_Y(d')$ est
ample. En combinant ces résultats avec Propriétés,
Lemme
\ref{properties-lemma-ample-power-ample} Nous concluons que $\mathcal{L}$
est suffisant comme souhaité.
\end{proof}

\begin{lemma}
\label{lemma-ample-over-affine}
\begin{reference}
\cite[II Corollaire 4.6.6]{EGA}
\end{reference}
Soit $f : X \to S$ un morphisme de schémas. Soit $\mathcal{L}$
un module $\mathcal{O}_X$ inversible. Supposons que $S$ soit
affine. Alors $\mathcal{L}$
est $f$-relativement amplement si et seulement si $\mathcal{L}$
est amplement sur $X$.
\end{lemma}

\begin{proof}
Immédiat du lemme \ref{lemma-characterize-relatively-ample} et des
définitions.
\end{proof}

\begin{lemma}
\label{lemma-quasi-affine-O-ample}
\begin{reference}
\cite[II Proposition 5.1.6]{EGA}
\end{reference}
Soit $f : X \to S$ un morphisme de schémas. Alors $f$ est quasi-affine si
et seulement si $\mathcal{O}_X$ est $f$-relativement ample.
\end{lemma}

\begin{proof}
Cela résulte de Propriétés, Lemme \ref{properties-lemma-quasi-affine-O-ample} et des
définitions.
\end{proof}

\begin{lemma}
\label{lemma-pullback-ample-tensor-relatively-ample}
Soit $f : X \to Y$ un morphisme de schémas, $\mathcal{M}$
un $\mathcal{O}_Y$-module inversible, et $\mathcal{L}$ un
$\mathcal{O}_X$-module inversible.
\begin{enumerate}
\item Si $\mathcal{L}$ est suffisant pour $f$
et que $\mathcal{M}$ est suffisant, alors $\mathcal{L} \otimes f^*\mathcal{M}^{\otimes a}$ est suffisant pour
$a \gg 0$.
\item Si $\mathcal{M}$ est ample et
$f$ quasi-affine, alors $f^*\mathcal{M}$ est ample.
\end{enumerate}
\end{lemma}

\begin{proof}
Supposons que $\mathcal{L}$ est $f$-ample et $\mathcal{M}$
ample. Par hypothèse $Y$
et $f$ sont quasi-compacts (voir Définition
\ref{definition-relatively-ample} et Propriétés, Définition \ref{properties-definition-ample}).
Donc $X$ est quasi-compact.
Par Propriétés, Lemme \ref{properties-lemma-ample-separated} le schéma $Y$
est séparé et par le lemme \ref{lemma-relatively-ample-separated} le morphisme
$f$ est séparé. Ainsi $X$ est
séparé par Schémas, Lemme \ref{schemes-lemma-separated-permanence}. Choisissez
$x \in X$. On peut choisir $m \geq 1$
et $t \in \Gamma(Y, \mathcal{M}^{\otimes m})$ tels que $Y_t$ soit affine et
$f(x) \in Y_t$. Puisque $\mathcal{L}$ se limite à un faisceau
amplement inversible sur $f^{-1}(Y_t) = X_{f^*t}$, nous
pouvons choisir $n \geq 1$ et $s \in \Gamma(X_{f^*t}, \mathcal{L}^{\otimes n})$ avec $x \in (X_{f^*t})_s$ avec
$(X_{f^*t})_s$ affines. Par Propriétés, Lemme
\ref{properties-lemma-invert-s-sections} partie (2) dont les hypothèses sont satisfaites
par ce qui précède, il existe
un entier $e \geq 1$ et une section
$s' \in \Gamma(X, \mathcal{L}^{\otimes n} \otimes f^*\mathcal{M}^{\otimes em})$ qui se restreint à $s(f^*t)^e$ sur $X_{f^*t}$. Pour
tout $b > 0$, considérez la section $s'' = s'(f^*t)^b$
de $\mathcal{L}^{\otimes n} \otimes f^*\mathcal{M}^{\otimes (e + b)m}$. Alors $X_{s''} = (X_{f^*t})_s$
est un ouvert affine de $X$ contenant $x$.
En choisissant $b$ de telle sorte que $n$
divise $e + b$, nous voyons
$\mathcal{L}^{\otimes n} \otimes f^*\mathcal{M}^{\otimes (e + b)m}$ est la $n$ième puissance de $\mathcal{L} \otimes f^*\mathcal{M}^{\otimes a}$
pour certains $a$ et nous pouvons obtenir n'importe quel
$a$ divisible par $m$ et assez grand. Puisque
$X$ est quasi-compact, un nombre fini de
ces ouverts affines couvrent $X$. Nous concluons
que pour certains $a$
suffisamment divisibles et assez grands le faisceau inversible
$\mathcal{L} \otimes f^*\mathcal{M}^{\otimes a}$ est amplement suffisant sur $X$.
D'un autre côté, nous savons que $\mathcal{M}^{\otimes c}$ (et donc
son retrait vers $X$) est généré globalement pour tous
les $c \gg 0$ par Propriétés, Proposition \ref{properties-proposition-characterize-ample}.
Ainsi $\mathcal{L} \otimes f^*\mathcal{M}^{\otimes a + c}$ est suffisant (Propriétés, Lemme \ref{properties-lemma-ample-tensor-globally-generated})
pour $c \gg 0$ et (1) est prouvé.

\medskip\noindent
La partie (2) découle du lemme \ref{lemma-quasi-affine-O-ample},
Propriétés, Lemme \ref{properties-lemma-ample-power-ample} et de la partie
(1).
\end{proof}

\begin{lemma}
\label{lemma-ample-composition}
Soient $g : Y \to S$ et $f : X \to Y$ des morphismes de schémas.
Soit $\mathcal{M}$ un module $\mathcal{O}_Y$ inversible.
Soit $\mathcal{L}$ un module $\mathcal{O}_X$ inversible. Si
$S$ est quasi-compact, $\mathcal{M}$ est $g$-ample
et $\mathcal{L}$ est
$f$-ample, alors $\mathcal{L} \otimes f^*\mathcal{M}^{\otimes a}$ est $g \circ f$-ample
pour $a \gg 0$.
\end{lemma}

\begin{proof}
Soit $S = \bigcup_{i = 1, \ldots, n} V_i$ un recouvrement ouvert affine
fini. Par le Lemme \ref{lemma-characterize-relatively-ample}
il suffit de
prouver que $\mathcal{L} \otimes f^*\mathcal{M}^{\otimes a}$ suffit
sur $(g \circ f)^{-1}(V_i)$ pour $i = 1, \ldots, n$.
Ainsi le lemme
découle du lemme \ref{lemma-pullback-ample-tensor-relatively-ample}.
\end{proof}

\begin{lemma}
\label{lemma-ample-base-change}
Soit $f : X \to S$ un morphisme de schémas. Soit
$\mathcal{L}$ un module $\mathcal{O}_X$ inversible. Soit $S' \to S$
un morphisme de schémas. Soit $f' : X' \to S'$
le changement de base de $f$ et notons $\mathcal{L}'$
le retrait de $\mathcal{L}$ vers $X'$. Si $\mathcal{L}$
est $f$-ample, alors $\mathcal{L}'$ est $f'$-ample.
\end{lemma}

\begin{proof}
Par le Lemme \ref{lemma-characterize-relatively-ample} il suffit de trouver un recouvrement
ouvert affine $S' = \bigcup U'_i$ tel que $\mathcal{L}'$
se restreint à un faisceau ample inversible
sur $(f')^{-1}(U_i')$ pour tout $i$. Nous pouvons
choisir le mappage $U'_i$ dans un $U_i \subset S$
ouvert affine. Dans ce cas le morphisme $(f')^{-1}(U'_i) \to f^{-1}(U_i)$
est affine comme un changement de base
du morphisme affine $U'_i \to U_i$ (Lemme
\ref{lemma-base-change-affine}). Ainsi $\mathcal{L}'|_{(f')^{-1}(U'_i)}$ est
amplement suffisant par le lemme \ref{lemma-pullback-ample-tensor-relatively-ample}.
\end{proof}

\begin{lemma}
\label{lemma-ample-permanence}
Soient $g : Y \to S$ et $f : X \to Y$ des morphismes de schémas. Soit
$\mathcal{L}$ un module $\mathcal{O}_X$ inversible. Si $\mathcal{L}$
est $g \circ f$-ample et $f$ est quasi-compact\footnote{Cela
s'ensuit si $g$ est quasi-séparé
par Schémas, Lemme \ref{schemes-lemma-quasi-compact-permanence}.} alors $\mathcal{L}$
est $f$-ample.
\end{lemma}

\begin{proof}
Supposons que $f$ soit quasi-compact et que $\mathcal{L}$
soit $g \circ f$-ample. Soit $U \subset S$
être un ouvert affine et soit
$V \subset Y$ être un ouvert affine avec
$g(V) \subset U$. Alors $\mathcal{L}|_{(g \circ f)^{-1}(U)}$
est suffisant sur $(g \circ f)^{-1}(U)$ par hypothèse.
Puisque $f^{-1}(V) \subset (g \circ f)^{-1}(U)$ on voit que $\mathcal{L}|_{f^{-1}(V)}$
est suffisant sur $f^{-1}(V)$ par Propriétés, Lemme
\ref{properties-lemma-ample-on-locally-closed}. À savoir, $f^{-1}(V) \to (g \circ f)^{-1}(U)$ est une immersion
ouverte quasi-compacte
par Schémas, Lemme \ref{schemes-lemma-quasi-compact-permanence} puisque
$(g \circ f)^{-1}(U)$ est séparé
(Propriétés, Lemme \ref{properties-lemma-ample-separated}) et $f^{-1}(V)$ est
quasi-compact (car $f$ est quasi-compact).
Ainsi Nous concluons que $\mathcal{L}$ est $f$-ample
par le lemme \ref{lemma-characterize-relatively-ample}.
\end{proof}







\section{Faisceaux très amples}
\label{section-very-ample}

\noindent
Rappelons que étant donné un faisceau quasi-cohérent
$\mathcal{E}$ sur un schéma $S$ le fibré projectif
{\it }
associé à $\mathcal{E}$ est le morphisme $\mathbf{P}(\mathcal{E}) \to S$, où
$\mathbf{P}(\mathcal{E}) = \underline{\text{Proj}}_S(\text{Sym}(\mathcal{E}))$,
voir Constructions, Définition \ref{constructions-definition-projective-bundle}.

\begin{definition}
\label{definition-very-ample}
Soit $f : X \to S$ un morphisme de schémas. Soit $\mathcal{L}$
un module $\mathcal{O}_X$ inversible. On dit que $\mathcal{L}$
est {\it relativement très ample} ou plus
précisément {\it $f$-relativement
très ample}, ou encore {\it
très ample sur $X/S$}, ou {\it
$f$-très ample} s'il existe un $\mathcal{O}_S$-module
$\mathcal{E}$ quasi-cohérent
et une immersion $i : X \to \mathbf{P}(\mathcal{E})$ sur $S$ tel que $\mathcal{L} \cong i^*\mathcal{O}_{\mathbf{P}(\mathcal{E})}(1)$.
\end{definition}



\noindent
Puisqu'il n'y a pas d'hypothèse de quasi-compacité dans cette définition, il n'est pas vrai en général
qu'un ample inversible à faisceau relativement très large soit un ample inversible à faisceau
relativement important.

\begin{lemma}
\label{lemma-ample-very-ample}
\begin{reference}
\cite[II, Proposition 4.6.2]{EGA}
\end{reference}
Soit $f : X \to S$ un morphisme de schémas. Soit
$\mathcal{L}$ un module $\mathcal{O}_X$ inversible. Si $f$
est quasi-compact et $\mathcal{L}$ est un ample inversible
relativement très gros, alors $\mathcal{L}$ est un ample inversible
relativement gros.
\end{lemma}

\begin{proof}
Par définition il existe un $\mathcal{O}_S$-module
$\mathcal{E}$ quasi-cohérent et une immersion
$i : X \to \mathbf{P}(\mathcal{E})$
sur $S$ telle que $\mathcal{L} \cong i^*\mathcal{O}_{\mathbf{P}(\mathcal{E})}(1)$.
Posons $\mathcal{A} = \text{Sym}(\mathcal{E})$, donc
$\mathbf{P}(\mathcal{E}) = \underline{\text{Proj}}_S(\mathcal{A})$ par définition. L'algèbre $\mathcal{O}_S$
graduée $\mathcal{A}$ est munie d'un
homomorphisme
$$
\psi :
\mathcal{A} \to
\bigoplus\nolimits_{n \geq 0}
\pi_*\mathcal{O}_{\mathbf{P}(\mathcal{E})}(n) \to
\bigoplus\nolimits_{n \geq 0}
f_*\mathcal{L}^{\otimes n}
$$
où la deuxième flèche utilise l'identification
$\mathcal{L} \cong i^*\mathcal{O}_{\mathbf{P}(\mathcal{E})}(1)$. Par adjointité de $f_*$
et $f^*$ on obtient un morphisme
$\psi : f^*\mathcal{A} \to \bigoplus_{n \geq 0}\mathcal{L}^{\otimes n}$. Nous omettons de vérifier que le
morphisme $r_{\mathcal{L}, \psi}$ associé à cette application
est exactement l'immersion $i$.
Le résultat découle
donc de la partie (6) du lemme \ref{lemma-characterize-relatively-ample}.
\end{proof}

\noindent
Pour formuler une réciproque correcte de ce lemme, demandons-nous
si, étant donné un faisceau inversible relativement ample
$\mathcal{L}$, il existe un entier $n \geq 1$
tel que $\mathcal{L}^{\otimes n}$ soit relativement très ample ? En général, c'est
faux. Il y a plusieurs choses qui empêchent que cela soit vrai :
\begin{enumerate}
\item Même si $S$ est affine, il peut arriver qu'aucun entier
fini $n$ ne fonctionne car $X \to S$ n'est
pas de type fini, voir Exemple \ref{example-not-finite-type-proj}.
\item Si la base n'est pas quasi-compacte, le résultat peut être faux
même lorsque $f$ est de type fini. En effet, étant
donné un corps $k$ il existe un schéma $X_d$ de type fini
sur $k$ avec un faisceau inversible ample $\mathcal{O}_{X_d}(1)$ de sorte que
la plus petite puissance tensorielle de $\mathcal{O}_{X_d}(1)$ qui est très ample soit la
$d$ième puissance. Voir Exemple \ref{example-not-bounded}.
Prendre pour $f$ la réunion disjointe des schémas $X_d$, envoyée
sur la réunion disjointe de copies de $\Spec(k)$, donne un exemple.
\end{enumerate}
Pour voir notre version de l'inverse, jetez un œil
dans le lemme \ref{lemma-finite-type-ample-very-ample} ci-dessous. Nous ferons quelques
travaux préliminaires avant de le prouver.

\begin{example}
\label{example-very-ample}
Soit $S$ un schéma. Soit $\mathcal{A}$ une algèbre $\mathcal{O}_S$
graduée quasi-cohérente générée par $\mathcal{A}_1$ sur $\mathcal{A}_0$.
Posons $X = \underline{\text{Proj}}_S(\mathcal{A})$. Dans ce cas $\mathcal{O}_X(1)$ est inversible
et très ample sur $X/S$. À savoir, le morphisme
associé à l'homomorphisme d'algèbres graduées $\mathcal{O}_S$
$$
\text{Sym}_{\mathcal{O}_X}^*(\mathcal{A}_1)
\longrightarrow
\mathcal{A}
$$
est une immersion fermée $X \to \mathbf{P}(\mathcal{A}_1)$
qui ramène $\mathcal{O}_{\mathbf{P}(\mathcal{A}_1)}(1)$ à $\mathcal{O}_X(1)$,
voir
Constructions, Lemme \ref{constructions-lemma-surjective-generated-degree-1-map-relative-proj}.
\end{example}

\begin{example}
\label{example-not-finite-type-proj}
Soit $k$ un corps.
Considérons l'algèbre $k$ graduée
$$
A = k[U, V, Z_1, Z_2, Z_3, \ldots]/I
\quad
\text{avec}
\quad
I = (U^2 - Z_1^2, U^4 - Z_2^2, U^6 - Z_3^2, \ldots)
$$
avec la graduation définie par $\deg(U) = \deg(V) = \deg(Z_1) = 1$ et $\deg(Z_d) = d$.
Notez que $X = \text{Proj}(A)$
est couvert par $D_{+}(U)$ et $D_{+}(V)$. Ainsi les faisceaux
$\mathcal{O}_X(n)$ sont tous inversibles et isomorphes
à $\mathcal{O}_X(1)^{\otimes n}$. En particulier $\mathcal{O}_X(1)$ est ample et
$f$-ample pour le morphisme $f : X \to \Spec(k)$. Nous
affirmons qu'aucune puissance de $\mathcal{O}_X(1)$
n'est $f$-relativement très ample. En effet, il est facile
de voir que $\Gamma(X, \mathcal{O}_X(n))$ est la composante homogène de degré $n$
de l'algèbre $A$. Donc si $\mathcal{O}_X(n)$ était très ample, alors
$X$ serait un sous-schéma fermé d'un espace projectif sur
$k$ et donc de type fini sur $k$. Par contre
$D_{+}(V)$ est le
spectre de $k[t, t_1, t_2, \ldots]/(t^2 - t_1^2, t^4 - t_2^2, t^6 - t_3^2, \ldots)$ qui n'est pas de type
fini sur $k$.
\end{example}

\begin{example}
\label{example-not-bounded}
Soit $k$ un corps infini. Soit $\lambda_1, \lambda_2, \lambda_3, \ldots$ des éléments distincts par paires
de $k^*$. (Ce n'est pas strictement nécessaire, et en fait l'exemple
fonctionne parfaitement même si tous les $\lambda_i$ sont égaux à $1$.)
Considérons
l'algèbre $k$ graduée
$$
A_d = k[U, V, Z]/I_d
\quad
\text{avec}
\quad
I_d = (Z^2 - \prod\nolimits_{i = 1}^{2d} (U - \lambda_i V)).
$$
avec la graduation définie par $\deg(U) = \deg(V) = 1$ et $\deg(Z) = d$. Ensuite $X_d = \text{Proj}(A_d)$ a un faisceau inversible ample
$\mathcal{O}_{X_d}(1)$. Nous affirmons que si $\mathcal{O}_{X_d}(n)$ est
très ample, alors $n \geq d$. La raison en est que $Z$
a le diplôme $d$, et donc $\Gamma(X_d, \mathcal{O}_{X_d}(n)) =
k[U, V]_n$
pour $n < d$. Détails omis.
\end{example}

\begin{lemma}
\label{lemma-relatively-very-ample-separated}
Soit $f : X \to S$ un morphisme de schémas. Soit $\mathcal{L}$
un faisceau inversible sur $X$. Si $\mathcal{L}$
est relativement très ample sur $X/S$ alors $f$
est séparé.
\end{lemma}

\begin{proof}
La séparation est locale sur la base
(voir Schémas, section \ref{schemes-section-separation-axioms}). Une immersion
est séparée
(voir Schémas, Lemme \ref{schemes-lemma-immersions-monomorphisms}). D'où le lemme s'ensuit
puisque localement $X$ a une immersion
dans le spectre homogène d'un anneau gradué qui est
séparé, voir Constructions, Lemme \ref{constructions-lemma-proj-separated}.
\end{proof}

\begin{lemma}
\label{lemma-relatively-very-ample}
Soit $f : X \to S$ un morphisme de schémas. Soit $\mathcal{L}$
un faisceau inversible sur $X$. Supposons que $f$
soit quasi-compact. Les conditions suivantes sont
équivalentes
\begin{enumerate}
\item $\mathcal{L}$ est relativement très ample sur $X/S$,
\item il existe un recouvrement ouvert $S = \bigcup V_j$ tel
que $\mathcal{L}|_{f^{-1}(V_j)}$ est relativement très ample sur $f^{-1}(V_j)/V_j$
pour tout $j$,
\item il existe un faisceau quasi-cohérent
d'algèbres $\mathcal{O}_S$ graduées
$\mathcal{A}$ générées en degré $1$ sur $\mathcal{O}_S$
et une application d'algèbres $\mathcal{O}_X$
graduées $\psi : f^*\mathcal{A} \to \bigoplus_{n \geq 0} \mathcal{L}^{\otimes n}$ telle que $f^*\mathcal{A}_1 \to \mathcal{L}$
est surjectif
et le morphisme associé $r_{\mathcal{L}, \psi} : X \to \underline{\text{Proj}}_S(\mathcal{A})$ est une immersion,
et
\item $f$ est quasi-séparé,
l'application canonique $\psi : f^*f_*\mathcal{L} \to \mathcal{L}$ est surjectif,
et l'application associée $r_{\mathcal{L}, \psi} : X \to \mathbf{P}(f_*\mathcal{L})$ est une
immersion.
\end{enumerate}
\end{lemma}

\begin{proof}
Il est clair que (1) implique (2). Il est également
clair que (4) implique (1) ; l'hypothèse de
quasi-séparation dans (4) est utilisée pour garantir que $f_*\mathcal{L}$
est quasi-cohérent via Schémas, Lemme \ref{schemes-lemma-push-forward-quasi-coherent}.

\medskip\noindent
Supposons (2). Nous
allons démontrer (4). Soit $S = \bigcup V_j$ un
recouvrement ouvert comme dans (2). Définissez
$X_j = f^{-1}(V_j)$ et $f_j : X_j \to V_j$ la restriction
de $f$. On voit que $f$ est
séparé par le lemme \ref{lemma-relatively-very-ample-separated} (car la séparation
est locale sur la base). Par hypothèse il existe un $\mathcal{O}_{V_j}$-module
$\mathcal{E}_j$ quasi-cohérent
et une immersion $i_j : X_j \to \mathbf{P}(\mathcal{E}_j)$ avec $\mathcal{L}|_{X_j} \cong i_j^*\mathcal{O}_{\mathbf{P}(\mathcal{E}_j)}(1)$. Le morphisme
$i_j$ correspond à une surjection
$f_j^*\mathcal{E}_j \to \mathcal{L}|_{X_j}$, voir Constructions
, section \ref{constructions-section-projective-bundle}. Cette application est adjointe
à une application $\mathcal{E}_j \to f_*\mathcal{L}|_{V_j}$ telle que
la composition
$$
f_j^*\mathcal{E}_j \to (f^*f_*\mathcal{L})|_{X_j} \to \mathcal{L}|_{X_j}
$$
est surjectif. Nous concluons
que $\psi : f^*f_*\mathcal{L} \to \mathcal{L}$ est surjectif. Soit $r_{\mathcal{L}, \psi} : X \to \mathbf{P}(f_*\mathcal{L})$
le morphisme associé. Il reste à montrer que
$r_{\mathcal{L}, \psi}$ est une immersion ; nous invitons le lecteur
à le prouver par lui-même. L'application $\mathcal{O}_{V_j}$-module
$\mathcal{E}_j \to f_*\mathcal{L}|_{V_j}$ détermine un homomorphisme sur les
algèbres symétriques, qui à son tour définit un
morphisme
$$
\mathbf{P}(f_*\mathcal{L}|_{V_j}) \supset U_j \longrightarrow
\mathbf{P}(\mathcal{E}_j)
$$
où $U_j$ est le sous-schéma
ouvert des Constructions, Lemme \ref{constructions-lemma-morphism-relative-proj}.
La compatibilité de $\psi$ avec $\mathcal{E}_j \to f_*\mathcal{L}|_{V_j}$
montre que $r_{\mathcal{L}, \psi}(X_j) \subset U_j$ et qu'il existe une
factorisation
$$
\xymatrix{
X_j \ar[r]^-{r_{\mathcal{L}, \psi}} & U_j \ar[r] &
\mathbf{P}(\mathcal{E}_j)
}
$$
Nous omettons la vérification.
Cela montre que $r_{\mathcal{L}, \psi}$ est une immersion.

\medskip\noindent
À ce stade, nous voyons que (1), (2) et (4) sont équivalents.
Clairement (4) implique (3). Supposons (3). Nous
allons démontrer (1). Soit $\mathcal{A}$ un faisceau quasi-cohérent de $\mathcal{O}_S$-algèbres
graduées générées en degré $1$ sur $\mathcal{O}_S$.
Considérons l'application des $\mathcal{O}_S$-algèbres graduées $\text{Sym}(\mathcal{A}_1) \to \mathcal{A}$. Ceci
est surjectif par hypothèse et induit donc une immersion fermée
$$
\underline{\text{Proj}}_S(\mathcal{A})
\longrightarrow
\mathbf{P}(\mathcal{A}_1)
$$
qui ramène $\mathcal{O}(1)$ à $\mathcal{O}(1)$,
voir Constructions,
Lemme \ref{constructions-lemma-surjective-generated-degree-1-map-relative-proj}. Il est donc clair
que (3) implique (1).
\end{proof}

\begin{lemma}
\label{lemma-very-ample-base-change}
Soit $f : X \to S$ un morphisme de schémas. Soit
$\mathcal{L}$ un module $\mathcal{O}_X$ inversible. Soit $S' \to S$
un morphisme de schémas. Soit $f' : X' \to S'$
le changement de base de $f$ et notons $\mathcal{L}'$
le retrait de $\mathcal{L}$ vers $X'$. Si $\mathcal{L}$
est $f$-très ample, alors $\mathcal{L}'$ est $f'$-très ample.
\end{lemma}

\begin{proof}
Par Définition \ref{definition-very-ample} il existe un $\mathcal{O}_S$-module
$\mathcal{E}$ quasi-cohérent et une
immersion $i : X \to \mathbf{P}(\mathcal{E})$ sur $S$
telle que $\mathcal{L} \cong i^*\mathcal{O}_{\mathbf{P}(\mathcal{E})}(1)$. Le changement de
base de $\mathbf{P}(\mathcal{E})$ vers $S'$
est le fibré projectif associé à l'image réciproque
$\mathcal{E}'$ de $\mathcal{E}$,
et l'image réciproque de $\mathcal{O}_{\mathbf{P}(\mathcal{E})}(1)$
est
$\mathcal{O}_{\mathbf{P}(\mathcal{E}')}(1)$, voir Constructions,
Lemme \ref{constructions-lemma-relative-proj-base-change}. Enfin, le changement
de base
d'une immersion est une
immersion (Schémas, Lemme \ref{schemes-lemma-base-change-immersion}).
\end{proof}







\section{Faisceaux amples et très amples relativement aux morphismes de type fini}
\label{section-ample-finite-type}

\noindent
En fait, la majeure partie du contenu de cette section concerne la notion de
projet de (quasi-)morphisme que nous n'avons pas encore défini.

\begin{lemma}
\label{lemma-very-ample-finite-type-over-affine}
Soit $f : X \to S$ un morphisme de schémas. Soit $\mathcal{L}$
un faisceau inversible sur $X$. Supposons
que
\begin{enumerate}
\item le faisceau inversible $\mathcal{L}$ est très ample sur $X/S$,
\item le morphisme $X \to S$ est de type fini, et
\item $S$ est affine.
\end{enumerate}
Alors il existe un $n \geq 0$
et une immersion $i : X \to \mathbf{P}^n_S$ sur
$S$ tel que $\mathcal{L} \cong i^*\mathcal{O}_{\mathbf{P}^n_S}(1)$.
\end{lemma}

\begin{proof}
Supposons (1), (2) et
(3). La condition (3) signifie $S = \Spec(R)$ pour
certains anneaux $R$.
La condition (1) signifie par définition
qu'il existe un $\mathcal{O}_S$-module $\mathcal{E}$ quasi-cohérent
et une immersion $\alpha : X \to \mathbf{P}(\mathcal{E})$ telle que $\mathcal{L} = \alpha^*\mathcal{O}_{\mathbf{P}(\mathcal{E})}(1)$. Écrivez $\mathcal{E} = \widetilde{M}$
pour certains $R$-module $M$.
On a donc
$$
\mathbf{P}(\mathcal{E}) = \text{Proj}(\text{Sym}_R(M)).
$$
Puisque $\alpha$ est une immersion, et puisque
la topologie de $\text{Proj}(\text{Sym}_R(M))$ est générée par le standard
ouvre $D_{+}(f)$, $f \in \text{Sym}_R^d(M)$, $d \geq 1$,
on peut trouver pour chaque
$x \in X$ un $f \in \text{Sym}_R^d(M)$, $d \geq 1$, avec $\alpha(x) \in D_{+}(f)$ tel
que
$$
\alpha|_{\alpha^{-1}(D_{+}(f))} : \alpha^{-1}(D_{+}(f)) \to D_{+}(f)
$$
soit une immersion fermée.
La condition (2) implique que $X$ est quasi-compact.
On peut donc trouver
une collection finie d'éléments $f_j \in \text{Sym}_R^{d_j}(M)$,
$d_j \geq 1$ telle que pour chaque $f = f_j$ l'application
affichée ci-dessus est une immersion fermée et telle que
$\alpha(X) \subset \bigcup D_{+}(f_j)$. Écrivez $U_j = \alpha^{-1}(D_{+}(f_j))$. Notez que $U_j$ est affine
comme sous-schéma fermé du schéma affine $D_{+}(f_j)$.
Écrivez $U_j = \Spec(A_j)$. La condition (2) implique également
que $A_j$ est de type
fini sur $R$, voir Lemme \ref{lemma-locally-finite-type-characterize}. Choisissez
un nombre fini de $x_{j, k} \in A_j$ que
génère $A_j$ en tant qu'algèbre $R$. Puisque $\alpha|_{U_j}$ est une
immersion fermée on voit que $x_{j, k}$ est l'image d'un élément
$$
f_{j, k}/f_j^{e_{j, k}} \in \text{Sym}_R(M)_{(f_j)}
=
\Gamma(D_{+}(f_j), \mathcal{O}_{\text{Proj}(\text{Sym}_R(M))}).
$$
Enfin, choisir $n \geq 1$ et les éléments $y_0, \ldots, y_n \in M$ tels que chacun des polynômes
$f_j, f_{j, k} \in \text{Sym}_R(M)$ soit un polynôme dans les éléments $y_t$ à coefficients
en $R$. Considérons l'homomorphisme
gradué d'anneaux
$$
\psi : R[Y_0, \ldots, Y_n] \longrightarrow \text{Sym}_R(M),
\quad Y_i \longmapsto y_i.
$$
Désignons $F_j$, $F_{j, k}$ les éléments de $R[Y_0, \ldots, Y_n]$
tels que $\psi(F_j) = f_j$ et $\psi(F_{j, k}) = f_{j, k}$.
Par Constructions, Lemme \ref{constructions-lemma-morphism-proj} on obtient un
sous-schéma ouvert
$$
U(\psi) \subset \text{Proj}(\text{Sym}_R(M))
$$
et un morphisme
$r_\psi : U(\psi) \to \mathbf{P}^n_R$. Ce morphisme satisfait
à $r_\psi^{-1}(D_{+}(F_j)) = D_{+}(f_j)$, et nous voyons donc que
$\alpha(X) \subset U(\psi)$. De plus, il est clair
que
$$
i = r_\psi \circ \alpha : X \longrightarrow \mathbf{P}^n_R
$$
est encore une immersion
puisque $i^\sharp(F_{j, k}/F_j^{e_{j, k}}) = x_{j, k} \in
A_j = \Gamma(U_j, \mathcal{O}_X)$ par construction.
De plus, le morphisme $r_\psi$ est équipé d'une application
$\theta : r_\psi^*\mathcal{O}_{\mathbf{P}^n_R}(1)
\to \mathcal{O}_{\text{Proj}(\text{Sym}_R(M))}(1)|_{U(\psi)}$ qui est dans ce
cas un isomorphisme (pour la construction
$\theta$ voir le lemme cité ci-dessus
; certains détails omis). Puisque
l'application d'origine $\alpha$
était
supposée avoir la propriété que $\mathcal{L} = \alpha^*\mathcal{O}_{\text{Proj}(\text{Sym}_R(M))}(1)$
on obtient le résultat.
\end{proof}

\begin{lemma}
\label{lemma-quasi-affine-finite-type-over-S}
Soit $\pi : X \to S$ un morphisme de schémas. Supposons
que $X$ soit quasi-affine et que $\pi$ soit localement de
type fini. Il existe alors $n \geq 0$ et une immersion $i : X \to \mathbf{A}^n_S$ sur
$S$.
\end{lemma}

\begin{proof}
Soit $A = \Gamma(X, \mathcal{O}_X)$. Par hypothèse $X$ est quasi-compact
et s'identifie à un sous-schéma ouvert de $\Spec(A)$,
voir Propriétés, Lemme \ref{properties-lemma-quasi-affine}. De plus,
l'ensemble des ouvertures $X_f$, pour celles $f \in A$
telles que $X_f$ est affine, constitue une base
pour la topologie de $X$, voir la démonstration
de Propriétés, Lemme \ref{properties-lemma-quasi-affine}. On peut donc trouver un
nombre fini de $f_j \in A$, $j = 1, \ldots, m$ tel que $X = \bigcup X_{f_j}$,
et tel que $\pi(X_{f_j}) \subset V_j$ pour certains
$V_j \subset S$ ouverts affines. Par le Lemme
\ref{lemma-locally-finite-type-characterize} les homomorphismes d'anneaux $\mathcal{O}(V_j) \to \mathcal{O}(X_{f_j}) = A_{f_j}$ sont
de type fini. Ainsi, nous pouvons choisir $a_1, \ldots, a_N \in A$ tel que
les éléments $a_1, \ldots, a_N, 1/f_j$ génèrent $A_{f_j}$
sur $\mathcal{O}(V_j)$ pour chaque $j$. Prenons $n = m + N$ et
soit
$$
i : X \longrightarrow \mathbf{A}^n_S
$$
le morphisme donné par les sections
globales $f_1, \ldots, f_m, a_1, \ldots, a_N$ du faisceau structural de $X$.
Soit $D(x_j) \subset \mathbf{A}^n_S$ le sous-schéma ouvert où la $j$ème
fonction de coordonnées
est non nulle. Alors pour $1 \leq j \leq m$ nous avons
$i^{-1}(D(x_j)) = X_{f_j}$ et le morphisme induit $X_{f_j} \to D(x_j)$ se factorise
par l'ouvert affine $\Spec(\mathcal{O}(V_j)[x_1, \ldots, x_n, 1/x_j])$ de $D(x_j)$.
Puisque l'homomorphisme
d'anneaux $\mathcal{O}(V_j)[x_1, \ldots, x_n, 1/x_j] \to A_{f_j}$ est surjectif
par construction Nous concluons que $i^{-1}(D(x_j)) \to D(x_j)$ est une
immersion comme voulu.
\end{proof}

\begin{lemma}
\label{lemma-quasi-projective-finite-type-over-S}
Soit $f : X \to S$ un morphisme de schémas. Soit $\mathcal{L}$
un faisceau inversible sur $X$. Supposons
que
\begin{enumerate}
\item le faisceau inversible $\mathcal{L}$ est suffisant sur $X$, et que
\item le morphisme $X \to S$ est localement de type fini.
\end{enumerate}
Alors il existe un $d_0 \geq 1$ tel que pour tout
$d \geq d_0$ il existe un
$n \geq 0$ et une immersion $i : X \to \mathbf{P}^n_S$
sur $S$ tel que $\mathcal{L}^{\otimes d} \cong i^*\mathcal{O}_{\mathbf{P}^n_S}(1)$.
\end{lemma}

\begin{proof}
Soit
$A = \Gamma_*(X, \mathcal{L}) =
\bigoplus_{d \geq 0} \Gamma(X, \mathcal{L}^{\otimes d})$. Par Propriétés,
la Proposition \ref{properties-proposition-characterize-ample} l'ensemble des ouvertes affines $X_a$ avec $a \in A_{+}$ homogène forme
une base pour la topologie de $X$. Par conséquent,
nous pouvons trouver un nombre fini
d'éléments de ce type $a_0, \ldots, a_n \in A_{+}$ tels que
\begin{enumerate}
\item nous avons $X = \bigcup_{i = 0, \ldots, n} X_{a_i}$,
\item chaque $X_{a_i}$ est affine, et
\item chaque application $X_{a_i}$ dans un $V_i \subset S$ ouvert affine.
\end{enumerate}
Par le Lemme \ref{lemma-locally-finite-type-characterize} on voit que les homomorphismes
d'anneaux $\mathcal{O}_S(V_i) \to \mathcal{O}_X(X_{a_i})$
sont de type fini. Nous pouvons donc
trouver un nombre fini d'éléments
$f_{ij} \in \mathcal{O}_X(X_{a_i})$, $j = 1, \ldots, n_i$ qui génèrent $\mathcal{O}_X(X_{a_i})$ comme une
$\mathcal{O}_S(V_i)$-algèbre. Par Propriétés, Lemme \ref{properties-lemma-invert-s-sections}, nous
pouvons écrire chaque $f_{ij}$ comme $a_{ij}/a_i^{e_{ij}}$
pour certains
$a_{ij} \in A_{+}$ homogènes. Soit $N$
un entier positif qui est un multiple commun de tous
les degrés des éléments $a_i$, $a_{ij}$.
Considérez les éléments
$$
a_i^{N/\deg(a_i)}, \ a_{ij}a_i^{(N/\deg(a_i)) - e_{ij}} \in A_N.
$$
Par construction ceux-ci génèrent le faisceau inversible
$\mathcal{L}^{\otimes N}$ sur $X$. Ils donnent donc lieu à
un morphisme
$$
j : X \longrightarrow
\mathbf{P}_S^{m}
\quad
\text{avec } m = n + \sum n_i
$$
sur $S$, voir Constructions, Lemme \ref{constructions-lemma-projective-space} et
Définition \ref{constructions-definition-projective-space}. De plus, $j^*\mathcal{O}_{\mathbf{P}_S}(1) = \mathcal{L}^{\otimes N}$. On
nomme les coordonnées homogènes $T_0, \ldots, T_n, T_{ij}$ au lieu
de $T_0, \ldots, T_m$. Pour $i = 0, \ldots, n$, nous avons $j^{-1}(D_{+}(T_i)) = X_{a_i}$.
De plus, en
retirant l'élément $T_{ij}/T_i$ via $j^\sharp$
on obtient l'élément $f_{ij} \in \mathcal{O}_X(X_{a_i})$. D'où le morphisme
$j$ restreint à $X_{a_i}$
donne une immersion fermée de $X_{a_i}$
dans l'ouvert affine $D_{+}(T_i) \cap \mathbf{P}^m_{V_i}$ de
$\mathbf{P}^N_S$. Nous concluons donc que le
morphisme $j$ est une immersion.
Cela implique que le lemme est valable pour certains $d$ et $n$, ce qui est suffisant
dans pratiquement toutes les applications.

\medskip\noindent
Cela prouve que pour un $d_2 \geq 1$
(à savoir $d_2 = N$ ci-dessus), certains
$m \geq 0$ il existe une certaine immersion $j : X \to \mathbf{P}^m_S$
donnée par les sections globales $s'_0, \ldots, s'_m \in \Gamma(X, \mathcal{L}^{\otimes d_2})$. Par
Propriétés, Proposition \ref{properties-proposition-characterize-ample} nous savons qu'il existe
un entier $d_1$ tel
que $\mathcal{L}^{\otimes d}$ est généré globalement pour tous les
$d \geq d_1$. Posons $d_0 = d_1 + d_2$. Nous affirmons que
le lemme est valable avec cette valeur de $d_0$. À savoir, étant donné un entier $d \geq d_0$,
nous pouvons choisir $s''_1, \ldots, s''_t
\in \Gamma(X, \mathcal{L}^{\otimes d - d_2})$ qui génère
$\mathcal{L}^{\otimes d - d_2}$ sur $X$. Définissez $k = (m + 1)t$ et
désignez $s_0, \ldots, s_k$ l'ensemble des sections
$s'_\alpha s''_\beta$, $\alpha = 0, \ldots, m$, $\beta = 1, \ldots, t$.
Ceux-ci génèrent $\mathcal{L}^{\otimes d}$ sur $X$ et
définissent donc un morphisme
$$
i : X \longrightarrow \mathbf{P}^{k - 1}_S
$$
tel que $i^*\mathcal{O}_{\mathbf{P}^n_S}(1) \cong \mathcal{L}^{\otimes d}$. Pour voir que $i$ est une immersion,
observez que $i$ est la composition
$$
X \longrightarrow \mathbf{P}^m_S \times_S \mathbf{P}^{t - 1}_S
\longrightarrow \mathbf{P}^{k - 1}_S
$$
où le premier morphisme est $(j, j')$ avec
$j'$ donné par $s''_1, \ldots, s''_t$
et le deuxième morphisme
est le plongement de Segre (Constructions, Lemme
\ref{constructions-lemma-segre-embedding}). Puisque $j$
est une immersion, $(j, j')$ l'est
aussi (appliquez Lemme \ref{lemma-immersion-permanence} à $X \to \mathbf{P}^m_S \times_S \mathbf{P}^{t - 1}_S
\to \mathbf{P}^m_S$).
Ainsi $i$ est une composition d'immersions et donc
une immersion (Schémas, Lemme \ref{schemes-lemma-composition-immersion}).
\end{proof}

\begin{lemma}
\label{lemma-finite-type-over-affine-ample-very-ample}
Soit $f : X \to S$ un morphisme de schémas. Soit $\mathcal{L}$
un module $\mathcal{O}_X$ inversible. Supposons que $S$
affine et $f$ de type fini. Les conditions
suivantes sont équivalentes
\begin{enumerate}
\item $\mathcal{L}$ est suffisant sur $X$,
\item $\mathcal{L}$ est $f$-ample,
\item $\mathcal{L}^{\otimes d}$ est $f$-très suffisant pour certains $d \geq 1$,
\item $\mathcal{L}^{\otimes d}$ est $f$-très suffisant pour tous les $d \gg 1$,
\item pour certains $d \geq 1$ il existe
$n \geq 1$ et une immersion
$i : X \to \mathbf{P}^n_S$ telle que $\mathcal{L}^{\otimes d} \cong i^*\mathcal{O}_{\mathbf{P}^n_S}(1)$, et
\item pour tout $d \gg 1$ il existe $n \geq 1$
et une immersion
$i : X \to \mathbf{P}^n_S$ telle que $\mathcal{L}^{\otimes d} \cong i^*\mathcal{O}_{\mathbf{P}^n_S}(1)$.
\end{enumerate}
\end{lemma}

\begin{proof}
L'équivalence de (1) et (2) est Lemme \ref{lemma-ample-over-affine}. L'implication
(2) $\Rightarrow$ (6) est Lemme
\ref{lemma-quasi-projective-finite-type-over-S}. Trivialement (6) implique (5).
Comme $\mathbf{P}^n_S$
est un fibré projectif sur $S$ (voir Constructions,
Lemme \ref{constructions-lemma-projective-space-bundle}) on voit que (5) implique
(3) et (6)
implique (4) de la définition d'un faisceau relativement
très ample.
Trivialement (4) implique (3). Pour finir il
faut montrer que (3) implique (2) qui découle du lemme
\ref{lemma-ample-very-ample} et Lemme \ref{lemma-ample-power-ample}.
\end{proof}

\begin{lemma}
\label{lemma-finite-type-ample-very-ample}
Soit $f : X \to S$ un morphisme de schémas. Soit $\mathcal{L}$
un module $\mathcal{O}_X$ inversible. Supposons $S$ quasi-compact
et $f$ de type fini. Les conditions
suivantes sont équivalentes
\begin{enumerate}
\item $\mathcal{L}$ est $f$-ample,
\item $\mathcal{L}^{\otimes d}$ est $f$-très suffisant pour certains $d \geq 1$,
\item $\mathcal{L}^{\otimes d}$ est $f$-très suffisant pour tous les $d \gg 1$.
\end{enumerate}
\end{lemma}

\begin{proof}
Trivialement (3) implique (2). Lemme \ref{lemma-ample-very-ample} garantit que (2) implique
(1) puisqu'un morphisme de type fini est quasi-compact
par définition. Supposons que $\mathcal{L}$ soit $f$-ample.
Choisissons un recouvrement fini par des ouverts affines $S = V_1 \cup \ldots \cup V_m$. Écrivons
$X_i = f^{-1}(V_i)$. Par le Lemme \ref{lemma-finite-type-over-affine-ample-very-ample} ci-dessus on voit qu'il existe
un $d_0$ tel que $\mathcal{L}^{\otimes d}$ est relativement
très ample sur $X_i/V_i$ pour tout $d \geq d_0$. Nous concluons
donc que (1) implique (3) par le lemme \ref{lemma-relatively-very-ample}.
\end{proof}

\noindent
Les deux lemmes suivants fournissent les caractérisations les plus utilisées
et les plus utiles des faisceaux inversibles relativement très amples
et relativement amples lorsque le morphisme est de type fini.

\begin{lemma}
\label{lemma-characterize-very-ample-on-finite-type}
Soit $f : X \to S$ un morphisme de schémas. Soit $\mathcal{L}$
un faisceau inversible sur $X$. Supposons que $f$
soit de type fini. Les conditions
suivantes sont équivalentes :
\begin{enumerate}
\item $\mathcal{L}$ est $f$-relativement très ample, et
\item il existe un recouvrement ouvert $S = \bigcup V_j$, pour chaque
$j$ un entier $n_j$, et des immersions
$$
i_j :
X_j = f^{-1}(V_j) = V_j \times_S X
\longrightarrow
\mathbf{P}^{n_j}_{V_j}
$$
sur $V_j$
tel que $\mathcal{L}|_{X_j} \cong i_j^*\mathcal{O}_{\mathbf{P}^{n_j}_{V_j}}(1)$.
\end{enumerate}
\end{lemma}

\begin{proof}
On voit que (1) implique (2) en prenant un recouvrement ouvert
affine de $S$ et en appliquant le Lemme \ref{lemma-very-ample-finite-type-over-affine} à chacune
des restrictions de $f$
et $\mathcal{L}$. On voit que (2) implique
(1) par le lemme \ref{lemma-relatively-very-ample}.
\end{proof}

\begin{lemma}
\label{lemma-characterize-ample-on-finite-type}
Soit $f : X \to S$ un morphisme de schémas. Soit $\mathcal{L}$
un faisceau inversible sur $X$. Supposons que $f$
soit de type fini. Les conditions
suivantes sont équivalentes :
\begin{enumerate}
\item $\mathcal{L}$ est $f$-relativement ample, et
\item il existe un recouvrement ouvert $S = \bigcup V_j$, pour chaque
$j$ un entier $d_j \geq 1$, $n_j \geq 0$,
et des immersions
$$
i_j :
X_j = f^{-1}(V_j) = V_j \times_S X
\longrightarrow
\mathbf{P}^{n_j}_{V_j}
$$
sur $V_j$ tel que
$\mathcal{L}^{\otimes d_j}|_{X_j} \cong
i_j^*\mathcal{O}_{\mathbf{P}^{n_j}_{V_j}}(1)$.
\end{enumerate}
\end{lemma}

\begin{proof}
On voit que (1) implique (2) en prenant un recouvrement ouvert
affine de $S$ et en appliquant le Lemme \ref{lemma-finite-type-over-affine-ample-very-ample} à
chacune des restrictions de $f$
et $\mathcal{L}$. On voit que (2)
implique (1) par le lemme \ref{lemma-characterize-relatively-ample}.
\end{proof}

\begin{lemma}
\label{lemma-invertible-add-enough-ample-very-ample}
Soit $f : X \to S$ un morphisme de schémas. Soit
$\mathcal{N}$, $\mathcal{L}$ des modules $\mathcal{O}_X$ inversibles. Supposons
que $S$ est quasi-compact, $f$ est de type fini
et $\mathcal{L}$ est
$f$-ample. Alors $\mathcal{N} \otimes_{\mathcal{O}_X} \mathcal{L}^{\otimes d}$ est $f$-très suffisant
pour tous les $d \gg 1$.
\end{lemma}

\begin{proof}
Par le Lemme \ref{lemma-characterize-very-ample-on-finite-type} on se réduit au cas
$S$ est affine. En combinant
Lemme \ref{lemma-finite-type-over-affine-ample-very-ample} et Propriétés, Proposition
\ref{properties-proposition-characterize-ample}, nous pouvons trouver un entier
$d_0$ tel que $\mathcal{N} \otimes \mathcal{L}^{\otimes d_0}$
est généré globalement. Choisissez
les sections globales $s_0, \ldots, s_n$
de $\mathcal{N} \otimes \mathcal{L}^{\otimes d_0}$ qui le génèrent. Ceci détermine
un morphisme $j : X \to \mathbf{P}^n_S$
sur $S$. Par le Lemme
\ref{lemma-finite-type-over-affine-ample-very-ample} on peut aussi choisir un entier
$d_1$ tel que pour tout $d \geq d_1$
il existe des sections $t_{d, 0}, \ldots, t_{d, n(d)}$ de $\mathcal{L}^{\otimes d}$
qui le génèrent et définissent une
immersion
$$
j_d = \varphi_{\mathcal{L}^{\otimes d}, t_{d, 0}, \ldots, t_{d, n(d)}} :
X
\longrightarrow
\mathbf{P}^{n(d)}_S
$$
sur $S$. Alors pour $d \geq d_0 + d_1$ on peut considérer le
morphisme
$$
\varphi_{\mathcal{N} \otimes \mathcal{L}^{\otimes d},
s_j \otimes t_{d - d_0, i}} :
X
\longrightarrow
\mathbf{P}^{(n + 1)(n(d - d_0) + 1) - 1}_S
$$
Ce morphisme est une immersion car c'est la composition
$$
X \to \mathbf{P}^n_S \times_S \mathbf{P}^{n(d - d_0)}_S
\to \mathbf{P}^{(n + 1)(n(d - d_0) + 1) - 1}_S
$$
où le premier morphisme est $(j, j_{d - d_0})$ et
le second est le plongement
de Segre (Constructions, Lemme \ref{constructions-lemma-segre-embedding}). Puisque
$j$ est une immersion, $(j, j_{d - d_0})$
l'est aussi (appliquer le Lemme \ref{lemma-immersion-permanence}). Nous avons une composition
d'immersions et
donc une immersion (Schémas, Lemme \ref{schemes-lemma-composition-immersion}).
\end{proof}










\section{Morphismes quasi-projectifs}
\label{section-quasi-projective}

\noindent
La discussion dans la section précédente suggère les définitions suivantes. Nous reprenons
notre définition du quasi-projectif de \cite{EGA}. La version avec la lettre
``H'' est la définition dans \cite{H}.

\begin{definition}
\label{definition-quasi-projective}
\begin{reference}
\cite[II, Définition 5.3.1]{EGA} et \cite[page 103]{H}
\end{reference}
Soit $f : X \to S$ un morphisme de schémas.
\begin{enumerate}
\item On dit que $f$ est {\it quasi-projectif} si $f$
est de type fini et qu'il existe un $f$-relativement ample module inversible $\mathcal{O}_X$.
\item On dit que $f$ est {\it H-quasi-projectif}
s'il existe une immersion quasi-compacte $X \to \mathbf{P}^n_S$ sur $S$
pour un certain $n$.\footnote{Ce n'est pas exactement la
même définition que celle dans Hartshorne. À savoir, la définition dans Hartshorne
(8e impression corrigée, 1997) est que $f$ devrait être la composition d'une
immersion ouverte suivie d'un H-morphisme projectif (voir Définition \ref{definition-projective}),
ce qui n'implique pas que
$f$ soit quasi-compact. Voir le Lemme \ref{lemma-H-quasi-projective-open-H-projective} pour
l'implication dans l'autre sens.}
\item On dit que $f$ est {\it localement quasi-projectif}
s'il existe un recouvrement ouvert $S = \bigcup V_j$ tel que chaque $f^{-1}(V_j) \to V_j$
soit quasi-projectif.
\end{enumerate}
\end{definition}

\noindent
Comme cette définition le suggère, la propriété d'être quasi-projectif
n'est pas locale sur $S$. À un stade ultérieur, nous
pourrons en dire plus sur la catégorie des schémas
quasi-projectifs, voir Compléments sur les morphismes, section \ref{more-morphisms-section-quasi-projective}.

\begin{lemma}
\label{lemma-base-change-quasi-projective}
Un changement de base d'un morphisme quasi-projectif est quasi-projectif.
\end{lemma}

\begin{proof}
Cela résulte
des Lemme \ref{lemma-base-change-finite-type}
et \ref{lemma-ample-base-change}.
\end{proof}

\begin{lemma}
\label{lemma-composition-quasi-projective}
Soient $f : X \to Y$ et $g : Y \to S$ des morphismes de schémas. Si $S$
est quasi-compact et $f$ et $g$ sont quasi-projectifs,
alors $g \circ f$ est quasi-projectif.
\end{lemma}

\begin{proof}
Cela résulte
des Lemme \ref{lemma-composition-finite-type}
et \ref{lemma-ample-composition}.
\end{proof}

\begin{lemma}
\label{lemma-quasi-projective-properties}
Soit $f : X \to S$ un morphisme de schémas. Si $f$ est quasi-projectif, ou
H-quasi-projectif ou localement quasi-projectif, alors $f$
est séparé de type fini.
\end{lemma}

\begin{proof}
Omis.
\end{proof}

\begin{lemma}
\label{lemma-H-quasi-projective-quasi-projective}
Un H-morphisme quasi-projectif est quasi-projectif.
\end{lemma}

\begin{proof}
Omis.
\end{proof}

\begin{lemma}
\label{lemma-characterize-locally-quasi-projective}
Soit $f : X \to S$ un morphisme de schémas. Les conditions
suivantes sont équivalentes :
\begin{enumerate}
\item Le morphisme $f$ est localement quasi-projectif.
\item Il existe un recouvrement ouvert $S = \bigcup V_j$ tel
que chaque $f^{-1}(V_j) \to V_j$ soit H-quasi-projectif.
\end{enumerate}
\end{lemma}

\begin{proof}
D'après le Lemme \ref{lemma-H-quasi-projective-quasi-projective} nous voyons que (2) implique (1).
Supposons (1). La question est locale sur
$S$ et par conséquent nous pouvons supposer que $S$ est affine,
$X$ de type fini sur
$S$ et $\mathcal{L}$ est un faisceau inversible relativement ample
sur $X/S$. D'après le Lemme \ref{lemma-finite-type-over-affine-ample-very-ample} nous pouvons supposer
que $\mathcal{L}$ est ample sur $X$. D'après
le Lemme \ref{lemma-quasi-projective-finite-type-over-S} nous voyons qu'il existe une immersion de $X$ dans
un espace projectif sur $S$,
c'est-à-dire que $X$ est H-quasi-projectif sur $S$ comme
souhaité.
\end{proof}

\begin{lemma}
\label{lemma-quasi-affine-finite-type-quasi-projective}
\begin{reference}
\cite[II, Proposition 5.3.4 (i)]{EGA}
\end{reference}
Un morphisme quasi-affine de type fini est quasi-projectif.
\end{lemma}

\begin{proof}
Cela résulte du Lemme \ref{lemma-quasi-affine-O-ample}.
\end{proof}

\begin{lemma}
\label{lemma-quasi-projective-permanence}
Soient $g : Y \to S$ et $f : X \to Y$ des morphismes de schémas. Si $g \circ f$
est quasi-projectif et $f$ est
quasi-compact\footnote{Cela découle si $g$ est quasi-séparé
d'après Schémas, Lemme \ref{schemes-lemma-quasi-compact-permanence}.}, alors $f$
est quasi-projectif.
\end{lemma}

\begin{proof}
Observez que $f$ est de
type fini d'après le Lemme \ref{lemma-permanence-finite-type}.
Ainsi, le lemme découle du Lemme \ref{lemma-ample-permanence} et des
définitions.
\end{proof}





\section{Morphismes propres}
\label{section-proper}

\noindent
La notion de morphisme propre joue un rôle important en géométrie algébrique.
Un exemple important de morphisme propre sera le morphisme
de structure $\mathbf{P}^n_S \to S$ de l'espace projectif $n$, et
ceci est en fait l'exemple motivant qui a conduit à la définition.

\begin{definition}
\label{definition-proper}
Soit $f : X \to S$ un morphisme de schémas. Nous disons que $f$
est {\it propre} si $f$ est séparé, de type fini
et universellement fermé.
\end{definition}

\noindent
Le morphisme de la droite affine avec zéro doublé vers la droite affine est de type
fini et universellement fermé, donc la condition de séparation est nécessaire
dans la définition ci-dessus. Dans le
reste de cette section, nous prouvons certaines des propriétés de base des
morphismes propres et des morphismes universellement fermés.

\begin{lemma}
\label{lemma-universally-closed-local-on-the-base}
Soit $f : X \to S$ un morphisme de schémas. Les conditions
suivantes sont équivalentes :
\begin{enumerate}
\item Le morphisme $f$ est universellement fermé.
\item Il existe un recouvrement ouvert $S = \bigcup V_j$ tel que $f^{-1}(V_j) \to V_j$
soit universellement fermé pour tous les indices $j$.
\end{enumerate}
\end{lemma}

\begin{proof}
Ceci est clair d'après la définition.
\end{proof}

\begin{lemma}
\label{lemma-proper-local-on-the-base}
Soit $f : X \to S$ un morphisme de schémas. Les conditions
suivantes sont équivalentes :
\begin{enumerate}
\item Le morphisme $f$ est propre.
\item Il existe un recouvrement ouvert $S = \bigcup V_j$ tel que
$f^{-1}(V_j) \to V_j$ soit propre pour tous les indices $j$.
\end{enumerate}
\end{lemma}

\begin{proof}
Omis.
\end{proof}

\begin{lemma}
\label{lemma-composition-proper}
La composition de morphismes propres est propre. Il en
va de même pour les morphismes universellement fermés.
\end{lemma}

\begin{proof}
Une composition de morphismes fermés est
fermée. Si $X \to Y \to Z$ sont des morphismes universellement
fermés et $Z' \to Z$ est un morphisme
quelconque, alors nous voyons que $Z' \times_Z X = (Z' \times_Z Y) \times_Y X  \to Z' \times_Z Y$ est fermé
et $Z' \times_Z Y \to Z'$ est fermé. D'où le résultat pour
les morphismes universellement fermés. Nous
avons vu que ``séparé'' et ``de
type fini'' sont
conservés sous les compositions (Schémas, Lemme \ref{schemes-lemma-separated-permanence}
et Lemme \ref{lemma-composition-finite-type}). D'où le résultat pour
les morphismes propres.
\end{proof}

\begin{lemma}
\label{lemma-base-change-proper}
Le changement de base d'un morphisme propre est propre. Il
en est de même pour les morphismes universellement fermés.
\end{lemma}

\begin{proof}
Ceci est vrai par définition pour les morphismes universellement
fermés. C'est vrai pour les
morphismes séparés (Schémas, Lemme \ref{schemes-lemma-separated-permanence}). C'est
vrai pour les morphismes de type fini
(Lemme \ref{lemma-base-change-finite-type}). Par conséquent, c'est
vrai pour les morphismes propres.
\end{proof}

\begin{lemma}
\label{lemma-closed-immersion-proper}
Une immersion fermée est propre, donc a fortiori universellement fermée.
\end{lemma}

\begin{proof}
Le changement de base d'une immersion fermée est une
immersion fermée (Schémas, Lemme \ref{schemes-lemma-base-change-immersion}). Par conséquent,
elle est universellement
fermée. Une immersion
fermée est séparée (Schémas, Lemme \ref{schemes-lemma-immersions-monomorphisms}). Une
immersion fermée est de type
fini (Lemme \ref{lemma-immersion-locally-finite-type}). Par conséquent, une
immersion fermée est propre.
\end{proof}

\begin{lemma}
\label{lemma-image-proper-scheme-closed}
Supposons donné un diagramme commutatif de schémas
$$
\xymatrix{
X \ar[rr] \ar[rd] & &
Y \ar[ld] \\
& S &
}
$$
avec $Y$ séparé sur $S$.
\begin{enumerate}
\item Si $X \to S$ est universellement fermé, alors le morphisme $X \to Y$
est universellement fermé.
\item Si $X$ est propre sur $S$, alors le morphisme $X \to Y$ est propre.
\end{enumerate}
En particulier, dans les deux cas, l'image de $X$ dans $Y$ est fermée.
\end{lemma}

\begin{proof}
Supposons que $X \to S$ soit universellement fermé (resp. \ propre).
Nous factorisons le morphisme comme
$X \to X \times_S Y \to Y$. Le premier morphisme est une immersion
fermée, voir Schémas, Lemme \ref{schemes-lemma-semi-diagonal}. Par
conséquent, le premier morphisme est propre (Lemme \ref{lemma-closed-immersion-proper}). La projection $X \times_S Y \to Y$
est le changement de base d'un morphisme universellement
fermé (resp. \ propre) et donc universellement
fermé (resp. \ propre), voir Lemme \ref{lemma-base-change-proper}. Ainsi $X \to Y$
est universellement fermé (resp. \ propre) comme la composition
de morphismes universellement fermés (resp.
\ propres) (Lemme \ref{lemma-composition-proper}).
\end{proof}

\noindent
La démonstration du lemme suivant est due à Bjorn Poonen ; voir
\href{https://mathoverflow.net/questions/23337/is-a-universally-closed-morphism-of-schemes-quasi-compact/23528#23528}{cette page}.

\begin{lemma}
\label{lemma-universally-closed-quasi-compact}
\begin{reference}
Due à Bjorn Poonen.
\end{reference}
Un morphisme universellement fermé de schémas est quasi-compact.
\end{lemma}

\begin{proof}
Soit $f : X \to S$ un morphisme. Supposons que $f$ n'est pas quasi-compact.
Notre objectif est de montrer que $f$ n'est pas universellement
fermé. D'après Schémas, Lemme \ref{schemes-lemma-quasi-compact-affine} il existe
un ouvert affine $V \subset S$ tel que $f^{-1}(V)$ n'est pas quasi-compact.
Pour atteindre notre objectif il suffit de montrer que $f^{-1}(V) \to V$
n'est pas universellement fermé, donc nous pouvons supposer que $S = \Spec(A)$
pour un certain anneau $A$.

\medskip\noindent
Écrivons $X = \bigcup_{i \in I} X_i$ où les $X_i$ sont des sous-schémas
ouverts affines de $X$. Soit
$T = \Spec(A[y_i ; i \in I])$. Soit $T_i = D(y_i) \subset T$. Soit $Z$
l'ensemble fermé $(X \times_S T) - \bigcup_{i \in I} (X_i \times_S T_i)$. Il suffit de prouver que
l'image $f_T(Z)$ de $Z$ sous $f_T : X \times_S T \to T$ n'est pas
fermée.

\medskip\noindent Il
existe un point $s \in S$ tel qu'il n'existe aucun voisinage
$U$ de $s$ dans $S$ tel que $X_U$ soit quasi-compact.
Sinon, nous pourrions recouvrir $S$ avec un nombre fini
de tels $U$ et Schémas, le Lemme \ref{schemes-lemma-quasi-compact-affine} impliquerait
que $f$ est quasi-compact. Fixons un tel $s \in S$.

\medskip\noindent
D'abord, nous vérifions que $f_T(Z_s) \ne T_s$. Soit $t \in T$ le point situé
au-dessus de $s$ avec $\kappa(t) = \kappa(s)$ tel que $y_i = 1$ dans
$\kappa(t)$ pour tous
$i$. Alors $t \in T_i$ pour tous $i$, et la fibre de
$Z_s \to T_s$ au-dessus de $t$ est isomorphe à $(X - \bigcup_{i \in I} X_i)_s$, qui est
vide. Ainsi $t \in T_s - f_T(Z_s)$.

\medskip\noindent Supposons
que $f_T(Z)$ soit fermé dans $T$. Alors il existe
un élément $g \in A[y_i; i \in I]$ avec $f_T(Z) \subset V(g)$ mais $t \not \in V(g)$. Par conséquent, l'image de
$g$ dans $\kappa(t)$ est non nulle. En particulier, un certain coefficient
de $g$ a une image non nulle dans $\kappa(s)$. Par conséquent, ce coefficient
est inversible sur un certain voisinage $U$ de $s$. Soit $J$
l'ensemble fini de $j \in I$ tel que $y_j$ apparaisse dans $g$.
Puisque $X_U$ n'est pas quasi-compact, nous pouvons choisir un point $x \in X - \bigcup_{j \in J} X_j$
situé au-dessus de certain $u \in U$. Puisque $g$ a un coefficient
qui est inversible sur $U$, nous pouvons trouver un point $t' \in T$
situé au-dessus de $u$ tel que $t' \not \in V(g)$ et $t' \in V(y_i)$
pour tous $i \notin J$. Ceci est vrai parce que $V(y_i; i \in I, i \not\in J) = \Spec(A[y_j; j\in J])$.
et l'ensemble des points de ce schéma situés au-dessus de $u$
est bijectif avec $\Spec(\kappa(u)[y_j; j \in J])$. En d'autres termes $t' \notin T_i$ pour
chaque $i \notin J$. D'après
Schémas, Lemme \ref{schemes-lemma-points-fibre-product}, nous pouvons trouver
un point $z$ de $X \times_S T$ se projetant sur $x \in X$ et sur
$t' \in T$. Puisque $x \not \in X_j$ pour $j \in J$ et $t' \not \in T_i$ pour $i \in I \setminus J$,
nous voyons que $z \in Z$. D'autre part, $f_T(z) = t' \not \in V(g)$,
ce qui contredit $f_T(Z) \subset V(g)$. Ainsi, l'hypothèse ``$f_T(Z)$
clôt'' est fausse et nous concluons en effet que $f_T$
n'est pas fermée, comme souhaité.
\end{proof}

\noindent
Le lemme suivant dit que l'image d'un schéma propre (dans un schéma séparé
de type fini sur la base) est propre.

\begin{lemma}
\label{lemma-image-proper-is-proper}
Soit $S$ un schéma. Soit $f : X \to Y$ un morphisme de schémas au-dessus de $S$. Si $X$ est
universellement fermé au-dessus de $S$ et que $f$ est surjectif, alors $Y$
est universellement fermé au-dessus de $S$. En particulier, si également $Y$ est séparé
et localement de type fini sur $S$, alors $Y$ est propre sur $S$.
\end{lemma}

\begin{proof}
Supposons que $X$ soit universellement fermé
et que $f$ soit surjectif. Notons $p : X \to S$,
$q : Y \to S$ les morphismes de structure. Soit $S' \to S$
un morphisme de schémas. Le changement
de base $f' : X_{S'} \to Y_{S'}$ est surjectif (Lemme \ref{lemma-base-change-surjective}),
et le changement de base $p' : X_{S'} \to S'$
est fermé. Si $T \subset Y_{S'}$ est fermé, alors $(f')^{-1}(T) \subset X_{S'}$
est fermé, donc $p'((f')^{-1}(T)) = q'(T)$ est fermé. Ainsi $q'$
est fermé. Cela prouve la première affirmation.
Donc $Y \to S$ est quasi-compact
par le Lemme \ref{lemma-universally-closed-quasi-compact} et par conséquent
$Y \to S$ est propre par définition
si de plus $Y \to S$ est localement de type fini et séparé.
\end{proof}

\begin{lemma}
\label{lemma-scheme-theoretic-image-is-proper}
Supposons donné un diagramme commutatif de schémas
$$
\xymatrix{
X \ar[rr]_h \ar[rd]_f & & Y \ar[ld]^g \\
& S
}
$$
Supposons
\begin{enumerate}
\item $X \to S$ est un morphisme universellement fermé (par exemple propre), et
\item $Y \to S$ est séparé et localement de type fini.
\end{enumerate}
Alors l'image schématique $Z \subset Y$ de $h$ est propre
sur $S$ et $X \to Z$ est surjectif.
\end{lemma}

\begin{proof}
L'image schématique de $h$ est construite
dans la section \ref{section-scheme-theoretic-image}. Puisque $f$ est
quasi-compact (Lemme \ref{lemma-universally-closed-quasi-compact}), nous
trouvons que $h$ est
quasi-compact (Schémas, Lemme \ref{schemes-lemma-quasi-compact-permanence}). Par
conséquent, $h(X) \subset Z$
est dense (Lemme \ref{lemma-quasi-compact-scheme-theoretic-image}). D'autre part, $h(X)$
est fermé dans $Y$
(Lemme \ref{lemma-image-proper-scheme-closed}), donc $X \to Z$
est surjectif. Ainsi,
$Z \to S$ est un morphisme propre (Lemme \ref{lemma-image-proper-is-proper}).
\end{proof}

\noindent
Le but d'un schéma séparé sous un morphisme surjectif
universellement fermé est séparé.

\begin{lemma}
\label{lemma-image-universally-closed-separated}
Soit $S$ un schéma. Soit $f : X \to Y$ un morphisme surjectif universellement fermé
de schémas sur $S$.
\begin{enumerate}
\item Si $X$ est quasi-séparé, alors $Y$ est quasi-séparé.
\item Si $X$ est séparé, alors $Y$ est séparé.
\item Si $X$ est quasi-séparé sur $S$, alors $Y$ est quasi-séparé sur $S$.
\item Si $X$ est séparé sur $S$, alors $Y$ est séparé sur $S$.
\end{enumerate}
\end{lemma}

\begin{proof}
Les parties (1) et (2) sont une conséquence de (3)
et (4) pour $S = \Spec(\mathbf{Z})$
(voir Schémas, Définition \ref{schemes-definition-separated}). Considérons
le diagramme commutatif
$$
\xymatrix{
X \ar[d] \ar[rr]_{\Delta_{X/S}} & & X \times_S X \ar[d] \\
Y \ar[rr]^{\Delta_{Y/S}} & & Y \times_S Y
}
$$
La flèche verticale de gauche est surjective (c’est-à-dire universellement
surjective). La flèche verticale de droite est universellement
fermée comme composition
des morphismes universellement fermés $X \times_S X \to X \times_S Y \to Y \times_S Y$. Par conséquent,
elle est également
quasi-compacte, voir Lemme \ref{lemma-universally-closed-quasi-compact}.

\medskip\noindent
Supposons que $X$ soit quasi-séparé sur $S$, c'est-à-dire
que $\Delta_{X/S}$ soit quasi-compact. Si $V \subset Y \times_S Y$ est
un ouvert quasi-compact, alors $V \times_{Y \times_S Y} X \to \Delta_{Y/S}^{-1}(V)$ est surjectif et $V \times_{Y \times_S Y} X$
est quasi-compact d'après nos remarques ci-dessus. Nous concluons que
$\Delta_{Y/S}$ est quasi-compact, c'est-à-dire que $Y$ est quasi-séparé
sur $S$.

\medskip\noindent
Supposons que $X$ soit séparé sur $S$, c'est-à-dire que
$\Delta_{X/S}$ soit une immersion fermée. Alors $X \to Y \times_S Y$
est fermé comme composition de morphismes fermés.
Puisque $X \to Y$ est surjectif, il s'ensuit que $\Delta_{Y/S}(Y)$ est fermé dans
$Y \times_S Y$. Ainsi $Y$ est séparé sur $S$ d'après
la discussion suivant Schémas, Définition \ref{schemes-definition-separated}.
\end{proof}








\section{Critères valuatifs}
\label{section-valuative-criteria}

\noindent
Nous avons déjà discuté du critère valuatif pour la fermeture universelle
et pour la séparabilité dans
les schémas, sections \ref{schemes-section-valuative-criterion-universal-closedness} et \ref{schemes-section-valuative-separatedness}. Dans cette
section, nous discuterons de quelques
conséquences et variantes. Dans Limites, section \ref{limits-section-Noetherian-valuative-criterion},
nous montrerons qu'il suffit de considérer les anneaux
de valuation discrète lorsqu'on travaille avec
des schémas localement noethériens et des morphismes
de type fini.

\begin{lemma}[Critère valuatif pour la propreté]
\label{lemma-characterize-proper}
\begin{reference}
\cite[II Théorème 7.3.8]{EGA}
\end{reference}
Soit $S$ un schéma. Soit $f : X \to Y$ un morphisme de schémas sur $S$.
Supposons que $f$ soit de type fini et quasi-séparé. Alors les
conditions suivantes sont équivalentes
\begin{enumerate}
\item $f$ est propre,
\item $f$ satisfait le critère
valuatif (Schémas, Définition \ref{schemes-definition-valuative-criterion}),
\item donné tout diagramme solide commutatif
$$
\xymatrix{
\Spec(K) \ar[r] \ar[d] & X \ar[d] \\
\Spec(A) \ar[r] \ar@{-->}[ru] & Y
}
$$
où $A$ est un anneau de valuation avec corps de fractions $K$, il existe une
flèche pointillée unique rendant le diagramme commutatif.
\end{enumerate}
\end{lemma}

\begin{proof}
La partie (3) est une reformulation de
(2). Ainsi,
le lemme est une conséquence formelle de Schémas,
Proposition \ref{schemes-proposition-characterize-universally-closed} et Lemme
\ref{schemes-lemma-separated-implies-valuative}, Lemme \ref{schemes-lemma-valuative-criterion-separatedness}, et des
définitions.
\end{proof}

\noindent
On n'a généralement pas besoin de considérer tous les diagrammes possibles
lorsqu'on teste le critère valuatif. Nous appellerons un critère valuatif
comme dans le lemme suivant un ``critère valuatif raffiné''.

\begin{lemma}
\label{lemma-refined-valuative-criterion-universally-closed}
Soient $f : X \to S$ et $h : U \to X$ des morphismes de schémas. Supposons que $f$
et $h$ soient quasi-compacts et que $h(U)$ soit dense dans $X$. Si,
pour tout diagramme commutatif solide donné
$$
\xymatrix{
\Spec(K) \ar[r] \ar[d] & U \ar[r]^h & X \ar[d]^f \\
\Spec(A) \ar[rr] \ar@{-->}[rru] & & S
}
$$
où $A$ est un anneau de valuation avec corps des fractions $K$, il
existe un unique trait pointillé rendant le diagramme commutatif, alors $f$
est universellement fermé. Si de plus $f$ est quasi-séparé, alors $f$
est séparé.
\end{lemma}

\begin{proof}
Pour prouver que $f$ est universellement fermé, nous allons vérifier
la partie existence du critère valuatif pour
$f$, ce qui suffit d'après Schémas, Proposition \ref{schemes-proposition-characterize-universally-closed}. Pour ce faire,
considérons un diagramme commutatif
$$
\xymatrix{
\Spec(K) \ar[r] \ar[d] & X \ar[d] \\
\Spec(A) \ar[r] & S
}
$$
où $A$ est un anneau de valuation et $K$ est le corps des fractions
de $A$. Notez que, puisque les anneaux de valuation et
les corps sont réduits, nous pouvons remplacer $U$, $X$
et $S$ par leurs réductions respectives d'après Schémas,
Lemme \ref{schemes-lemma-map-into-reduction}. Dans ce cas, l'hypothèse selon laquelle
$h(U)$ est dense signifie que l'image schématique de $h : U \to X$
est $X$, voir Lemme \ref{lemma-scheme-theoretic-image-reduced}. Nous pouvons
aussi remplacer $S$ par un ouvert affine par lequel
le morphisme $\Spec(A) \to S$ se factorise. Ainsi, nous pouvons
supposer que $S = \Spec(R)$.

\medskip\noindent
Soit $\Spec(B) \subset X$ un ouvert affine par lequel le
morphisme $\Spec(K) \to X$ se factorise. Choisissons une algèbre
polynomiale $P$ sur $B$ et une surjection de
$B$-algèbres $P \to K$. Alors $\Spec(P) \to X$
est plat. Par conséquent, l'image schématique du morphisme
$U \times_X \Spec(P) \to \Spec(P)$ est $\Spec(P)$ par le Lemme \ref{lemma-flat-base-change-scheme-theoretic-image}.
D'après le Lemme \ref{lemma-reach-points-scheme-theoretic-image} nous pouvons trouver
un diagramme commutatif
$$
\xymatrix{
\Spec(K') \ar[r] \ar[d] & U \times_X \Spec(P) \ar[d] \\
\Spec(A') \ar[r] & \Spec(P)
}
$$
où $A'$ est un anneau de valuation et $K'$ est le corps des fractions de $A'$,
tel que le point fermé de $\Spec(A')$ s'envoie
dans $\Spec(K) \subset \Spec(P)$. En d'autres termes,
il existe un homomorphisme de $B$-algèbres $\varphi : K \to A'/\mathfrak m_{A'}$.
Choisissons un anneau de valuation $A'' \subset A'/\mathfrak m_{A'}$ dominant
$\varphi(A)$ et ayant pour corps des fractions
$K'' = A'/\mathfrak m_{A'}$ (Algèbre, Lemme \ref{algebra-lemma-dominate}). On pose
$$
C = \{\lambda \in A' \mid \lambda \bmod \mathfrak m_{A'} \in A''\}.
$$
qui est un anneau de valuation
par l'Algèbre, Lemme \ref{algebra-lemma-stack-valuation-rings}. Comme $C$ est une $R$-algèbre
avec corps des fractions $K'$, nous obtenons un diagramme
commutatif
$$
\xymatrix{
\Spec(K') \ar[r] \ar[d] & U \ar[r] & X \ar[d] \\
\Spec(C) \ar[rr] \ar@{-->}[rru] & & S
}
$$
comme dans l'énoncé du lemme. Ainsi, une flèche pointillée
s'insérant dans le diagramme comme indiqué. D'après
l'hypothèse d'unicité du lemme, la composition $\Spec(A') \to \Spec(C) \to X$
coïncide avec le morphisme donné $\Spec(A') \to \Spec(P) \to \Spec(B) \subset X$. Par
conséquent, la restriction du morphisme au spectre
de $C/\mathfrak m_{A'} = A''$ induit le morphisme donné
$\Spec(K'') = \Spec(A'/\mathfrak m_{A'}) \to \Spec(K) \to X$. Soit $x \in X$ l'image du point fermé
de $\Spec(A'') \to X$. L'image de l'homomorphisme d'anneaux induit
$\mathcal{O}_{X, x} \to A''$ est un sous-anneau local contenu dans
$K \subset K''$. Comme $A$ est maximal pour
la relation de domination dans $K$ et comme
$A \subset A''$, nous avons $A = K \cap A''$. Nous concluons
que $\mathcal{O}_{X, x} \to A''$ se factorise par $A \subset A''$. De cette manière,
nous obtenons la flèche désirée $\Spec(A) \to X$.

\medskip\noindent
Enfin, supposons que $f$ soit quasi-séparé. Alors $\Delta : X \to X \times_S X$ est quasi-compact.
Étant donné un diagramme solide
$$
\xymatrix{
\Spec(K) \ar[r] \ar[d] & U \ar[r]^h & X \ar[d]^\Delta \\
\Spec(A) \ar[rr] \ar@{-->}[rru] & & X \times_S X
}
$$
où $A$ est un anneau de valuation avec corps de fractions $K$,
il existe un unique flèche pointillée rendant le diagramme commutatif.
À savoir, la flèche horizontale inférieure est la même chose qu'une paire
de morphismes $\Spec(A) \to X$ qui peut servir de flèche pointillée dans le diagramme
du lemme. Ainsi, l'unicité requise montre que la flèche horizontale
inférieure se factorise par $\Delta$. Par conséquent,
nous pouvons appliquer le résultat que nous
venons de démontrer à $\Delta : X \to X \times_S X$ et $h : U \to X$ et conclure que $\Delta$
est universellement fermé. Clairement, cela signifie que $f$
est séparé.
\end{proof}

\begin{remark}
\label{remark-check-val-on-open}
L'hypothèse sur l'unicité des flèches
pointillées dans le Lemme \ref{lemma-refined-valuative-criterion-universally-closed} est nécessaire
(détails omis). Bien sûr, l'unicité est garantie
si $f$
est séparé (Schémas, Lemme \ref{schemes-lemma-separated-implies-valuative}).
\end{remark}

\begin{lemma}
\label{lemma-morphism-defined-local-ring}
Soit $S$ un schéma. Soient $X$, $Y$ des schémas sur $S$.
Soient $s \in S$ et $x \in X$, $y \in Y$ des points au-dessus de $s$.
\begin{enumerate}
\item Soient $f, g : X \to Y$ des morphismes sur $S$
tels que
$f(x) = g(x) = y$ et $f^\sharp_x = g^\sharp_x : \mathcal{O}_{Y, y} \to \mathcal{O}_{X, x}$. Alors il existe un voisinage
ouvert $U \subset X$ contenant $f|_U = g|_U$
dans les cas suivants
\begin{enumerate}
\item $Y$ est localement de type fini sur $S$,
\item $X$ est intègre,
\item $X$ est localement noethérien, ou
\item $X$ est réduit avec un nombre fini de composantes irréductibles.
\end{enumerate}
\item Soit $\varphi : \mathcal{O}_{Y, y} \to \mathcal{O}_{X, x}$ un homomorphisme local
de $\mathcal{O}_{S, s}$-algèbres. Alors il existe un voisinage ouvert $U \subset X$
de $x$ et un morphisme $f : U \to Y$ envoyant $x$
vers $y$ avec $f^\sharp_x = \varphi$ dans les cas suivants
\begin{enumerate}
\item $Y$ est localement de présentation finie sur $S$,
\item $Y$ est localement de type fini et $X$ est intègre,
\item $Y$ est localement de type fini et $X$ est localement noethérien, ou
\item $Y$ est localement de type fini et $X$ est réduit avec un nombre
fini de composantes irréductibles.
\end{enumerate}
\end{enumerate}
\end{lemma}

\begin{proof}
Démonstration de (1). Nous pouvons remplacer $X$, $Y$, $S$ par des voisinages
ouverts affines appropriés de $x$, $y$, $s$ et réduire au problème algébrique
suivant : étant donné un anneau $R$, deux applications d'algèbres $R$ $\varphi, \psi : B \to A$
telles que
\begin{enumerate}
\item $R \to B$ est de type fini, ou $A$ est un domaine, ou $A$
est noethérien, ou $A$ est réduit et a un nombre fini d'idéaux premiers minimaux,
\item les deux applications $B \to A_\mathfrak p$ sont les mêmes pour
un certain premier $\mathfrak p \subset A$,
\end{enumerate}
montrez que $\varphi, \psi$ définit la même application $B \to A_g$
pour un $g \in A$, $g \not \in \mathfrak p$ approprié. Si $R \to B$ est de type
fini, soit $t_1, \ldots, t_m \in B$ des générateurs de $B$ comme $R$-algèbre.
Pour chaque $j$ nous
pouvons trouver $g_j \in A$, $g_j \not \in \mathfrak p$
tels que $\varphi(t_j)$ et $\psi(t_j)$ aient la même image dans
$A_{g_j}$. Alors nous définissons $g = \prod g_j$.
Dans les autres cas (si $A$ est un domaine, de
type de Noether ou réduit avec un nombre fini de premiers minimaux),
nous pouvons trouver un $g \in A$, $g \not \in \mathfrak p$ tels que
$A_g \subset A_\mathfrak p$. Voir Algèbre, Lemme \ref{algebra-lemma-subring-of-local-ring}. Ainsi les applications
$B \to A_g$ sont égales comme souhaité.

\medskip\noindent
Démonstration de (2). Pour ce faire, nous pouvons remplacer $X$, $Y$ et $S$
par des ouverts affines appropriés. Disons $X = \Spec(A)$,
$Y = \Spec(B)$ et $S = \Spec(R)$. Soit $\mathfrak p \subset A$ l'idéal premier correspondant
à $x$. Soit $\mathfrak q \subset B$ l'idéal premier correspondant à $y$.
Alors $\varphi$ est un homomorphisme local de $R$-algèbres
$\varphi : B_\mathfrak q \to A_\mathfrak p$. Si $R \to B$
est un homomorphisme d'anneaux de présentation finie, alors il existe
un $g \in A \setminus \mathfrak p$ et un homomorphisme de $R$-algèbres $B \to A_g$ tel que
$$
\xymatrix{
B_\mathfrak q \ar[r]_\varphi & A_\mathfrak p \\
B \ar[u] \ar[r] & A_g \ar[u]
}
$$
commute, voir
Algèbre, Lemmes \ref{algebra-lemma-characterize-finite-presentation} et \ref{algebra-lemma-localization-colimit}. Le morphisme induit
$\Spec(A_g) \to \Spec(B)$ fonctionne. Si $B$
est de type fini sur $R$, soit $t_1, \ldots, t_m \in B$
des générateurs de $B$ comme une $R$-algèbre.
Alors nous pouvons choisir $g_j \in A$, $g_j \not \in \mathfrak p$
de telle sorte que $\varphi(t_j) \in \Im(A_{g_j} \to A_\mathfrak p)$.
Ainsi, après avoir remplacé $A$ par $A[1/\prod g_j]$,
nous pouvons supposer que $B$ s'envoie
dans l'image de $A \to A_\mathfrak p$. Si nous pouvons
trouver un $g \in A$, $g \not \in \mathfrak p$ de telle
sorte que $A_g \to A_\mathfrak p$ soit injectif, alors
nous obtiendrons l'application de $R$-algèbre
désirée $B \to A_g$. Une nouvelle application du
Lemme \ref{algebra-lemma-subring-of-local-ring} d'Algèbre achève donc la démonstration.
\end{proof}

\begin{lemma}
\label{lemma-extend-across}
Soit $S$ un schéma. Soient $X$, $Y$ des schémas sur $S$. Soit
$x \in X$. Soit $U \subset X$ un ouvert et soit $f : U \to Y$ un morphisme sur $S$.
Supposons
\begin{enumerate}
\item $x$ est dans l'adhérence de $U$,
\item $X$ est réduit avec un nombre fini de composantes irréductibles ou $X$
est noethérien,
\item $\mathcal{O}_{X, x}$ est un anneau de valuation,
\item $Y \to S$ est propre
\end{enumerate}
Alors il existe un ouvert $U \subset U' \subset X$ contenant $x$ et un
$S$-morphisme $f' : U' \to Y$ prolongeant $f$.
\end{lemma}

\begin{proof}
Il est sans danger de remplacer $X$ par un voisinage ouvert de
$x$ dans $X$ (petit détail
omis). Par les Propriétés, Lemme \ref{properties-lemma-ring-affine-open-injective-local-ring}, nous pouvons supposer
que $X$ est affine
avec $\Gamma(X, \mathcal{O}_X) \subset \mathcal{O}_{X, x}$. En particulier $X$ est intègre
avec un point générique unique $\xi$ dont le corps résiduel
est le corps des fractions $K$ de l'anneau
de valuation $\mathcal{O}_{X, x}$. Puisque
$x$ est dans l'adhérence de $U$, nous voyons
que $U$ n'est pas vide, donc $U$ contient $\xi$.
Ainsi, par le critère valuatif de propreté
(Lemme \ref{lemma-characterize-proper}), il existe un morphisme $t : \Spec(\mathcal{O}_{X, x}) \to Y$
s'insérant dans un diagramme commutatif
$$
\xymatrix{
\Spec(K) \ar[d]_\xi \ar[r] & \Spec(\mathcal{O}_{X, x}) \ar[d]_t \\
U \ar[r]^f & Y
}
$$
de morphismes de schémas sur $S$. En
appliquant le Lemme \ref{lemma-morphism-defined-local-ring} avec $y = t(x)$
et $\varphi = t^\sharp_x$, nous obtenons un voisinage ouvert $V \subset X$
de $x$ et un morphisme $g : V \to Y$ sur $S$ qui
envoie $x$ vers $y$ et tel que $g^\sharp_x = t^\sharp_x$. Comme $Y \to S$
est séparé, l'égaliseur $E$ de $f|_{U \cap V}$ et $g|_{U \cap V}$
est un sous-schéma fermé de $U \cap V$, voir Schémas, Lemme
\ref{schemes-lemma-where-are-they-equal}. Puisque $f$ et $g$ déterminent
le même morphisme $\Spec(K) \to Y$ par construction, nous voyons
que $E$ contient le point générique du schéma intègre
$U \cap V$. Ainsi $E = U \cap V$ et Nous concluons que $f$
et $g$ se recollent pour donner un morphisme $U' = U \cup V \to Y$
comme souhaité.
\end{proof}





\section{Morphismes projectifs}
\label{section-projective}

\noindent
Nous utiliserons la définition d'un morphisme projectif de \cite{EGA}. La version
de la définition avec le ``H'' est celle de
\cite{H}. Les définitions résultantes sont différentes. Les deux sont utiles.

\begin{definition}
\label{definition-projective}
Soit $f : X \to S$ un morphisme de schémas.
\begin{enumerate}
\item On dit que $f$ est {\it projectif}
si $X$ est isomorphe en tant que $S$-schéma à
un sous-schéma fermé d'un fibré projectif
$\mathbf{P}(\mathcal{E})$ pour un certain $\mathcal{O}_S$-module quasi-cohérent de type fini $\mathcal{E}$.
\item On dit que $f$ est {\it H-projectif} s'il existe un entier
$n$ et une immersion fermée $X \to \mathbf{P}^n_S$ sur $S$.
\item On dit que $f$ est {\it localement projectif}
s'il existe un recouvrement ouvert $S = \bigcup U_i$ tel que chaque $f^{-1}(U_i) \to U_i$ soit
projectif.
\end{enumerate}
\end{definition}

\noindent
Comme prévu, un morphisme projectif est quasi-projectif,
voir le Lemme \ref{lemma-projective-quasi-projective}. Inversement,
les morphismes quasi-projectifs sont souvent des compositions
d'immersions ouvertes et de morphismes
projectifs, voir le Lemme \ref{lemma-quasi-projective-open-projective}. Pour
un aperçu des propriétés des morphismes projectifs
sur une base quasi-projective,
voir Compléments sur les morphismes, section \ref{more-morphisms-section-projective}.

\begin{example}
\label{example-projective}
Soit $S$ un schéma. Soit $\mathcal{A}$ une $\mathcal{O}_S$-algèbre
quasi-cohérente graduée générée par $\mathcal{A}_1$ sur $\mathcal{A}_0$.
Supposons de plus que $\mathcal{A}_1$ soit de type fini
sur $\mathcal{O}_S$. Posons $X = \underline{\text{Proj}}_S(\mathcal{A})$. Dans ce cas, $X \to S$
est projectif. À savoir, le morphisme
associé à l'application de $\mathcal{O}_S$-algèbres graduées
$$
\text{Sym}_{\mathcal{O}_X}^*(\mathcal{A}_1)
\longrightarrow
\mathcal{A}
$$
est une immersion
fermée,
voir Constructions, Lemme \ref{constructions-lemma-surjective-generated-degree-1-map-relative-proj}.
\end{example}
 
\begin{lemma}
\label{lemma-H-projective}
Un H-morphisme projectif est H-quasi-projectif. Un
H-morphisme projectif est projectif.
\end{lemma}

\begin{proof}
La première affirmation est immédiate d'après les
définitions. La seconde tient car $\mathbf{P}^n_S$ est un fibré projectif
sur $S$, voir Constructions, Lemme \ref{constructions-lemma-projective-space-bundle}.
\end{proof}

\begin{lemma}
\label{lemma-characterize-locally-projective}
Soit $f : X \to S$ un morphisme de schémas. Les conditions
suivantes sont équivalentes :
\begin{enumerate}
\item Le morphisme $f$ est localement projectif.
\item Il existe un recouvrement ouvert $S = \bigcup U_i$ tel
que chaque $f^{-1}(U_i) \to U_i$ soit H-projectif.
\end{enumerate}
\end{lemma}

\begin{proof}
D'après le Lemme \ref{lemma-H-projective}, nous voyons que (2) implique (1). Supposons (1).
Pour chaque point $s \in S$ nous pouvons trouver $\Spec(R) = U \subset S$ un
voisinage ouvert affine de $s$ tel que $X_U$ soit isomorphe
à un sous-schéma fermé de $\mathbf{P}(\mathcal{E})$ pour un certain faisceau quasi-cohérent
de type fini de $\mathcal{O}_U$-modules $\mathcal{E}$. Écrivons $\mathcal{E} = \widetilde{M}$
pour un certain module $R$ de type fini $M$
(voir Propriétés,
Lemme \ref{properties-lemma-finite-type-module}). Choisissons des générateurs $x_0, \ldots, x_n \in M$
de $M$ en tant que $R$-module. Considérons l'application
surjective de $R$-algèbres graduées
$$
R[X_0, \ldots, X_n] \longrightarrow \text{Sym}_R(M).
$$
Selon Constructions,
Lemme \ref{constructions-lemma-surjective-graded-rings-map-proj} le morphisme
correspondant
$$
\mathbf{P}(\mathcal{E}) \to \mathbf{P}^n_R
$$
est une immersion fermée. Par conséquent, nous concluons que $f^{-1}(U)$ est isomorphe à
un sous-schéma fermé de $\mathbf{P}^n_U$ (en tant que schéma sur $U$). En d'autres
termes : (2) est vérifié.
\end{proof}

\begin{lemma}
\label{lemma-locally-projective-proper}
Un morphisme localement projectif est propre.
\end{lemma}

\begin{proof}
Soit $f : X \to S$ localement projectif.
Pour montrer que $f$ est propre, nous
pouvons travailler localement sur la
base, voir le Lemme \ref{lemma-proper-local-on-the-base}. Par conséquent,
d'après le Lemme \ref{lemma-characterize-locally-projective} ci-dessus, nous pouvons supposer
qu'il existe une immersion
fermée $X \to \mathbf{P}^n_S$. D'après les Lemmes \ref{lemma-composition-proper}
et \ref{lemma-closed-immersion-proper}, il suffit de prouver
que $\mathbf{P}^n_S \to S$ est propre.
Puisque $\mathbf{P}^n_S \to S$ est le changement de base de
$\mathbf{P}^n_{\mathbf{Z}} \to \Spec(\mathbf{Z})$, il suffit de montrer que $\mathbf{P}^n_{\mathbf{Z}} \to \Spec(\mathbf{Z})$
est propre, voir le Lemme \ref{lemma-base-change-proper}.
Par Constructions, Lemme \ref{constructions-lemma-proj-separated} le schéma $\mathbf{P}^n_{\mathbf{Z}}$
est séparé. Par Constructions,
Lemme \ref{constructions-lemma-proj-quasi-compact} le schéma $\mathbf{P}^n_{\mathbf{Z}}$ est quasi-compact.
Il est clair que $\mathbf{P}^n_{\mathbf{Z}} \to \Spec(\mathbf{Z})$
est localement de type fini puisque $\mathbf{P}^n_{\mathbf{Z}}$
est recouvert par les ouverts affines $D_{+}(X_i)$,
chacun d'eux étant le spectre de l'algèbre
de type fini $\mathbf{Z}$.
$$
\mathbf{Z}[X_0/X_i, \ldots, X_n/X_i].
$$
Enfin, nous devons
montrer que $\mathbf{P}^n_{\mathbf{Z}} \to \Spec(\mathbf{Z})$ est universellement
fermé. Cela
résulte de Constructions, Lemme \ref{constructions-lemma-proj-valuative-criterion}
et du critère valuatif
(voir Schémas, Proposition \ref{schemes-proposition-characterize-universally-closed}).
\end{proof}

\begin{lemma}
\label{lemma-proper-ample-locally-projective}
Soit $f : X \to S$ un morphisme propre de schémas. S'il existe un $f$-faisceau
inversible ample sur $X$, alors $f$ est localement projectif.
\end{lemma}

\begin{proof}
S'il existe un $f$-faisceau inversible ample, alors
nous pouvons localement sur $S$ trouver une
immersion $i : X \to \mathbf{P}^n_S$, voir Lemme \ref{lemma-finite-type-over-affine-ample-very-ample}. Comme $X \to S$
est propre, le morphisme $i$ est une immersion
fermée, voir Lemme \ref{lemma-image-proper-scheme-closed}.
\end{proof}

\begin{lemma}
\label{lemma-H-projective-composition}
Une composition de H-morphismes projectifs est H-projective.
\end{lemma}

\begin{proof}
Supposons que $X \to Y$ et $Y \to Z$ soient H-projectifs.
Alors il existe des immersions fermées $X \to \mathbf{P}^n_Y$ sur
$Y$, et $Y \to \mathbf{P}^m_Z$ sur $Z$. Considérons
le diagramme suivant
$$
\xymatrix{
X \ar[r] \ar[d] &
\mathbf{P}^n_Y \ar[r] \ar[dl] &
\mathbf{P}^n_{\mathbf{P}^m_Z} \ar[dl] \ar@{=}[r] &
\mathbf{P}^n_Z \times_Z \mathbf{P}^m_Z \ar[r] &
\mathbf{P}^{nm + n + m}_Z \ar[ddllll] \\
Y \ar[r] \ar[d] & \mathbf{P}^m_Z \ar[dl] & \\
Z & &
}
$$
Ici, la flèche horizontale supérieure la plus à droite est l'immersion
de Segre, voir Constructions, Lemme \ref{constructions-lemma-segre-embedding}. Le diagramme
identifie $X$
comme un sous-schéma fermé de $\mathbf{P}^{nm + n + m}_Z$ comme souhaité.
\end{proof}

\begin{lemma}
\label{lemma-H-projective-base-change}
Un changement de base d'un H-morphisme projectif est H-projectif.
\end{lemma}

\begin{proof}
C'est vrai parce que le changement de base d'un espace
projectif sur un schéma est un espace projectif, et le fait
que le changement de base d'une immersion fermée est
une immersion fermée, voir Schémas, Lemme \ref{schemes-lemma-base-change-immersion}.
\end{proof}

\begin{lemma}
\label{lemma-base-change-projective}
Un changement de base d'un morphisme projectif (localement) est projectif (localement).
\end{lemma}

\begin{proof}
Ceci est vrai parce que le changement de base d'un fibré
projectif sur un schéma est un fibré projectif,
le tiré en arrière d'un $\mathcal{O}$-module de type
fini est de type fini (Modules, Lemme \ref{modules-lemma-pullback-finite-type})
et le fait que le changement
de base d'une immersion fermée est une immersion
fermée, voir Schémas, Lemme \ref{schemes-lemma-base-change-immersion}. Certains
détails sont omis.
\end{proof}

\begin{lemma}
\label{lemma-projective-quasi-projective}
Un morphisme projectif est quasi-projectif.
\end{lemma}

\begin{proof}
Soit $f : X \to S$ un morphisme projectif. Choisissez une immersion fermée
$i : X \to \mathbf{P}(\mathcal{E})$ où $\mathcal{E}$ est un $\mathcal{O}_S$-module quasi-cohérent et
de type fini. Alors $\mathcal{L} = i^*\mathcal{O}_{\mathbf{P}(\mathcal{E})}(1)$ est
$f$-très ample. Puisque $f$ est propre (Lemme \ref{lemma-locally-projective-proper}), il
est quasi-compact. Par conséquent, le Lemme \ref{lemma-ample-very-ample} implique
que $\mathcal{L}$ est $f$-ample. Puisque $f$ est propre,
il est de type fini. Ainsi, nous avons vérifié que toutes les propriétés
définissant un quasi-projectif sont respectées et nous
avons gagné.
\end{proof}

\begin{lemma}
\label{lemma-H-quasi-projective-open-H-projective}
Soit $f : X \to S$ un H-morphisme quasi-projectif. Alors
$f$ se factorise en $X \to X' \to S$ où $X \to X'$ est une immersion
ouverte et $X' \to S$ est H-projectif.
\end{lemma}

\begin{proof}
Par définition, nous pouvons factoriser $f$ comme
une immersion quasi-compacte $i : X \to \mathbf{P}^n_S$ suivie de la projection $\mathbf{P}^n_S \to S$.
D'après le Lemme \ref{lemma-quasi-compact-immersion}, il existe un sous-schéma fermé
$X' \subset \mathbf{P}^n_S$ tel que $i$ se factorise par une immersion
ouverte $X \to X'$. Le lemme s'ensuit.
\end{proof}

\begin{lemma}
\label{lemma-quasi-projective-open-projective}
Soit $f : X \to S$ un morphisme quasi-projectif avec $S$ quasi-compact
et quasi-séparé. Alors $f$ se factorise en $X \to X' \to S$ où $X \to X'$ est
une immersion ouverte et $X' \to S$ est projectif.
\end{lemma}

\begin{proof}
Soit $\mathcal{L}$ $f$-ample. Comme $f$ est de
type fini et $S$ est quasi-compact, $\mathcal{L}^{\otimes n}$ est $f$-très
ample pour un certain $n > 0$,
voir Lemme \ref{lemma-finite-type-ample-very-ample}. Remplacez $\mathcal{L}$
par $\mathcal{L}^{\otimes n}$. Écrivez $\mathcal{F} = f_*\mathcal{L}$. C'est un
$\mathcal{O}_S$-module quasi-cohérent
d'après Schémas, Lemme \ref{schemes-lemma-push-forward-quasi-coherent} (les
morphismes quasi-projectifs sont quasi-compacts
et séparés, voir Lemme \ref{lemma-quasi-projective-properties}). D'après Propriétés,
Lemme \ref{properties-lemma-directed-colimit-finite-presentation}, nous pouvons trouver un ensemble dirigé
$I$ et un système de modules
quasi-cohérents de type fini $\mathcal{O}_S$ $\mathcal{E}_i$
sur $I$ tels que $\mathcal{F} = \colim \mathcal{E}_i$. Considérez
les compositions
$\psi_i : f^*\mathcal{E}_i \to f^*\mathcal{F} \to \mathcal{L}$. Choisissez un recouvrement ouvert
affine fini $S = \bigcup_{j = 1, \ldots, m} V_j$. Pour chaque $j$, nous pouvons
choisir des sections
$$
s_{j, 0}, \ldots, s_{j, n_j} \in
\Gamma(f^{-1}(V_j), \mathcal{L}) = f_*\mathcal{L}(V_j) = \mathcal{F}(V_j)
$$
qui génèrent $\mathcal{L}$ sur $f^{-1}V_j$ et définissent une immersion
$$
f^{-1}V_j \longrightarrow \mathbf{P}^{n_j}_{V_j},
$$
voir Lemme \ref{lemma-very-ample-finite-type-over-affine}. Choisissons $i$ de sorte
qu'il existe des sections $e_{j, t} \in \mathcal{E}_i(V_j)$ dont les images soient $s_{j, t}$
dans $\mathcal{F}$ pour tous $j = 1, \ldots, m$ et $t = 1, \ldots, n_j$. Alors
l'application $\psi_i$ est surjective puisque les
sections $f^*e_{j, t}$ ont la même image que les sections $s_{j, t}$
qui génèrent $\mathcal{L}|_{f^{-1}V_j}$. D'où nous obtenons un morphisme
$$
r_{\mathcal{L}, \psi_i} : X \longrightarrow \mathbf{P}(\mathcal{E}_i)
$$
sur $S$ tel que sur $V_j$ nous avons une factorisation
$$
f^{-1}V_j \to \mathbf{P}(\mathcal{E}_i)|_{V_j} \to \mathbf{P}^{n_j}_{V_j}
$$
de l'immersion donnée ci-dessus. Il s'ensuit que $r_{\mathcal{L}, \psi_i}|_{V_j}$ est
une immersion, voir Lemme \ref{lemma-immersion-permanence}. Puisque
$S = \bigcup V_j$ Nous concluons que $r_{\mathcal{L}, \psi_i}$ est une
immersion.
Notez que $r_{\mathcal{L}, \psi_i}$ est quasi-compact car
$X \to S$ est quasi-compact et $\mathbf{P}(\mathcal{E}_i) \to S$ est séparé (voir
Schémas, Lemme \ref{schemes-lemma-quasi-compact-permanence}). Par le Lemme \ref{lemma-quasi-compact-immersion}
il existe un sous-schéma fermé $X' \subset \mathbf{P}(\mathcal{E}_i)$ tel que $i$
se factorise par une immersion ouverte $X \to X'$.
Alors $X' \to S$ est projectif par définition
et on obtient le résultat.
\end{proof}

\begin{lemma}
\label{lemma-projective-is-quasi-projective-proper}
Soit $S$ un schéma quasi-compact et quasi-séparé.
Soit $f : X \to S$ un morphisme de schémas. Alors
\begin{enumerate}
\item $f$ est projectif si et seulement si $f$ est quasi-projectif et propre, et
\item $f$ est H-projectif si et seulement si $f$ est H-quasi-projectif et propre.
\end{enumerate}
\end{lemma}

\begin{proof}
Si $f$ est projectif, alors $f$ est quasi-projectif
d'après le Lemme \ref{lemma-projective-quasi-projective} et propre d'après le
Lemme \ref{lemma-locally-projective-proper}. Réciproquement, si $X \to S$ est quasi-projectif
et propre, alors nous pouvons choisir une immersion
ouverte $X \to X'$ avec $X' \to S$ projectif d'après le
Lemme \ref{lemma-quasi-projective-open-projective}. Puisque $X \to S$ est propre, nous
voyons que $X$ est fermé dans $X'$ (Lemme \ref{lemma-image-proper-scheme-closed}),
c'est-à-dire que $X \to X'$ est une immersion
fermée (et ouverte). Comme $X'$ est
isomorphe à un sous-schéma fermé d'un fibré projectif sur $S$
(Définition \ref{definition-projective}), nous voyons que la même
chose est vraie pour $X$, c'est-à-dire
$X \to S$ est un morphisme projectif. Cela
prouve (1). La démonstration de (2) est
la même, sauf qu'elle utilise
les Lemmes \ref{lemma-H-projective} et \ref{lemma-H-quasi-projective-open-H-projective}.
\end{proof}

\begin{lemma}
\label{lemma-composition-projective}
Soient $f : X \to Y$ et $g : Y \to S$ des morphismes de schémas. Si $S$
est quasi-compact et quasi-séparé et que $f$ et $g$ sont projectifs,
alors $g \circ f$ est projectif.
\end{lemma}

\begin{proof}
D'après les Lemmes \ref{lemma-projective-quasi-projective} et \ref{lemma-locally-projective-proper},
nous voyons que $f$
et $g$ sont quasi-projectifs et
propres. D'après les Lemmes \ref{lemma-composition-proper}
et \ref{lemma-composition-quasi-projective}, nous voyons
que $g \circ f$ est propre et quasi-projectif.
Ainsi, $g \circ f$
est projectif d'après le Lemme \ref{lemma-projective-is-quasi-projective-proper}.
\end{proof}

\begin{lemma}
\label{lemma-projective-permanence}
Soient $g : Y \to S$ et $f : X \to Y$ des morphismes de schémas. Si $g \circ f$
est projectif et $g$ est séparé, alors $f$
est projectif.
\end{lemma}

\begin{proof}
Choisissez une immersion fermée $X \to \mathbf{P}(\mathcal{E})$ où $\mathcal{E}$ est un $\mathcal{O}_S$-module
quasi-cohérent de type fini. Alors nous
obtenons un morphisme $X \to \mathbf{P}(\mathcal{E}) \times_S Y$. Ce morphisme est
une immersion fermée parce qu'il est la composition
$$
X \to X \times_S Y \to \mathbf{P}(\mathcal{E}) \times_S Y
$$
où le premier morphisme est une immersion fermée
d'après Schémas, Lemme \ref{schemes-lemma-semi-diagonal} (et le
fait que $g$ est séparé) et
le second comme le changement de base d'une immersion
fermée. Enfin, le produit fibré $\mathbf{P}(\mathcal{E}) \times_S Y$ est isomorphe
à $\mathbf{P}(g^*\mathcal{E})$ et le tiré en arrière préserve les
modules quasi-cohérents de type fini.
\end{proof}

\begin{lemma}
\label{lemma-projective-over-quasi-projective-is-H-projective}
Soit $S$ un schéma qui admet un faisceau inversible ample. Alors
\begin{enumerate}
\item tout morphisme projectif $X \to S$ est H-projectif, et
\item tout morphisme quasi-projectif $X \to S$ est H-quasi-projectif.
\end{enumerate}
\end{lemma}

\begin{proof}
Les hypothèses sur $S$ impliquent que $S$ est
quasi-compact et séparé, voir Propriétés, Définition
\ref{properties-definition-ample} et Lemme \ref{properties-lemma-ample-immersion-into-proj} et Constructions,
Lemme \ref{constructions-lemma-proj-separated}. Ainsi, le Lemme \ref{lemma-quasi-projective-open-projective}
s'applique et nous voyons que (1) implique
(2). Soit $\mathcal{E}$ un $\mathcal{O}_S$-module
quasi-cohérent de type fini. D'après notre définition
de morphismes projectifs, il suffit de montrer que
$\mathbf{P}(\mathcal{E}) \to S$ est H-projectif. Si $\mathcal{E}$
est engendré par un nombre fini de sections globales,
alors la surjection correspondante $\mathcal{O}_S^{\oplus n} \to \mathcal{E}$ induit
une immersion fermée
$$
\mathbf{P}(\mathcal{E}) \longrightarrow
\mathbf{P}(\mathcal{O}_S^{\oplus n}) = \mathbf{P}^{n - 1}_S
$$
comme souhaité. En général, soit $\mathcal{L}$ un faisceau
inversible ample sur $S$. D'après Propriétés,
Proposition \ref{properties-proposition-characterize-ample},
il existe un entier $n$ tel que
$\mathcal{E} \otimes_{\mathcal{O}_S} \mathcal{L}^{\otimes n}$ soit globalement engendré
par un nombre
fini de sections. Puisque $\mathbf{P}(\mathcal{E}) =
\mathbf{P}(\mathcal{E} \otimes_{\mathcal{O}_S} \mathcal{L}^{\otimes n})$ d'après Constructions,
Lemme \ref{constructions-lemma-twisting-and-proj}, cela termine la démonstration.
\end{proof}

\begin{lemma}
\label{lemma-proper-ample-is-proj}
Soit $f : X \to S$ un morphisme universellement fermé. Soit
$\mathcal{L}$ un $f$-ample faisceau inversible $\mathcal{O}_X$. Alors
le morphisme canonique
$$
r : X
\longrightarrow
\underline{\text{Proj}}_S
\left(
\bigoplus\nolimits_{d \geq 0} f_*\mathcal{L}^{\otimes d}
\right)
$$
du Lemme \ref{lemma-characterize-relatively-ample} est un isomorphisme.
\end{lemma}

\begin{proof}
Remarquez que $f$ est quasi-compact parce que l'existence
d'un $f$-module inversible ample oblige $f$
à être quasi-compact. D'après le lemme cité, le morphisme
$r$ est une immersion ouverte. D'autre
part, l'image de $r$ est fermée
par le Lemme \ref{lemma-image-proper-scheme-closed} (la cible de $r$ est séparée
au-dessus de $S$ d'après Constructions, Lemme
\ref{constructions-lemma-relative-proj-separated}). Enfin, l'image de $r$
est dense par Propriétés, Lemme \ref{properties-lemma-ample-immersion-into-proj} (ici nous utilisons
également qu'il a été montré dans la démonstration
du Lemme \ref{lemma-characterize-relatively-ample} que le morphisme $r$
sur des ouverts affines de $S$
est donné par le morphisme canonique de
Propriétés, Lemme \ref{properties-lemma-map-into-proj}). Ainsi nous concluons que
$r$ est une immersion ouverte surjective, c'est-à-dire un
isomorphisme.
\end{proof}

\begin{lemma}
\label{lemma-proper-ample-delete-affine}
Soit $f : X \to S$ un morphisme universellement fermé. Soit $\mathcal{L}$
un $f$-ample module inversible $\mathcal{O}_X$.
Soit $s \in \Gamma(X, \mathcal{L})$. Alors $X_s \to S$ est un morphisme affine.
\end{lemma}

\begin{proof}
La question est locale sur $S$ (Lemme \ref{lemma-characterize-affine}) donc
nous pouvons supposer que $S$ est affine. D'après le Lemme
\ref{lemma-proper-ample-is-proj}, nous pouvons écrire $X = \text{Proj}(A)$ où $A$
est un anneau gradué et $s$
correspond à $f \in A_1$ et $X_s = D_+(f)$ (Propriétés, Lemme \ref{properties-lemma-map-into-proj})
ce qui prouve le lemme par construction de $\text{Proj}(A)$,
voir Constructions, section \ref{constructions-section-proj}.
\end{proof}











\section{Morphismes entiers et finis}
\label{section-integral}

\noindent
Rappelons qu'un homomorphisme d'anneaux $R \to A$ est dit {\it
entier} si chaque élément de $A$ satisfait à une équation
monique avec des coefficients dans $R$. Rappelez-vous qu'un homomorphisme
d'anneaux $R \to A$ est dit {\it fini} si
$A$ est fini en tant que $R$-module. Voir Algèbre, Définition \ref{algebra-definition-integral-ring-map}.

\begin{definition}
\label{definition-integral}
Soit $f : X \to S$ un morphisme de schémas.
\begin{enumerate}
\item On dit que $f$ est {\it entier}
si $f$ est affine et si pour tout ouvert affine
$\Spec(R) = V \subset S$ avec image réciproque $\Spec(A) = f^{-1}(V) \subset X$, l'homomorphisme
d'anneaux associé $R \to A$ est entier.
\item On dit que $f$ est {\it fini}
si $f$ est affine et si pour tout ouvert affine
$\Spec(R) = V \subset S$ avec image réciproque $\Spec(A) = f^{-1}(V) \subset X$ l'homomorphisme
d'anneaux associé $R \to A$ est fini.
\end{enumerate}
\end{definition}

\noindent
Il est clair que les morphismes entiers/finis sont séparés
et quasi-compacts. Il est également clair qu'un morphisme
fini est un morphisme de type fini. La plupart des lemmes de
cette section sont complètement
standards. Mais notez le lemme amusant \ref{lemma-integral-universally-closed}
à la fin de la section.

\begin{lemma}
\label{lemma-integral-local}
Soit $f : X \to S$ un morphisme de schémas. Les conditions
suivantes sont équivalentes :
\begin{enumerate}
\item Le morphisme $f$ est entier.
\item Il existe un recouvrement ouvert affine $S = \bigcup U_i$
tel que chaque $f^{-1}(U_i)$
soit affine et $\mathcal{O}_S(U_i) \to \mathcal{O}_X(f^{-1}(U_i))$ soit entier.
\item Il existe un recouvrement ouvert $S = \bigcup U_i$
tel que chaque $f^{-1}(U_i) \to U_i$ soit entier.
\end{enumerate}
De plus, si $f$ est entier alors pour tout sous-schéma
ouvert $U \subset S$ le morphisme $f : f^{-1}(U) \to U$ est entier.
\end{lemma}

\begin{proof}
Voir Algèbre, Lemme \ref{algebra-lemma-integral-local}. Certains détails
sont omis.
\end{proof}

\begin{lemma}
\label{lemma-finite-local}
Soit $f : X \to S$ un morphisme de schémas. Les conditions
suivantes sont équivalentes :
\begin{enumerate}
\item Le morphisme $f$ est fini.
\item Il existe un recouvrement ouvert affine $S = \bigcup U_i$
tel que chaque $f^{-1}(U_i)$
soit affine et $\mathcal{O}_S(U_i) \to \mathcal{O}_X(f^{-1}(U_i))$ soit fini.
\item Il existe un recouvrement ouvert $S = \bigcup U_i$
tel que chaque $f^{-1}(U_i) \to U_i$ soit fini.
\end{enumerate}
De plus, si $f$ est fini alors pour chaque sous-schéma
ouvert $U \subset S$ le morphisme $f : f^{-1}(U) \to U$ est fini.
\end{lemma}

\begin{proof}
Voir Algèbre, Lemme \ref{algebra-lemma-integral-local}. Certains détails
sont omis.
\end{proof}

\begin{lemma}
\label{lemma-finite-integral}
Un morphisme fini est entier.
Un morphisme entier qui est localement de type fini est fini.
\end{lemma}

\begin{proof}
Voir Algèbre, Lemme \ref{algebra-lemma-finite-is-integral}
et Lemme \ref{algebra-lemma-characterize-finite-in-terms-of-integral}.
\end{proof}

\begin{lemma}
\label{lemma-composition-finite}
Une composition de morphismes finis est finie. Il en
est de même pour les morphismes entiers.
\end{lemma}

\begin{proof}
Voir Algèbre, Lemme \ref{algebra-lemma-finite-transitive}
et \ref{algebra-lemma-integral-transitive}.
\end{proof}

\begin{lemma}
\label{lemma-base-change-finite}
Un changement de base d'un morphisme fini est fini. Il
en est de même pour les morphismes entiers.
\end{lemma}

\begin{proof}
Voir Algèbre, Lemme \ref{algebra-lemma-base-change-integral}.
\end{proof}

\begin{lemma}
\label{lemma-integral-universally-closed}
\begin{slogan}
entier $=$ affine $+$ universellement fermé
\end{slogan}
Soit $f : X \to S$ un morphisme de schémas. Les conditions
suivantes sont équivalentes
\begin{enumerate}
\item $f$ est entier, et
\item $f$ est affine et universellement fermé.
\end{enumerate}
\end{lemma}

\begin{proof}
Supposons (1). Un morphisme entier est affine par
définition. Un changement de base d'un morphisme entier est entier,
donc pour prouver (2) il suffit de montrer
qu'un morphisme entier est fermé. Cela résulte de Algèbre,
Lemmes \ref{algebra-lemma-integral-going-up} et \ref{algebra-lemma-going-up-closed}.

\medskip\noindent Supposons
(2). Nous pouvons supposer $f$ est
le morphisme $f : \Spec(A) \to \Spec(R)$ provenant d'un homomorphisme
d'anneaux $R \to A$. Soit $a$ un élément de $A$.
Nous devons montrer que $a$ est entier sur
$R$, c'est-à-dire que dans le noyau $I$ de l'application
$R[x] \to A$ envoyant $x$
à $a$ il existe un polynôme unitaire. Considérons
l'anneau $B = A[x]/(ax -1)$ et soit $J$ le noyau de la composition
$R[x]\to A[x] \to B$. Si $f\in J$ il existe $q\in A[x]$
tel que $f = (ax-1)q$ dans $A[x]$ donc si $f = \sum_i f_ix^i$ et $q = \sum_iq_ix^i$,
pour tout $i \geq 0$ nous avons $f_i = aq_{i-1} - q_i$. Pour $n \geq \deg q + 1$ le polynôme
$$
\sum\nolimits_{i \geq 0} f_i x^{n - i} =
\sum\nolimits_{i \geq 0} (a q_{i - 1} - q_i) x^{n - i} =
(a - x) \sum\nolimits_{i \geq 0} q_i x^{n - i - 1}
$$
est clairement dans $I$ ; si $f_0 = 1$ ce polynôme est aussi unitaire, nous
sommes donc réduits à prouver que $J$ contient un polynôme avec un terme constant $1$.
Nous le faisons en prouvant que $\Spec(R[x]/(J + (x))$ est vide.


\medskip\noindent
Puisque $f$ est universellement fermé, le changement
de base $\Spec(A[x]) \to \Spec(R[x])$ est fermé. Par conséquent, l'image du sous-ensemble
fermé $\Spec(B) \subset \Spec(A[x])$ est le sous-ensemble
fermé $\Spec(R[x]/J) \subset \Spec(R[x])$, voir Exemple \ref{example-scheme-theoretic-image}
et Lemme \ref{lemma-quasi-compact-scheme-theoretic-image}. En particulier, $\Spec(B) \to \Spec(R[x]/J)$
est surjectif. Considérons le diagramme suivant
où chaque carré est un produit fibré :
$$
\xymatrix{
\Spec(B) \ar@{->>}[r]^g &
\Spec(R[x]/J)  \ar[r] &
\Spec(R[x])\\
\emptyset \ar[u] \ar[r] &
\Spec(R[x]/(J + (x)))\ar[u] \ar[r] &
\Spec(R) \ar[u]^0
}
$$
Le coin inférieur gauche est vide car il s'agit du spectre de $R\otimes_{R[x]} B$
où l'application $R[x]\to B$ envoie $x$ sur un élément inversible
et $R[x]\to R$ envoie $x$ sur $0$. Comme $g$
est surjectif, cela implique que $\Spec(R[x]/(J + (x)))$ est vide, comme nous voulions
le montrer.
\end{proof}

\begin{lemma}
\label{lemma-integral-fibres}
Soit $f : X \to S$ un morphisme entier. Alors chaque
point de $X$ est fermé dans sa fibre.
\end{lemma}

\begin{proof}
Voir Algèbre, Lemme \ref{algebra-lemma-integral-no-inclusion}.
\end{proof}

\begin{lemma}
\label{lemma-integral-dimension}
Soit $f : X \to Y$ un morphisme entier. Alors $\dim(X) \leq \dim(Y)$. Si $f$
est surjectif alors $\dim(X) = \dim(Y)$.
\end{lemma}

\begin{proof}
Puisque la dimension de $X$ et $Y$ est le supremum
des dimensions des éléments d'un recouvrement ouvert affine,
nous pouvons supposer que $Y$
et $X$ sont affines. L'inégalité
découle de l'Algèbre, Lemme
\ref{algebra-lemma-integral-dim-up}. L'égalité découle ensuite de l'Algèbre,
Lemmes \ref{algebra-lemma-dimension-going-up} et \ref{algebra-lemma-integral-going-up}.
\end{proof}

\begin{lemma}
\label{lemma-finite-quasi-finite}
Un morphisme fini est quasi-fini.
\end{lemma}

\begin{proof}
Ceci est impliqué par l'Algèbre, Lemme \ref{algebra-lemma-quasi-finite} et Lemme
\ref{lemma-quasi-finite-locally-quasi-compact}. Alternativement, tous les points
des fibres sont des points fermés par le Lemme
\ref{lemma-integral-fibres} (et le fait qu'un morphisme fini est entier)
et utiliser le Lemme
\ref{lemma-quasi-finite-at-point-characterize} (3) pour voir que $f$ est quasi-fini
en $x$ pour tous $x \in X$.
\end{proof}

\begin{lemma}
\label{lemma-finite-proper}
Soit $f : X \to S$ un morphisme de schémas. Les conditions suivantes sont équivalentes
\begin{enumerate}
\item $f$ est fini, et
\item $f$ est affine et propre.
\end{enumerate}
\end{lemma}

\begin{proof}
Cela découle formellement
du Lemme \ref{lemma-integral-universally-closed}, du fait qu'un morphisme
fini est entier et séparé, du fait qu'un
morphisme propre est la même chose qu'un morphisme
de type fini, séparé, universellement
fermé, et du fait qu'un morphisme entier de type
fini est fini (Lemme \ref{lemma-finite-integral}).
\end{proof}

\begin{lemma}
\label{lemma-closed-immersion-finite}
Une immersion fermée est finie (et a fortiori entière).
\end{lemma}

\begin{proof}
Vrai car une immersion fermée est affine
(Lemme \ref{lemma-closed-immersion-affine}) et un homomorphisme
surjectif d'anneaux est fini et entier.
\end{proof}

\begin{lemma}
\label{lemma-finite-union-finite}
Soient $X_i \to Y$, $i = 1, \ldots, n$ des morphismes finis de
schémas. Alors $X_1 \amalg \ldots \amalg X_n \to Y$ est fini aussi.
\end{lemma}

\begin{proof}
Cela découle du fait algébrique que si $R \to A_i$, $i = 1, \ldots, n$ sont
des homomorphismes finis d'anneaux, alors $R \to A_1 \times \ldots \times A_n$ est fini aussi.
\end{proof}

\begin{lemma}
\label{lemma-finite-permanence}
Soient $f : X \to Y$ et $g : Y \to Z$ des morphismes.
\begin{enumerate}
\item Si $g \circ f$ est fini et $g$ séparé, alors $f$ est fini.
\item Si $g \circ f$ est entier et $g$ séparé, alors $f$ est entier.
\end{enumerate}
\end{lemma}

\begin{proof}
Supposons que $g \circ f$ soit fini (resp.\ entier) et que $g$ soit
séparé. Le changement de base $X \times_Z Y \to Y$ est fini (resp.\
entier) d'après le lemme
\ref{lemma-base-change-finite}. Le morphisme $X \to X \times_Z Y$
est une immersion fermée car $Y \to Z$ est séparé,
voir Schémas, Lemme \ref{schemes-lemma-section-immersion}. Une immersion
fermée est finie (resp. \  entière),
voir Lemme \ref{lemma-closed-immersion-finite}. La composition de morphismes
finis (resp. \  entiers) est finie (resp.
\  entière),
voir Lemme \ref{lemma-composition-finite}. Ainsi nous avons gagné.
\end{proof}

\begin{lemma}
\label{lemma-finite-monomorphism-closed}
Soit $f : X \to Y$ un morphisme de schémas. Si $f$
est fini et un monomorphisme, alors $f$ est une immersion fermée.
\end{lemma}

\begin{proof}
Cela se réduit
à Algèbre, Lemme \ref{algebra-lemma-finite-epimorphism-surjective}.
\end{proof}

\begin{lemma}
\label{lemma-finite-projective}
Un morphisme fini est projectif.
\end{lemma}

\begin{proof}
Soit $f : X \to S$ un morphisme fini. Alors $f_*\mathcal{O}_X$ est
un $\mathcal{O}_S$-module quasi-cohérent
(Lemme \ref{lemma-affine-equivalence-algebras}) de type fini (d'après
notre
définition des morphismes finis et
Propriétés, Lemme \ref{properties-lemma-finite-type-module}). Nous affirmons qu'il
existe une immersion fermée
$$
\sigma :
X
\longrightarrow
\mathbf{P}(f_*\mathcal{O}_X) =
\underline{\text{Proj}}_S(\text{Sym}^*_{\mathcal{O}_S}(f_*\mathcal{O}_X))
$$
sur $S$, ce qui termine
la démonstration. À savoir, nous prenons $\sigma$ comme étant
le morphisme qui correspond (via Constructions, Lemme \ref{constructions-lemma-apply-relative})
à la surjection
$$
f^*f_*\mathcal{O}_X \longrightarrow \mathcal{O}_X
$$
provenant de l'application d'adjonction $f^*f_* \to \text{id}$. Alors $\sigma$
est une immersion fermée. En effet, localement sur les
affines de $S$, nous pouvons écrire $X = \Spec(A)$ et $S = \Spec(R)$.
Puisque $X$ est fini sur $S$ nous pouvons choisir
$a_1, \ldots, a_n \in A$ générant $A$ en tant que $R$-module.
Alors l'homomorphisme de $R$-algèbres $R[T_0, T_1, \ldots, T_n] \to \text{Sym}_R^*(A)$ envoyant $T_0$
sur $1$ et $T_i$ sur $a_i$ pour $1 \leq i \leq n$
est surjective. D'où $\mathbf{P}(f_*\mathcal{O}_X)$ est
un sous-schéma fermé de $\mathbf{P}^n_S$ selon Constructions,
Lemme \ref{constructions-lemma-surjective-graded-rings-map-proj}. Ainsi, il suffit de prouver
que le morphisme induit $X \to \mathbf{P}^n_S$ est une
immersion fermée (voir par exemple, Lemme \ref{lemma-closed-immersion-ideals}).
Le lecteur vérifie que $X \to \mathbf{P}^n_S$ a une image contenue
dans le $D_+(T_0) \cong \mathbf{A}^n_S$ ouvert et que $X \to \mathbf{A}^n_S$ correspond à l'application
d'algèbre $R$ surjective
$R[x_1, \ldots, x_n] \to \text{Sym}_R^*(A)$ envoyant $x_i$ vers $a_i$.
D'où $X \to \mathbf{P}^n_S$ est une immersion (comme une composition
d'une immersion fermée et d'une
immersion ouverte). Comme $X$ est fini sur
$S$, l'image de cette immersion
est fermée (cela utilise le Lemme
\ref{lemma-finite-proper}, le Lemme \ref{lemma-image-proper-scheme-closed} et Constructions, Lemme \ref{constructions-lemma-projective-space-separated}).
Nous concluons par Schémas, Lemme \ref{schemes-lemma-immersion-when-closed}.
\end{proof}







\section{Homéomorphismes universels}
\label{section-universal-homeomorphisms}

\noindent
La définition suivante est vraiment superflue puisque un homéomorphisme
universel n'est vraiment qu'un morphisme entier,
universellement injectif
et surjectif, voir le Lemme \ref{lemma-universal-homeomorphism}.

\begin{definition}
\label{definition-universal-homeomorphism}
Un morphisme $f : X \to Y$ de schémas est appelé un {\it homéomorphisme universel}
si le changement de base $f' : Y' \times_Y X \to Y'$ est un homéomorphisme pour
tout morphisme $Y' \to Y$.
\end{definition}

\noindent
Tout d'abord, nous énonçons les lemmes obligatoires.

\begin{lemma}
\label{lemma-base-change-universal-homeomorphism}
Le changement de base d'un homéomorphisme universel de schémas
par tout morphisme de schémas est un homéomorphisme universel.
\end{lemma}

\begin{proof}
Cela découle immédiatement de la définition.
\end{proof}

\begin{lemma}
\label{lemma-composition-universal-homeomorphism}
La composition d'une paire d'homéomorphismes universels de schémas
est un homéomorphisme universel.
\end{lemma}

\begin{proof}
Omis.
\end{proof}

\noindent
Le lemme simple suivant est la clé pour caractériser les homéomorphismes
universels.

\begin{lemma}
\label{lemma-homeomorphism-affine}
Soit $f : X \to Y$ un morphisme de schémas. Si $f$ est un homéomorphisme sur un sous-ensemble
fermé de $Y$, alors $f$ est affine.
\end{lemma}

\begin{proof}
Soit $y \in Y$ un point. Si $y \not \in f(X)$, alors il existe un voisinage
affine de $y$ qui est disjoint de $f(X)$.
Si $y \in f(X)$, soit $x \in X$ l'unique point de $X$ s'envoyant
sur $y$. Soit $y \in V$ un voisinage
ouvert affine. Soit $U \subset X$ un voisinage ouvert affine de $x$
qui s'envoie dans $V$. Puisque $f(U) \subset V \cap f(X)$ est ouvert dans la
topologie induite par notre hypothèse
sur $f$, nous pouvons choisir un $h \in \Gamma(V, \mathcal{O}_Y)$
tel que $y \in D(h)$ et $D(h) \cap f(X) \subset f(U)$. Notons $h' \in \Gamma(U, \mathcal{O}_X)$ la restriction
de $f^\sharp(h)$ à $U$. Alors nous voyons que $D(h') \subset U$
est égal à $f^{-1}(D(h))$. En d'autres termes, chaque point de $Y$
a un voisinage ouvert dont l'image réciproque est affine.
Ainsi $f$ est
affine, voir Lemme \ref{lemma-characterize-affine}.
\end{proof}

\begin{lemma}
\label{lemma-universal-homeomorphism}
Soit $f : X \to Y$ un morphisme de schémas. Les conditions suivantes sont équivalentes
:
\begin{enumerate}
\item $f$ est un homéomorphisme universel, et
\item $f$ est entier, universellement injectif et surjectif.
\end{enumerate}
\end{lemma}

\begin{proof}
Supposons que $f$ soit un homéomorphisme universel.
D'après le Lemme \ref{lemma-homeomorphism-affine}, nous
voyons que $f$ est affine. Puisque $f$ est clairement universellement
fermé, nous voyons que $f$
est entier d'après le Lemme \ref{lemma-integral-universally-closed}. Il
est aussi clair que $f$ est universellement injectif et surjectif.

\medskip\noindent
Supposons que $f$ soit entier, universellement injectif et
surjectif. D'après le Lemme \ref{lemma-integral-universally-closed}, $f$
est universellement fermé. Puisqu'il est aussi universellement bijectif
(voir le Lemme \ref{lemma-base-change-surjective}), nous voyons
que c'est un homéomorphisme universel.
\end{proof}

\begin{lemma}
\label{lemma-reduction-universal-homeomorphism}
Soit $X$ un schéma. L'immersion fermée canonique $X_{red} \to X$
(voir Schémas, Définition \ref{schemes-definition-reduced-induced-scheme}) est un homéomorphisme
universel.
\end{lemma}

\begin{proof}
Omis.
\end{proof}

\begin{lemma}
\label{lemma-check-closed-infinitesimally}
Soient $f : X \to S$ et $S' \to S$ des morphismes de schémas.
Supposons
\begin{enumerate}
\item $S' \to S$ est une immersion fermée,
\item $S' \to S$ est bijective sur les points,
\item $X \times_S S' \to S'$ est une immersion fermée, et
\item $X \to S$ est de type fini ou $S' \to S$ est de présentation finie.
\end{enumerate}
Alors $f : X \to S$ est une immersion fermée.
\end{lemma}

\begin{proof}
Les hypothèses (1) et (2) impliquent que $S' \to S$ est un homéomorphisme universel
(par exemple parce que $S_{red} = S'_{red}$ et en utilisant
le Lemme \ref{lemma-reduction-universal-homeomorphism}). Par conséquent, (3) implique
que $X \to S$ est un homéomorphisme sur un sous-ensemble
fermé de $S$. Ensuite, $X \to S$
est affine d'après le Lemme \ref{lemma-homeomorphism-affine}.
Soit $U \subset S$ un ouvert affine, disons $U = \Spec(A)$. Alors $S' = \Spec(A/I)$ par (1) pour
un idéal nilpotent localement $I$ par (2). Comme $f$ est affine,
on voit que $f^{-1}(U) = \Spec(B)$.
L'hypothèse (4) nous dit que $B$ est une $A$-algèbre
de type fini (Lemme \ref{lemma-locally-finite-type-characterize}) ou que
$I$ est de type fini
(Lemme \ref{lemma-closed-immersion-finite-presentation}). L'hypothèse (3) est que $A/I \to B/IB$
est surjectif. D'après l'Algèbre, Lemme
\ref{algebra-lemma-surjective-mod-locally-nilpotent} si $A \to B$ est de type fini ou Algèbre,
Lemme \ref{algebra-lemma-NAK} si $I$
est de type fini et donc nilpotent, nous déduisons
que $A \to B$ est surjectif. Cela signifie
que $f$ est une immersion fermée,
voir Lemme \ref{lemma-closed-immersion}.
\end{proof}

\begin{lemma}
\label{lemma-permanence-universal-homeomorphism}
Soit $f : X \to Z$ la composition de deux morphismes $g : X \to Y$
et $h : Y \to Z$. Si deux des
morphismes $\{f, g, h\}$ sont des homéomorphismes universels, le troisième
morphisme l'est aussi.
\end{lemma}

\begin{proof}
Si $g$ et $h$ sont tous deux des homéomorphismes universels,
$f$ l'est aussi d'après le Lemme \ref{lemma-composition-universal-homeomorphism}.

\medskip\noindent Supposons
que $f$ et $g$ soient des homéomorphismes universels. Nous voulons montrer
que $h$ l'est aussi. Maintenant, faisons un changement de base du diagramme le long d'un
morphisme arbitraire $\alpha : Z' \to Z$ de schémas, nous obtenons le diagramme suivant avec tous les
carrés cartésiens :
$$
\xymatrix{
X' \ar[r]^{g'} \ar[d] & Y' \ar[r]^{h'} \ar[d] & Z' \ar[d] \\
X \ar[r]^{g} & Y \ar[r]^{h} & Z.
}
$$
Notre hypothèse implique que la composition $f'= h' \circ g' : X' \to Z'$ et $g' : X' \to Y'$ est un
homéomorphisme, donc $h'$ l'est également. Cela termine la
démonstration que $h$ est un homéomorphisme universel.

\medskip\noindent Enfin,
supposons que $f$ et $h$ soient
des homéomorphismes universels. Nous voulons montrer que $g$ est un homéomorphisme
universel. Soit $\beta : Y' \to Y$ un morphisme de schémas quelconque. Nous
obtenons le diagramme suivant avec tous les carrés cartésiens :
$$
\xymatrix{
X' \ar[r]^{g'} \ar[d] & Y' \ar[d]^{\gamma} & \\
X'' \ar[r]^{g''} \ar[d] & Y'' \ar[r]^{h''} \ar[d] & Y' \ar[d]^{h \circ \beta} \\
X \ar[r]^{g} & Y \ar[r]^{h} & Z.
}
$$
Ici, le morphisme $\gamma : Y' \to Y''$ est défini par la propriété universelle des produits
fibrés et les deux morphismes $id_{Y'} : Y' \to Y'$
et $\beta : Y' \to Y$. Nous allons prouver que
$g'$ est un homéomorphisme. Comme la
propriété d'être un homéomorphisme a la propriété 2-sur-3, nous voyons que
$g''$ est un homéomorphisme.
En regardant le carré du haut, il suffit de prouver que $\gamma$ est
un homéomorphisme universel. Comme $h''$ est un homéomorphisme,
nous voyons que c'est un morphisme affine d'après le
Lemme \ref{lemma-homeomorphism-affine} et a fortiori séparé (Lemme \ref{lemma-affine-separated}). Comme
$h'' \circ \gamma$ est l'identité, nous voyons que $\gamma$ est une
immersion fermée par
Schémas, Lemme \ref{schemes-lemma-section-immersion}. Puisque $h''$ est
bijectif, il s'ensuit que $\gamma$ est une immersion
fermée bijective et donc un homéomorphisme universel (par exemple
par la caractérisation dans le Lemme \ref{lemma-universal-homeomorphism}) comme désiré.
\end{proof}




\section{Homéomorphismes universels de schémas affines}
\label{section-universal-homeomorphisms-affine}

\noindent
Dans cette section, nous caractérisons les homéomorphismes universels des schémas affines.

\begin{lemma}
\label{lemma-subalgebra-inherits}
Soit $A \to B$ un homomorphisme d'anneaux tel que le morphisme de schémas
induit $f : \Spec(B) \to \Spec(A)$ soit un homéomorphisme universel, resp. \ 
un homéomorphisme universel induisant des isomorphismes sur les corps résiduels, resp.
\  universellement
fermé, resp. \  universellement fermé et universellement
injectif. Alors, pour toute $A$-sous-algèbre
$B' \subset B$, la même chose est vraie pour $f' : \Spec(B') \to \Spec(A)$.
\end{lemma}

\begin{proof}
Si $f$ est universellement fermé, alors
$B$ est entier sur $A$ d'après le Lemme \ref{lemma-integral-universally-closed}. Par conséquent, $B'$
est entier sur $A$ et $f'$
est universellement fermé (d'après le même lemme).
Cela prouve le cas où $f$ est universellement fermé.

\medskip\noindent
En continuant, nous voyons que $B$
est entier sur $B'$ (Algèbre, Lemme \ref{algebra-lemma-integral-permanence})
ce qui implique que $\Spec(B) \to \Spec(B')$ est surjectif
(Algèbre, Lemme \ref{algebra-lemma-integral-overring-surjective}). Ainsi, si $A \to B$ induit
des extensions purement inséparables des corps résiduels, alors il
en est de même pour $A \to B'$. Cela prouve le cas
où $f$ est universellement fermé et universellement
injectif, voir Lemme \ref{lemma-universally-injective}.

\medskip\noindent
Le cas où $f$ est un homéomorphisme universel découle des
remarques ci-dessus, du Lemme \ref{lemma-universal-homeomorphism}, et de l'observation
évidente que si $f$ est surjectif, alors $f'$ l'est aussi.

\medskip\noindent Si $A \to B$
induit des isomorphismes sur les corps résiduels, alors il en
est de même pour $A \to B'$ (voir l'argument au deuxième paragraphe).
De cette façon, nous voyons que le lemme est valable dans le cas restant.
\end{proof}

\begin{lemma}
\label{lemma-colimit-inherits}
Soit $A$ un anneau. Soit $B = \colim B_\lambda$ une colimite filtrante d'$A$-algèbres.
Si chaque $f_\lambda : \Spec(B_\lambda) \to \Spec(A)$ est un homéomorphisme universel, respectivement
\  un homéomorphisme
universel induisant des isomorphismes sur les corps résiduels, respectivement
\  universellement
fermé, respectivement \  universellement fermé et
universellement injectif, alors il en va de même pour $f : \Spec(B) \to \Spec(A)$.
\end{lemma}

\begin{proof}
Si $f_\lambda$ est universellement fermé, alors $B_\lambda$
est entier sur $A$ d'après
le Lemme \ref{lemma-integral-universally-closed}. Par conséquent, $B$
est entier sur $A$ et $f$
est universellement fermé (d'après le même lemme).
Cela prouve le cas où chaque $f_\lambda$ est universellement fermé.

\medskip\noindent
Pour un idéal premier $\mathfrak q \subset B$ au-dessus de $\mathfrak p \subset A$,
notons $\mathfrak q_\lambda \subset B_\lambda$ l'image réciproque. Alors $\kappa(\mathfrak q) = \colim \kappa(\mathfrak q_\lambda)$.
Ainsi, si $A \to B_\lambda$ induit des extensions
purement inséparables des corps résiduels, il
en est de même pour $A \to B$. Cela prouve le cas où $f_\lambda$
est universellement fermé et universellement injectif,
voir Lemme \ref{lemma-universally-injective}.

\medskip\noindent
Le cas où $f$ est un homéomorphisme universel découle
des remarques ci-dessus et du Lemme \ref{lemma-universal-homeomorphism} combiné au fait
que les idéaux premiers dans $B$ sont la même chose que les
suites compatibles d'idéaux premiers dans tous les $B_\lambda$.

\medskip\noindent Si $A \to B_\lambda$
induit des isomorphismes sur les corps résiduels, alors il en est de même
pour $A \to B$ (voir l'argument dans le deuxième paragraphe).
De cette manière, nous voyons que le lemme est valable dans le cas restant.
\end{proof}

\begin{lemma}
\label{lemma-special-elements-and-localization}
Soit $A \subset B$ une extension d'anneaux. Soit $S \subset A$ un
sous-ensemble multiplicatif. Soient $n \geq 1$
et $b_i \in B$ pour $1 \leq i \leq n$. Tout $x \in S^{-1}B$ tel que
$$
x \not \in S^{-1}A\text{ et } b_i x^i \in S^{-1}A\text{ pour }i = 1, \ldots, n
$$
est égal à $s^{-1}y$ avec $s \in S$ et $y \in B$ tels que
$$
y \not \in A\text{ et } b_i y^i \in A\text{ pour }i = 1, \ldots, n
$$
\end{lemma}

\begin{proof}
Omis. Astuce : clarifier les dénominateurs.
\end{proof}

\begin{lemma}
\label{lemma-nth-and-nplusone-implies-square-and-cube}
Soit $A \subset B$ une extension d'anneaux. S'il existe
$b \in B$, $b \not \in A$ et un entier $n \geq 2$
avec $b^n \in A$ et $b^{n + 1} \in A$, alors il existe
un $b' \in B$, $b' \not \in A$
avec $(b')^2 \in A$ et $(b')^3 \in A$.
\end{lemma}

\begin{proof}
Soient $b$ et $n$ comme dans
le lemme. Alors toutes les puissances suffisamment grandes
de $b$ sont dans $A$. À savoir, $(b^n)^k(b^{n + 1})^i = b^{(k + i)n + i}$
ce qui implique que toute puissance $b^m$ avec $m \geq n^2$
est dans $A$. Ainsi, si $i \geq 1$ est le plus
grand entier tel que $b^i \not \in A$, alors $(b^i)^2 \in A$ et $(b^i)^3 \in A$.
\end{proof}

\begin{lemma}
\label{lemma-square-and-cube}
Soit $A \subset B$ une extension d'anneau telle que $\Spec(B) \to \Spec(A)$ soit un homéomorphisme
universel induisant des isomorphismes sur les corps résiduels.
Si $A \not = B$, alors il existe un $b \in B$, $b \not \in A$ avec
$b^2 \in A$ et $b^3 \in A$.
\end{lemma}

\begin{proof}
Rappelons que l'extension $A \subset B$ est entière
(Lemme \ref{lemma-integral-universally-closed}). D'après le Lemme
\ref{lemma-subalgebra-inherits}, nous pouvons supposer
que $B$ est généré par un seul élément
sur $A$. Ainsi $B$
est fini sur $A$ (Algèbre, Lemme \ref{algebra-lemma-characterize-finite-in-terms-of-integral}). Par
conséquent, le support de $B/A$
en tant que module $A$
est fermé et non vide (Algèbre, Lemmes \ref{algebra-lemma-support-closed} et \ref{algebra-lemma-support-zero}).
Soit $\mathfrak p \subset A$ un idéal premier
minimal du support. Après avoir remplacé
$A \subset B$ par $A_\mathfrak p \subset B_\mathfrak p$ (permis par
Lemme \ref{lemma-special-elements-and-localization}) nous pouvons supposer que $(A, \mathfrak m)$
est un anneau local, que $B$ est fini
sur $A$, et que $B/A$ a pour support $\{\mathfrak m\}$ en
tant que $A$-module. Puisque $B/A$
est un module fini, nous voyons que $I = \text{Ann}_A(B/A)$ satisfait $\mathfrak m = \sqrt{I}$
(Algèbre, Lemme \ref{algebra-lemma-support-closed}). Soit $\mathfrak m' \subset B$
l'idéal premier unique au-dessus de $\mathfrak m$. Comme
$\Spec(B) \to \Spec(A)$ est un homéomorphisme, nous constatons que $\mathfrak m' = \sqrt{IB}$.
Pour $f \in \mathfrak m'$ choisissons
$n \geq 1$ de telle sorte que $f^n \in IB$. Alors
aussi $f^{n + 1} \in IB$. Puisque $IB \subset A$
par notre choix de $I$ Nous concluons que
$f^n, f^{n + 1} \in A$. En utilisant
le Lemme \ref{lemma-nth-and-nplusone-implies-square-and-cube}, nous concluons que notre lemme
est vrai si $\mathfrak m' \not \subset A$. Cependant, si $\mathfrak m' \subset A$,
alors $\mathfrak m' = \mathfrak m$ et nous concluons que $A = B$ puisque les
corps résiduels sont également isomorphes
par hypothèse. Cette contradiction termine la
démonstration.
\end{proof}

\begin{lemma}
\label{lemma-pth-power-and-multiple}
Soit $A \subset B$ une extension d'anneau telle que $\Spec(B) \to \Spec(A)$ soit un homéomorphisme
universel. Si
$A \not = B$, alors soit il existe un $b \in B$, $b \not \in A$ avec $b^2 \in A$
et $b^3 \in A$, soit il existe un nombre premier $p$
et un $b \in B$, $b \not \in A$ avec $pb \in A$ et $b^p \in A$.
\end{lemma}

\begin{proof}
L'argument est presque exactement le même que dans la démonstration
du Lemme \ref{lemma-square-and-cube}, mais nous écrivons tout pour être
sûrs que cela fonctionne.

\medskip\noindent
Rappelons que l'extension $A \subset B$ est
entière (Lemme \ref{lemma-integral-universally-closed}). D'après le Lemme
\ref{lemma-subalgebra-inherits}, nous pouvons supposer
que $B$ est engendré par un seul élément
sur $A$. Par conséquent,
$B$ est fini sur $A$ (Algèbre, Lemme \ref{algebra-lemma-characterize-finite-in-terms-of-integral}). Ainsi,
le support de $B/A$ en tant que
$A$-module est fermé et
non vide (Algèbre, Lemmes \ref{algebra-lemma-support-closed} et \ref{algebra-lemma-support-zero}). Soit $\mathfrak p \subset A$
un idéal premier minimal du support.
Après avoir remplacé $A \subset B$ par
$A_\mathfrak p \subset B_\mathfrak p$ (permis par le Lemme \ref{lemma-special-elements-and-localization}) nous
pouvons supposer que $(A, \mathfrak m)$ est un anneau
local, que $B$ est fini sur $A$,
et que $B/A$ a pour support $\{\mathfrak m\}$ en tant que
$A$-module. Puisque $B/A$ est
un module fini, nous voyons que $I = \text{Ann}_A(B/A)$ satisfait $\mathfrak m = \sqrt{I}$
(Algèbre, Lemme \ref{algebra-lemma-support-closed}). Soit $\mathfrak m' \subset B$
l'idéal premier unique au-dessus de $\mathfrak m$.
Comme $\Spec(B) \to \Spec(A)$ est un homéomorphisme, nous trouvons que
$\mathfrak m' = \sqrt{IB}$. Pour $f \in \mathfrak m'$ choisissez
$n \geq 1$ de telle sorte que $f^n \in IB$.
Alors également $f^{n + 1} \in IB$.
Puisque $IB \subset A$ par notre choix de $I$ Nous
concluons que $f^n, f^{n + 1} \in A$.
En utilisant le Lemme \ref{lemma-nth-and-nplusone-implies-square-and-cube} nous concluons
que notre lemme est vrai si $\mathfrak m' \not \subset A$. Si $\mathfrak m' \subset A$,
alors $\mathfrak m' = \mathfrak m$. Puisque $A \not = B$ nous concluons
que l'application $\kappa = A/\mathfrak m \to B/\mathfrak m' = \kappa'$ du corps
résiduel ne peut pas être un isomorphisme.
Par le Lemme \ref{lemma-universally-injective} Nous concluons
que la caractéristique de $\kappa$ est un
nombre premier $p$ et que l'extension $\kappa'/\kappa$
est purement inséparable. Choisissez $b \in B$
dont l'image dans $\kappa'$ est un élément
non contenu dans $\kappa$ mais dont la puissance $p$
est dans $\kappa$. Alors $b \not \in A$, $b^p \in A$,
et $pb \in A$ (parce que $pb \in \mathfrak m' = \mathfrak m \subset A$) comme
souhaité.
\end{proof}

\begin{proposition}
\label{proposition-universal-homeomorphism-equal-residue-fields}
Soit $A \subset B$ une extension d'anneau. Les conditions suivantes sont équivalentes
\begin{enumerate}
\item $\Spec(B) \to \Spec(A)$ est un homéomorphisme universel induisant des isomorphismes
sur les corps résiduels, et
\item chaque sous-ensemble fini $E \subset B$ est contenu dans une extension
$$
A[b_1, \ldots, b_n] \subset B
$$
telle que $b_i^2, b_i^3 \in A[b_1, \ldots, b_{i - 1}]$ pour $i = 1, \ldots, n$.
\end{enumerate}
\end{proposition}

\begin{proof}
Supposons (1). En utilisant la récurrence transfinie, nous construisons
pour chaque ordinal $\alpha$ un $A$-sous-algèbre $B_\alpha \subset B$
comme suit. Posons $B_0 = A$. Si $\alpha$ est un ordinal limite,
alors nous posons $B_\alpha = \colim_{\beta < \alpha} B_\beta$.
Si $\alpha = \beta + 1$, alors soit $B_\beta = B$ auquel cas nous
posons $B_\alpha = B_\beta$ ou $B_\beta \not = B$, auquel cas nous appliquons le Lemme
\ref{lemma-square-and-cube} pour choisir un $b_\alpha \in B$, $b_\alpha \not \in B_\beta$
avec $b_\alpha^2, b_\alpha^3 \in B_\beta$ et nous posons $B_\alpha = B_\beta[b_\alpha] \subset B$.
Il est clair que $B = \colim B_\alpha$ (en fait
$B = B_\alpha$ pour un certain ordinal $\alpha$ comme on
le voit en regardant les cardinalités). Nous prouverons,
par induction transfinie, que (2) tient pour
$A \to B_\alpha$ pour chaque ordinal $\alpha$. Cela est
clair pour $\alpha = 0$. Supposons que l'énoncé soit
vrai pour chaque $\beta < \alpha$ et soit $E \subset B_\alpha$
un sous-ensemble fini.
Si $\alpha$ est un ordinal limite, alors
$B_\alpha = \bigcup_{\beta < \alpha} B_\beta$ et nous voyons que $E \subset B_\beta$ pour
un certain $\beta < \alpha$ ce qui prouve le résultat
dans ce cas. Si $\alpha = \beta + 1$, alors $B_\alpha = B_\beta[b_\alpha]$.
Ainsi tout $e \in E$ peut être écrit comme
un polynôme $e = \sum d_{e, i}b_\alpha^i$ avec $d_{e, i} \in B_\beta$.
Soit $D \subset B_\beta$ l'ensemble $D = \{d_{e, i}\} \cup \{b_\alpha^2, b_\alpha^3\}$. Par hypothèse
d'induction
il existe une $A$-sous-algèbre $A[b_1, \ldots, b_n] \subset B_\beta$ comme dans
l'énoncé du lemme contenant $D$.
Alors $A[b_1, \ldots, b_n, b_\alpha] \subset B_\alpha$ est une $A$-sous-algèbre
de $B_\alpha$ comme dans l'énoncé du lemme
contenant $E$.

\medskip\noindent
Supposons (2). Écrivons $B = \colim B_\lambda$ comme le colimite de ses $A$-sous-algèbres
finies.
D'après le Lemme \ref{lemma-colimit-inherits}, il suffit de montrer que $\Spec(B_\lambda) \to \Spec(A)$
est un homéomorphisme universel induisant
des isomorphismes sur les corps résiduels. Les compositions de
morphismes universellement fermés sont universellement fermées et il
en va de même pour les morphismes qui induisent des isomorphismes sur les corps résiduels.
Ainsi, il suffit de montrer que si $A \subset B$ et $B$ sont
générés par un seul élément $b$ avec $b^2, b^3 \in A$, alors
(1) tient. Une telle extension est intégrale et donc $\Spec(B) \to \Spec(A)$ est universellement
fermé et surjectif.
(Lemme \ref{lemma-integral-universally-closed} et Algèbre, Lemme \ref{algebra-lemma-integral-overring-surjective}).
Notez que $(b^2)^3 = (b^3)^2$ dans $A$. Pour
tout homomorphisme d'anneaux
$\varphi : A \to K$ vers un corps $K$, nous voyons qu'il
existe un $\lambda \in K$ avec $\varphi(b^2) = \lambda^2$ et $\varphi(b^3) = \lambda^3$.
À savoir, $\lambda = 0$ si $\varphi(b^2) = 0$ et $\lambda = \varphi(b^3)/\varphi(b^2)$ sinon.
Ainsi, $B \otimes_A K$ est un quotient de $K[x]/(x^2 - \lambda^2, x^3 - \lambda^3)$.
Cet anneau a exactement un premier avec corps résiduel
$K$. Cela implique que $\Spec(B) \to \Spec(A)$
est bijectif et induit des isomorphismes sur les
corps résiduels. Combiné avec la clôture universelle,
cela montre que (1) est vrai, voir lemmes \ref{lemma-universal-homeomorphism}
et \ref{lemma-universally-injective}.
\end{proof}

\begin{proposition}
\label{proposition-universal-homeomorphism}
Soit $A \subset B$ une extension d'anneau. Les conditions suivantes sont équivalentes
\begin{enumerate}
\item $\Spec(B) \to \Spec(A)$ est un homéomorphisme universel, et
\item chaque sous-ensemble fini $E \subset B$ est contenu dans une extension
$$
A[b_1, \ldots, b_n] \subset B
$$
telle que pour $i = 1, \ldots, n$ nous avons
\begin{enumerate}
\item $b_i^2, b_i^3 \in A[b_1, \ldots, b_{i - 1}]$, ou
\item il existe un nombre premier
$p$ avec $pb_i, b_i^p \in A[b_1, \ldots, b_{i - 1}]$.
\end{enumerate}
\end{enumerate}
\end{proposition}

\begin{proof}
La démonstration est exactement la même
que la démonstration de la Proposition \ref{proposition-universal-homeomorphism-equal-residue-fields} sauf pour
les changements suivants :
\begin{enumerate}
\item Utilisez le Lemme \ref{lemma-pth-power-and-multiple} au lieu du Lemme
\ref{lemma-square-and-cube} ce qui signifie que pour chaque ordinal
successeur $\alpha = \beta + 1$ nous avons
soit $b_\alpha^2, b_\alpha^3 \in B_\beta$ ou nous avons un nombre premier
$p$ et $pb_\alpha, b_\alpha^p \in B_\beta$.
\item Si $\alpha$ est un ordinal successeur,
alors prenez $D = \{d_{e, i}\} \cup \{b_\alpha^2, b_\alpha^3\}$ ou
prenez $D = \{d_{e, i}\} \cup \{pb_\alpha, b_\alpha^p\}$ selon le cas dans lequel
$\alpha$ se trouve.
\item Dans la démonstration de (2) $\Rightarrow$ (1), nous devons également considérer
le cas où $B$ est généré sur $A$ par un seul élément
$b$ avec $pb, b^p \in B$ pour un certain nombre premier $p$.
Ici, $A \subset B$ induit un homéomorphisme universel sur les spectres
par exemple d'après Algèbre, Lemme \ref{algebra-lemma-p-ring-map}.
\end{enumerate}
Cela termine la démonstration.
\end{proof}

\begin{lemma}
\label{lemma-universal-homeo-iso-if-invert-p}
Soit $p$ un nombre premier. Soit $A \to B$ un homomorphisme d'anneaux
qui induit un isomorphisme $A[1/p] \to B[1/p]$ (par exemple si
$p$ est nilpotent dans $A$). Les conditions
suivantes sont équivalentes
\begin{enumerate}
\item $\Spec(B) \to \Spec(A)$ est un homéomorphisme universel, et
\item le noyau de $A \to B$ est un idéal localement nilpotent et pour tout $b \in B$
il existe une puissance $p$-ième $q$ avec $qb$ et $b^q$
dans l'image de $A \to B$.
\end{enumerate}
\end{lemma}

\begin{proof}
Si (2) est vrai, alors (1) est vrai par Algèbre, Lemme \ref{algebra-lemma-p-ring-map}. Supposons
(1). Alors le noyau de $A \to B$ consiste en des
éléments nilpotents par Algèbre, Lemme \ref{algebra-lemma-image-dense-generic-points}. Ainsi, nous
pouvons remplacer $A$ par l'image de $A \to B$ et
supposer que $A \subset B$. Par Algèbre, Lemme \ref{algebra-lemma-help-with-powers}
l'ensemble
$$
B' = \{b \in B \mid p^nb, b^{p^n} \in A\text{ pour un certain }n \geq 0\}
$$
est une $A$-sous-algèbre de $B$ (le fait qu'elle soit fermée
sous les produits est trivial). Nous
devons montrer $B' = B$. Sinon, d'après le Lemme
\ref{lemma-pth-power-and-multiple}, il existe un $b \in B$, $b \not \in B'$ avec soit
$b^2, b^3 \in B'$, soit il existe un nombre premier $\ell$
avec $\ell b, b^\ell \in B'$. Nous allons
montrer que les deux cas mènent à une contradiction, prouvant ainsi le lemme.

\medskip\noindent Puisque
$A[1/p] = B[1/p]$, nous pouvons choisir une $p$-puissance $q$ telle que $qb \in A$.

\medskip\noindent Si
$b^2, b^3 \in B'$, alors aussi $b^q \in B'$. Par définition de $B'$, nous
trouvons que $(b^q)^{q'} \in A$ pour une certaine $p$-puissance
$q'$. Alors $qq'b, b^{qq'} \in A$, d'où $b \in B'$ ce qui est une contradiction.

\medskip\noindent
Supposons maintenant qu'il existe un nombre premier $\ell$ avec $\ell b, b^\ell \in B'$.
Si $\ell \not = p$ alors $\ell b \in B'$ et $qb \in A \subset B'$ impliquent $b \in B'$, une
contradiction. Ainsi $\ell = p$ et $b^p \in B'$ et nous obtenons
une contradiction exactement comme avant.
\end{proof}

\begin{lemma}
\label{lemma-make-universal-homeo}
Soit $A$ un anneau. Soit $x, y \in A$.
\begin{enumerate}
\item Si $x^3 = y^2$ est dans $A$, alors $A \to B = A[t]/(t^2 - x, t^3 - y)$ induit des bijections
sur les corps résiduels et un homéomorphisme
universel sur les spectres.
\item S'il existe un nombre premier $p$ tel que $p^px = y^p$ dans $A$,
alors $A \to B = A[t]/(t^p - x, pt - y)$ induit un homéomorphisme universel
sur les spectres.
\end{enumerate}
\end{lemma}

\begin{proof}
Nous utiliserons le critère du Lemme \ref{lemma-universal-homeomorphism} pour vérifier cela. Dans les deux
cas, l'homomorphisme d'anneaux est entier. Il suffit donc de montrer que, étant
donné un corps $k$ et un homomorphisme d'anneaux $\varphi : A \to k$,
l'algèbre $k$-$B \otimes_A k$ possède un idéal premier
unique dont le corps résiduel est égal à $k$ dans
le cas (1) et purement inséparable sur $k$
dans le cas (2). Voir le Lemme \ref{lemma-universally-injective}.

\medskip\noindent
Dans le cas (1), posons $\lambda = 0$ si
$\varphi(x) = 0$ et posons $\lambda = \varphi(y)/\varphi(x)$ sinon.
Alors $B = k[t]/(t^2 - \lambda^2, t^3 - \lambda^2)$. Ainsi le résultat
est clair.

\medskip\noindent Dans
le cas (2), si la caractéristique de $k$ est $p$,
alors nous obtenons $\varphi(y) = 0$ et $B = k[t]/(t^p - \varphi(x))$ qui est une algèbre
locale artinienne $k$ dont le corps résiduel est soit
$k$, soit une extension purement inséparable de degré $p$
de $k$. Si la caractéristique de $k$ n'est pas
$p$, alors en posant $\lambda = \varphi(y)/p$ nous voyons $B = k[t]/(t - \lambda) = k$ et
nous concluons également.
\end{proof}

\begin{lemma}
\label{lemma-universal-homeo-limit}
Soit $A \to B$ un homomorphisme d'anneaux.
\begin{enumerate}
\item Si $A \to B$ induit un homéomorphisme universel sur les spectres, alors
$B = \colim B_i$ est une colimite filtrante d'algèbres $A$ de présentation finie $B_i$
telle que $A \to B_i$ induit un homéomorphisme universel sur les spectres.
\item Si $A \to B$ induit des isomorphismes sur les corps résiduels et un
homéomorphisme universel sur les spectres, alors $B = \colim B_i$ est une colimite
filtrante d'algèbres $A$ de présentation finie $B_i$ telles que $A \to B_i$
induit des isomorphismes sur les corps résiduels et un homéomorphisme
universel sur les spectres.
\end{enumerate}
\end{lemma}

\begin{proof}
Démonstration de (1). Nous utiliserons le critère
de Algèbre, Lemme \ref{algebra-lemma-when-colimit}. Soit
$A \to C$ de présentation finie et soit $\varphi : C \to B$ un homomorphisme
d'$A$-algèbres. Soit $B' = \varphi(C) \subset B$ l'image.
Alors $A \to B'$ induit un homéomorphisme universel sur
les spectres par le Lemme \ref{lemma-subalgebra-inherits}.
Par Algèbre, Lemme \ref{algebra-lemma-ring-colimit-fp} nous pouvons
écrire $B' = \colim_{i \in I} B_i$ avec $A \to B_i$
de présentation finie et des homomorphismes
de transition
surjectives. Par Algèbre, Lemme \ref{algebra-lemma-characterize-finite-presentation} nous pouvons choisir
un indice $0 \in I$
et une factorisation $C \to B_0 \to B'$ de l'application $C \to B'$. Nous affirmons
que $\Spec(B_i) \to \Spec(A)$ est un homéomorphisme universel pour $i$ suffisamment
grand. Cette affirmation termine la démonstration de (1).

\medskip\noindent
Démonstration de l'affirmation. D'après le Lemme \ref{lemma-reduction-universal-homeomorphism} l'homomorphisme
d'anneaux $A_{red} \to B'_{red}$ induit un homéomorphisme
universel sur les spectres. Ainsi $A_{red} \subset B'_{red}$
par l'Algèbre, Lemme \ref{algebra-lemma-image-dense-generic-points}. En posant
$A' = \Im(A \to B')$ nous avons des surjections $A \to A' \to A_{red}$ induisant
des bijections $\Spec(A_{red}) = \Spec(A') = \Spec(A)$. Ainsi $A' \subset B'$ induit
un homéomorphisme universel sur les spectres. D'après
la Proposition \ref{proposition-universal-homeomorphism} et le fait que $B'$
est de type fini sur $A'$, nous
pouvons trouver $n$ et $b'_1, \ldots, b'_n \in B'$
tels que $B' = A'[b'_1, \ldots, b'_n]$
et tel que pour $j = 1, \ldots, n$ nous avons
\begin{enumerate}
\item $(b'_j)^2, (b'_j)^3 \in A'[b'_1, \ldots, b'_{j - 1}]$, ou
\item il existe un nombre premier
$p$ avec $pb'_j, (b'_j)^p \in A'[b'_1, \ldots, b'_{j - 1}]$.
\end{enumerate}
Choisissez $b_1, \ldots, b_n \in B_0$ en relevant $b'_1, \ldots, b'_n$. Pour $i \geq 0$ notez
$b_{j, i}$ l'image de $b_j$ dans $B_i$. Pour un $i$
suffisamment grand nous aurons pour $j = 1, \ldots, n$
\begin{enumerate}
\item $b_{j, i}^2, b_{j, i}^3 \in A_i[b_{1, i}, \ldots, b_{j - 1, i}]$, ou
\item il existe un nombre
premier $p$ avec $pb_{j, i}, b_{j, i}^p \in A_i[b_{1, i}, \ldots, b_{j - 1, i}]$.
\end{enumerate}
Voici $A_i \subset B_i$, l'image de $A \to B_i$. Observez que $A \to A_i$ est un
homomorphisme surjectif d'anneaux dont le noyau est un idéal
localement nilpotent. En augmentant encore $i$ si
nécessaire, nous pouvons supposer que $B_i$ est engendré
par $b_1, \ldots, b_n$ sur $A_i$,
en d'autres termes $B_i = A_i[b_1, \ldots, b_n]$. D'après
l'Algèbre, les Lemmes \ref{algebra-lemma-p-ring-map} et \ref{algebra-lemma-2-3-ring-map}
Nous concluons que $A \to A_i \to A_i[b_1] \to \ldots \to A_i[b_1, \ldots, b_n] = B_i$ induit des homéomorphismes universels
sur les spectres. Cela termine la démonstration de
l'assertion.

\medskip\noindent La
démonstration de (2) est exactement la même.
\end{proof}







\section{Normalisation faible absolue et semi-normalisation}
\label{section-seminormalization}

\noindent
Motivés par les résultats prouvés dans la section précédente,
nous donnons la définition suivante.

\begin{definition}
\label{definition-seminormal-ring}
Soit $A$ un anneau.
\begin{enumerate}
\item On dit que $A$ est {\it semi-normal} si pour
tout $x, y \in A$ avec $x^3 = y^2$ il existe un unique $a \in A$
avec $x = a^2$ et $y = a^3$.
\item On dit que $A$ est {\it absolument faiblement
normal} si (a) $A$
est semi-normal et (b) pour tout nombre premier $p$ et $x, y \in A$ avec
$p^px = y^p$ il existe un unique $a \in A$ avec $x = a^p$ et $y = pa$.
\end{enumerate}
\end{definition}

\noindent
Une observation amusante, voir \cite{Costa},
est que dans la définition des anneaux semi-normaux, il suffit\footnote{Soit
$A$ un anneau tel que pour tout $x, y \in A$ avec
$x^3 = y^2$, il existe un $a \in A$ avec $x = a^2$ et $y = a^3$.
Alors $A$ est réduit : si $x^2 = 0$, alors $x^2 = x^3$
et donc il existe un $a$ tel que $x = a^3$ et $x = a^2$. Puis
$x = a^3 = ax = a^4 = x^2 = 0$. Enfin, si $a_1^2 = a_2^2$ et $a_1^3 = a_2^3$ pour $a_1, a_2$ dans un
anneau réduit, alors $(a_1 - a_2)^3 = a_1^3 - 3a_1^2a_2 +3a_1a_2^2 - a_2^3 =
(1 - 3 + 3 - 1)a_1^3 = 0$ et donc $a_1 = a_2$.}
d'assumer l'existence de $a$. Les schémas absolument faiblement normaux ont été définis dans \cite[Appendix
B]{rydh_descent}.

\begin{lemma}
\label{lemma-seminormal-local-property}
Être semi-normal ou être absolument faiblement
normal est une propriété locale
des anneaux, voir Propriétés, Définition \ref{properties-definition-property-local}.
\end{lemma}

\begin{proof}
Supposons que $A$ soit semi-normal et $f \in A$. Soit $x', y' \in A_f$ avec
$(x')^3 = (y')^2$. Écrivons $x' = x/f^{2n}$ et $y' = y/f^{3n}$ pour certains $n \geq 0$ et
$x, y \in A$. Après avoir remplacé $x, y$ par $f^{2m}x, f^{3m}y$ et $n$
par $n + m$, nous voyons que $x^3 = y^2$ dans $A$.
Puis nous trouvons un unique $a \in A$ avec $x = a^2$ et
$y = a^3$. En posant $a' = a/f^n$, nous obtenons $x' = (a')^2$
et $y' = (a')^3$ comme souhaité. L'unicité de $a'$ découle
de l'unicité de $a$. De la même manière
exactement, le lecteur montre que si $A$ est absolument
faiblement normal, alors $A_f$
est absolument faiblement normal.

\medskip\noindent Supposons
que $A$ soit un anneau et que $f_1, \ldots, f_n \in A$ génère l'idéal unité.
Supposons que $A_{f_i}$ soit semi-normal pour chaque
$i$. Soit $x, y \in A$ avec $x^3 = y^2$. Pour chaque $i$ nous
trouvons un unique $a_i \in A_{f_i}$ avec $x = a_i^2$ et $y = a_i^3$ dans $A_{f_i}$.
Par unicité et le résultat du premier paragraphe (qui nous dit que
$A_{f_if_j}$ est semi-normal) nous voyons que $a_i$ et
$a_j$ s'envoient sur le même élément de $A_{f_if_j}$. Par l'algèbre,
Lemme \ref{algebra-lemma-standard-covering} nous trouvons un unique $a \in A$ se
projetant sur $a_i$ dans
$A_{f_i}$ pour tous les $i$. Ensuite $x = a^2$
et $y = a^3$ par le même principe. Clairement, ce $a$ est unique.
Ainsi $A$ est semi-normal. Si nous supposons que $A_{f_i}$ est absolument faiblement normal,
alors le même argument exact montre que $A$ est absolument faiblement normal.
\end{proof}

\noindent
Ensuite, nous définissons les schémas semi-normaux et les schémas absolument faiblement normaux.

\begin{definition}
\label{definition-seminormal}
Soit $X$ un schéma.
\begin{enumerate}
\item Nous disons que $X$ est {\it semi-normal}
si chaque $x \in X$ possède un voisinage ouvert affine $\Spec(R) = U \subset X$
tel que l'anneau $R$ soit semi-normal.
\item Nous disons que $X$ est {\it absolument faiblement normal}
si chaque $x \in X$ possède un voisinage ouvert affine $\Spec(R) = U \subset X$
tel que l'anneau $R$ soit absolument faiblement normal.
\end{enumerate}
\end{definition}

\noindent
Voici le lemme obligatoire.

\begin{lemma}
\label{lemma-locally-seminormal}
Soit $X$ un schéma. Les conditions suivantes sont équivalentes :
\begin{enumerate}
\item Le schéma $X$ est semi-normal.
\item Pour tout ouvert affine $U \subset X$, l'anneau $\mathcal{O}_X(U)$ est
semi-normal.
\item Il existe un recouvrement ouvert affine $X = \bigcup U_i$ tel que
chaque $\mathcal{O}_X(U_i)$ soit semi-normal.
\item Il existe un recouvrement ouvert $X = \bigcup X_j$ tel que chaque
sous-schéma ouvert $X_j$ soit semi-normal.
\end{enumerate}
De plus, si $X$ est semi-normal, alors
tout sous-schéma ouvert est semi-normal. Les mêmes affirmations
sont vraies lorsque ``semi-normal'' est remplacé par ``absolument
faiblement normal''.
\end{lemma}

\begin{proof}
Combiner les Propriétés, Lemme \ref{properties-lemma-locally-P}
et Lemme \ref{lemma-seminormal-local-property}.
\end{proof}

\begin{lemma}
\label{lemma-seminormal-reduced}
Un schéma ou anneau semi-normal est réduit. A fortiori, il en est de même
pour les schémas ou anneaux absolument faiblement normaux.
\end{lemma}

\begin{proof}
Soit $A$ un anneau.
Si $a \in A$ est non nul mais $a^2 = 0$, alors $a^2 = 0^2$ et $a^3 = 0^3$
et donc $A$ ne sont pas semi-normaux.
\end{proof}

\begin{lemma}
\label{lemma-seminormalization-ring}
Soit $A$ un anneau.
\begin{enumerate}
\item La catégorie des homomorphismes d'anneaux $A \to B$ induisant
un homéomorphisme universel sur les spectres a un objet final $A \to A^{awn}$.
\item Étant donné $A \to B$ dans la catégorie de (1) l'application résultante
$B \to A^{awn}$ est un isomorphisme si et seulement si $B$ est
absolument faiblement normal.
\item La catégorie des homomorphismes d'anneaux $A \to B$ induisant des isomorphismes
sur les corps résiduels et un homéomorphisme universel sur les spectres a un
objet final $A \to A^{sn}$.
\item Étant donné $A \to B$ dans la catégorie de (3) l'application résultante
$B \to A^{sn}$ est un isomorphisme si et seulement si $B$ est semi-normal.
\end{enumerate}
Pour tout homomorphisme d'anneaux $\varphi : A \to A'$ il
existe des applications uniques $\varphi^{awn} : A^{awn} \to (A')^{awn}$
et $\varphi^{sn} : A^{sn} \to (A')^{sn}$ compatibles avec $\varphi$.
\end{lemma}

\begin{proof}
Nous prouvons (1) et (2) et nous omettons la démonstration de (3) et (4) ainsi
que l'énoncé final. Considérons la catégorie des $A$-algèbres de la forme
$$
B = A[x_1, \ldots, x_n]/J
$$
où $J$ est un idéal de type fini tel que $A \to B$ définit un homéomorphisme
universel sur les spectres. Nous affirmons que cette catégorie
est dirigée (Catégories, Définition \ref{categories-definition-directed}). À savoir, étant donné
$$
B = A[x_1, \ldots, x_n]/J
\quad\text{et}\quad
B' = A[x_1, \ldots, x_{n'}]/J'
$$
alors nous pouvons considérer
$$
B'' = A[x_1, \ldots, x_{n + n'}]/J''
$$
où $J''$ est généré par les éléments de $J$ et
les éléments $f(x_{n + 1}, \ldots, x_{n + n'})$ où $f \in J'$. Ensuite,
nous avons des homomorphismes d'algèbre $A$ $B \to B''$
et $B' \to B''$ qui induisent un isomorphisme
$B \otimes_A B' \to B''$.
Il découle des Lemme \ref{lemma-base-change-universal-homeomorphism} et \ref{lemma-composition-universal-homeomorphism}
que $\Spec(B'') \to \Spec(A)$ est un homéomorphisme universel
et donc $A \to B''$ est dans notre catégorie. Enfin,
étant donné $\varphi, \varphi' : B \to B'$ dans notre
catégorie avec $B$ comme affiché
ci-dessus, nous considérons alors
le quotient $B''$ de $B'$
par l'idéal engendré par $\varphi(x_i) - \varphi'(x_i)$, $i = 1, \ldots, n$. Puisque $\Spec(B') = \Spec(B)$ nous
voyons que $\Spec(B'') \to \Spec(B')$ est une immersion bijective fermée, donc
un homéomorphisme universel. Ainsi $B''$ est dans notre
catégorie et $\varphi, \varphi'$ sont égalisés par $B' \to B''$. Ceci complète
la démonstration de notre affirmation. Nous posons
$$
A^{awn} = \colim B
$$
où la colimite est sur la catégorie précédemment décrite. Observez que
$A \to A^{awn}$ induit un homéomorphisme universel sur les spectres selon
le Lemme \ref{lemma-colimit-inherits} (c'est ici que nous utilisons que la catégorie est
dirigée).

\medskip\noindent Étant
donné un homomorphisme d'anneaux $A \to B$ de présentation finie induisant
un homéomorphisme universel sur les spectres, nous obtenons une application canonique
$B \to A^{awn}$ par la construction même de $A^{awn}$. Puisque chaque
$A \to B$ comme dans (1) est une colimite filtrante de $A \to B$ comme
dans (1) de présentation finie (Lemme \ref{lemma-universal-homeo-limit}), nous voyons
que $A \to A^{awn}$ est final dans la catégorie (1).

\medskip\noindent
Soient $x, y \in A^{awn}$ des éléments tels que $x^3 = y^2$.
Alors $A^{awn} \to A^{awn}[t]/(t^2 - x, t^3 - y)$ induit un homéomorphisme universel
sur les spectres par le Lemme \ref{lemma-make-universal-homeo}. Ainsi $A \to A^{awn}[t]/(t^2 - x, t^3 - y)$ est dans
la catégorie (1) et nous obtenons une unique application
d'$A$-algèbre $A^{awn}[t]/(t^2 - x, t^3 - y) \to A^{awn}$.
L'image $a \in A^{awn}$ de $t$ est donc
l'élément unique tel que $a^2 = x$ et $a^3 = y$
dans $A^{awn}$. De la même manière, donné un
premier $p$ et $x, y \in A^{awn}$ avec $p^px = y^p$, nous trouvons
un unique $a \in A^{awn}$ avec $a^p = x$ et $pq = y$.
Ainsi, $A^{awn}$ est absolument faiblement normal
par définition.

\medskip\noindent
Enfin, soit $A \to B$ dans la catégorie (1) avec $B$ absolument
faiblement normal. Puisque $A^{awn} \to B^{awn}$ induit un homéomorphisme
universel sur les spectres et puisque $A^{awn}$ est réduit
(Lemme \ref{lemma-seminormal-reduced}),
nous trouvons $A^{awn} \subset B^{awn}$ (voir Algèbre, Lemme \ref{algebra-lemma-image-dense-generic-points}).
Si cette inclusion n'est pas une égalité,
alors le Lemme \ref{lemma-pth-power-and-multiple} implique
qu'il existe un élément $b \in B^{awn}$, $b \not \in A^{awn}$ tel
que soit $b^2, b^3 \in A^{awn}$ soit $pb, b^p \in A^{awn}$ pour un certain
nombre premier $p$. Cependant, par l'existence et l'unicité
dans la Définition \ref{definition-seminormal-ring}, cela force $b \in A^{awn}$
et donc nous obtenons la contradiction qui termine la démonstration.
\end{proof}

\begin{lemma}
\label{lemma-seminormalization}
Soit $X$ un schéma.
\begin{enumerate}
\item La catégorie des homéomorphismes universels $Y \to X$
a un objet initial $X^{awn} \to X$.
\item Étant donné $Y \to X$ dans la catégorie de (1) le morphisme résultant $X^{awn} \to Y$
est un isomorphisme si et seulement si $Y$ est absolument
faiblement normal.
\item La catégorie des homéomorphismes universels $Y \to X$ qui induisent
des isomorphismes sur les corps résiduels a un objet initial
$X^{sn} \to X$.
\item Étant donné $Y \to X$ dans la catégorie de (3) le morphisme résultant $X^{sn} \to Y$
est un isomorphisme si et seulement si $Y$ est semi-normal.
\end{enumerate}
Pour tout morphisme $h : X' \to X$ de schémas, il existe des
morphismes uniques $h^{awn} : (X')^{awn} \to X^{awn}$ et $h^{sn} : (X')^{sn} \to X^{sn}$ compatibles avec
$h$.
\end{lemma}

\begin{proof}
Nous allons prouver (1) et (2) et omettre la démonstration
de (3) et (4). Soit $h : X' \to X$ un morphisme de schémas. Si (1)
vaut pour $X$ et $X'$, alors $X' \times_X X^{awn} \to X'$ est un homéomorphisme
universel et nous obtenons donc un morphisme unique $(X')^{awn} \to X' \times_X X^{awn}$
au-dessus de $X'$ par la propriété universelle de $(X')^{awn} \to X'$.
Composé avec la projection $X' \times_X X^{awn} \to X^{awn}$, nous obtenons $h^{awn}$.
Si de plus (2) vaut pour $X$ et $X'$ et que $h$ est
une immersion ouverte, alors $X' \times_X X^{awn}$ est
absolument faiblement normal (Lemme \ref{lemma-locally-seminormal}) et nous
en déduisons que $(X')^{awn} \to X' \times_X X^{awn}$ est un isomorphisme.

\medskip\noindent
Rappelez-vous que tout homéomorphisme universel
est affine, voir Lemme \ref{lemma-homeomorphism-affine}. Ainsi, si $X$
est affine, alors (1) et (2) en découlent
immédiatement d'après le Lemme \ref{lemma-seminormalization-ring}.
Soit $X$ un schéma et soit $\mathcal{B}$ l'ensemble
des ouverts affines de $X$. Pour chaque $U \in \mathcal{B}$
nous obtenons $U^{awn} \to U$ et pour $V \subset U$, $V, U \in \mathcal{B}$
nous obtenons un isomorphisme canonique $\rho_{V, U} : V^{awn} \to V \times_U U^{awn}$ selon la discussion
du paragraphe précédent. Ainsi, par recollement
relatif (Constructions, Lemme \ref{constructions-lemma-relative-glueing}), nous obtenons
un morphisme $X^{awn} \to X$ qui se restreint
à $U^{awn}$ sur $U$ de manière compatible avec le
$\rho_{V, U}$. Ensuite, soit $Y \to X$ un homéomorphisme
universel. Alors $U \times_X Y \to U$ est un homéomorphisme universel pour $U \in \mathcal{B}$
et nous obtenons un morphisme unique $g_U : U^{awn} \to U \times_X Y$ sur $U$. Ces $g_U$
sont compatibles avec les morphismes $\rho_{V, U}$ ; détails omis.
Par conséquent, il existe un morphisme unique $g : X^{awn} \to Y$ sur
$X$ concordant avec $g_U$ sur
$U$, voir Constructions, Remarque \ref{constructions-remark-relative-glueing-functorial}. Cela prouve (1) pour $X$.
La partie (2) s'ensuit car elle tient localement sur les affines.
\end{proof}

\begin{definition}
\label{definition-seminormalization}
Soit $X$ un schéma.
\begin{enumerate}
\item Le morphisme $X^{sn} \to X$ construit dans le Lemme \ref{lemma-seminormalization}
est la {\it semi-normalisation}
de $X$.
\item Le morphisme $X^{awn} \to X$ construit dans le Lemme \ref{lemma-seminormalization}
est la {\it normalisation
faible absolue} de $X$.
\end{enumerate}
\end{definition}

\noindent
Pour être sûr, la semi-normalisation $X^{sn}$ de $X$ est un schéma
semi-normal et la normalisation faible absolue $X^{awn}$ est un
schéma absolument faiblement normal.
De plus, pour tout morphisme $h : Y \to X$ de schémas, nous obtenons un
diagramme commutatif canonique
$$
\xymatrix{
Y^{awn} \ar[d]^{h^{awn}} \ar[r] &
Y^{sn} \ar[d]^{h^{sn}} \ar[r] &
Y \ar[d]^h \\
X^{awn} \ar[r] &
X^{sn} \ar[r] &
X
}
$$
de schémas ; les flèches $h^{sn}$ et $h^{awn}$ sont les uniques compatibles
avec $h$.

\begin{lemma}
\label{lemma-characterize-seminormal}
Soit $X$ un schéma. Les conditions suivantes sont équivalentes
\begin{enumerate}
\item $X$ est semi-normal,
\item $X$ est égal à sa propre semi-normalisation, c'est-à-dire que le morphisme
$X^{sn} \to X$ est un isomorphisme,
\item si $\pi : Y \to X$ est un homéomorphisme universel induisant des isomorphismes sur les corps
résiduels avec $Y$ réduit, alors $\pi$ est un isomorphisme.
\end{enumerate}
\end{lemma}

\begin{proof}
L'équivalence de (1) et (2) est claire d'après
le Lemme \ref{lemma-seminormalization}. Si (3) est vrai, alors $X^{sn} \to X$ est un isomorphisme
et nous voyons que (2) est vrai.

\medskip\noindent
Supposons que (2) vaut et soit $\pi : Y \to X$ un homéomorphisme universel
induisant des isomorphismes sur les corps résiduels avec $Y$ réduit.
Alors il existe une factorisation $X \to Y \to X$ de $\text{id}_X$
par le Lemme \ref{lemma-seminormalization}. Alors $X \to Y$ est une immersion fermée
(par Schémas, Lemme \ref{schemes-lemma-section-immersion} et le fait que $\pi$
est séparé, par exemple par le Lemme \ref{lemma-universally-injective-separated}).
Puisque $X \to Y$ est également une bijection sur les points, la caractéristique
d'être réduit de $Y$ montre qu'il doit s'agir
d'un isomorphisme. Cela termine la démonstration.
\end{proof}

\begin{lemma}
\label{lemma-characterize-absolutely-weakly-normal}
Soit $X$ un schéma. Les conditions suivantes sont équivalentes
\begin{enumerate}
\item $X$ est faiblement normal absolu,
\item $X$ est égal à sa propre normalisation faible absolue, c’est-à-dire que le morphisme
$X^{awn} \to X$ est un isomorphisme,
\item si $\pi : Y \to X$ est un homéomorphisme universel avec $Y$ réduit, alors $\pi$
est un isomorphisme.
\end{enumerate}
\end{lemma}

\begin{proof}
Cela se démontre exactement de la même
manière que le Lemme \ref{lemma-characterize-seminormal}.
\end{proof}





\section{Morphismes finis localement libres}
\label{section-finite-locally-free}

\noindent
Dans de nombreux articles, les auteurs utilisent des morphismes finis plats alors qu’ils
veulent réellement dire des morphismes finis localement libres. La raison en
est que si la base est localement noethérienne alors c’est la même chose. Mais
en général ce n’est pas le cas, voir Exercices, Exercice \ref{exercises-exercise-flat-not-projective}.

\begin{definition}
\label{definition-finite-locally-free}
Soit $f : X \to S$ un morphisme de schémas. On dit que $f$
est {\it fini localement libre} si $f$
est affine et que $f_*\mathcal{O}_X$ est un $\mathcal{O}_S$-module fini
localement libre. Dans ce cas, on dit que $f$ a {\it
rang} ou {\it degré}
$d$ si le faisceau $f_*\mathcal{O}_X$ est fini localement libre de degré $d$.
\end{definition}

\noindent
Remarquez que si $f : X \to S$ est fini localement libre alors $S$ est l'union disjointe
de ouverts et de sous-schémas fermés $S_d$ tels que $f^{-1}(S_d) \to S_d$ soit fini
localement libre de degré $d$.

\begin{lemma}
\label{lemma-finite-flat}
Soit $f : X \to S$ un morphisme de schémas. Les conditions
suivantes sont équivalentes :
\begin{enumerate}
\item $f$ est fini localement libre,
\item $f$ est fini, plat et localement de présentation finie.
\end{enumerate}
Si $S$ est localement noethérien, cela équivaut également à
\begin{enumerate}
\item[(3)] $f$ est fini et plat.
\end{enumerate}
\end{lemma}

\begin{proof}
Soit $V \subset S$ ouvert affine. Dans les trois cas, le morphisme
est affine, donc $f^{-1}(V)$ est affine. Ainsi, nous pouvons
écrire $V = \Spec(R)$ et $f^{-1}(V) = \Spec(A)$ pour une certaine $R$-algèbre $A$.
Supposons (1). Cela
signifie que nous pouvons couvrir $S$ par des ouverts affines
$V = \Spec(R)$ tels que $A$ est libre de type fini en tant que $R$-module.
Alors $R \to A$ est de présentation
finie par l’Algèbre, Lemme \ref{algebra-lemma-finite-finite-type}. Ainsi (2) est vérifié. Réciproquement,
supposons (2). Pour chaque ouvert
affine $V = \Spec(R)$ de $S$, le homomorphisme d’anneaux $R \to A$
est fini et de présentation finie et $A$ est plat en tant que
$R$-module. Par
Algèbre, Lemme \ref{algebra-lemma-finite-finitely-presented-extension} nous voyons que $A$ est de présentation
finie en tant que $R$-module.
Ainsi Algèbre, Lemme \ref{algebra-lemma-finite-projective} implique que $A$
est fini localement libre. Ainsi (1)
est vrai. Le cas noethérien découle du fait qu'un module
fini sur un anneau noethérien est un module
de présentation finie, voir Algèbre, Lemme \ref{algebra-lemma-Noetherian-finite-type-is-finite-presentation}.
\end{proof}

\begin{lemma}
\label{lemma-composition-finite-locally-free}
Une composition de morphismes finis localement libres est finie localement libre.
\end{lemma}

\begin{proof}
Omis.
\end{proof}

\begin{lemma}
\label{lemma-base-change-finite-locally-free}
Un changement de base d'un morphisme fini localement libre est fini localement libre.
\end{lemma}

\begin{proof}
Omis.
\end{proof}

\begin{lemma}
\label{lemma-finite-locally-free}
Soit $f : X \to S$ un morphisme fini localement libre de schémas. Il
existe une décomposition en union disjointe
$S = \coprod_{d \geq 0} S_d$ par des ouverts et sous-schémas fermés telle qu'en
posant $X_d = f^{-1}(S_d)$, les restrictions $f|_{X_d}$ soient
des morphismes finis localement libres $X_d \to S_d$ de degré
$d$.
\end{lemma}

\begin{proof}
Cela est vrai car un faisceau fini localement libre a localement
un rang bien défini. Détails omis.
\end{proof}

\begin{lemma}
\label{lemma-massage-finite}
Soit $f : Y \to X$ un morphisme fini avec $X$ affine.
Il existe un diagramme
$$
\xymatrix{
Z' \ar[rd] &
Y' \ar[l]^i \ar[d] \ar[r] &
Y \ar[d] \\
 & X' \ar[r] & X
}
$$
où
\begin{enumerate}
\item $Y' \to Y$ et $X' \to X$ sont surjectifs, finis et localement libres,
\item $Y' = X' \times_X Y$,
\item $i : Y' \to Z'$ est une immersion fermée,
\item $Z' \to X'$ est fini localement libre, et
\item $Z' = \bigcup_{j = 1, \ldots, m} Z'_j$ est une union finie (au sens des ensembles) de sous-schémas
fermés, chacun desquels se projette isomorphiquement sur
$X'$.
\end{enumerate}
\end{lemma}

\begin{proof}
Écrivez $X = \Spec(A)$ et $Y = \Spec(B)$. Voir aussi Plus
sur l'Algèbre, section \ref{more-algebra-section-descent-flatness-integral}. Soit $x_1, \ldots, x_n \in B$ des générateurs de $B$
sur $A$. Pour chaque $i$, nous pouvons choisir
un polynôme unitaire $P_i(T) \in A[T]$ tel que $P(x_i) = 0$ soit dans
$B$. D'après l'Algèbre,
Lemme \ref{algebra-lemma-adjoin-roots} (appliqué $n$ fois) il
existe une extension d'anneaux fini localement libre $A \subset A'$
telle que chaque $P_i$ se décompose complètement :
$$
P_i(T) = \prod\nolimits_{k = 1, \ldots, d_i} (T - \alpha_{ik})
$$
pour certains $\alpha_{ik} \in A'$. Posons
$$
C = A'[T_1, \ldots, T_n]/(P_1(T_1), \ldots, P_n(T_n))
$$
et $B' = A' \otimes_A B$. L'application $C \to B'$, $T_i \mapsto 1 \otimes x_i$ est une surjection d'$A'$-algèbres.
En posant $X' = \Spec(A')$, $Y' = \Spec(B')$ et $Z' = \Spec(C)$,
nous voyons que (1) -- (4) sont vérifiées. La partie
(5) est vérifiée parce que, du point de vue des ensembles,
$\Spec(C)$ est l'union des sous-schémas fermés
coupés par les idéaux
$$
(T_1 - \alpha_{1k_1}, T_2 - \alpha_{2k_2}, \ldots, T_n - \alpha_{nk_n})
$$
pour tout $1 \leq k_i \leq d_i$.
\end{proof}

\noindent
Le lemme suivant est énoncé dans la généralité
correcte dans le Lemme \ref{lemma-image-nowhere-dense-quasi-finite} ci-dessous.

\begin{lemma}
\label{lemma-image-nowhere-dense-finite}
Soit $f : Y \to X$ un morphisme fini de schémas. Soit $T \subset Y$
une partie fermée nulle part dense de $Y$. Alors $f(T) \subset X$
est une partie fermée nulle part dense de $X$.
\end{lemma}

\begin{proof}
D'après le Lemme \ref{lemma-finite-proper} nous savons que $f(T) \subset X$ est fermé. Soit $X = \bigcup X_i$
un recouvrement affine. Comme $T$ est
nulle part dense dans $Y$, nous voyons que $T \cap f^{-1}(X_i)$ est également
nulle part dense dans $f^{-1}(X_i)$. Ainsi, si nous pouvons prouver le théorème dans
le cas affine, alors nous voyons que $f(T) \cap X_i$ est nulle part dense.
Cela implique alors que $T$ est nulle part
dense dans $X$ selon la Topologie, Lemme \ref{topology-lemma-nowhere-dense-local}.

\medskip\noindent Supposons
que $X$ soit affine. Choisissez un diagramme
$$
\xymatrix{
Z' \ar[rd] &
Y' \ar[l]^i \ar[d]^{f'} \ar[r]_a &
Y \ar[d]^f \\
 & X' \ar[r]^b & X
}
$$
comme dans le Lemme \ref{lemma-massage-finite}. Les morphismes $a$,
$b$ sont ouverts puisqu'ils
sont finis localement libres (Lemmes \ref{lemma-finite-flat} et \ref{lemma-fppf-open}).
Par conséquent, $T' = a^{-1}(T)$
est nulle part dense, voir Topologie, Lemme \ref{topology-lemma-open-inverse-image-closed-nowhere-dense}. Le morphisme
$b$ est surjectif et ouvert.
Par conséquent, si nous pouvons prouver que
$f'(T') = b^{-1}(f(T))$ est nulle part dense, alors $f(T)$
est nulle part dense, voir Topologie, Lemme \ref{topology-lemma-open-inverse-image-closed-nowhere-dense}. Comme
$i$ est une immersion
fermée, d'après Topologie, Lemme \ref{topology-lemma-closed-image-nowhere-dense} nous voyons
que $i(T') \subset Z'$ est fermée et nulle part dense.
Ainsi, nous avons réduit le problème au cas discuté dans
le paragraphe suivant.

\medskip\noindent Supposons
que $Y = \bigcup_{i = 1, \ldots, n} Y_i$ soit une union finie de parties fermées, chacune s'envoyant
isomorphiquement sur $X$. Considérons $T_i = Y_i \cap T$. Si chacun
des $T_i$ est nulle part dense dans $Y_i$, alors chacun des $f(T_i)$
est nulle part dense dans $X$ car $Y_i \to X$ est un isomorphisme. Par conséquent,
$f(T) = f(T_i)$ est une union finie de sous-ensembles fermés nulle part
denses de $X$ et nous avons
gagné, voir Topologie, Lemme \ref{topology-lemma-nowhere-dense}. Supposons que ce ne
soit pas le cas, disons que $T_1$ contient un $V \subset Y_1$ ouvert non
vide. Nous allons montrer que cela conduit à une contradiction.
Considérons $Y_2 \cap V \subset V$. Ceci est soit un sous-ensemble
fermé propre, soit égal à $V$. Dans le premier cas, nous remplaçons
$V$ par $V \setminus V \cap Y_2$, donc $V \subset T_1$ est ouvert dans $Y_1$ et ne rencontre
pas $Y_2$. Dans le second cas, nous avons
$V \subset Y_1 \cap Y_2$ est ouvert à la fois dans $Y_1$ et $Y_2$.
Répétez séquentiellement avec $i = 3, \ldots, n$. Le résultat est une décomposition
en union disjointe
$$
\{1, \ldots, n\} = I_1 \amalg I_2, \quad 1 \in I_1
$$
et un $V$ ouvert de $Y_1$ contenu dans $T_1$ tel que
$V \subset Y_i$ pour $i \in I_1$ et $V \cap Y_i = \emptyset$ pour $i \in I_2$. Posons $U = f(V)$.
C'est un ouvert de $X$ puisque $f|_{Y_1} : Y_1 \to X$ est un
isomorphisme. Alors
$$
f^{-1}(U) = V\ \amalg\ \bigcup\nolimits_{i \in I_2} (Y_i \cap f^{-1}(U))
$$
Comme $\bigcup_{i \in I_2} Y_i$ est fermé, cela implique que $V \subset f^{-1}(U)$ est
ouvert, donc $V \subset Y$ est ouvert. Cela contredit l'hypothèse
que $T$ est quelque part dense dans $Y$, comme souhaité.
\end{proof}







\section{Applications rationnelles}
\label{section-rational-maps}

\noindent
Soit $X$ un schéma. Remarquez que si $U$, $V$ sont des ouverts denses dans $X$,
alors il en est de même de $U \cap V$.

\begin{definition}
\label{definition-rational-map}
Soit $X$, $Y$ des schémas.
\begin{enumerate}
\item Soit $f : U \to Y$, $g : V \to Y$ des morphismes de schémas définis sur des ouverts denses $U$,
$V$ de $X$. On dit que $f$ est {\it
équivalent} à $g$ si $f|_W = g|_W$ pour un certain $W \subset U \cap V$ ouvert
dense dans $X$.
\item Une {\it application rationnelle de $X$
vers $Y$} est une classe d'équivalence pour la relation d'équivalence définie en (1).
\item Si $X$, $Y$ sont des schémas sur un schéma de base $S$, nous disons qu'une
application rationnelle de $X$ vers $Y$ est une {\it $S$-application
rationnelle de $X$ vers $Y$} s'il existe un représentant $f : U \to Y$ de la classe
d'équivalence qui est un $S$-morphisme.
\end{enumerate}
\end{definition}

\noindent
Nous disons que deux morphismes $f$, $g$ comme dans (1) de la définition définissent
la même application rationnelle au lieu de dire qu'ils sont équivalents. Dans certains cas,
les applications rationnelles sont déterminées par des applications sur les anneaux locaux
aux points génériques.

\begin{lemma}
\label{lemma-rational-map-finite-presentation}
Soit $S$ un schéma. Soient $X$ et $Y$ des schémas sur $S$. Supposons
que $X$ ait un nombre fini de composantes irréductibles avec des
points génériques $x_1, \ldots, x_n$. Soit $s_i \in S$ l'image de $x_i$. Considérons
l'application
$$
\left\{
\begin{matrix}
S\text{-applications rationnelles} \\
\text{de }X\text{ vers }Y
\end{matrix}
\right\}
\longrightarrow
\left\{
\begin{matrix}
(y_1, \varphi_1, \ldots, y_n, \varphi_n)\text{ où }
y_i \in Y\text{ est au-dessus de }s_i\text{ et}\\
\varphi_i : \mathcal{O}_{Y, y_i} \to \mathcal{O}_{X, x_i}
\text{ est un homomorphisme local de }\mathcal{O}_{S, s_i}\text{-algèbres}
\end{matrix}
\right\}
$$
qui envoie $f : U \to Y$ sur le $2n$-uplet
avec $y_i = f(x_i)$ et $\varphi_i = f^\sharp_{x_i}$. Alors
\begin{enumerate}
\item Si $Y \to S$ est localement de type fini, alors l'application est injective.
\item Si $Y \to S$ est localement de présentation finie, alors l'application est bijective.
\item Si $Y \to S$ est localement de type fini et que $X$ est réduit, alors
l'application est bijective.
\end{enumerate}
\end{lemma}

\begin{proof}
Remarquez que tout ouvert dense de $X$ contient les points $x_i$,
donc la construction a du sens. Pour prouver (1)
ou (2), nous pouvons remplacer $X$ par tout ouvert dense. Ainsi, si
$Z_1, \ldots, Z_n$ sont les composantes irréductibles de $X$, alors
nous pouvons remplacer $X$ par $X \setminus \bigcup_{i \not = j} Z_i \cap Z_j$. Après avoir
fait cela, $X$ est l'union disjointe de ses composantes irréductibles
(considérées comme ouvertes et sous-schémas fermés). Ensuite,
à la fois le côté droit et le côté gauche de la flèche sont des produits
sur les composantes irréductibles et nous réduisons au cas où $X$
est irréductible.

\medskip\noindent Supposons
que $X$ est irréductible avec le point générique $x$ situé
au-dessus de $s \in S$. La partie
(1) découle de la partie (1) du Lemme \ref{lemma-morphism-defined-local-ring}. Les
parties (2) et (3) découlent de la partie (2 du même lemme.
\end{proof}

\begin{definition}
\label{definition-rational-function}
Soit $X$ un schéma. Une fonction rationnelle {\it sur $X$} est une
application rationnelle de $X$ vers $\mathbf{A}^1_{\mathbf{Z}}$.
\end{definition}

\noindent
Voir Constructions, Définition \ref{constructions-definition-affine-n-space} pour la définition de la droite affine
$\mathbf{A}^1$. Soit $X$ un schéma sur $S$. Pour tout ouvert
$U \subset X$, un morphisme $U \to \mathbf{A}^1_{\mathbf{Z}}$ est le
même qu'un morphisme $U \to \mathbf{A}^1_S$ sur $S$. Par conséquent,
une fonction rationnelle est également la même qu'une
$S$-application rationnelle de $X$ vers $\mathbf{A}^1_S$.

\medskip\noindent
Rappelons que nous avons l'identification
canonique $\Mor(T, \mathbf{A}^1_{\mathbf{Z}}) = \Gamma(T, \mathcal{O}_T)$ pour tout schéma $T$,
voir Schémas, Exemple \ref{schemes-example-global-sections}. Ainsi $\mathbf{A}^1_{\mathbf{Z}}$ est
un objet anneau dans la catégorie des schémas.
Plus précisément, les morphismes
\begin{eqnarray*}
+ : \mathbf{A}^1_{\mathbf{Z}} \times \mathbf{A}^1_{\mathbf{Z}}
& \longrightarrow
& \mathbf{A}^1_{\mathbf{Z}}
\\ (f, g) & \longmapsto & f + g \\ *
: \mathbf{A}^1_{\mathbf{Z}} \times \mathbf{A}^1_{\mathbf{Z}}
& \longrightarrow
& \mathbf{A}^1_{\mathbf{Z}}
\\ (f, g) & \longmapsto & fg
\end{eqnarray*}
satisfaire tous les axiomes de l'addition et de la multiplication dans un anneau (commutatif
avec $1$ comme toujours). Par conséquent, l'ensemble des applications
rationnelles dans $\mathbf{A}^1_{\mathbf{Z}}$ possède également une structure d'anneau naturelle.

\begin{definition}
\label{definition-ring-of-rational-functions}
Soit $X$ un schéma. L'anneau {\it des fonctions rationnelles sur $X$}
est l'anneau $R(X)$ dont les éléments sont des fonctions rationnelles avec
l'addition et la multiplication telles que décrites ci-dessus.
\end{definition}

\noindent
Pour les schémas avec un nombre fini de composantes irréductibles, nous pouvons le calculer.

\begin{lemma}
\label{lemma-integral-scheme-rational-functions}
Soit $X$ un schéma avec un nombre fini de composantes irréductibles
$X_1, \ldots, X_n$. Si $\eta_i \in X_i$ est le point générique, alors
$$
R(X) = \mathcal{O}_{X, \eta_1} \times \ldots \times \mathcal{O}_{X, \eta_n}
$$
Si $X$ est réduit, cela équivaut à $\prod \kappa(\eta_i)$.
Si $X$ est intègre, alors $R(X) = \mathcal{O}_{X, \eta} = \kappa(\eta)$ est un
corps.
\end{lemma}

\begin{proof}
Soit $U \subset X$ un ouvert dense. Alors
$U_i = (U \cap X_i) \setminus (\bigcup_{j \not = i} X_j)$ est non vide et ouvert car il contient
$\eta_i$, contenu dans $X_i$, et $\bigcup U_i \subset U \subset X$
est dense. Ainsi, l'identification
dans le lemme provient de la chaîne d'égalités
\begin{align*}
R(X)
& =
\colim_{U \subset X\text{ ouvert dense}} \Mor(U, \mathbf{A}^1_\mathbf{Z}) \\
& =
\colim_{U \subset X\text{ ouvert dense}} \mathcal{O}_X(U) \\
& =
\colim_{\eta_i \in U_i \subset X\text{ ouvert}} \prod \mathcal{O}_X(U_i) \\
& =
\prod \colim_{\eta_i \in U_i \subset X\text{ ouvert}} \mathcal{O}_X(U_i) \\
& =
\prod \mathcal{O}_{X, \eta_i}
\end{align*}
où la deuxième égalité
est Schémas, Exemple \ref{schemes-example-global-sections}.
L'énoncé final
découle d'Algèbre, Lemme \ref{algebra-lemma-minimal-prime-reduced-ring}.
\end{proof}

\begin{definition}
\label{definition-function-field}
Soit $X$ un schéma intègre. La fonction {\it
corps}, ou le {\it corps des fonctions rationnelles }
de $X$ est le corps $R(X)$.
\end{definition}

\noindent
Nous pouvons occasionnellement indiquer ce corps $k(X)$ au lieu de
$R(X)$. Nous pouvons utiliser la notion de corps de fonctions
pour élucider la condition de séparation sur
un schéma intègre. Notez que d'après le Lemme \ref{lemma-integral-scheme-rational-functions} sur un schéma
intègre, chaque anneau local $\mathcal{O}_{X, x}$ peut être considéré comme un sous-anneau
local de $R(X)$.

\begin{lemma}
\label{lemma-distinct-local-rings}
Soit $X$ un schéma intègre séparé.
Soient $Z_1$ et $Z_2$ des sous-ensembles fermés irréductibles
distincts de $X$. Soit
$\eta_i$ le point générique
de $Z_i$. Si $Z_1 \not\subset Z_2$, alors $\mathcal{O}_{X, \eta_1} \not \subset \mathcal{O}_{X, \eta_2}$ comme sous-anneaux
de $R(X)$. En
particulier, si $Z_1 = \{x\}$ consiste en un point fermé $x$,
il existe une fonction régulière dans un voisinage de
$x$ qui n'est pas dans $\mathcal{O}_{X, \eta_{2}}$.
\end{lemma}

\begin{proof}
Observez d'abord que sous l'hypothèse que $X$ est séparé,
il existe une application unique
de schémas $\Spec(\mathcal{O}_{X, \eta_2}) \to X$ sur $X$ telle
que la composition
$$
\Spec(R(X)) \longrightarrow
\Spec(\mathcal{O}_{X, \eta_2}) \longrightarrow X
$$
soit l'application canonique $\Spec(R(X)) \to X$.
À savoir, il existe l'application
canonique $can : \Spec(\mathcal{O}_{X, \eta_2}) \to X$, voir Schémas,
Équation (\ref{schemes-equation-canonical-morphism}). Étant donné un second morphisme $a$
vers $X$, nous avons que $a$ coïncide
avec $can$ sur le
point générique de $\Spec(\mathcal{O}_{X, \eta_2})$ par hypothèse.
Maintenant, le fait que $X$ soit séparé
garantit que le sous-ensemble
dans $\Spec(\mathcal{O}_{X, \eta_2})$ où ces deux applications
coïncident est fermé, voir Schémas, Lemme
\ref{schemes-lemma-where-are-they-equal}. Ainsi $a = can$ sur tout $\Spec(\mathcal{O}_{X, \eta_2})$.

\medskip\noindent
Supposons $Z_1 \not \subset Z_2$ et supposons au contraire
que $\mathcal{O}_{X, \eta_{1}} \subset \mathcal{O}_{X, \eta_{2}}$ comme sous-anneaux de $R(X)$.
Alors nous obtiendrions un second morphisme
$$
\Spec(\mathcal{O}_{X, \eta_{2}}) \longrightarrow
\Spec(\mathcal{O}_{X, \eta_{1}}) \longrightarrow
X.
$$
D'après ce qui précède, cette composition devrait être égale
à $can$. Cela implique que $\eta_2$ se spécialise
en $\eta_1$ (voir Schémas, Lemme \ref{schemes-lemma-specialize-points}).
Mais cela contredit notre hypothèse $Z_1 \not \subset Z_2$.
\end{proof}

\begin{definition}
\label{definition-domain-of-definition}
Soit $\varphi$ une application rationnelle entre deux schémas $X$ et $Y$. Nous disons
que $\varphi$ est {\it défini en un point $x \in X$} s'il
existe un représentant $(U, f)$ de $\varphi$ avec $x \in U$. Le {\it
domaine de définition} de $\varphi$ est l'ensemble de tous les
points où $\varphi$ est défini.
\end{definition}

\noindent
Avec cette définition, il n'est pas vrai en général que $\varphi$
possède un représentant défini sur tout le domaine de définition.

\begin{lemma}
\label{lemma-rational-map-from-reduced-to-separated}
Soient $X$ et $Y$ des schémas. Supposons $X$ réduit et $Y$ séparé.
Soit $\varphi$ une application rationnelle de $X$ vers $Y$ avec domaine de
définition $U \subset X$. Alors il existe un unique morphisme $f : U \to Y$ représentant $\varphi$.
Si $X$ et $Y$ sont des schémas sur un schéma séparé $S$ et
si $\varphi$ est une application rationnelle $S$, alors $f$ est un morphisme
sur $S$.
\end{lemma}

\begin{proof}
Soient $(V, g)$ et $(V', g')$ des représentants de $\varphi$.
Alors $g, g'$ coïncident sur un sous-schéma ouvert dense
$W \subset V \cap V'$. D'un autre côté, l'égaliseur $E$ de $g|_{V \cap V'}$ et $g'|_{V \cap V'}$
est un sous-schéma fermé de $V \cap V'$ (Schémas,
Lemme \ref{schemes-lemma-where-are-they-equal}). Maintenant $W \subset E$ implique que
$E = V \cap V'$ d'un point de vue des ensembles. Comme $V \cap V'$
est réduit, nous concluons $E = V \cap V'$ d'un point de vue schématique,
c'est-à-dire, $g|_{V \cap V'} = g'|_{V \cap V'}$. Il s'ensuit que nous pouvons recoller
les représentants $g : V \to Y$ de $\varphi$ en un morphisme $f : U \to Y$,
voir Schémas, Lemme \ref{schemes-lemma-glue}. Nous
omettons la démonstrations de l'énoncé final.
\end{proof}

\noindent
En général, il n'est pas logique de composer des applications rationnelles. La raison est
que l'image d'un représentant de la première application rationnelle peut avoir une intersection
vide avec le domaine de définition de la seconde. Cependant, si nous supposons
que nos schémas sont irréductibles et que nous considérons des applications rationnelles
dominantes, alors nous pouvons composer des applications rationnelles.

\begin{definition}
\label{definition-dominant-rational}
Soient $X$ et $Y$ des schémas irréductibles. Une application rationnelle de $X$ vers
$Y$ est appelée {\it dominante} si tout représentant $f : U \to Y$ est un morphisme
de schémas dominant.
\end{definition}

\noindent
D'après le Lemme \ref{lemma-dominant-finite-number-irreducible-components}, il est équivalent d'exiger que le point générique
$\eta \in X$ s'envoie sur le point générique $\xi$ de $Y$,
c'est-à-dire $f(\eta) = \xi$ pour tout représentant $f : U \to Y$.
Nous pouvons composer une application rationnelle dominante $\varphi$ entre
des schémas irréductibles $X$ et $Y$ avec une application rationnelle
arbitraire $\psi$ de $Y$ vers
$Z$. À savoir, choisir des représentants $f : U \to Y$
avec $U \subset X$ ouvert dense et $g : V \to Z$ avec $V \subset Y$
ouvert dense. Alors $W = f^{-1}(V) \subset X$ est ouvert non vide (car il contient
le point générique de $X$) et nous laissons $\psi \circ \varphi$
être la classe d'équivalence de $g \circ f|_W : W \to Z$. Nous omettons la vérification
que cela est bien défini.

\medskip\noindent De cette
manière, nous obtenons une catégorie dont les objets sont des schémas irréductibles et dont les morphismes
sont des applications rationnelles dominantes. Étant donné un schéma de base $S$, nous pouvons
de manière similaire définir une catégorie dont les objets sont des schémas irréductibles au-dessus de $S$
et dont les morphismes sont des $S$-applications rationnelles dominantes.

\begin{definition}
\label{definition-birational-schemes}
Soient $X$ et $Y$ des schémas irréductibles.
\begin{enumerate}
\item Nous disons que $X$ et $Y$ sont {\it birationnels} si $X$ et $Y$
sont isomorphes dans la catégorie des schémas irréductibles et des applications rationnelles dominantes.
\item Supposons que $X$ et $Y$ sont des schémas sur un schéma de base $S$. On dit
que $X$ et $Y$ sont {\it $S$-birationnels} si $X$
et $Y$ sont isomorphes dans la catégorie des schémas irréductibles sur $S$ et des applications
rationnelles dominantes $S$.
\end{enumerate}
\end{definition}

\noindent
Si $X$ et $Y$ sont des schémas irréductibles birationnels, alors l'ensemble des applications
rationnelles de $X$ vers $Z$ est bijectif avec l'ensemble des applications rationnelles
de $Y$ vers $Z$ pour tous les schémas $Z$ (fonctoriellement en $Z$).
Pour des schémas irréductibles ``généraux'', ceci
n'est qu'une définition possible. Une autre serait de demander que $X$ et $Y$
aient des anneaux de fonctions rationnelles isomorphes.
Pour les variétés, ces conditions sont équivalentes, voir Lemme \ref{lemma-criterion-birational-finite-presentation}.

\begin{lemma}
\label{lemma-birational-integral}
Soient $X$ et $Y$ des schémas irréductibles.
\begin{enumerate}
\item Les schémas $X$ et $Y$ sont birationnels si et seulement si ils ont des
ouverts non vides isomorphes.
\item Supposons que $X$ et $Y$ soient des schémas sur un schéma de base $S$.
Alors $X$ et $Y$ sont $S$-birationnels si et seulement s'il
existe des ouverts non vides $U \subset X$ et $V \subset Y$ qui sont $S$-isomorphes.
\end{enumerate}
\end{lemma}

\begin{proof}
Supposons que $X$ et $Y$ soient birationnels. Soient $f : U \to Y$ et $g : V \to X$
définissant des applications rationnelles dominantes inverses de $X$ vers $Y$
et de $Y$ vers $X$. Nous pouvons supposer $V$ affine. Nous pouvons remplacer
$U$ par un ouvert affine de $f^{-1}(V)$. Comme $g \circ f$ est l'identité
en tant qu'application rationnelle dominante, nous voyons que la composition $U \to V \to X$
est l'identité sur un ouvert dense de $U$. Ainsi, après avoir remplacé $U$
par un ouvert affine plus petit, nous pouvons supposer que $U \to V \to X$ est
l'inclusion de $U$ dans
$X$. Il s'ensuit que $U \to V$ est une immersion (appliquer Schémas,
Lemme \ref{schemes-lemma-section-immersion} à $U \to g^{-1}(U) \to U$).
Cependant, en inversant les rôles de $U$ et $V$ et en refaisant
l'argument ci-dessus, nous voyons qu'il existe un ouvert affine non vide $V' \subset V$
tel que l'inclusion se factorise en $V' \to U \to V$. Alors $V' \to U$ est nécessairement
une immersion ouverte. En effet, $V' \to f^{-1}(V') \to V'$ sont des monomorphismes (Schémas,
Lemme
\ref{schemes-lemma-immersions-monomorphisms}) se composant pour donner l'identité, donc des isomorphismes.
Ainsi $V'$ est isomorphe à un ouvert à
la fois de $X$ et $Y$. Dans le cas des $S$-applications
rationnelles, le même argument exact fonctionne.
\end{proof}

\begin{remark}
\label{remark-generalize-category}
Voici une généralisation de la catégorie des schémas irréductibles et des applications
rationnelles dominantes. Pour un schéma $X$, notons $X^0$ l'ensemble
des points $x \in X$ avec $\dim(\mathcal{O}_{X, x}) = 0$, en d'autres termes, $X^0$ est l'ensemble
des points génériques des composantes irréductibles de $X$. Ensuite,
nous pouvons considérer la catégorie avec
\begin{enumerate}
\item objets étant les schémas $X$ tels que chaque ouvert quasi-compact possède
un nombre fini de composantes irréductibles, et
\item morphismes de $X$ vers $Y$ étant des applications rationnelles $f : U \to Y$
de $X$ vers $Y$ telles que $f(U^0) = Y^0$.
\end{enumerate}
Si $U \subset X$ est un ouvert dense d'un schéma, alors $U^0 \subset X^0$
n'a pas besoin d'être une égalité, mais si $X$ est un objet
de notre catégorie, alors c'est le cas. Ainsi, étant
donnés deux morphismes dans notre catégorie, la composition est
bien définie et un morphisme dans notre catégorie.
\end{remark}

\begin{remark}
\label{remark-pseudo-morphisms}
Il existe une variante de la Définition \ref{definition-rational-map} où nous considérons
uniquement les morphismes $U \to Y$ définis sur des sous-schémas
ouverts schématiquement denses $U \subset X$. Nous utilisons le
Lemme \ref{lemma-intersection-scheme-theoretically-dense} pour voir que nous obtenons une relation d'équivalence.
Une classe d'équivalence de ceux-ci est appelée
un {\it pseudo-morphisme de $X$
à $Y$}. Si $X$
est réduit, les deux notions coïncident.
\end{remark}






\section{Morphismes birationnels}
\label{section-birational}

\noindent
Vous êtes peut-être habitué à la notion d'application birationnelle
de variétés ayant la propriété d'être un isomorphisme sur
un ouvert de la cible. Cependant, en général, un morphisme
birationnel peut ne pas être un isomorphisme sur
un ouvert non vide, voir Exemple \ref{example-birational-not-iso-over-open}. Voici
la définition formelle.

\begin{definition}
\label{definition-birational}
\begin{reference}
\cite[(2.2.9)]{EGA1}
\end{reference}
Soient $X$, $Y$ des schémas. Supposons que $X$ et $Y$ aient un nombre
fini de composantes irréductibles. On dit qu'un morphisme $f : X \to Y$ est {\it
birationnel} si
\begin{enumerate}
\item $f$ induit une bijection entre l'ensemble des points génériques des
composantes irréductibles de $X$ et l'ensemble des points génériques
des composantes irréductibles de $Y$, et
\item pour chaque point générique $\eta \in X$ d'une composante irréductible
de $X$,
l'homomorphisme d'anneaux locaux $\mathcal{O}_{Y, f(\eta)} \to \mathcal{O}_{X, \eta}$ est un
isomorphisme.
\end{enumerate}
\end{definition}

\noindent
Nous verrons ci-dessous que les fibres d'un morphisme birationnel au-dessus
des points génériques sont des singletons. De plus, nous verrons que dans la
plupart des cas rencontrés en pratique, l'existence d'un morphisme birationnel entre des schémas irréductibles
$X$ et $Y$ implique que $X$ et
$Y$ sont des schémas birationnels.

\begin{lemma}
\label{lemma-birational-dominant}
Soit $f : X \to Y$ un morphisme de schémas ayant un nombre fini de composantes
irréductibles. Si $f$ est birationnel alors $f$
est dominant.
\end{lemma}

\begin{proof}
Cela découle du Lemme \ref{lemma-generic-points-in-image-dominant} et de la
définition.
\end{proof}

\begin{lemma}
\label{lemma-birational-generic-fibres}
Soit $f : X \to Y$ un morphisme birationnel de schémas ayant un
nombre fini de composantes irréductibles. Si $y \in Y$
est le point générique d'une composante
irréductible, alors le changement de base $X \times_Y \Spec(\mathcal{O}_{Y, y}) \to \Spec(\mathcal{O}_{Y, y})$ est un
isomorphisme.
\end{lemma}

\begin{proof}
Nous pouvons supposer que $Y = \Spec(B)$ est affine et irréductible.
Alors $X$ est irréductible aussi. Si nous prouvons
le résultat pour tout ouvert affine non vide $U \subset X$,
alors le résultat vaut pour $X$ (petit argument
omis). Par conséquent, nous
pouvons supposer que $X$ est affine aussi, disons $X = \Spec(A)$.
Soit $y \in Y$ correspondant au premier idéal minimal $\mathfrak q \subset B$.
D'après l'hypothèse, $A$ a un unique premier idéal minimal $\mathfrak p$
au-dessus de $\mathfrak q$ et $B_\mathfrak q \to A_\mathfrak p$ est un isomorphisme. Il
s'ensuit que $A_\mathfrak q \to \kappa(\mathfrak p)$ est surjectif, donc $\mathfrak p A_\mathfrak q$ est un idéal maximal.
D'autre part, $\mathfrak p A_\mathfrak q$ est le seul premier idéal minimal de $A_\mathfrak q$.
Nous concluons que $\mathfrak p A_\mathfrak q$ est le seul idéal premier de
$A_\mathfrak q$ et que $A_\mathfrak q = A_\mathfrak p$. Puisque
$A_\mathfrak q = A \otimes_B B_\mathfrak q$, le lemme s'ensuit.
\end{proof}

\begin{example}
\label{example-birational-not-iso-over-open}
Voici deux exemples de morphismes birationnels qui ne sont pas des isomorphismes
sur aucun ouvert de la cible.

\medskip\noindent
Premier exemple. Soit $k$ un corps infini. Soit
$A = k[x]$. Soit $B = k[x, \{y_{\alpha}\}_{\alpha \in k}]/
((x-\alpha)y_\alpha, y_\alpha y_\beta)$. Il existe
une inclusion $A \subset B$ et une rétraction $B \to A$ définissant tous les $y_\alpha$
égaux à zéro. Le morphisme $\Spec(A) \to \Spec(B)$
ainsi que le morphisme $\Spec(B) \to \Spec(A)$
sont birationnels mais ne sont pas des
isomorphismes sur aucun ouvert.

\medskip\noindent Second
exemple. Soit $A$ un domaine. Soit $S \subset A$ un sous-ensemble multiplicatif
ne contenant pas
$0$. Avec $B = S^{-1}A$, le morphisme $f : \Spec(B) \to \Spec(A)$ est birationnel.
S'il existe un ouvert $U$ de $\Spec(A)$ tel
que $f^{-1}(U) \to U$ soit un isomorphisme, alors il existe un $a \in A$ tel
que chaque élément de $S$ devienne inversible dans la localisation
principale $A_a$. En prenant $A = \mathbf{Z}$ et $S$,
l'ensemble des entiers impairs donne un contre-exemple.
\end{example}

\begin{lemma}
\label{lemma-birational-birational}
Soit $f : X \to Y$ un morphisme birationnel de schémas ayant un nombre fini de composantes
irréductibles sur un schéma de base $S$. Supposons que l'une des
conditions suivantes soit remplie
\begin{enumerate}
\item $f$ est localement de type fini et $Y$ réduit,
\item $f$ est localement de présentation finie.
\end{enumerate}
Alors il existe des ouverts denses $U \subset X$ et $V \subset Y$ tels que $f(U) \subset V$
et $f|_U : U \to V$ est un isomorphisme. En particulier, si $X$ et
$Y$ sont irréductibles, alors $X$ et $Y$
sont $S$-birationnels.
\end{lemma}

\begin{proof}
Il y a une réduction immédiate au cas où $X$ et $Y$ sont irréductibles
que nous omettons. De plus, après avoir encore restreint, nous pouvons
supposer que $X$ et $Y$ sont affines,
disons $X = \Spec(A)$ et $Y = \Spec(B)$. Par hypothèse, $A$, resp.\ $B$
a un unique idéal premier minimal $\mathfrak p$, resp.\ $\mathfrak q$,
le premier $\mathfrak p$
est au-dessus de $\mathfrak q$, et $B_\mathfrak q = A_\mathfrak p$.
Par le Lemme \ref{lemma-birational-generic-fibres} nous avons $B_\mathfrak q = A_\mathfrak q = A_\mathfrak p$.

\medskip\noindent
Supposons que $B \to A$ soit de type fini, disons $A = B[x_1, \ldots, x_n]$.
Il existe un $b_i \in B$ et $g_i \in B \setminus \mathfrak q$ tels que $b_i/g_i$ s'envoie
à l'image de $x_i$ dans $A_\mathfrak q$. Par conséquent,
$b_i  - g_ix_i$ s'envoie sur zéro dans $A_{g_i'}$ pour
un certain $g_i' \in B \setminus \mathfrak q$. En posant $g = \prod g_i g'_i$ nous voyons que
$B_g \to A_g$ est surjectif. Si de plus $Y$ est réduit, alors
l'application $B_g \to B_\mathfrak q$ est injective et donc $B_g \to A_g$
est également injective. Cela prouve le cas (1).

\medskip\noindent
Démonstration de (2). Par l'argument donné dans le paragraphe précédent,
nous pouvons supposer que $B \to A$ est surjectif. Comme $f$
est localement de présentation finie, le noyau $J \subset B$
est un idéal de type fini. Disons $J = (b_1, \ldots, b_r)$. Puisque
$B_\mathfrak q = A_\mathfrak q$, il existe $g_i \in B \setminus \mathfrak q$ tels que $g_i b_i = 0$.
En posant $g = \prod g_i$, nous voyons que $B_g \to A_g$ est un
isomorphisme.
\end{proof}

\begin{lemma}
\label{lemma-criterion-birational-finite-presentation}
Soit $S$ un schéma. Soient $X$ et $Y$ des schémas irréductibles
localement de présentation finie sur $S$. Soient $x \in X$ et $y \in Y$
les points génériques. Les conditions suivantes sont équivalentes
\begin{enumerate}
\item $X$ et $Y$ sont $S$-birationnels,
\item il existe des ouverts non vides de $X$ et $Y$
qui sont $S$-isomorphes, et
\item $x$ et $y$ s'envoient sur le même point
$s$ de $S$ et $\mathcal{O}_{X, x}$ et $\mathcal{O}_{Y, y}$ sont isomorphes
en tant qu'algèbres $\mathcal{O}_{S, s}$.
\end{enumerate}
\end{lemma}

\begin{proof}
Nous avons vu l'équivalence de (1) et (2) dans
le Lemme \ref{lemma-birational-integral}. Il est immédiat
que (2) implique (3). Pour finir,
nous supposons que (3) est vrai et nous
prouvons (1). Par le Lemme \ref{lemma-rational-map-finite-presentation}, il existe une
application rationnelle $f : U \to Y$
qui envoie $x \in U$ sur $y$
et induit l'isomorphisme donné $\mathcal{O}_{Y, y} \cong \mathcal{O}_{X, x}$. Ainsi $f$ est
un morphisme birationnel et induit donc un isomorphisme
sur des ouverts non
vides par le Lemme \ref{lemma-birational-birational}. Cela termine
la démonstration.
\end{proof}

\begin{lemma}
\label{lemma-common-open}
Soit $S$ un schéma. Soient $X$ et $Y$ des schémas intègres localement
de type fini sur $S$. Soient $x \in X$ et $y \in Y$ les points génériques. Les conditions
suivantes sont équivalentes
\begin{enumerate}
\item $X$ et $Y$ sont $S$-birationnels,
\item il existe des ouverts non vides de $X$ et $Y$ qui sont $S$-isomorphes, et
\item $x$ et $y$ ont la même image
$s \in S$ et $\kappa(x) \cong \kappa(y)$ en tant qu'extensions $\kappa(s)$.
\end{enumerate}
\end{lemma}

\begin{proof}
Nous avons vu l'équivalence de (1) et
(2) dans le Lemme \ref{lemma-birational-integral}. Il
est immédiat que (2) implique
(3). Pour finir, nous supposons que
(3) est vrai et nous prouvons (1). Observez
que $\mathcal{O}_{X, x} = \kappa(x)$ et $\mathcal{O}_{Y, y} = \kappa(y)$
par l'Algèbre, Lemme \ref{algebra-lemma-minimal-prime-reduced-ring}. Par le Lemme
\ref{lemma-rational-map-finite-presentation} il existe une application rationnelle
$f : U \to Y$ qui envoie $x \in U$
sur $y$ et induit
l'isomorphisme donné $\mathcal{O}_{Y, y} \cong \mathcal{O}_{X, x}$. Ainsi $f$ est
un morphisme birationnel et induit donc
un isomorphisme sur des ouverts
non vides par le Lemme \ref{lemma-birational-birational}. Cela termine
la démonstration.
\end{proof}







\section{Morphismes génériquement finis}
\label{section-generically-finite}

\noindent
Dans cette section, nous caractérisons les applications entre schémas
qui sont localement de type fini et qui sont ``génériquement fini'' d'une
certaine manière.

\begin{lemma}
\label{lemma-generically-finite}
Soient $X$, $Y$
des schémas. Soit $f : X \to Y$ localement de type fini. Soit
$\eta \in Y$ un point générique d'une composante irréductible de $Y$.
Les conditions suivantes sont équivalentes :
\begin{enumerate}
\item l'ensemble $f^{-1}(\{\eta\})$ est fini,
\item il existe des ouverts affines $U_i \subset X$,
$i = 1, \ldots, n$ et $V \subset Y$ avec $f(U_i) \subset V$,
$\eta \in V$ et $f^{-1}(\{\eta\}) \subset \bigcup U_i$ tels
que chaque $f|_{U_i} : U_i \to V$ soit fini.
\end{enumerate}
Si $f$ est quasi-séparé, alors ceci est également équivalent à
\begin{enumerate}
\item[(3)] il existe des ouverts affines
$V \subset Y$ et $U \subset X$ avec $f(U) \subset V$,
$\eta \in V$ et $f^{-1}(\{\eta\}) \subset U$ tels
que $f|_U : U \to V$ soit fini.
\end{enumerate}
Si $f$ est quasi-compact et quasi-séparé, alors
ceux-ci sont également équivalents à
\begin{enumerate}
\item[(4)] il existe un ouvert affine $V \subset Y$, $\eta \in V$
tel que $f^{-1}(V) \to V$ soit fini.
\end{enumerate}
\end{lemma}

\begin{proof}
La question est locale sur la base. Par conséquent, nous pouvons remplacer $Y$ par
un voisinage affine de $\eta$, et nous pouvons et faisons l'hypothèse tout au long de
la démonstration ci-dessous que $Y$ est affine, disons $Y = \Spec(R)$.

\medskip\noindent
Il est clair que (2) implique (1).
Supposons que $f^{-1}(\{\eta\}) = \{\xi_1, \ldots, \xi_n\}$ soit fini. Choisissez des ouverts affines
$U_i \subset X$ avec $\xi_i \in U_i$. Par Algèbre, Lemme \ref{algebra-lemma-generically-finite}
nous voyons qu'après avoir remplacé $Y$ par
un ouvert principal dans $Y$, chacun
des morphismes $U_i \to Y$ est fini. En d'autres
termes, (2) est vrai.

\medskip\noindent Il
est clair que (3) implique (1). Supposons que $f$ soit quasi-séparé et
(1). Écrivons $f^{-1}(\{\eta\}) = \{\xi_1, \ldots, \xi_n\}$. Il n'y a pas de spécialisations parmi les $\xi_i$
d'après le Lemme \ref{lemma-finite-fibre}. Puisque chaque $\xi_i$ s'envoie sur le
point générique $\eta$ d'une composante
irréductible de $Y$, il ne peut y avoir de spécialisation non triviale
$\xi \leadsto \xi_i$ dans $X$ (puisque $\xi$ s'enverrait
également sur $\eta$). Nous concluons
que chaque $\xi_i$ est un point générique d'une
composante irréductible de $X$.
Puisque $Y$ est affine et $f$ quasi-séparé, nous voyons que
$X$ est quasi-séparé. D'après les Propriétés, Lemme \ref{properties-lemma-maximal-points-affine}.
nous pouvons trouver un ouvert affine $U \subset X$ contenant chaque $\xi_i$.
Par l'Algèbre, Lemme \ref{algebra-lemma-generically-finite} nous voyons qu'après avoir remplacé
$Y$ par un ouvert principal dans $Y$
les morphismes $U \to Y$ sont finis. En
autres termes (3) est vrai.

\medskip\noindent Il
est clair que (4) implique toutes les conditions (1) -- (3) sans autre hypothèse
sur $f$. Supposons que $f$ soit quasi-compact et quasi-séparé. Nous
devons montrer que les conditions équivalentes (1) -- (3) impliquent
(4). Soient $U$, $V$ comme dans (3). Remplacez $Y$ par $V$.
Puisque $f$ est quasi-compact et $Y$ est quasi-compact (étant
affine), nous voyons que $X$ est quasi-compact. Par conséquent,
$Z = X \setminus U$ est quasi-compact, donc le morphisme $f|_Z : Z \to Y$ est quasi-compact.
Par la construction de $Z$,
nous voyons que $\eta \not \in f(Z)$. Par conséquent, d'après le Lemme
\ref{lemma-quasi-compact-generic-point-not-in-image}, nous voyons qu'il existe un
ouvert affine voisinage $V'$ de $\eta$ dans $Y$ tel que $f^{-1}(V') \cap Z = \emptyset$.
Alors nous avons $f^{-1}(V') \subset U$ et cela signifie
que $f^{-1}(V') \to V'$ est fini.
\end{proof}

\begin{example}
\label{example-counter-generically-finite}
Soit $A = \prod_{n \in \mathbf{N}} \mathbf{F}_2$. Chaque élément de $A$ est idempotent.
Ainsi, tout idéal premier est maximal avec corps résiduel $\mathbf{F}_2$.
Ainsi, la topologie sur $X = \Spec(A)$
est totalement déconnectée et quasi-compacte. Les applications
de projection $A \to \mathbf{F}_2$ définissent des points ouverts de $\Spec(A)$. Il
ne peut pas être le cas que tous les points de $X$ soient ouverts
puisque $X$ est quasi-compact. Soit $x \in X$ un point fermé
qui n'est pas ouvert. Alors nous pouvons former un schéma $Y$ constitué
de deux copies de $X$ collées le long de $X \setminus \{x\}$. En
d'autres termes, il s'agit de $X$
avec $x$ doublé, comparer Schémas, Exemple \ref{schemes-example-affine-space-zero-doubled}. Le
morphisme
$f : Y \to X$ est quasi-compact, de type fini et a des fibres finies mais n'est
pas quasi-séparé. Le point
$x \in X$ est un point générique d'une composante irréductible de $X$
(puisque $X$ est totalement déconnecté). Mais les propriétés (3)
et (4) du Lemme \ref{lemma-generically-finite} ne sont pas vérifiées. La raison est que pour tout
voisinage ouvert $x \in U \subset X$, l'image réciproque $f^{-1}(U)$ n'est pas affine
parce que les fonctions sur $f^{-1}(U)$ ne peuvent pas séparer les deux
points situés au-dessus de $x$ (démonstration omise ; c'est un
bel exercice). Ainsi, la condition que $f$ soit quasi-séparé est nécessaire
dans les parties (3) et (4) du lemme.
\end{example}

\begin{remark}
\label{remark-quasi-finite-finite-over-dense-open}
Une alternative
au Lemme \ref{lemma-generically-finite} est
l'énoncé qu'un morphisme quasi-fini est fini
sur un ouvert dense de la cible. Cela
sera démontré
dans Compléments sur les morphismes, Lemme \ref{more-morphisms-lemma-quasi-finite-finite-over-dense-open}.
\end{remark}

\begin{lemma}
\label{lemma-quasi-finiteness-over-generic-point}
Soient $X$, $Y$ des schémas. Soit $f : X \to Y$ localement de type fini. Soient $X^0$,
resp. \ $Y^0$ l'ensemble des points génériques des composantes irréductibles
de $X$, resp. \ $Y$. Soit $\eta \in Y^0$. Les conditions suivantes sont
équivalentes
\begin{enumerate}
\item $f^{-1}(\{\eta\}) \subset X^0$,
\item $f$ est quasi-fini en tous les points au-dessus de $\eta$,
\item $f$ est quasi-fini en tous $\xi \in X^0$ au-dessus de $\eta$.
\end{enumerate}
\end{lemma}

\begin{proof}
La condition (1) implique qu'il n'y a pas de spécialisations parmi les points
de la fibre $X_\eta$. Par conséquent, (2)
est vrai par le Lemme \ref{lemma-quasi-finite-at-point-characterize}. L'implication (2) $\Rightarrow$
(3) est immédiate. Puisque $\eta$ est un point
générique de $Y$, les points génériques de $X_\eta$ sont des points
génériques de $X$. Par conséquent,
(3) et le Lemme \ref{lemma-quasi-finite-at-point-characterize} impliquent que les points génériques
de $X_\eta$ sont également fermés. Ainsi, tous les
points de $X_\eta$ sont génériques et nous voyons que (1) est vrai.
\end{proof}

\begin{lemma}
\label{lemma-finite-over-dense-open}
Soient $X$, $Y$ des schémas. Soit $f : X \to Y$ localement de type fini. Soient $X^0$,
resp. \ $Y^0$ l'ensemble des points génériques des composantes irréductibles de
$X$, resp. \ $Y$. Supposons
\begin{enumerate}
\item $X^0$ et $Y^0$ sont finis et $f^{-1}(Y^0) = X^0$,
\item soit $f$ est quasi-compact, soit $f$ est séparé.
\end{enumerate}
Alors il existe un ouvert dense $V \subset Y$
tel que $f^{-1}(V) \to V$ soit fini.
\end{lemma}

\begin{proof}
Puisque $Y$ a un nombre fini de composantes irréductibles, nous pouvons trouver un ouvert
dense qui est une union disjointe de ses composantes irréductibles. Ainsi, nous pouvons supposer
que $Y$ est un affine irréductible avec un point générique $\eta$. Alors la fibre au-dessus
de $\eta$ est finie car $X^0$ est fini.

\medskip\noindent
Supposons que $f$ soit séparé et $Y$ affine irréductible.
Choisissez $V \subset Y$ et $U \subset X$ comme dans le Lemme \ref{lemma-generically-finite} partie
(3). Puisque $f|_U : U \to V$ est fini, nous voyons que $U \subset f^{-1}(V)$
est fermé ainsi qu'ouvert (Lemmes \ref{lemma-image-proper-scheme-closed} et \ref{lemma-finite-proper}). Ainsi $f^{-1}(V) = U \amalg W$
pour un certain sous-schéma ouvert $W$ de $X$. Cependant,
puisque $U$ contient tous les points génériques de $X$,
nous concluons que $W = \emptyset$ comme souhaité.

\medskip\noindent
Supposons que $f$ soit quasi-compact et $Y$ affine
irréductible. Alors $X$ est quasi-compact, donc il existe
un sous-schéma ouvert
dense $U \subset X$ qui est séparé (Propriétés, Lemme \ref{properties-lemma-quasi-compact-dense-open-separated}). Comme
l'ensemble des points génériques $X^0$ est fini, nous voyons que $X^0 \subset U$.
Ainsi $\eta \not \in f(X \setminus U)$. Comme $X \setminus U \to Y$ est quasi-compact, nous
concluons qu'il existe un ouvert non vide $V \subset Y$ tel que $f^{-1}(V) \subset U$,
voir Lemme \ref{lemma-quasi-compact-dominant}. Après avoir remplacé $X$ par $f^{-1}(V)$
et $Y$ par $V$, nous nous réduisons au cas
séparé que nous avons traité dans le paragraphe précédent.
\end{proof}

\begin{lemma}
\label{lemma-birational-isomorphism-over-dense-open}
Soient $X$, $Y$ des schémas. Soit $f : X \to Y$ un morphisme birationnel
entre des schémas qui ont un nombre fini de composantes irréductibles.
Supposons
\begin{enumerate}
\item soit $f$ est quasi-compact, soit $f$ est séparé, et
\item soit $f$ est localement de type fini et $Y$ est réduit, soit $f$
est localement de présentation finie.
\end{enumerate}
Alors, il existe un ouvert dense $V \subset Y$
tel que $f^{-1}(V) \to V$ soit un isomorphisme.
\end{lemma}

\begin{proof}
Par le Lemme \ref{lemma-finite-over-dense-open} nous pouvons supposer que $f$ est fini. Comme $Y$
a un nombre fini de composantes irréductibles, nous pouvons trouver un ouvert dense
qui est une union disjointe de ses composantes irréductibles. Ainsi nous pouvons
supposer que $Y$ est irréductible. Par le Lemme \ref{lemma-birational-birational} nous trouvons un ouvert
non vide $U \subset X$ tel que $f|_U : U \to Y$ est une immersion ouverte. Après avoir retiré
le sous-ensemble fermé (comme $f$ est fini)
$f(X \setminus U)$ de $Y$ nous voyons que $f$ est un isomorphisme.
\end{proof}

\begin{lemma}
\label{lemma-finite-degree}
Soient $X$, $Y$ des schémas intègres.
Soit $f : X \to Y$ localement de type fini. Supposons que
$f$ est dominant. Les conditions
suivantes sont équivalentes :
\begin{enumerate}
\item l'extension $R(Y) \subset R(X)$ a un degré de transcendance
$0$,
\item l'extension $R(Y) \subset R(X)$ est finie,
\item il existe des ouverts affines non vides $U \subset X$
et $V \subset Y$ tels que $f(U) \subset V$
et $f|_U : U \to V$ soit fini, et
\item le point générique de $X$ est le seul point de $X$ qui s'envoie
sur le point générique de $Y$.
\end{enumerate}
Si $f$ est séparé ou si $f$ est quasi-compact, alors cela équivaut
également à
\begin{enumerate}
\item[(5)] il existe un ouvert affine non vide $V \subset Y$ tel
que $f^{-1}(V) \to V$ soit fini.
\end{enumerate}
\end{lemma}

\begin{proof}
Choisissez des ouverts affines $\Spec(A) = U \subset X$ et $\Spec(R) = V \subset Y$
tels que $f(U) \subset V$. Alors $R$ et $A$
sont des anneaux intègres par définition. L'homomorphisme
d'anneaux $R \to A$ est
de type fini (Lemme \ref{lemma-locally-finite-type-characterize}). D'après le Lemme \ref{lemma-dominant-finite-number-irreducible-components},
le point générique de $X$ s'envoie sur le point générique
de $Y$, donc $R \to A$ est injectif. Soit
$K = R(Y)$ le corps des fractions de $R$ et $L = R(X)$ le
corps des fractions de $A$. Alors $L/K$ est une
extension de corps de type fini. Ainsi, nous voyons que (1)
est équivalent à (2).

\medskip\noindent
Supposons que (2) soit vrai. Soient $x_1, \ldots, x_n \in A$ des générateurs de
$A$ sur $R$. D'après l'hypothèse, il existe des polynômes
non nuls $P_i(X) \in R[X]$ tels que $P_i(x_i) = 0$. Soit $f_i \in R$ le coefficient
principal de $P_i$. Nous concluons alors
que $R_{f_1 \ldots f_n} \to A_{f_1 \ldots f_n}$ est fini, c'est-à-dire que (3) est vrai. Notez que (3) implique
(2). Nous voyons donc maintenant que (1), (2) et (3) sont tous
équivalents.

\medskip\noindent
Soit $\eta$ le point générique de $X$, et soit $\eta' \in Y$
le point générique de $Y$.
Supposons (4). Alors $\dim_\eta(X_{\eta'}) = 0$ et nous voyons que $R(X) = \kappa(\eta)$
a un degré de transcendance $0$ sur $R(Y) = \kappa(\eta')$
par le Lemme \ref{lemma-dimension-fibre-at-a-point}. En d'autres
termes, (1) est vrai. Supposons les conditions équivalentes (1),
(2) et (3). Supposons que $x \in X$ soit un point se
mappant sur $\eta'$. Comme $x$
est une spécialisation de $\eta$, cela donne des inclusions
$R(Y) \subset \mathcal{O}_{X, x} \subset R(X)$, ce qui implique que $\mathcal{O}_{X, x}$
est un corps, voir Algèbre, Lemme \ref{algebra-lemma-integral-over-field}.
Ainsi $x = \eta$. Nous voyons donc que (1) -- (4) sont
tous équivalents.

\medskip\noindent
Il est clair que (5) implique (3) sans hypothèses supplémentaires
sur $f$. Il reste à prouver que si $f$ est soit
séparé soit quasi-compact, alors les conditions équivalentes (1)
-- (4) impliquent (5). Cela résulte du lemme \ref{lemma-finite-over-dense-open}.
\end{proof}

\begin{definition}
\label{definition-degree}
Soient $X$ et $Y$ des
schémas intègres. Soit $f : X \to Y$ localement de type
fini et dominant. Supposons $[R(X) : R(Y)] < \infty$, ou toute autre des conditions
équivalentes (1) -- (4) du Lemme \ref{lemma-finite-degree}.
Alors l'entier positif
$$
\deg(X/Y) = [R(X) : R(Y)]
$$
est appelé le degré {\it de $X$ sur $Y$}.
\end{definition}

\noindent
Il est possible d’étendre cette notion à un morphisme $f : X \to Y$
si (a) $Y$ est intègre avec point générique $\eta$,
(b) $f$ est localement de type fini, et (c) $f^{-1}(\{\eta\})$ est fini. Dans ce
cas, nous pouvons définir
$$
\deg(X/Y)
=
\sum\nolimits_{\xi \in X, \ f(\xi) = \eta}
\dim_{R(Y)} (\mathcal{O}_{X, \xi}).
$$
À savoir, étant donné que $R(Y) = \kappa(\eta) = \mathcal{O}_{Y, \eta}$ (Lemme \ref{lemma-integral-scheme-rational-functions})
les dimensions ci-dessus sont finies par
le Lemme \ref{lemma-generically-finite} ci-dessus.
Cependant, pour la plupart des
applications, la définition donnée ci-dessus est
la bonne.

\begin{lemma}
\label{lemma-degree-composition}
Soient $X$, $Y$, $Z$
des schémas intègres. Soient $f : X \to Y$ et $g : Y \to Z$ des morphismes
dominants localement de type fini. Supposons que $[R(X) : R(Y)] < \infty$
et $[R(Y) : R(Z)] < \infty$. Alors
$$
\deg(X/Z) = \deg(X/Y) \deg(Y/Z).
$$
\end{lemma}

\begin{proof}
Cela provient de la multiplicativité des degrés
dans les tours d'extensions
finies de corps, voir Corps, Lemme \ref{fields-lemma-multiplicativity-degrees}.
\end{proof}

\begin{remark}
\label{remark-definition-generically-finite}
Soit $f : X \to Y$ un morphisme de schémas qui est localement de type fini. Il y a (au moins) deux
propriétés que nous pourrions utiliser pour définir les morphismes {\it
génériquement finis}. Celles-ci correspondent à la question de savoir
si vous voulez que la propriété soit locale sur la source ou locale sur le but :
\begin{enumerate}
\item (Locale sur le but ; suggéré par Ravi Vakil.) Supposons
que tout ouvert quasi-compact de $Y$ ait un
nombre fini de composantes irréductibles (par exemple si $Y$ est localement
noethérien). L'exigence est que l'image réciproque de chaque point générique
soit finie, voir Lemme \ref{lemma-generically-finite}.
\item (Local sur la source.) L'exigence est qu'il existe un ouvert
dense $U \subset X$ tel que $U \to Y$ soit localement quasi-fini.
\end{enumerate}
Dans le cas (1), l'exigence peut être formulée sans la condition auxiliaire sur $Y$,
mais ne donne probablement pas la bonne notion pour les schémas en général.
La propriété (2) telle que formulée n'implique pas que les fibres au-dessus des points génériques
soient finies ; cependant, si $f$ est quasi-compact
et que $Y$ est comme dans (1), alors c'est le cas.
\end{remark}

\begin{definition}
\label{definition-modification}
Soit $X$ un schéma intègre. Une modification {\it de $X$}
est un morphisme birationnel propre $f : X' \to X$ avec $X'$
intègre.
\end{definition}

\noindent
Soit $f : X' \to X$ une modification comme dans la définition. D'après
le Lemme \ref{lemma-finite-degree} il existe un $U \subset X$ non vide tel que $f^{-1}(U) \to U$ soit
fini. Par platitude générique (Proposition \ref{proposition-generic-flatness})
nous pouvons supposer que $f^{-1}(U) \to U$ est plat et
de présentation finie. Ainsi $f^{-1}(U) \to U$ est fini localement libre (Lemme
\ref{lemma-finite-flat}). Puisque $f$ est birationnel, le degré de $X'$ sur
$X$ est $1$. Par conséquent, $f^{-1}(U) \to U$ est fini
localement libre de degré $1$, autrement dit c'est un isomorphisme.
Ainsi nous pouvons {\it redéfinir}
une modification pour être un morphisme propre $f : X' \to X$ de schémas
intègres tel que $f^{-1}(U) \to U$ soit un isomorphisme pour un ouvert non vide
$U \subset X$.

\begin{definition}
\label{definition-alteration}
\begin{reference}
\cite[Définition 2.20]{alterations}
\end{reference}
Soit $X$ un schéma intègre. Une {\it altération
de $X$} est un morphisme propre dominant $f : Y \to X$ avec $Y$
intègre tel que $f^{-1}(U) \to U$ soit fini pour un certain ouvert non vide $U \subset X$.
\end{definition}

\noindent
Ceci est la définition donnée dans \cite{alterations}, sauf que ici nous ne supposons
pas que $X$ et $Y$ soient noethériens. En argumentant comme
ci-dessus, nous voyons qu'une altération est un morphisme propre dominant $f : Y \to X$
de schémas intègres qui induit une extension finie de
corps de fonctions, c’est-à-dire, telle que les conditions équivalentes
du Lemme \ref{lemma-finite-degree} soient vérifiées.



\section{La formule de dimension}
\label{section-dimension-formula}

\noindent
Pour les morphismes entre schémas noethériens, nous pouvons dire un peu plus
sur les dimensions des anneaux locaux. Voici un résultat important (et pas si
difficile à prouver). Rappelons que $R(X)$ désigne le corps des fonctions d'un
schéma intègre $X$.

\begin{lemma}
\label{lemma-dimension-formula}
Soit $S$ un schéma.
Soit $f : X \to S$ un morphisme de schémas. Soit
$x \in X$, et posons $s = f(x)$.
Supposons
\begin{enumerate}
\item $S$ localement noethérien,
\item $f$ localement de type fini,
\item $X$ et $S$ intègres, et
\item $f$ dominant.
\end{enumerate}
Nous avons
\begin{equation}
\label{equation-dimension-formula}
\dim(\mathcal{O}_{X, x})
\leq
\dim(\mathcal{O}_{S, s}) + \text{trdeg}_{R(S)}R(X)
- \text{trdeg}_{\kappa(s)} \kappa(x).
\end{equation}
De plus, l'égalité est vérifiée si $S$ est universellement caténaire.
\end{lemma}

\begin{proof}
L'énoncé algébrique correspondant
est Algèbre, Lemme \ref{algebra-lemma-dimension-formula}.
\end{proof}

\begin{lemma}
\label{lemma-dimension-formula-general}
Soit $S$ un schéma. Soit $f : X \to S$ un morphisme de schémas. Soit $x \in X$,
et posons $s = f(x)$. Supposons que $S$ soit localement noethérien et que $f$
soit localement de type fini. Nous
avons
\begin{equation}
\label{equation-dimension-formula-general}
\dim(\mathcal{O}_{X, x})
\leq
\dim(\mathcal{O}_{S, s}) + E - \text{trdeg}_{\kappa(s)} \kappa(x).
\end{equation}
où $E$ est le maximum de $\text{trdeg}_{\kappa(f(\xi))}(\kappa(\xi))$ lorsque $\xi$ parcourt les
points génériques des composantes irréductibles de $X$ contenant
$x$.
\end{lemma}

\begin{proof}
Soient $X_1, \ldots, X_n$ les composantes irréductibles de $X$
contenant $x$ munies de leur
structure de schéma induite réduite. Celles-ci
correspondent aux idéaux premiers minimaux
$\mathfrak q_i$ de $\mathcal{O}_{X, x}$ et il y en
a donc un nombre fini (Schémas, Lemme \ref{schemes-lemma-specialize-points} et
Algèbre,
Lemme \ref{algebra-lemma-Noetherian-irreducible-components}). Alors $\dim(\mathcal{O}_{X, x}) = \max \dim(\mathcal{O}_{X, x}/\mathfrak q_i)
= \max \dim(\mathcal{O}_{X_i, x})$. Les $\xi$
apparaissant dans la définition de $E$ sont exactement les
points génériques $\xi_i \in X_i$. Soit $Z_i = \overline{\{f(\xi_i)\}} \subset S$ muni de la
structure de schéma induite réduite.
La composition $X_i \to X \to S$ se factorise
par $Z_i$ (Schémas, Lemme \ref{schemes-lemma-map-into-reduction}).
Ainsi nous pouvons appliquer
la formule de dimension (Lemme \ref{lemma-dimension-formula})
pour voir que $\dim(\mathcal{O}_{X_i, x}) \leq \dim(\mathcal{O}_{Z_i, x}) +
\text{trdeg}_{\kappa(f(\xi))}(\kappa(\xi)) -
\text{trdeg}_{\kappa(s)} \kappa(x)$.
En combinant le tout, nous obtenons le lemme.
\end{proof}

\noindent
Une application est la construction d'une fonction de dimension
sur tout schéma de type fini sur un schéma universellement
caténaire doté d'une fonction de dimension.
Pour la définition des fonctions
de dimension, voir Topologie, Définition \ref{topology-definition-dimension-function}.

\begin{lemma}
\label{lemma-dimension-function-propagates}
Soit $S$ un schéma localement noethérien et universellement caténaire.
Soit $\delta : S \to \mathbf{Z}$ une fonction de dimension. Soit $f : X \to S$
un morphisme de schémas. Supposons $f$
localement de type fini. Alors
l'application
\begin{align*}
\delta = \delta_{X/S} : X & \longrightarrow \mathbf{Z} \\
x & \longmapsto \delta(f(x)) + \text{trdeg}_{\kappa(f(x))} \kappa(x)
\end{align*}
est une fonction de dimension sur $X$.
\end{lemma}

\begin{proof}
Soit $f : X \to S$ localement de type fini.
Soient $x \leadsto y$, $x \not = y$ une spécialisation dans
$X$. Nous devons montrer que $\delta_{X/S}(x) > \delta_{X/S}(y)$ et
que $\delta_{X/S}(x) = \delta_{X/S}(y) + 1$ si $y$ est une spécialisation
immédiate de $x$.

\medskip\noindent Choisissez
un ouvert affine $V \subset S$ contenant l'image de $y$
et choisissons un ouvert affine $U \subset X$ dont l'image est contenue dans $V$
et contenant $y$. Nous pouvons clairement remplacer
$X$ par $U$ et $S$ par $V$.
Ainsi nous pouvons supposer que $X = \Spec(A)$ et $S = \Spec(R)$
et que $f$ est donné par un homomorphisme d'anneaux
$R \to A$. L'anneau $R$ est universellement
caténaire (Lemme \ref{lemma-universally-catenary-local}) et l'application
$R \to A$ est de type fini (Lemme \ref{lemma-locally-finite-type-characterize}).

\medskip\noindent
Soit $\mathfrak q \subset A$ l'idéal premier correspondant au point $x$
et soit $\mathfrak p \subset R$ l'idéal premier correspondant à $f(x)$.
La restriction $\delta'$ de $\delta$ à $S' = \Spec(R/\mathfrak p) \subset S$ est une fonction
de dimension. L'anneau $R/\mathfrak p$ est universellement
caténaire. La restriction de $\delta_{X/S}$ à $X' = \Spec(A/\mathfrak q)$ est
clairement égale à la fonction $\delta_{X'/S'}$ construite à l'aide de
la fonction de dimension $\delta'$. Par conséquent, on peut
supposer, en plus de ce qui précède, que $R \subset A$ sont des
domaines, en d'autres termes que $X$ et $S$ sont
des schémas intègres, et que $x$ est le point générique
de $X$ et $f(x)$ est le point générique de $S$.

\medskip\noindent
Notez que $\mathcal{O}_{X, x} = R(X)$ et que puisque $x \leadsto y$,
$x \not = y$, le spectre de $\mathcal{O}_{X, y}$ a au
moins
deux points (Schémas, Lemme \ref{schemes-lemma-specialize-points})
donc $\dim(\mathcal{O}_{X, y}) > 0$. Si $y$
est une spécialisation immédiate
de $x$, alors $\Spec(\mathcal{O}_{X, y}) = \{x, y\}$
et $\dim(\mathcal{O}_{X, y}) = 1$.

\medskip\noindent Écrivez
$s = f(x)$ et $t = f(y)$. Nous calculons
\begin{align*}
\delta_{X/S}(x) - \delta_{X/S}(y)
& =
\delta(s) + \text{trdeg}_{\kappa(s)} \kappa(x)
- \delta(t) - \text{trdeg}_{\kappa(t)} \kappa(y) \\
& =
\delta(s) - \delta(t) +
\text{trdeg}_{R(S)} R(X) - \text{trdeg}_{\kappa(t)} \kappa(y) \\
& =
\delta(s) - \delta(t) + \dim(\mathcal{O}_{X, y})
- \dim(\mathcal{O}_{S, t})
\end{align*}
où nous utilisons l'égalité dans (\ref{equation-dimension-formula}) à
la dernière étape. Puisque $\delta$ est une fonction
de dimension sur le schéma $S$
et $s \in S$ est le point générique, la différence $\delta(s) - \delta(t)$
est égale à $\text{codim}(\overline{\{t\}}, S)$ par Topologie, Lemme
\ref{topology-lemma-dimension-function-catenary}. C'est égal à $\dim(\mathcal{O}_{S, t})$
par Propriétés,
Lemme \ref{properties-lemma-codimension-local-ring}. Ainsi Nous
concluons que
$$
\delta_{X/S}(x) - \delta_{X/S}(y) = \dim(\mathcal{O}_{X, y})
$$
et le lemme découle de ce que nous avons dit ci-dessus à propos de $\dim(\mathcal{O}_{X, y})$.
\end{proof}

\noindent
Une autre application de la formule de la dimension est que la dimension ne change pas sous
``altérations'' (à définir plus tard).

\begin{lemma}
\label{lemma-alteration-dimension}
Soit $f : X \to Y$ un morphisme de schémas. Supposons que
\begin{enumerate}
\item $Y$ soit localement noethérien,
\item $X$ et $Y$ soient des schémas intègres,
\item $f$ soit dominant, et
\item $f$ soit localement de type fini.
\end{enumerate}
Alors nous avons
$$
\dim(X) \leq \dim(Y) + \text{trdeg}_{R(Y)} R(X).
$$
Si $f$ est fermé\footnote{Par exemple si $f$ est propre,
voir Définition \ref{definition-proper}.} alors l'égalité est vérifiée.
\end{lemma}

\begin{proof}
Soit $f : X \to Y$ comme dans le
lemme. Soit $\xi_0 \leadsto \xi_1 \leadsto \ldots \leadsto \xi_e$ une suite de spécialisations
dans $X$. Posons $x = \xi_e$ et
$y = f(x)$. Remarquons que $e \leq \dim(\mathcal{O}_{X, x})$ tel que les
spécialisations données apparaissent dans
le spectre de $\mathcal{O}_{X, x}$, voir Schémas, Lemme \ref{schemes-lemma-specialize-points}. Par la
formule de dimension, Lemme \ref{lemma-dimension-formula}, nous voyons
que
\begin{align*}
e & \leq \dim(\mathcal{O}_{X, x}) \\
& \leq \dim(\mathcal{O}_{Y, y}) +
\text{trdeg}_{R(Y)} R(X) - \text{trdeg}_{\kappa(y)} \kappa(x) \\
& \leq \dim(\mathcal{O}_{Y, y}) + \text{trdeg}_{R(Y)} R(X)
\end{align*}
Par conséquent, nous concluons que $e \leq \dim(Y) + \text{trdeg}_{R(Y)} R(X)$ comme souhaité.

\medskip\noindent
Ensuite, supposons que $f$
soit également fermé. Supposons que $\overline{\xi}_0 \leadsto \overline{\xi}_1 \leadsto \ldots
\leadsto \overline{\xi}_d$ soit
une suite de spécialisations dans $Y$. Nous voulons montrer que $\dim(X) \geq d + r$. Nous pouvons
supposer que $\overline{\xi}_0 = \eta$ est le point générique de $Y$.
La fibre générique $X_\eta$ est un schéma localement de
type fini sur $\kappa(\eta) = R(Y)$. Elle est non vide puisque $f$
est dominant. Ainsi, d'après le Lemme \ref{lemma-ubiquity-Jacobson-schemes}, c'est
un Schéma de Jacobson. Donc, d'après le Lemme
\ref{lemma-jacobson-finite-type-points}, nous pouvons trouver un point fermé $\xi_0 \in X_\eta$
et l'extension $\kappa(\eta) \subset \kappa(\xi_0)$ est
une extension finie. Notez
que $\mathcal{O}_{X, \xi_0} = \mathcal{O}_{X_\eta, \xi_0}$ parce que $\eta$ est le point générique
de $Y$. Ainsi, nous voyons que
$\dim(\mathcal{O}_{X, \xi_0}) = r$ par le Lemme \ref{lemma-dimension-formula} appliqué au schéma $X_\eta$
sur le schéma universellement caténaire $\Spec(\kappa(\eta))$
(voir Lemme \ref{lemma-ubiquity-uc}) et le point $\xi_0$. Cela
signifie que nous pouvons trouver $\xi_{-r} \leadsto \ldots \leadsto \xi_{-1} \leadsto \xi_0$
dans $X$. D'autre part, comme $f$
est des spécialisations fermées, on peut relever
le long de $f$,
voir Topologie, Lemme \ref{topology-lemma-closed-open-map-specialization}. Ainsi, comme $\xi_0$
se situe
au-dessus de $\eta = \overline{\xi}_0$, nous pouvons trouver des spécialisations
$\xi_0 \leadsto \xi_1 \leadsto \ldots \leadsto \xi_d$ au-dessus de $\overline{\xi}_0 \leadsto \overline{\xi}_1 \leadsto \ldots
\leadsto \overline{\xi}_d$. En
d'autres termes, nous avons
$$
\xi_{-r} \leadsto \ldots \leadsto \xi_{-1} \leadsto \xi_0
\leadsto \xi_1 \leadsto \ldots \leadsto \xi_d
$$
ce qui signifie que $\dim(X) \geq d + r$ comme souhaité.
\end{proof}

\begin{lemma}
\label{lemma-alteration-dimension-general}
Soit $f : X \to Y$ un morphisme de schémas. Supposons que $Y$ soit
localement noethérien et que $f$ soit localement de type fini.
Alors
$$
\dim(X) \leq \dim(Y) + E
$$
où $E$ est le supremum de $\text{trdeg}_{\kappa(f(\xi))}(\kappa(\xi))$ lorsque $\xi$ parcourt
les points génériques des composantes irréductibles de
$X$.
\end{lemma}

\begin{proof}
Conséquence immédiate du Lemme \ref{lemma-dimension-formula-general} et
des Propriétés, Lemme \ref{properties-lemma-dimension}.
\end{proof}








\section{Normalisation relative}
\label{section-normalization-X-in-Y}

\noindent
Dans cette section, nous construisons la normalisation {\it d'un schéma dans un autre }.

\begin{lemma}
\label{lemma-integral-closure}
Soit $X$ un schéma. Soit $\mathcal{A}$ un faisceau quasi-cohérent
de $\mathcal{O}_X$-algèbres. Le sous-faisceau $\mathcal{A}' \subset \mathcal{A}$ défini
par la règle
$$
U \longmapsto \{f \in \mathcal{A}(U) \mid
f_x \in \mathcal{A}_x \text{ entier sur } \mathcal{O}_{X, x}
\text{ pour tout }x \in U\}
$$
est une $\mathcal{O}_X$-algèbre quasi-cohérente, le germe $\mathcal{A}'_x$
est la clôture intégrale de $\mathcal{O}_{X, x}$ dans $\mathcal{A}_x$,
et pour tout ouvert affine $U \subset X$
l'anneau $\mathcal{A}'(U) \subset \mathcal{A}(U)$ est la clôture
intégrale de $\mathcal{O}_X(U)$ dans $\mathcal{A}(U)$.
\end{lemma}

\begin{proof}
Il s'agit d'un sous-faisceau par la nature locale des conditions.
C'est une $\mathcal{O}_X$-algèbre
d'après Algèbre, Lemme \ref{algebra-lemma-integral-closure-is-ring}. Soit $U \subset X$ un
ouvert affine. Soit $U = \Spec(R)$ et supposons que $\mathcal{A}$
soit le faisceau quasi-cohérent associé à la $R$-algèbre
$A$. Alors, d'après Algèbre,
Lemme \ref{algebra-lemma-integral-closure-stalks}, la valeur de $\mathcal{A}'$ sur $U$
est donnée par la clôture intégrale $A'$ de
$R$ dans $A$. Cela prouve la dernière assertion
du lemme. Pour prouver que $\mathcal{A}'$ est quasi-cohérent,
il suffit de montrer que $\mathcal{A}'(D(f)) = A'_f$. Cela résulte du fait
que la clôture intégrale et la localisation commutent, voir
Algèbre, Lemme \ref{algebra-lemma-integral-closure-localize}. Le même fait montre
que les fibres sont comme annoncé.
\end{proof}

\begin{definition}
\label{definition-integral-closure}
Soit $X$ un schéma. Soit $\mathcal{A}$ un faisceau quasi-cohérent
d'$\mathcal{O}_X$-algèbres. La clôture {\it intégrale
de $\mathcal{O}_X$ dans $\mathcal{A}$} est la $\mathcal{O}_X$-sous-algèbre
quasi-cohérente $\mathcal{A}' \subset \mathcal{A}$ construite
dans le Lemme \ref{lemma-integral-closure} ci-dessus.
\end{definition}

\noindent
Dans le cadre de la définition ci-dessus, nous pouvons considérer le morphisme
des spectres relatifs
$$
\xymatrix{
Y = \underline{\Spec}_X(\mathcal{A}) \ar[rr] \ar[rd] & &
X' = \underline{\Spec}_X(\mathcal{A}') \ar[ld] \\
& X &
}
$$
voir Lemme \ref{lemma-affine-equivalence-algebras}. Le schéma $X' \to X$ sera
la normalisation de $X$ dans le schéma $Y$. Voici un
contexte légèrement plus général. Supposons que nous
ayons un quasi-compact et un quasi-morphisme séparé
$f : Y \to X$ de schémas. Dans ce
cas, le faisceau de $\mathcal{O}_X$-algèbres $f_*\mathcal{O}_Y$ est quasi-cohérent,
voir Schémas, Lemme \ref{schemes-lemma-push-forward-quasi-coherent}. En prenant la clôture
intégrale $\mathcal{O}' \subset f_*\mathcal{O}_Y$ nous obtenons un faisceau quasi-cohérent
de $\mathcal{O}_X$-algèbres dont le spectre relatif est
la normalisation de $X$ dans $Y$. Voici la définition
formelle.

\begin{definition}
\label{definition-normalization-X-in-Y}
Soit $f : Y \to X$ un morphisme de schémas quasi-compact et quasi-séparé. Soit $\mathcal{O}'$ la clôture
intégrale de $\mathcal{O}_X$ dans $f_*\mathcal{O}_Y$. La normalisation {\it
de $X$ dans $Y$} est le schéma\footnote{
Le schéma $X'$ n'a pas besoin d'être normal, par exemple si $Y = X$
et $f = \text{id}_X$, alors $X' = X$.}
$$
\nu : X' = \underline{\Spec}_X(\mathcal{O}') \to X
$$
au-dessus de $X$. Il est muni d'une factorisation naturelle
$$
Y \xrightarrow{f'} X' \xrightarrow{\nu} X
$$
du morphisme initial $f$.
\end{definition}

\noindent
La factorisation est la composition du morphisme
canonique $Y \to \underline{\Spec}_X(f_*\mathcal{O}_Y)$ (voir Constructions,
Lemme \ref{constructions-lemma-canonical-morphism})
et du morphisme des spectres
relatifs provenant de l'application d'inclusion $\mathcal{O}' \to f_*\mathcal{O}_Y$.
Nous pouvons caractériser la normalisation
comme suit.

\begin{lemma}
\label{lemma-characterize-normalization}
Soit $f : Y \to X$ un morphisme de schémas quasi-compact et quasi-séparé. La factorisation
$f = \nu \circ f'$, où $\nu : X' \to X$ est la normalisation de $X$
dans $Y$, est caractérisée par les deux propriétés
suivantes :
\begin{enumerate}
\item le morphisme $\nu$ est entier, et
\item pour toute factorisation $f = \pi \circ g$, avec $\pi : Z \to X$ entier,
il existe un diagramme commutatif
$$
\xymatrix{
Y \ar[d]_{f'} \ar[r]_g & Z \ar[d]^\pi \\
X' \ar[ru]^h \ar[r]^\nu & X
}
$$
pour un morphisme unique $h : X' \to Z$.
\end{enumerate}
De plus, le morphisme $f' : Y \to X'$ est dominant et dans (2) le morphisme
$h : X' \to Z$ est la normalisation de $Z$ dans $Y$.
\end{lemma}

\begin{proof}
Soit $\mathcal{O}' \subset f_*\mathcal{O}_Y$ la clôture intégrale de $\mathcal{O}_X$ comme dans
la définition \ref{definition-normalization-X-in-Y}. Le morphisme $\nu$ est entier
par construction, ce qui prouve (1). Supposons donnée
une factorisation $f = \pi \circ g$ avec $\pi : Z \to X$ entier comme
dans (2). Par la Définition \ref{definition-integral} $\pi$
est affine, et donc $Z$ est le
spectre relatif d'un faisceau quasi-cohérent d'algèbres $\mathcal{O}_X$
$\mathcal{B}$. Le morphisme $g : Y \to Z$ correspond à un homomorphisme
d'$\mathcal{O}_X$-algèbres $\chi : \mathcal{B} \to f_*\mathcal{O}_Y$. Comme l'algèbre $\mathcal{B}(U)$ est entière
sur $\mathcal{O}_X(U)$ pour chaque ouvert affine $U \subset X$
(par la Définition \ref{definition-integral}), nous
voyons à partir du Lemme \ref{lemma-integral-closure}
cela $\chi(\mathcal{B}) \subset \mathcal{O}'$. Par la fonctorialité
du spectre relatif Lemme \ref{lemma-affine-equivalence-algebras},
cela nous fournit un morphisme
unique $h : X' \to Z$. Nous omettons
la vérification que le diagramme commute.

\medskip\noindent
Il est clair que (1) et (2) caractérisent la factorisation $f = \nu \circ f'$ puisqu'elle
la caractérise comme un objet initial dans une catégorie.

\medskip\noindent D'après
la propriété universelle en (2), nous voyons que $f'$ ne se factorise pas par
un sous-schéma fermé propre de $X'$. Ainsi, l'image schématique de $f'$
est $X'$. Puisque $f'$ est quasi-compact (d'après
Schémas, Lemme \ref{schemes-lemma-quasi-compact-permanence} et le fait que $\nu$ est séparé en tant que morphisme
affine), nous voyons que $f'(Y)$ est dense dans $X'$. Par conséquent,
$f'$ est dominant.

\medskip\noindent
Observez que $g$ est quasi-compact et quasi-séparé
d'après Schémas, Lemme \ref{schemes-lemma-compose-after-separated} et \ref{schemes-lemma-quasi-compact-permanence}.
Ainsi, la dernière affirmation du lemme a du sens. Le morphisme
$h$ dans
(2) est entier d'après le Lemme \ref{lemma-finite-permanence}. Étant donné une factorisation
$g = \pi' \circ g'$ avec $\pi' : Z' \to Z$ entier, nous obtenons
une factorisation $f = (\pi \circ \pi') \circ g'$ et nous obtenons un morphisme
$h' : X' \to Z'$. L'unicité implique que $\pi' \circ h' = h$. Par
conséquent, la caractérisation (1), (2) s'applique au
morphisme $h : X' \to Z$ ce qui donne l'affirmation finale du lemme.
\end{proof}

\begin{lemma}
\label{lemma-functoriality-normalization}
Soit
$$
\xymatrix{
Y_2 \ar[d]_{f_2} \ar[r] & Y_1 \ar[d]^{f_1} \\
X_2 \ar[r] & X_1
}
$$
soit un diagramme commutatif de morphismes de schémas.
Supposons $f_1$, $f_2$ quasi-compacts et
quasi-séparés. Soient $f_i = \nu_i \circ f_i'$, $i = 1, 2$
les factorisations canoniques, où $\nu_i : X_i' \to X_i$ est la normalisation
de $X_i$ dans $Y_i$. Alors il existe un unique flèche
$X'_2 \to X'_1$ s'insérant dans un diagramme
commutatif
$$
\xymatrix{
Y_2 \ar[d]_{f_2'} \ar[r] & Y_1 \ar[d]^{f_1'} \\
X_2' \ar[d]_{\nu_2} \ar[r] & X_1' \ar[d]^{\nu_1} \\
X_2 \ar[r] & X_1
}
$$
\end{lemma}

\begin{proof}
D'après les Lemme \ref{lemma-characterize-normalization} (1) et \ref{lemma-base-change-finite},
le changement de base
$X_2 \times_{X_1} X'_1 \to X_2$ est entier. Notons que $f_2$
se factorise par ce morphisme. Ainsi on
obtient un morphisme unique
$X'_2 \to X_2 \times_{X_1} X'_1$ d'après le Lemme
\ref{lemma-characterize-normalization} (2). Cela donne la flèche
$X'_2 \to X'_1$ s'insérant dans le diagramme
commutatif et l'unicité s'ensuit également.
\end{proof}

\begin{lemma}
\label{lemma-normalization-localization}
Soit $f : Y \to X$ un morphisme quasi-compact et quasi-séparé de schémas. Soit $U \subset X$ un
sous-schéma ouvert et posons $V = f^{-1}(U)$. Alors la normalisation de $U$
dans $V$ est l'image réciproque de $U$ dans la normalisation
de $X$ dans $Y$.
\end{lemma}

\begin{proof}
Clair d'après la construction.
\end{proof}

\begin{lemma}
\label{lemma-normalization-is-normalization}
Soit $f : Y \to X$ un morphisme quasi-compact et quasi-séparé de schémas. Soit $X'$ la normalisation
de $X$ dans $Y$. Alors la normalisation de $X'$ dans $Y$
est $X'$.
\end{lemma}

\begin{proof}
Si $Y \to X'' \to X'$ est la normalisation de $X'$ dans $Y$, alors
nous pouvons appliquer le Lemme \ref{lemma-characterize-normalization} à la composition
$X'' \to X$ pour obtenir un morphisme canonique $h : X' \to X''$ sur
$X$. Nous omettons la vérification que les morphismes $h$
et $X'' \to X'$ sont mutuellement inverses (en utilisant l'unicité
de la factorisation dans le lemme).
\end{proof}

\begin{lemma}
\label{lemma-normalization-in-reduced}
Soit $f : Y \to X$ un morphisme de schémas quasi-compact et quasi-séparé. Soit $X' \to X$ la normalisation
de $X$ dans $Y$. Si $Y$ est réduit, alors $X'$ l'est
également.
\end{lemma}

\begin{proof}
Cela résulte du fait qu'un sous-anneau d'un anneau réduit est réduit. Quelques
détails sont omis.
\end{proof}

\begin{lemma}
\label{lemma-normalization-generic}
Soit $f : Y \to X$ un morphisme de schémas quasi-compact et quasi-séparé. Soit $X' \to X$
la normalisation de $X$ dans $Y$. Chaque point générique d'une composante
irréductible de $X'$ est l'image d'un point générique d'une composante
irréductible de $Y$.
\end{lemma}

\begin{proof}
D'après le Lemme \ref{lemma-normalization-localization}, on peut supposer que $X = \Spec(A)$ est
affine. Choisissons un recouvrement fini par des ouverts affines $Y = \bigcup \Spec(B_i)$.
Ensuite $X' = \Spec(A')$ et les morphismes $\Spec(B_i) \to Y \to X'$ définissent
conjointement une application $A$-algèbre
injective $A' \to \prod B_i$.
Ainsi, le lemme découle de l'Algèbre, Lemme \ref{algebra-lemma-injective-minimal-primes-in-image}.
\end{proof}

\begin{lemma}
\label{lemma-normalization-in-disjoint-union}
Soit $f : Y \to X$ un morphisme de schémas quasi-compact et quasi-séparé. Supposons que
$Y = Y_1 \amalg Y_2$ soit une réunion disjointe de deux schémas. Écrivons $f_i = f|_{Y_i}$.
Soit $X_i'$ la normalisation de $X$ dans $Y_i$. Alors $X_1' \amalg X_2'$
est la normalisation de $X$ dans $Y$.
\end{lemma}

\begin{proof}
En termes de clôtures intégrales, cela correspond au fait suivant : Soit $A \to B$ un
homomorphisme d'anneaux. Supposons que $B = B_1 \times B_2$. Soit $A_i'$ la clôture
intégrale de $A$ dans $B_i$. Alors $A_1' \times A_2'$ est la clôture
intégrale de $A$ dans $B$. La raison pour laquelle
cela fonctionne est que les éléments $(1, 0)$ et $(0, 1)$ de $B$ sont idempotents
et donc entiers sur $A$. Ainsi, la clôture intégrale $A'$ de $A$
dans $B$ est un produit et il n'est pas difficile de voir que les se factorise
par sont les clôtures intégrales $A'_i$ comme décrit ci-dessus (certains détails
omis).
\end{proof}

\begin{lemma}
\label{lemma-normalization-in-universally-closed}
Soit $f : X \to S$ un morphisme de schémas quasi-compact, quasi-séparé
et universellement fermé. Alors
$f_*\mathcal{O}_X$ est entier sur $\mathcal{O}_S$. En d'autres termes, la normalisation
de $S$ dans $X$ est égale à la factorisation
$$
X \longrightarrow \underline{\Spec}_S(f_*\mathcal{O}_X)
\longrightarrow S
$$
de Constructions, Lemme \ref{constructions-lemma-canonical-morphism}.
\end{lemma}

\begin{proof}
La question est locale sur $S$, donc nous pouvons supposer que $S = \Spec(R)$
est affine. Soit $h \in \Gamma(X, \mathcal{O}_X)$. Nous devons montrer que $h$ satisfait
une équation unitaire sur $R$. Considérons $h$ comme un morphisme
comme dans le diagramme commutatif suivant
$$
\xymatrix{
X \ar[rr]_h \ar[rd]_f & & \mathbf{A}^1_S \ar[ld] \\
& S &
}
$$
Soit $Z \subset \mathbf{A}^1_S$ l'image schématique de $h$, voir
Définition \ref{definition-scheme-theoretic-image}. Le morphisme $h$
est quasi-compact car $f$ est quasi-compact
et $\mathbf{A}^1_S \to S$ est séparé, voir
Schémas, Lemme \ref{schemes-lemma-quasi-compact-permanence}. D'après le Lemme \ref{lemma-quasi-compact-scheme-theoretic-image},
le morphisme $X \to Z$ est dominant. D'après
le Lemme \ref{lemma-image-proper-scheme-closed}, le
morphisme $X \to Z$ est fermé. Donc $h(X) = Z$
(au sens des ensembles). Ainsi nous pouvons utiliser
le Lemme
\ref{lemma-image-proper-is-proper} pour conclure que $Z \to S$
est universellement fermé (et même propre).
Puisque $Z \subset \mathbf{A}^1_S$, nous voyons que $Z \to S$ est affine
et propre, donc entier d'après le Lemme \ref{lemma-integral-universally-closed}. En écrivant $\mathbf{A}^1_S = \Spec(R[T])$
Nous concluons que l'idéal $I \subset R[T]$ de
$Z$ contient un polynôme unitaire $P(T) \in R[T]$.
Donc $P(h) = 0$ et nous avons terminé.
\end{proof}

\begin{lemma}
\label{lemma-normalization-in-integral}
Soit $f : Y \to X$ un morphisme entier. Alors la normalisation
de $X$ dans $Y$ est égale à $Y$.
\end{lemma}

\begin{proof}
D'après le Lemme \ref{lemma-integral-universally-closed} ceci est un cas particulier
du Lemme \ref{lemma-normalization-in-universally-closed}.
\end{proof}

\begin{lemma}
\label{lemma-normal-normalization}
Soit $f : Y \to X$ un morphisme de schémas quasi-compact et quasi-séparé. Soit
$X'$ la normalisation de $X$ dans $Y$. Supposons
\begin{enumerate}
\item $Y$ est un schéma normal,
\item les ouverts quasi-compacts de $Y$ ont un nombre fini de composantes irréductibles.
\end{enumerate}
Alors $X'$ est une réunion disjointe de schémas normaux intègres. De plus,
le morphisme $Y \to X'$ est dominant et induit une bijection entre les
composantes irréductibles.
\end{lemma}

\begin{proof}
Soit $U \subset X$ un ouvert affine. Considérons
l'image réciproque $U'$ de
$U$ dans $X'$. Posons $V = f^{-1}(U)$. D'après le Lemme
\ref{lemma-normalization-localization}, le diagramme $V \to U' \to U$ est la normalisation de
$U$ dans $V$. Disons $U = \Spec(A)$. Alors $V$ est
quasi-compact, et possède donc un nombre
fini de composantes irréductibles par hypothèse. Ainsi $V = \coprod_{i = 1, \ldots n} V_i$
est une union
finie disjointe de schémas intègres normaux d'après les Propriétés, Lemme
\ref{properties-lemma-normal-locally-finite-nr-irreducibles}. D'après le Lemme \ref{lemma-normalization-in-disjoint-union}, nous
voyons que $U' = \coprod_{i = 1, \ldots, n} U_i'$, où $U'_i$ est la
normalisation de $U$ dans $V_i$.
Par les propriétés, Lemme \ref{properties-lemma-normal-integral-sections} nous voyons que
$B_i = \Gamma(V_i, \mathcal{O}_{V_i})$ est un domaine normal. Notez que $U_i' = \Spec(A_i')$,
où $A_i' \subset B_i$ est la clôture intégrale de
$A$ dans $B_i$, voir
Lemme \ref{lemma-integral-closure}. Par l'algèbre,
Lemme \ref{algebra-lemma-integral-closure-in-normal} nous voyons que $A_i' \subset B_i$
est un domaine normal. Ainsi $U' = \coprod U_i'$
est une union finie de schémas intègres normaux et est
donc normal.

\medskip\noindent
Comme $X'$ a un recouvrement ouvert par les schémas $U'$,
nous concluons d'après Propriétés, Lemme \ref{properties-lemma-locally-normal} que $X'$ est normal. D'autre
part, chaque $U'$ est une réunion finie disjointe de schémas
irréductibles, donc chaque ouvert quasi-compact de $X'$ a un nombre fini de
composantes irréductibles (par un argument topologique que nous omettons).
Ainsi $X'$ est une réunion disjointe
de schémas intègres normaux d'après Propriétés, Lemme \ref{properties-lemma-normal-locally-finite-nr-irreducibles}. Il est clair d'après
la description de $X'$ ci-dessus que $Y \to X'$ est dominant
et induit une bijection sur les composantes irréductibles $V \to U'$
pour chaque ouvert affine $U \subset X$. La bijection des composantes
irréductibles pour le morphisme $Y \to X'$
résulte de cela par un argument topologique (omis).
\end{proof}

\begin{lemma}
\label{lemma-nagata-normalization-finite-general}
Soit $f : X \to S$ un morphisme. Supposons que
\begin{enumerate}
\item $S$ soit un schéma de Nagata,
\item $f$ soit quasi-compact et quasi-séparé,
\item les ouverts quasi-compacts de $X$ ont un nombre fini de composantes irréductibles,
\item si $x \in X$ est un point générique d'une composante irréductible,
alors l'extension de corps $\kappa(x)/\kappa(f(x))$ est de type
fini, et
\item $X$ est réduit.
\end{enumerate}
Alors la normalisation $\nu : S' \to S$ de $S$ dans $X$ est finie.
\end{lemma}

\begin{proof}
Il y a une réduction immédiate au cas $S = \Spec(R)$ où $R$
est un anneau de Nagata par l'hypothèse (1). Nous devons montrer
que la clôture intégrale $A$ de $R$ dans $\Gamma(X, \mathcal{O}_X)$
est finie sur $R$. Étant donné que $f$ est quasi-compact
par l'hypothèse (2), nous pouvons écrire $X = \bigcup_{i = 1, \ldots, n} U_i$ avec
chaque $U_i$ affine. Disons $U_i = \Spec(B_i)$. Chaque $B_i$
est réduit par l'hypothèse (5) et
possède un nombre fini de premiers minimaux
$\mathfrak q_{i1}, \ldots, \mathfrak q_{im_i}$ par l'hypothèse
(3) et l'Algèbre, Lemme \ref{algebra-lemma-irreducible}. Nous
avons
$$
\Gamma(X, \mathcal{O}_X) \subset B_1 \times \ldots \times B_n
\subset
\prod\nolimits_{i = 1, \ldots, n}
\prod\nolimits_{j = 1, \ldots, m_i} (B_i)_{\mathfrak q_{ij}}
$$
la deuxième inclusion
par l'algèbre, Lemme \ref{algebra-lemma-reduced-ring-sub-product-fields}. Nous avons $\kappa(\mathfrak q_{ij}) = (B_i)_{\mathfrak q_{ij}}$
par l'algèbre, Lemme \ref{algebra-lemma-minimal-prime-reduced-ring}. Par conséquent,
la clôture intégrale $A$ de $R$ dans
$\Gamma(X, \mathcal{O}_X)$ est contenue dans le produit des clôtures
intégrales $A_{ij}$ de $R$ dans $\kappa(\mathfrak q_{ij})$.
Comme $R$ est noethérien, il suffit
de montrer que $A_{ij}$ est un $R$-module fini
pour chaque $i, j$. Soit $\mathfrak p_{ij} \subset R$ l'image de $\mathfrak q_{ij}$.
Comme $\kappa(\mathfrak q_{ij})/\kappa(\mathfrak p_{ij})$ est une extension
de corps de type fini par l'hypothèse (4), nous
voyons que $R \to \kappa(\mathfrak q_{ij})$ est essentiellement de type fini.
Ainsi $R \to A_{ij}$ est fini
par l'algèbre, Lemme \ref{algebra-lemma-nagata-in-reduced-finite-type-finite-integral-closure}.
\end{proof}

\begin{lemma}
\label{lemma-nagata-normalization-finite}
Soit $f : X \to S$ un morphisme. Supposons que
\begin{enumerate}
\item $S$ est un schéma de Nagata,
\item $f$ est de type fini,
\item $X$ est réduit.
\end{enumerate}
Alors la normalisation $\nu : S' \to S$ de $S$ dans $X$ est finie.
\end{lemma}

\begin{proof}
Ceci est un cas particulier
du Lemme \ref{lemma-nagata-normalization-finite-general}. À savoir, (2) est vérifié
car le morphisme de type fini $f$ est quasi-compact par
définition et quasi-séparé par
le Lemme \ref{lemma-finite-type-Noetherian-quasi-separated}. La condition (3) est vérifiée parce
que $X$ est localement noethérien par le
Lemme \ref{lemma-finite-type-noetherian}. Enfin, la condition (4) est vérifiée car
un morphisme de type fini induit des extensions de corps résiduels
de type fini.
\end{proof}

\begin{lemma}
\label{lemma-relative-normalization-normal-codim-1}
Soit $f : Y \to X$ un morphisme de schémas de type fini avec $Y$ réduit et $X$
Nagata. Soit $X'$ la normalisation de $X$ dans $Y$.
Soit $x' \in X'$ un point tel que
\begin{enumerate}
\item $\dim(\mathcal{O}_{X', x'}) = 1$, et
\item la fibre de $Y \to X'$ au-dessus de $x'$ est vide.
\end{enumerate}
Alors $\mathcal{O}_{X', x'}$ est un anneau de valuation discret.
\end{lemma}

\begin{proof}
Nous pouvons remplacer $X$ par un voisinage affine de l'image de
$x'$. Ainsi, nous pouvons supposer $X = \Spec(A)$
avec $A$ Nagata. Par le Lemme \ref{lemma-nagata-normalization-finite} le
morphisme $X' \to X$ est fini. Ainsi, nous pouvons écrire
$X' = \Spec(A')$ pour une $A$-algèbre finie $A'$.
Par le Lemme \ref{lemma-normalization-is-normalization} après avoir remplacé
$X$ par $X'$
nous revenons au cas décrit dans le paragraphe suivant.

\medskip\noindent
Le cas $X = X' = \Spec(A)$ avec $A$ noethérien. Soit $\mathfrak p \subset A$
l'idéal premier correspondant à notre point
$x'$. Choisissez $g \in \mathfrak p$ qui n'est contenu dans
aucun idéal premier minimal de $A$ (utilisez l'évitement des
idéaux premiers et le fait que $A$ possède un nombre
fini d'idéaux premiers minimaux, voir Algèbre, Lemmes
\ref{algebra-lemma-silly} et \ref{algebra-lemma-Noetherian-irreducible-components}). Posons $Z = f^{-1}V(g) \subset Y$ ; c'est
un sous-schéma fermé de $Y$. Alors $f(Z)$ ne contient aucun
point générique par choix de $g$ et ne contient pas $x'$
car $x'$ n'est pas dans l'image de $f$. L'adhérence
de $f(Z)$ est l'ensemble des spécialisations des points
de $f(Z)$ par le lemme \ref{lemma-reach-points-scheme-theoretic-image}. Ainsi, l'adhérence de $f(Z)$
ne contient pas $x'$ car la condition $\dim(\mathcal{O}_{X', x'}) = 1$ implique
seulement que les points génériques de $X = X'$ se spécialisent
en $x'$. En d'autres termes, après avoir remplacé
$X$ par un voisinage ouvert affine de $x'$,
nous pouvons supposer que $f^{-1}V(g) = \emptyset$. Ainsi $g$ s'envoie
sur une fonction globale inversible sur $Y$ et nous obtenons
une factorisation
$$
A \to A_g \to \Gamma(Y, \mathcal{O}_Y)
$$
Puisque $X = X'$, cela implique que $A$ est égal
à la clôture intégrale de
$A$ dans $A_g$. Par l'algèbre, Lemme \ref{algebra-lemma-integral-closure-localize}
Nous concluons que $A_\mathfrak p$ est la clôture intégrale
de $A_\mathfrak p$ dans $A_\mathfrak p[1/g]$. Par
notre choix de $g$, puisque $\dim(A_\mathfrak p) = 1$ et
puisque $A$ est réduit, nous voyons que $A_\mathfrak p[1/g]$
est un produit fini de corps (le produit des corps résiduels
des idéaux premiers minimaux contenus dans $\mathfrak p$). Par conséquent
$A_\mathfrak p$ est normal (Algèbre, Lemme \ref{algebra-lemma-characterize-reduced-ring-normal}) et la démonstration
est terminée. Certains détails sont omis.
\end{proof}






\section{Normalisation}
\label{section-normalization}

\noindent
Ensuite, nous arrivons à la normalisation d'un schéma $X$.
Nous ne le définissons/construisons que lorsque $X$ a localement un nombre fini de
composantes irréductibles. Soit $X$ un schéma tel que chaque ouvert quasi-compact
possède un nombre fini de composantes irréductibles.
Soit $X^{(0)} \subset X$ l'ensemble des points génériques des composantes irréductibles de $X$.
Soit
\begin{equation}
\label{equation-generic-points}
f :
Y = \coprod\nolimits_{\eta \in X^{(0)}} \Spec(\kappa(\eta))
\longrightarrow
X
\end{equation}
l'inclusion des points génériques dans $X$ à
l'aide des applications canoniques de Schémas, section \ref{schemes-section-points}.
Notez que ce morphisme est quasi-compact par
hypothèse et quasi-séparé car $Y$
est séparé (voir Schémas, section \ref{schemes-section-separation-axioms}).

\begin{definition}
\label{definition-normalization}
Soit $X$ un schéma tel que tout ouvert quasi-compact ait un nombre fini de
composantes irréductibles. Nous définissons la normalisation {\it
} de $X$ comme le morphisme
$$
\nu : X^\nu \longrightarrow X
$$
qui est la normalisation de $X$ dans le morphisme
$f : Y \to X$ (\ref{equation-generic-points}) construit ci-dessus.
\end{definition}

\noindent
Tout schéma localement noethérien possède un ensemble localement fini de composantes
irréductibles et la définition s'y applique. Habituellement,
la normalisation est définie uniquement pour les schémas réduits.
Avec la définition ci-dessus, la normalisation de $X$ est la même
que la normalisation de la réduction $X_{red}$ de $X$.

\begin{lemma}
\label{lemma-normalization-reduced}
Soit $X$ un schéma tel que tout ouvert quasi-compact ait
un nombre fini de composantes irréductibles. Le morphisme de normalisation
$\nu$ se factorise par la réduction $X_{red}$ et $X^\nu \to X_{red}$
est la normalisation de $X_{red}$.
\end{lemma}

\begin{proof}
Soit $f : Y \to X$ le morphisme (\ref{equation-generic-points}). Nous obtenons
une factorisation $Y \to X_{red} \to X$ de $f$ provenant
des Schémas, Lemme \ref{schemes-lemma-map-into-reduction}. D'après le
Lemme \ref{lemma-characterize-normalization}, nous obtenons un morphisme canonique $X^\nu \to X_{red}$
et que $X^\nu$
est la normalisation de $X_{red}$ dans $Y$.
Le lemme s'ensuit puisque $Y \to X_{red}$ est identique au
morphisme (\ref{equation-generic-points}) construit pour $X_{red}$.
\end{proof}

\noindent
Si $X$ est réduit, alors la normalisation de
$X$ est la même que le spectre relatif de l'adhérence
intégrale de $\mathcal{O}_X$ dans le faisceau
de fonctions méromorphes $\mathcal{K}_X$ (voir Diviseurs, section
\ref{divisors-section-meromorphic-functions}). À savoir, $\mathcal{K}_X = f_*\mathcal{O}_Y$ dans ce cas,
voir Diviseurs, Lemme \ref{divisors-lemma-reduced-finite-irreducible} et sa démonstration.
Nous décrivons cela ici explicitement.

\begin{lemma}
\label{lemma-description-normalization}
Soit $X$ un schéma réduit tel que chaque ouvert quasi-compact
possède un nombre fini de composantes irréductibles. Soit $\Spec(A) = U \subset X$
un ouvert affine. Alors
\begin{enumerate}
\item $A$ possède un nombre fini
de premiers minimaux $\mathfrak q_1, \ldots, \mathfrak q_t$,
\item l'anneau total des fractions
$Q(A)$ de $A$ est $Q(A/\mathfrak q_1) \times \ldots \times Q(A/\mathfrak q_t)$,
\item la clôture intégrale $A'$ de $A$ dans $Q(A)$
est le produit des clôtures intégrales des domaines $A/\mathfrak q_i$
dans le corps $Q(A/\mathfrak q_i)$, et
\item $\nu^{-1}(U)$ est identifié au spectre de $A'$ où $\nu : X^\nu \to X$
est le morphisme de normalisation.
\end{enumerate}
\end{lemma}

\begin{proof}
Les idéaux premiers minimaux correspondent
aux composantes irréductibles
(Algèbre, Lemme
\ref{algebra-lemma-irreducible}), donc nous avons (1) par hypothèse.
Ensuite $(0) = \mathfrak q_1 \cap \ldots \cap \mathfrak q_t$ parce que
$A$
est réduit (Algèbre, Lemme \ref{algebra-lemma-Zariski-topology}).
Ensuite nous avons $Q(A) = \prod A_{\mathfrak q_i} = \prod \kappa(\mathfrak q_i)$ par Algèbre,
Lemmes \ref{algebra-lemma-total-ring-fractions-no-embedded-points} et \ref{algebra-lemma-minimal-prime-reduced-ring}.
Cela prouve (2).
La partie (3) découle de Algèbre, Lemme \ref{algebra-lemma-characterize-reduced-ring-normal},
ou Lemme \ref{lemma-normalization-in-disjoint-union}.
La partie (4) est vraie parce qu'il est clair que $f^{-1}(U) \to U$ est le morphisme
$$
\Spec\left(\prod \kappa(\mathfrak q_i)\right)
\longrightarrow
\Spec(A)
$$
où $f : Y \to X$ est le morphisme (\ref{equation-generic-points}).
\end{proof}

\begin{lemma}
\label{lemma-stalk-normalization}
Soit $X$ un schéma tel que tout ouvert quasi-compact ait un nombre
fini de composantes irréductibles. Soit $\nu : X^\nu \to X$ la normalisation
de $X$. Soit $x \in X$. Alors les suivants sont
canoniquement isomorphes en tant que $\mathcal{O}_{X, x}$-algèbres
\begin{enumerate}
\item le germe $(\nu_*\mathcal{O}_{X^\nu})_x$,
\item la clôture intégrale de $\mathcal{O}_{X, x}$ dans
l'anneau total des fractions de $(\mathcal{O}_{X, x})_{red}$,
\item la clôture intégrale de $\mathcal{O}_{X, x}$ dans le produit
des corps résiduels des idéaux premiers minimaux de $\mathcal{O}_{X, x}$
(et il y en a un nombre fini).
\end{enumerate}
\end{lemma}

\begin{proof}
Après avoir remplacé $X$ par un ouvert affine
voisinage de $x$, nous pouvons supposer que $X$
a un nombre fini de composantes irréductibles
et que $x$ est contenu dans chacune d'elles.
Alors le germe $(\nu_*\mathcal{O}_{X^\nu})_x$ est l'adhérence entière de $A = \mathcal{O}_{X, x}$
dans le produit $L$ des corps résiduels
des idéaux premiers minimaux de $A$. Cela résulte
de la construction de la normalisation
et du Lemme \ref{lemma-integral-closure}.
Alternativement, vous pouvez utiliser
le Lemme \ref{lemma-description-normalization} et le fait que la normalisation
commute avec la localisation (Algèbre, Lemme \ref{algebra-lemma-integral-closure-localize}).
Puisque $A_{red}$ a un nombre fini de premiers minimaux
(parce qu'ils correspondent exactement aux points génériques
des composantes irréductibles de $X$ passant
par $x$), nous voyons que $L$ est l'anneau total
des fractions de $A_{red}$ (Algèbre, Lemme \ref{algebra-lemma-total-ring-fractions-no-embedded-points}). Ainsi, notre anneau est
également la clôture intégrale de $A$ dans l'anneau total
des fractions de $A_{red}$.
\end{proof}

\begin{lemma}
\label{lemma-normalization-normal}
Soit $X$ un schéma tel que tout ouvert quasi-compact ait un
nombre fini de composantes irréductibles.
\begin{enumerate}
\item La normalisation $X^\nu$ est une union disjointe de schémas intègres normaux.
\item Le morphisme $\nu : X^\nu \to X$ est entier, surjectif, et induit une
bijection sur les composantes irréductibles.
\item Pour tout morphisme entier $\alpha : X' \to X$ tel que pour $U \subset X$
ouvert quasi-compact l'image réciproque $\alpha^{-1}(U)$ ait un nombre fini
de composantes irréductibles et $\alpha|_{\alpha^{-1}(U)} : \alpha^{-1}(U) \to U$
soit birationnel\footnote{Cette formulation maladroite est
nécessaire car nous n'avons défini ce que signifie qu'un
morphisme soit birationnel que si la source et la cible ont
un nombre fini de composantes irréductibles. Il suffit
que $X'_{red} \to X_{red}$ satisfasse la condition.}
il existe une factorisation
$X^\nu \to X' \to X$ et $X^\nu \to X'$ est la normalisation de $X'$.
\item Pour tout morphisme $Z \to X$ avec $Z$ un schéma normal
tel que chaque composante irréductible de $Z$ domine une composante
irréductible de $X$ il existe une unique factorisation $Z \to X^\nu \to X$.
\end{enumerate}
\end{lemma}

\begin{proof}
Soit $f : Y \to X$ comme dans (\ref{equation-generic-points}). Le schéma $X^\nu$
est une réunion disjointe de schémas intègres normaux
car $Y$ est normal et tout ouvert affine de
$Y$ possède un nombre
fini de composantes irréductibles, voir le Lemme
\ref{lemma-normal-normalization}. Cela prouve (1). Alternativement,
on peut déduire (1) des
Lemmes \ref{lemma-normalization-reduced} et \ref{lemma-description-normalization}.

\medskip\noindent
Le morphisme $\nu$ est entier d'après le Lemme \ref{lemma-characterize-normalization}. D'après le Lemme
\ref{lemma-normal-normalization}, le morphisme $Y \to X^\nu$ induit
une bijection sur les composantes irréductibles, et par construction
de $Y$ cela implique que $X^\nu \to X$ induit une bijection
sur les composantes irréductibles. Par construction $f : Y \to X$
est dominant, donc également $\nu$ est dominant. Puisqu'un morphisme
entier est fermé (Lemme \ref{lemma-integral-universally-closed}), cela implique que $\nu$
est surjectif. Cela prouve (2).

\medskip\noindent
Supposons que $\alpha : X' \to X$ soit comme dans (3). Il est clair
que $X'$ satisfait les hypothèses sous lesquelles
la normalisation est définie. Soit $f' : Y' \to X'$
le morphisme (\ref{equation-generic-points}) construit à partir de $X'$.
Comme $\alpha$ est localement birationnel, il
est clair que $Y' = Y$ et $f = \alpha \circ f'$.
Par conséquent, la factorisation $X^\nu \to X' \to X$
existe et $X^\nu \to X'$ est la normalisation
de $X'$ d'après le Lemme \ref{lemma-characterize-normalization}. Cela prouve (3).

\medskip\noindent
Soit $g : Z \to X$ un morphisme dont le domaine est un schéma normal
et tel que chaque composante irréductible domine une composante irréductible
de $X$. D'après le Lemme \ref{lemma-normalization-reduced} nous avons
$X^\nu = X_{red}^\nu$ et d'après Schémas, Lemme
\ref{schemes-lemma-map-into-reduction} $Z \to X$ se factorise par $X_{red}$.
Par conséquent, nous pouvons remplacer $X$ par $X_{red}$
et supposer que $X$ est réduit. De plus, comme la factorisation
est unique, il suffit de la construire localement
sur $Z$. Soient $W \subset Z$ et $U \subset X$ des
ouverts affines tels que $g(W) \subset U$. Écrivons $U = \Spec(A)$
et $W = \Spec(B)$, avec $g|_W$ donné par $\varphi : A \to B$.
Nous utiliserons librement les résultats du Lemme \ref{lemma-description-normalization}. Soient $\mathfrak p_1, \ldots, \mathfrak p_t$ les
premiers minimaux de $A$. Comme $Z$ est normal, nous voyons
que $B$ est un anneau normal, en particulier
réduit. De plus, d'après l'hypothèse, pour tout premier
minimal $\mathfrak q \subset B$, nous avons que $\varphi^{-1}(\mathfrak q)$ est un premier minimal de
$A$. Ainsi, si $x \in A$ est un non-diviseur de zéro, c'est-à-dire $x \not \in \bigcup \mathfrak p_i$,
alors $\varphi(x)$ est un non-diviseur de zéro dans $B$. Ainsi, nous
obtenons un homomorphisme canonique d'anneaux $Q(A) \to Q(B)$. Comme $B$
est normal, il est égal à sa clôture intégrale dans
$Q(B)$ (voir Algèbre, Lemme \ref{algebra-lemma-normal-ring-integrally-closed}). Ainsi, nous voyons
que la clôture intégrale $A' \subset Q(A)$ de $A$ s'envoie dans
$B$ par l'homomorphisme canonique $Q(A) \to Q(B)$.
Puisque $\nu^{-1}(U) = \Spec(A')$, cela donne la factorisation canonique
$W \to \nu^{-1}(U) \to U$ de $\nu|_W$. Nous omettons
la vérification qu'elle est unique.
\end{proof}

\begin{lemma}
\label{lemma-normalization-in-terms-of-components}
Soit $X$ un schéma tel que chaque ouvert quasi-compact
ait un nombre fini de composantes irréductibles. Soient $Z_i \subset X$,
$i \in I$ les composantes irréductibles de $X$
munies de la structure réduite induite. Soit $Z_i^\nu \to Z_i$ la normalisation.
Alors $\coprod_{i \in I} Z_i^\nu \to X$ est la normalisation de $X$.
\end{lemma}

\begin{proof}
Nous pouvons supposer que $X$ est réduit, voir Lemme \ref{lemma-normalization-reduced}.
Ensuite, le lemme découle soit de la description locale
dans le lemme \ref{lemma-description-normalization} soit du Lemme
\ref{lemma-normalization-normal} partie (3) car $\coprod Z_i \to X$ est entier et
localement birationnel (puisque $X$ est réduit et a localement
un nombre fini de composantes irréductibles).
\end{proof}

\begin{lemma}
\label{lemma-normalization-birational}
Soit $X$ un schéma réduit avec un nombre fini de composantes irréductibles.
Alors le morphisme de normalisation $X^\nu \to X$ est birationnel.
\end{lemma}

\begin{proof}
La normalisation induit une bijection des composantes irréductibles
par le Lemme \ref{lemma-normalization-normal}. Soit $\eta \in X$ un point générique
d'une composante irréductible de $X$ et soit
$\eta^\nu \in X^\nu$ le point générique de la composante irréductible correspondante
de $X^\nu$. Alors $\eta^\nu \mapsto \eta$ et pour terminer la démonstration,
nous devons montrer que $\mathcal{O}_{X, \eta} \to \mathcal{O}_{X^\nu, \eta^\nu}$ est
un isomorphisme, voir la Définition \ref{definition-birational}. Comme
$X$ et $X^\nu$ sont réduits, nous voyons que les
deux anneaux locaux sont égaux
à leur corps résiduel (Algèbre, Lemme \ref{algebra-lemma-minimal-prime-reduced-ring}).
D'autre part, par la construction de la normalisation
comme la normalisation de $X$ dans $Y = \coprod \Spec(\kappa(\eta))$
nous voyons que
nous avons $\kappa(\eta) \subset \kappa(\eta^\nu) \subset \kappa(\eta)$ et la démonstration
est complète.
\end{proof}

\begin{lemma}
\label{lemma-finite-birational-over-normal}
Un morphisme fini (ou même entier) birationnel $f : X \to Y$ de schémas
intègres avec $Y$ normal est un isomorphisme.
\end{lemma}

\begin{proof}
Soit $V \subset Y$ un ouvert affine avec
image réciproque $U \subset X$ qui est également un ouvert affine.
Comme $f$ est un morphisme birationnel de schémas intègres, l'homomorphisme
$\mathcal{O}_Y(V) \to \mathcal{O}_X(U)$ est un homomorphisme injectif de domaines intègres qui
induit un isomorphisme de corps de fractions. Comme $Y$
est normal, l'anneau $\mathcal{O}_Y(V)$ est intégralement clos dans le corps des
fractions. Comme $f$ est fini (ou entier), chaque élément
de $\mathcal{O}_X(U)$ est entier sur $\mathcal{O}_Y(V)$. Nous concluons
que $\mathcal{O}_Y(V) = \mathcal{O}_X(U)$. Cela prouve que $f$ est un isomorphisme
comme souhaité.
\end{proof}

\begin{lemma}
\label{lemma-normalization-trivial}
Soit $X$ un schéma avec localement un nombre fini de composantes irréductibles.
Le morphisme de normalisation $\nu : X^\nu \to X$ est un isomorphisme si et seulement
si $X$ est normal.
\end{lemma}

\begin{proof}
Si $\nu$ est un isomorphisme, alors $X$
est normal d'après le Lemme
\ref{lemma-normalization-normal}. Réciproquement,
supposons que $X$ soit normal. D'après
le Lemme \ref{lemma-normalization-localization} et Propriétés, Lemme \ref{properties-lemma-normal-locally-finite-nr-irreducibles}, nous pouvons
supposer que $X$
est intègre. D'après le Lemme \ref{lemma-normalization-normal}, le morphisme
$\nu$ est intègre et $X^\nu$ a une unique
composante irréductible
(donc c'est un schéma intègre). Nous concluons d'après
les Lemme \ref{lemma-normalization-birational} et \ref{lemma-finite-birational-over-normal}.
\end{proof}

\begin{lemma}
\label{lemma-Japanese-normalization}
Soit $X$ un schéma intègre, japonais.
La normalisation $\nu : X^\nu \to X$ est un morphisme fini.
\end{lemma}

\begin{proof}
Cela découle de la définition
(Propriétés, Définition \ref{properties-definition-nagata}) et du Lemme
\ref{lemma-description-normalization}. À savoir, dans ce cas, le lemme dit que
$\nu^{-1}(\Spec(A))$ est le spectre de la clôture intégrale
de $A$ dans son corps de fractions.
\end{proof}

\begin{lemma}
\label{lemma-nagata-normalization}
Soit $X$ un Schéma de
Nagata. La normalisation $\nu : X^\nu \to X$ est un morphisme fini.
\end{lemma}

\begin{proof}
Notez qu'un schéma de Nagata est localement noethérien,
ainsi la Définition \ref{definition-normalization}
s'applique. Le lemme est maintenant un
cas particulier du Lemme \ref{lemma-nagata-normalization-finite-general} mais nous pouvons
également le prouver directement
comme suit. Écrivons $X^\nu \to X$ comme
la composition $X^\nu \to X_{red} \to X$. Comme $X_{red} \to X$ est une immersion
fermée, elle est finie. Par conséquent, il suffit de
prouver le lemme pour un schéma de Nagata réduit
(d'après le Lemme \ref{lemma-composition-finite}). Soit $\Spec(A) = U \subset X$
un ouvert affine. D'après le Lemme \ref{lemma-description-normalization}
nous avons $\nu^{-1}(U) = \Spec(\prod A_i')$ où $A_i'$ est la clôture
intégrale de $A/\mathfrak q_i$ dans son corps des fractions. Comme $A$
est un anneau de Nagata (voir Propriétés, Lemme \ref{properties-lemma-locally-nagata}) chacune
des extensions d'anneaux
$A/\mathfrak q_i \subset A'_i$ sont finies. Donc $A \to \prod A'_i$ est un homomorphisme d'anneaux
fini et on obtient le résultat.
\end{proof}

\begin{lemma}
\label{lemma-normalization-geom-unibranch-univ-homeo}
Soit $X$ un schéma irréductible, géométriquement unibranche.
Le morphisme de normalisation $\nu : X^\nu \to X$ est un homéomorphisme
universel.
\end{lemma}

\begin{proof}
Nous devons montrer que $\nu$ est intègre, universellement
injectif et surjectif, voir Lemme \ref{lemma-universal-homeomorphism}. D'après
le Lemme \ref{lemma-normalization-normal}, le morphisme $\nu$ est
intègre. Soit $x \in X$ et posons $A = \mathcal{O}_{X, x}$. Puisque
$X$ est irréductible, nous voyons que
$A$ possède un seul idéal premier minimal $\mathfrak p$
et $A_{red} = A/\mathfrak p$. Par le Lemme \ref{lemma-stalk-normalization},
le germe $A' = (\nu_*\mathcal{O}_{X^\nu})_x$ est la clôture intégrale de $A$
dans le corps des fractions de
$A_{red}$. D'après Compléments d'algèbre, Définition \ref{more-algebra-definition-unibranch},
nous voyons que $A'$ possède un seul idéal
premier $\mathfrak m'$ au-dessus de $\mathfrak m_x \subset A$ et $\kappa(\mathfrak m')/\kappa(x)$.
est purement inséparable. Ainsi $\nu$ est bijectif (donc
surjectif) et universellement injectif par le Lemme \ref{lemma-universally-injective}.
\end{proof}







\section{Normalisation faible}
\label{section-weak-normalization}

\noindent
Nous définirons seulement la normalisation faible d'un schéma lorsque celui-ci
a localement un nombre fini de composantes irréductibles ; similaire
au cas de la normalisation.

\begin{lemma}
\label{lemma-relative-weak-normalization-algebra}
Soit $A \to B$ un homomorphisme d'anneaux induisant un morphisme
dominant $\Spec(B) \to \Spec(A)$ de spectres. Il existe une $A$-sous-algèbre
$B' \subset B$ telle que
\begin{enumerate}
\item $\Spec(B') \to \Spec(A)$ est un homéomorphisme universel,
\item étant donnée une factorisation $A \to C \to B$ telle que
$\Spec(C) \to \Spec(A)$ est un homéomorphisme universel, l'image de
$C \to B$ est contenue dans $B'$.
\end{enumerate}
\end{lemma}

\begin{proof}
Nous utiliserons Lemme \ref{lemma-reduction-universal-homeomorphism} sans autre mention.
Considérons le diagramme commutatif
$$
\xymatrix{
B \ar[r] & B_{red} \\
A \ar[u] \ar[r] & A_{red} \ar[u]
}
$$
Pour toute factorisation $A \to C \to B$ de $A \to B$ comme dans (2),
nous voyons que $A_{red} \to C_{red} \to B_{red}$ est une factorisation de
$A_{red} \to B_{red}$ comme dans (2). Il s'ensuit que si le lemme
est vrai pour $A_{red} \to B_{red}$ et produit la $A_{red}$-sous-algèbre $B'_{red} \subset B_{red}$,
alors en posant $B' \subset B$ égal à l'image
réciproque de $B'_{red}$ résout le lemme pour $A \to B$.
Cela nous ramène au cas discuté dans le paragraphe suivant.

\medskip\noindent
Supposons que $A$ et $B$ sont réduits.
Dans ce cas $A \subset B$ par Algèbre, Lemme \ref{algebra-lemma-image-dense-generic-points}.
Soit $A \to C \to B$ une factorisation comme dans
(2). Alors nous pouvons appliquer la Proposition \ref{proposition-universal-homeomorphism} à $A \subset C$
pour voir que chaque élément de $C$ est
contenu dans une extension $A[c_1, \ldots, c_n] \subset C$ telle que pour
$i = 1, \ldots, n$ nous avons
\begin{enumerate}
\item $c_i^2, c_i^3 \in A[c_1, \ldots, c_{i - 1}]$, ou
\item il existe un nombre premier
$p$ avec $pc_i, c_i^p \in A[c_1, \ldots, c_{i - 1}]$.
\end{enumerate}
Ainsi, la propriété (2) est vérifiée si nous définissons $B' \subset B$
comme le sous-ensemble des éléments $b \in B$ qui sont contenus dans
une extension $
A[b_1, \ldots, b_n] \subset B
$ telle
que (*) est vérifiée : pour $i = 1, \ldots, n$ nous avons
\begin{enumerate}
\item $b_i^2, b_i^3 \in A[b_1, \ldots, b_{i - 1}]$, ou
\item il existe un nombre premier
$p$ avec $pb_i, b_i^p \in A[b_1, \ldots, b_{i - 1}]$.
\end{enumerate}
Il n'y a que deux choses à vérifier : (a) $B'$ est une $A$-sous-algèbre,
et (b) $\Spec(B') \to \Spec(A)$ est un homéomorphisme
universel. La partie (a) découle du fait que étant donné
$n \geq 0$ et $b_1, \ldots, b_n \in B$ satisfaisant (*), et $m \geq 0$
et $b'_1, \ldots, b'_m \in B$ satisfaisant
(*), l'entier $n + m$ et $b_1, \ldots, b_n, b'_1, \ldots, b'_m \in B$ satisfait
également (*). Enfin, la partie (b)
tient par la Proposition \ref{proposition-universal-homeomorphism} et notre construction de
$B'$.
\end{proof}

\begin{lemma}
\label{lemma-relative-weak-normalization-localization}
Soit $A \to B$ un homomorphisme d'anneaux induisant un morphisme
dominant $\Spec(B) \to \Spec(A)$ des spectres. La formation de la $A$-sous-algèbre
$B' \subset B$ dans le Lemme \ref{lemma-relative-weak-normalization-algebra} commute avec la
localisation (voir la démonstration pour explication).
\end{lemma}

\begin{proof}
Soit $S \subset A$ un sous-ensemble multiplicatif. Alors $S^{-1}A \to S^{-1}B$
est un homomorphisme d'anneaux qui induit également un morphisme
dominant $\Spec(S^{-1}B) \to \Spec(S^{-1}A)$ (voir Lemmes \ref{lemma-base-change-dominant}
et \ref{lemma-generalizations-lift-flat}). Par conséquent,
le Lemme \ref{lemma-relative-weak-normalization-algebra} produit une $S^{-1}A$-sous-algèbre
$(S^{-1}B)' \subset S^{-1}B$. L'énoncé signifie
que $S^{-1}B' = (S^{-1}B)'$ en tant que $S^{-1}A$-sous-algèbres de
$S^{-1}B$.

\medskip\noindent
Pour voir que cela est vrai, nous utiliserons la construction
de $B'$ et $(S^{-1}B)'$
dans la démonstration du Lemme \ref{lemma-relative-weak-normalization-algebra}. Dans
la première étape, nous voyons que $B'$ est l'image
réciproque de la $A_{red}$-sous-algèbre $B'_{red} \subset B_{red}$ construite
pour l'homomorphisme d'anneaux $A_{red} \to B_{red}$ et de même pour
$(S^{-1}B)'$. En notant que $S^{-1}B_{red} = (S^{-1}B)_{red}$, cela nous ramène
au cas discuté dans le paragraphe suivant.

\medskip\noindent Si
$A$ et $B$ sont réduits, nous avons construit $B'$
comme l'union des sous-algèbres $A[b_1, \ldots, b_n]$ telles que pour $i = 1, \ldots, n$ nous
avons
\begin{enumerate}
\item $b_i^2, b_i^3 \in A[b_1, \ldots, b_{i - 1}]$, ou
\item il existe un nombre premier
$p$ avec $pb_i, b_i^p \in A[b_1, \ldots, b_{i - 1}]$.
\end{enumerate}
De même pour $(S^{-1}B)' \subset S^{-1}B$. Ainsi il est clair que l'image de $B' \to B \to S^{-1}B$
est contenue dans $(S^{-1}B)'$. Pour montrer que
l'application correspondante $S^{-1}B' \to (S^{-1}B)'$ est surjective,
on utilise le Lemme \ref{lemma-special-elements-and-localization} pour éliminer les dénominateurs
successivement ; nous omettons les détails.
\end{proof}

\begin{lemma}
\label{lemma-relative-seminormalization-algebra}
Soit $A \to B$ un homomorphisme d'anneaux induisant un morphisme
dominant $\Spec(B) \to \Spec(A)$ des spectres. Il existe un $A$-sous-algèbre
$B' \subset B$ tel que
\begin{enumerate}
\item $\Spec(B') \to \Spec(A)$ est un homéomorphisme universel induisant des isomorphismes
sur les corps résiduels,
\item donnée une factorisation $A \to C \to B$ telle que $\Spec(C) \to \Spec(A)$ est un homéomorphisme
universel induisant des isomorphismes sur les corps résiduels, l'image
de $C \to B$ est contenue dans $B'$.
\end{enumerate}
\end{lemma}

\begin{proof}
Cette démonstration est exactement
la même que celle du Lemme \ref{lemma-relative-weak-normalization-algebra}
sauf que nous utilisons la Proposition \ref{proposition-universal-homeomorphism-equal-residue-fields}
au lieu de la Proposition \ref{proposition-universal-homeomorphism}
\end{proof}

\begin{lemma}
\label{lemma-relative-seminormalization-localization}
Soit $A \to B$ un homomorphisme d'anneaux induisant un morphisme
dominant $\Spec(B) \to \Spec(A)$ de spectres. La formation de la $A$-sous-algèbre
$B' \subset B$ dans le Lemme \ref{lemma-relative-seminormalization-algebra} commute avec la
localisation (voir la démonstration pour explication).
\end{lemma}

\begin{proof}
La démonstration est la même
que la démonstration du Lemme \ref{lemma-relative-weak-normalization-localization}.
\end{proof}

\begin{lemma}
\label{lemma-relative-seminormalization}
Soit $f : Y \to X$ un morphisme de schémas quasi-compact, quasi-séparé
et dominant.
\begin{enumerate}
\item La catégorie des factorisations $Y \to X' \to X$ où $X' \to X$
est un homéomorphisme universel a un objet initial $Y \to X^{Y/wn} \to X$.
\item La catégorie des factorisations $Y \to X' \to X$ où $X' \to X$ est un
homéomorphisme universel induisant des isomorphismes sur les corps
résiduels a un objet initial $Y \to X^{Y/sn} \to X$.
\end{enumerate}
De plus, la formation de la factorisation $Y \to X^{Y/wn} \to X$ et $Y \to X^{Y/sn} \to X$ commute
avec le changement de base vers des sous-schémas ouverts de $X$.
\end{lemma}

\begin{proof}
Nous allons prouver (1) et omettre la démonstration de (2) ; de plus,
l'assertion finale découlera de la construction de la factorisation.
Nous utiliserons le Lemme \ref{lemma-universal-homeomorphism}
sans mention supplémentaire. Tout d'abord, soit $Y \to X^{Y/n} \to X$ la normalisation
de $X$ dans $Y$, voir Définition \ref{definition-normalization-X-in-Y}.
Pour $Y \to X' \to X$ comme en (1), nous obtenons un morphisme
unique $X^{Y/n} \to X'$ compatible avec les morphismes donnés,
voir Lemme \ref{lemma-characterize-normalization}. Ainsi, il suffit de prouver le
lemme avec $f$ remplacé par $X^{Y/n} \to X$. Cela nous ramène
au cas étudié dans le paragraphe suivant.

\medskip\noindent
Supposons que $f$ soit entier (le reste de la démonstration fonctionne
plus généralement si $f$ est affine). Soit $U = \Spec(A)$
un ouvert affine de $X$ et soit $V = f^{-1}(U) = \Spec(B)$ l'image
réciproque dans $Y$. Alors $A \to B$ est un homomorphisme
d'anneaux qui induit un morphisme dominant sur les spectres. Par le
Lemme \ref{lemma-relative-weak-normalization-algebra}, nous obtenons une $A$-sous-algèbre $B' \subset B$
telle qu'en posant $U^{V/wn} = \Spec(B')$, la factorisation $V \to U^{V/wn} \to U$
est initiale dans la catégorie des factorisations $V \to U' \to U$ où $U' \to U$
est un homéomorphisme universel.

\medskip\noindent
Si $U_1 \subset U_2 \subset X$ sont des ouverts affines, alors en posant $V_i = f^{-1}(U_i)$
nous obtenons un morphisme canonique
$$
\rho_{U_1}^{U_2} : U_1^{V_1/wn} \to U_1 \times_{U_2} U_2^{V_2/wn}
$$
sur $U_1$ par la propriété universelle de $U_1^{V_1/wn}$. Ces morphismes
satisfont une fonctorialité naturelle que nous laissons au lecteur
de formuler et de démontrer. De plus,
le morphisme $\rho_{U_1}^{U_2}$ est
un isomorphisme ; cela résulte du Lemme \ref{lemma-relative-weak-normalization-localization}
à condition que $U_1 \subset U_2$ soit un ouvert
principal et, dans le cas général, peut être
réduit à ce cas par la nature fonctorielle
de ces applications et Schémas, Lemme \ref{schemes-lemma-standard-open-two-affines}
(détails omis). Ainsi, par recollement
relatif (Constructions, Lemme \ref{constructions-lemma-relative-glueing}), nous obtenons
un morphisme $X^{Y/wn} \to X$ qui se restreint
à $U^{V/wn} \to U$ sur $U$ de manière compatible avec le $\rho_{U_1}^{U_2}$.
Bien sûr, les morphismes $V \to U^{V/wn}$ se recollent pour former un
morphisme $Y \to X^{Y/wn}$ (voir
Constructions, Remarque \ref{constructions-remark-relative-glueing-functorial}) et nous obtenons
notre factorisation $Y \to X^{Y/wn} \to X$ où le deuxième morphisme est un
homéomorphisme universel.

\medskip\noindent
Enfin, soit $Y \to X' \to X$ une factorisation comme en (1).
Avec $V \to U^{V/wn} \to U \subset X$ comme ci-dessus, nous obtenons
une factorisation $V \to U \times_X X' \to U$ où la deuxième flèche
est un homéomorphisme universel et nous obtenons
un morphisme unique $g_U : U^{V/wn} \to U \times_X X'$ au-dessus de $U$ par la propriété
universelle de $U^{V/wn}$. Ces
$g_U$ sont compatibles avec les morphismes $\rho_{U_1}^{U_2}$ ; détails
omis. Ainsi il existe un morphisme unique $g : X^{Y/wn} \to X'$ au-dessus
de $X$ concordant avec $g_U$
au-dessus de $U$, voir Constructions, Remarque \ref{constructions-remark-relative-glueing-functorial}.
Cela prouve que $Y \to X^{Y/wn} \to X$ est initial dans notre catégorie
et la démonstration est terminée.
\end{proof}

\begin{definition}
\label{definition-relative-seminormalization}
Soit $f : Y \to X$ un morphisme de schémas quasi-compact, quasi-séparé
et dominant.
\begin{enumerate}
\item La factorisation $Y \to X^{Y/sn} \to X$ construite dans le Lemme \ref{lemma-relative-seminormalization}
partie (2) est la {\it semi-normalisation
de $X$ dans $Y$}.
\item La factorisation $Y \to X^{Y/wn} \to X$ construite dans le Lemme \ref{lemma-relative-seminormalization}
partie (1) est la {\it normalisation
faible de $X$ dans $Y$}.
\end{enumerate}
\end{definition}

\noindent
Voici une façon de réinterpréter la semi-normalisation d'un schéma qui
localement a un nombre fini de composantes irréductibles.

\begin{lemma}
\label{lemma-seminormal-and-normalization}
Soit $X$ un schéma tel que chaque ouvert quasi-compact possède
un nombre fini de composantes irréductibles. Soit $\nu : X^\nu \to X$
la normalisation de $X$. Alors la semi-normalisation de $X$
dans $X^\nu$ est la semi-normalisation de $X$.
En formule : $X^{sn} = X^{X^\nu/sn}$.
\end{lemma}

\begin{proof}
Soit $f : Y \to X$ comme dans (\ref{equation-generic-points}) de sorte que $X^\nu$
soit la normalisation de $X$ dans $Y$.
La semi-normalisation $X^{sn} \to X$ de $X$ est l'objet initial dans la catégorie
des homéomorphismes universels $X' \to X$ induisant des isomorphismes
sur les corps résiduels. Puisque $Y$ est l'union disjointe des spectres
des corps résiduels aux points génériques des composantes irréductibles
de $X$, nous voyons que pour tout $X' \to X$ dans cette catégorie,
nous obtenons un relèvement canonique
$f' : Y \to X'$ de $f$. Alors, par le Lemme
\ref{lemma-characterize-normalization}, nous obtenons un morphisme canonique
$X^\nu \to X'$. D'où, à son tour, un morphisme canonique $X^{X^\nu/sn} \to X'$
par la propriété universelle de $X^{X^\nu/sn}$.
De cette façon, nous voyons que $X^{X^\nu/sn}$ satisfait la même
propriété universelle que $X^{sn}$ possède et nous concluons.
\end{proof}

\noindent
Le lemme \ref{lemma-seminormal-and-normalization} motive la définition suivante. Puisque
nous n'avons construit la normalisation que dans le cas
où $X$ a localement un nombre fini de composantes irréductibles,
nous nous limiterons également à ce cas pour la normalisation faible.

\begin{definition}
\label{definition-weak-normalization}
Soit $X$ un schéma tel que tout ouvert quasi-compact ait un nombre fini
de composantes irréductibles. Nous définissons la {\it
normalisation faible} de $X$ comme la normalisation faible
$$
X^\nu \longrightarrow X^{wn} \longrightarrow X
$$
de $X$ dans la normalisation $X^\nu$
de $X$ (Définition \ref{definition-normalization}).
En formule : $X^{wn} = X^{X^\nu/wn}$.
\end{definition}

\noindent
Combiné avec Lemme \ref{lemma-seminormal-and-normalization}, nous voyons que pour un schéma $X$
qui a localement un nombre fini de composantes irréductibles, il existe
des morphismes canoniques
$$
X^\nu \to X^{wn} \to X^{sn} \to X
$$
Ayant fait cette définition, nous pouvons dire ce que cela signifie pour un schéma d'être
faiblement normal (à condition qu'il ait localement un nombre fini de composantes
irréductibles).

\begin{definition}
\label{definition-weakly-normal}
Soit $X$ un schéma tel que chaque ouvert quasi-compact ait un
nombre fini de composantes irréductibles. Nous disons que $X$ est {\it
faiblement normal} si la normalisation faible
$X^{wn} \to X$ est un isomorphisme (Définition \ref{definition-weak-normalization}).
\end{definition}

\noindent
Il découle immédiatement des définitions que pour un schéma $X$ tel que chaque
ouvert quasi-compact possède un nombre fini de composantes irréductibles, nous
avons
$$
X\text{ normal} \Rightarrow
X\text{ faiblement normal} \Rightarrow
X\text{ semi-normal}
$$
Nous pouvons déterminer le sens de la normalité faible dans le cas affine comme suit.

\begin{lemma}
\label{lemma-affine-weakly-normal}
Soit $X = \Spec(A)$ un schéma affine qui a un nombre fini de composantes irréductibles.
Alors $X$ est faiblement normal si et seulement si
\begin{enumerate}
\item $A$ est semi-normal (Définition \ref{definition-seminormal-ring}),
\item pour un nombre premier $p$ et $z, w \in A$ tels que (a) $z$
est un non-diviseur de zéro, (b) $w^p$ est divisible par $z^p$ et (c) $pw$
est divisible par $z$, alors $w$ est divisible par $z$.
\end{enumerate}
\end{lemma}

\begin{proof}
Supposons que $X$ soit faiblement normal. Puisqu'un schéma faiblement normal
est semi-normal, on voit que (1) est vérifiée (par définition des schémas
faiblement normaux). En particulier, $A$ est réduit.
Soit $p, z, w$ comme dans (2). Choisissez $x, y \in A$ de sorte
que $z^p x = w^p$ et $zy = pw$. Alors
$p^p x = y^p$. L’homomorphisme d’anneaux $A \to C = A[t]/(t^p - x, pt - y)$ induit un homéomorphisme universel
sur les spectres. La normalisation $X^\nu$
de $X$ est le spectre de la clôture
intégrale $A'$ de $A$ dans l’anneau total
des fractions de $A$, voir Lemme \ref{lemma-description-normalization}.
Notons que $a = w/z \in A'$ puisque $a^p = x$.
Ainsi, nous avons un homomorphisme d'$A$-algèbres $A \to C \to A'$
envoyant $t$ sur $a$.
À ce stade, la propriété définissante $X = X^{wn} = X^{X^\nu/wn}$ d'être faiblement
normal nous dit que l'image de $C \to A'$ est contenue dans $A$. Ainsi,
nous trouvons $a \in A$ comme désiré.

\medskip\noindent
Inversement, supposons (1) et (2). Soit $A'$ comme dans le
paragraphe précédent. Nous devons montrer que $X^{X^\nu/wn} = X$. Par
construction dans la démonstration du Lemme \ref{lemma-relative-weak-normalization-algebra}, le
schéma $X^{X^\nu/wn}$ est le spectre du sous-anneau
de $A'$ qui est l'union des sous-anneaux $A[a_1, \ldots, a_n] \subset A'$ tels
que pour $i = 1, \ldots, n$ nous avons
\begin{enumerate}
\item[(a)] $a_i^2, a_i^3 \in A[a_1, \ldots, a_{i - 1}]$, ou
\item[(b)] il existe un nombre premier
$p$ avec $pa_i, a_i^p \in A[a_1, \ldots, a_{i - 1}]$.
\end{enumerate}
Nous pouvons alors utiliser (1) et (2) pour voir de manière inductive que $a_1, \ldots, a_n \in A$ ; nous
omettons les détails. Par conséquent, nous avons $X = X^{X^\nu/wn}$ et donc $X$
est faiblement normal.
\end{proof}

\noindent
Voici le lemme obligatoire.

\begin{lemma}
\label{lemma-locally-weakly-normal}
Soit $X$ un schéma tel que chaque ouvert quasi-compact possède un nombre
fini de composantes irréductibles. Les conditions suivantes sont équivalentes :
\begin{enumerate}
\item Le schéma $X$ est faiblement normal.
\item Pour chaque ouvert affine $U \subset X$, l'anneau $\mathcal{O}_X(U)$
satisfait les conditions (1) et (2 du Lemme \ref{lemma-affine-weakly-normal}.
\item Il existe un recouvrement ouvert affine $X = \bigcup U_i$ tel
que chaque anneau $\mathcal{O}_X(U_i)$
satisfasse aux conditions (1) et (2) du Lemme \ref{lemma-affine-weakly-normal}.
\item Il existe un recouvrement ouvert $X = \bigcup X_j$ tel que chaque
sous-schéma ouvert $X_j$ soit faiblement normal.
\end{enumerate}
De plus, si $X$ est faiblement normal, alors tout sous-schéma ouvert est faiblement normal.
\end{lemma}

\begin{proof}
La condition pour que $X$ soit faiblement normal est que
le morphisme $X^{wn} = X^{X^\nu/wn} \to X$ soit un isomorphisme. Puisque la construction de
$X^\nu \to X$ commute avec le changement de base vers les sous-schémas
ouverts et que la construction de $X^{X^\nu/wn}$ commute avec le
changement de base vers les sous-schémas ouverts de $X$ (Lemme \ref{lemma-relative-seminormalization}), le
lemme est clair.
\end{proof}










\section{Théorème principal de Zariski (version algébrique)}
\label{section-Zariski}

\noindent
Ceci est la version que vous pouvez prouver en utilisant uniquement
des méthodes algébriques. Avant de pouvoir prouver des versions
plus puissantes (pour les morphismes non affines), nous
devons développer plus d'outils. Voir Cohomologie
des
schémas, section \ref{coherent-section-applications-formal-functions}
et Compléments sur les morphismes, section \ref{more-morphisms-section-application-etale-neighbourhoods}.

\begin{theorem}[Version algébrique du Théorème principal de Zariski]
\label{theorem-main-theorem}
Soit $f : Y \to X$ un morphisme affine de schémas. Supposons que $f$
soit de type fini. Soit $X'$
la normalisation de $X$ dans $Y$. Illustration :
$$
\xymatrix{
Y \ar[rd]_f \ar[rr]_{f'} & & X' \ar[ld]^\nu \\
& X &
}
$$
Alors il existe un sous-schéma ouvert $U' \subset X'$ tel que
\begin{enumerate}
\item $(f')^{-1}(U') \to U'$ soit un isomorphisme, et
\item $(f')^{-1}(U') \subset Y$ est l'ensemble des points en lesquels $f$
est quasi-fini.
\end{enumerate}
\end{theorem}

\begin{proof}
Il y a une réduction immédiate au cas où $X$ et donc $Y$
sont affines. Disons $X = \Spec(R)$ et $Y = \Spec(A)$.
Alors $X' = \Spec(A')$, où $A'$ est la clôture
intégrale de $R$ dans $A$, voir Définitions
\ref{definition-integral-closure} et \ref{definition-normalization-X-in-Y}. Par Algèbre,
Théorème \ref{algebra-theorem-main-theorem} pour tout $y \in Y$ en
lequel $f$ est quasi-fini, il existe un ouvert $U'_y \subset X'$
tel que $(f')^{-1}(U'_y) \to U'_y$ soit un isomorphisme. Posons
$U' = \bigcup U'_y$ où $y \in Y$ parcourt tous les points où $f$
est quasi-fini. Il reste à montrer que $f$ est
quasi-fini en tous les points de $(f')^{-1}(U')$. Si
$y \in (f')^{-1}(U')$ avec image $x \in X$, alors nous voyons que
$Y_x \to X'_x$ est un isomorphisme dans un voisinage de $y$. Par conséquent,
il n'existe aucun point de $Y_x$ qui se spécialise en $y$,
puisque cela est vrai pour $f'(y)$ dans $X'_x$, voir Lemme
\ref{lemma-integral-fibres}. D'après le Lemme \ref{lemma-quasi-finite-at-point-characterize} partie (3), cela implique que
$f$ est quasi-fini en $y$.
\end{proof}

\noindent
Nous pouvons utiliser la version algébrique du Théorème principal de Zariski pour montrer
que l'ensemble des points où un morphisme est quasi-fini est ouvert.

\begin{lemma}
\label{lemma-quasi-finite-points-open}
\begin{slogan}
Le lieu localement quasi-fini d'un morphisme est ouvert
\end{slogan}
Soit $f : X \to S$ un morphisme de schémas. L'ensemble
des points de $X$ où $f$ est quasi-fini est un $U \subset X$
ouvert. Le morphisme induit $U \to S$ est localement quasi-fini.
\end{lemma}

\begin{proof}
Supposons que $f$ soit quasi-fini
en $x$. Soient $x \in U = \Spec(A) \subset X$, $V = \Spec(R) \subset S$ des ouverts
affines comme dans la Définition \ref{definition-quasi-finite}. D'après le Théorème
\ref{theorem-main-theorem} ci-dessus ou Algèbre, Lemme \ref{algebra-lemma-quasi-finite-open},
l'ensemble des idéaux premiers $\mathfrak q$ en lesquels
$R \to A$ est quasi-fini est ouvert dans $\Spec(A)$. Puisque
tous correspondent à des points de $X$ où $f$
est quasi-fini, on obtient la première affirmation. La seconde
affirmation est évidente.
\end{proof}

\noindent
Nous améliorerons le lemme suivant aux morphismes
quasi-finis séparés
plus tard,
voir Compléments sur les morphismes, Lemme \ref{more-morphisms-lemma-quasi-finite-separated-pass-through-finite}.

\begin{lemma}
\label{lemma-quasi-finite-affine}
Soit $f : Y \to X$ un morphisme de schémas.
Supposons
\begin{enumerate}
\item $X$ et $Y$ soient affines, et
\item $f$ est quasi-fini.
\end{enumerate}
Alors il existe un diagramme
$$
\xymatrix{
Y \ar[rd]_f \ar[rr]_j & & Z \ar[ld]^\pi \\
& X &
}
$$
avec $Z$ affine, $\pi$ fini et $j$ une immersion ouverte.
\end{lemma}

\begin{proof}
Ceci
est l'Algèbre, Lemme \ref{algebra-lemma-quasi-finite-open-integral-closure} reformulé dans
le langage des schémas.
\end{proof}

\begin{lemma}
\label{lemma-image-nowhere-dense-quasi-finite}
Soit $f : Y \to X$ un morphisme quasi-fini de schémas. Soit $T \subset Y$
une partie fermée nulle part dense de $Y$. Alors $f(T) \subset X$
est une partie nulle part dense de $X$.
\end{lemma}

\begin{proof}
Comme dans la démonstration du Lemme \ref{lemma-image-nowhere-dense-finite}, cela se réduit immédiatement
au cas où la base $X$ est affine. Dans ce cas,
$Y = \bigcup_{i = 1, \ldots, n} Y_i$ est une union finie d'ouverts affines (puisque $f$
est quasi-compact). Comme chaque $T \cap Y_i$ est nulle part dense,
et qu'une union finie de parties nulle part denses est elle-même
nulle part dense (voir
Topologie, Lemme \ref{topology-lemma-nowhere-dense}), il suffit de prouver
que l'image $f(T \cap Y_i)$ est nulle part dense dans $X$. Cela nous ramène
au cas où $X$ et $Y$ sont tous deux affines. À ce stade, nous
appliquons le Lemme \ref{lemma-quasi-finite-affine} ci-dessus pour obtenir un diagramme
$$
\xymatrix{
Y \ar[rd]_f \ar[rr]_j & & Z \ar[ld]^\pi \\
& X &
}
$$
avec $Z$ affine, $\pi$ fini et $j$
une immersion ouverte. Posons $\overline{T} = \overline{j(T)} \subset Z$.
Par la Topologie, Lemme \ref{topology-lemma-image-nowhere-dense-open} nous voyons
que $\overline{T}$ est nulle part dense
dans $Z$. Puisque $f(T) \subset \pi(\overline{T})$
le lemme découle du résultat correspondant dans
le cas fini, voir Lemme \ref{lemma-image-nowhere-dense-finite}.
\end{proof}




\section{Fibres universellement bornées}
\label{section-universally-bounded}

\noindent
Soit $X$ un schéma sur un corps $k$. Si $X$ est fini sur
$k$, alors $X = \Spec(A)$ où $A$ est une $k$-algèbre finie. Une
autre façon de dire cela est que $X$ est fini localement libre sur $\Spec(k)$,
voir Définition \ref{definition-finite-locally-free}. Par conséquent $X \to \Spec(k)$ a un
{\it degré} qui est un entier $d \geq 0$, à savoir $d = \dim_k(A)$.
Nous appelons parfois cela le {\it degré} du schéma (fini) $X$
sur $k$.

\begin{definition}
\label{definition-universally-bounded}
Soit $f : X \to Y$ un morphisme de schémas.
\begin{enumerate}
\item On dit que l'entier $n$ {\it bornes les degrés des fibres
de $f$} si pour
tout $y \in Y$ la fibre $X_y$ est un schéma fini sur $\kappa(y)$
dont le degré sur $\kappa(y)$ est $\leq n$.
\item On dit que les fibres {\it de $f$ sont universellement borné}\footnote{Ceci
est probablement une notation
non standard.} s'il existe un entier $n$ qui bornes les degrés des fibres de
$f$.
\end{enumerate}
\end{definition}

\noindent
Notez qu'en particulier le nombre de points dans une fibre est lui aussi borné par $n$.
(La réciproque n'est pas vraie, même si toutes les fibres sont des schémas réduits
finis.)

\begin{lemma}
\label{lemma-characterize-universally-bounded}
Soit $f : X \to Y$ un morphisme de schémas. Soit $n \geq 0$. Les conditions
suivantes sont équivalentes :
\begin{enumerate}
\item l'entier $n$ borne les degrés des fibres de $f$, et
\item pour tout morphisme $\Spec(k) \to Y$, où $k$ est un corps,
le produit fibré $X_k = \Spec(k) \times_Y X$ est fini sur $k$ de degré
$\leq n$.
\end{enumerate}
Dans ce cas, les fibres de $f$ sont universellement bornées et les schémas
$X_k$ ont au plus $n$ points. Plus
précisément, si $X_k = \{x_1, \ldots, x_t\}$, alors nous avons
$$
n \geq \sum\nolimits_{i = 1, \ldots, t} [\kappa(x_i) : k]
$$
\end{lemma}

\begin{proof}
L'implication (2) $\Rightarrow$ (1) est triviale. L'autre implication
tient parce que si l'image de $\Spec(k) \to Y$ est $y$,
alors $X_k = \Spec(k) \times_{\Spec(\kappa(y))} X_y$. Par définition, les fibres de $f$
étant universellement bornées signifie qu'il existe
un(e) $n$. Enfin, supposons que $X_k = \Spec(A)$.
Alors $\dim_k A = n$. Par conséquent, $A$ est artinien, tous les
idéaux premiers sont des idéaux maximaux $\mathfrak m_i$, et $A$
est le produit des localisations en ces idéaux maximaux.
Voir Algèbre, Lemme \ref{algebra-lemma-finite-dimensional-algebra} et \ref{algebra-lemma-artinian-finite-length}. Alors
$\mathfrak m_i$ correspond à $x_i$,
nous avons $A_{\mathfrak m_i} = \mathcal{O}_{X_k, x_i}$ et donc
il existe une surjection
$A \to  \bigoplus \kappa(\mathfrak m_i) = \bigoplus \kappa(x_i)$ ce qui implique l'inégalité
dans l'énoncé du lemme par algèbre
linéaire.
\end{proof}

\begin{lemma}
\label{lemma-finite-locally-free-universally-bounded}
Si $f$ est un morphisme fini localement libre de degré $d$, alors
$d$ borne le degré des fibres de $f$.
\end{lemma}

\begin{proof}
C'est vrai car tout changement de base de $f$
est fini localement libre de
degré $d$ (Lemme \ref{lemma-base-change-finite-locally-free}) et donc les fibres
de $f$ ont toutes le degré $d$.
\end{proof}

\begin{lemma}
\label{lemma-composition-universally-bounded}
Une composition de morphismes avec des fibres universellement bornées
est un morphisme avec des fibres universellement bornées. Plus
précisément, supposons que $n$ borne les degrés des fibres de $f : X \to Y$
et que $m$ borne les degrés de
$g : Y \to Z$. Alors $nm$ borne les degrés des fibres de $g \circ f : X \to Z$.
\end{lemma}

\begin{proof}
Soient $f : X \to Y$ et $g : Y \to Z$ des fibres universellement bornées. Disons
que $\deg(X_y/\kappa(y)) \leq n$ pour tout $y \in Y$, et que $\deg(Y_z/\kappa(z)) \leq m$ pour tout $z \in Z$.
Soit $z \in Z$ un point. Par hypothèse,
le schéma $Y_z$ est fini sur $\Spec(\kappa(z))$. En particulier,
l'espace topologique sous-jacent
de $Y_z$ est un ensemble fini discret. Les fibres du
morphisme $f_z : X_z \to Y_z$ sont les fibres de $f$ aux
points correspondants de $Y$, qui sont des ensembles finis
discrets d'après le raisonnement ci-dessus. Par conséquent, nous concluons
que l'espace topologique sous-jacent de $X_z$ est
également un ensemble fini discret. Ainsi, $X_z$ est un schéma affine
(c'est un bel exercice ; cela découle également par exemple de
Propriétés, Lemme \ref{properties-lemma-maximal-points-affine} appliqué à l'ensemble
de tous les points de $X_z$). Écrire $X_z = \Spec(A)$, $Y_z = \Spec(B)$,
et $k = \kappa(z)$. Ensuite $k \to B \to A$ et
nous savons que (a) $\dim_k(B) \leq m$, et (b) pour chaque idéal
maximal $\mathfrak m \subset B$ nous avons
$\dim_{\kappa(\mathfrak m)}(A/\mathfrak mA) \leq n$. Nous affirmons que cela implique
que $\dim_k(A) \leq nm$. Notez que $B$
est le produit de ses localisations $B_{\mathfrak m}$, par exemple
parce que $Y_z$ est une union disjointe de schémas
à points $1$, ou par l'Algèbre, Lemmes \ref{algebra-lemma-finite-dimensional-algebra}
et \ref{algebra-lemma-artinian-finite-length}. Ainsi, nous
voyons que
$\dim_k(B) = \sum_{\mathfrak m} \dim_k(B_{\mathfrak m})$ et $\dim_k(A) = \sum_{\mathfrak m} \dim_k(A_{\mathfrak m})$, où dans les deux cas
$\mathfrak m$ parcourt les idéaux maximaux de $B$
(et non de $A$). D'après ce
qui précède, et le lemme de Nakayama (Algèbre,
Lemme \ref{algebra-lemma-NAK}), nous
voyons que chaque $A_{\mathfrak m}$ est un quotient
de $B_{\mathfrak m}^{\oplus n}$ en tant que $B_{\mathfrak m}$-module.
Par conséquent $\dim_k(A_{\mathfrak m}) \leq n \dim_k(B_{\mathfrak m})$. En rassemblant le
tout, nous voyons que
$$
\dim_k(A) = \sum\nolimits_{\mathfrak m} \dim_ta	(A_{\mathfrak m})
\leq \sum\nolimits_{\mathfrak m} n \dim_k(B_{\mathfrak m})
= n \dim_k(B) \leq nm
$$
comme souhaité.
\end{proof}

\begin{lemma}
\label{lemma-base-change-universally-bounded}
Un changement de base d'un morphisme à fibres universellement bornées
est un morphisme à fibres universellement bornées. Plus précisément,
si $n$ borne les degrés des fibres de $f : X \to Y$ et que $Y' \to Y$
est un morphisme quelconque, alors les degrés des fibres du changement de
base $f' : Y' \times_Y X \to Y'$ sont également bornés par $n$.
\end{lemma}

\begin{proof}
Ceci est clair d'après
le résultat du Lemme \ref{lemma-characterize-universally-bounded}.
\end{proof}

\begin{lemma}
\label{lemma-descent-universally-bounded}
Soit $f : X \to Y$ un morphisme de schémas. Soit $Y' \to Y$
un morphisme de schémas, et soit $f' : X' = X_{Y'} \to Y'$ le
changement de base de $f$. Si $Y' \to Y$ est surjectif
et $f'$ a des fibres universellement bornées, alors $f$ a des
fibres universellement bornées. Plus précisément, si $n$ borne le degré
des fibres de $f'$, alors $n$ borne également les degrés des
fibres de $f$.
\end{lemma}

\begin{proof}
Soit $n \geq 0$ un entier bornant les degrés des fibres de $f'$. Nous affirmons que
$n$ convient également pour $f$. C’est-à-dire, si $y \in Y$ est un point,
choisissez un point $y' \in Y'$ situé au-dessus de $y$ et observez que
$$
X'_{y'} = \Spec(\kappa(y')) \times_{\Spec(\kappa(y))} X_y.
$$
Puisque $X'_{y'}$ est supposé fini de degré $\leq n$ sur $\kappa(y')$, il s'ensuit
que également $X_y$ est fini de degré $\leq n$ sur $\kappa(y)$. (Certains
détails sont omis.)
\end{proof}

\begin{lemma}
\label{lemma-immersion-universally-bounded}
Une immersion a des fibres universellement bornées.
\end{lemma}

\begin{proof}
L'entier $n = 1$ fonctionne dans la définition.
\end{proof}

\begin{lemma}
\label{lemma-etale-universally-bounded}
Soit $f : X \to Y$ un morphisme étale de schémas. Soit $n \geq 0$.
Les conditions suivantes sont équivalentes
\begin{enumerate}
\item l'entier $n$ borne les degrés des fibres,
\item pour tout corps $k$ et morphisme $\Spec(k) \to Y$ le
changement de base $X_k = \Spec(k) \times_Y X$ a au plus $n$ points, et
\item pour chaque $y \in Y$ et chaque clôture algébrique
séparable $\kappa(y) \subset \kappa(y)^{sep}$, le schéma $X_{\kappa(y)^{sep}}$
a au plus $n$ points.
\end{enumerate}
\end{lemma}

\begin{proof}
Cela résulte du
Lemme \ref{lemma-characterize-universally-bounded} et du fait que les fibres
$X_y$ sont des unions disjointes de spectres d'extensions
de corps finies et séparables de $\kappa(y)$,
voir Lemme \ref{lemma-etale-over-field}.
\end{proof}

\noindent
Avoir des fibres universellement bornées est une notion absolue et non une notion relative.
C'est pourquoi la condition dans le lemme suivant est que $X$ est quasi-compact,
et non que $f$ est quasi-compact.

\begin{lemma}
\label{lemma-locally-quasi-finite-qc-source-universally-bounded}
Soit $f : X \to Y$ un morphisme de schémas. Supposons
que
\begin{enumerate}
\item $f$ est localement quasi-fini, et
\item $X$ est quasi-compact.
\end{enumerate}
Alors $f$ a des fibres universellement bornées.
\end{lemma}

\begin{proof}
Puisque $X$ est quasi-compact, il existe un recouvrement
ouvert affine fini $X = \bigcup_{i = 1, \ldots, n} U_i$ et des ouverts affines $V_i \subset Y$,
$i = 1, \ldots, n$ tels que $f(U_i) \subset V_i$. En raison
de la nature locale de ``quasi-finitude
locale'' (voir le Lemme \ref{lemma-quasi-finite-at-point-characterize} partie (4)), nous
voyons que les morphismes $f|_{U_i} : U_i \to V_i$ sont localement
des morphismes quasi-finis d'affines, donc quasi-finis,
voir le Lemme \ref{lemma-quasi-finite-locally-quasi-compact}. Pour $y \in Y$,
il est clair que $X_y = \bigcup_{y \in V_i} (U_i)_y$ est un recouvrement ouvert. Il
suffit donc de prouver le lemme pour un morphisme
quasi-fini d'affines (à savoir, si $n_i$ fonctionne
pour le morphisme $f|_{U_i} : U_i \to V_i$, alors $\sum n_i$)
fonctionne pour $f$).

\medskip\noindent Supposons
que $f : X \to Y$ soit un morphisme quasi-fini de schémas affines.
D'après le Lemme \ref{lemma-quasi-finite-affine}, nous pouvons
trouver un diagramme
$$
\xymatrix{
X \ar[rd]_f \ar[rr]_j & & Z \ar[ld]^\pi \\
& Y &
}
$$
avec $Z$ affine, $\pi$ fini et $j$ une immersion
ouverte. Puisque $j$
a des fibres universellement bornées (Lemme
\ref{lemma-immersion-universally-bounded}), cela nous ramène à montrer que $\pi$ a
des fibres universellement bornées (Lemme \ref{lemma-composition-universally-bounded}).

\medskip\noindent Cela
nous ramène à un morphisme de la forme $\Spec(B) \to \Spec(A)$ où $A \to B$ est fini.
Supposons que $B$ soit engendré par $x_1, \ldots, x_n$ comme $A$-module,
de sorte que
$$
A^{\oplus n} \longrightarrow B, \quad
(a_1, \ldots, a_n) \longmapsto \sum a_i x_i
$$
est une application de module $A$ surjective. Pour tout premier $\mathfrak p \subset A$, cela
induit une application surjective entre les espaces vectoriels $\kappa(\mathfrak p)$
$$
\kappa(\mathfrak p)^{\oplus n} \longrightarrow B \otimes_A \kappa(\mathfrak p)
$$
En d'autres termes, l'entier $n$ intervient dans la définition
d'un morphisme à fibres universellement bornées.
\end{proof}

\begin{lemma}
\label{lemma-universally-bounded-permanence}
Considérons un diagramme commutatif de morphismes de schémas
$$
\xymatrix{
X \ar[rd]_g \ar[rr]_f & & Y \ar[ld]^h \\
& Z &
}
$$
Si $g$ a des fibres universellement bornées, et que $f$ est surjectif et plat,
alors $h$ a également des fibres universellement bornées. Plus précisément, si $n$
borne le degré des fibres de $g$, alors $n$ borne également
le degré des fibres de $h$.
\end{lemma}

\begin{proof}
Supposons que $g$ a des fibres universellement bornées, et que $f$
est surjectif et plat. Disons que le degré des fibres de $g$
est borné par $n \in \mathbf{N}$. Nous affirmons
que $n$ fonctionne également pour $h$. Soit
$z \in Z$. Considérons le morphisme de schémas $X_z \to Y_z$. Il est
plat et surjectif. Par hypothèse $X_z$ est un schéma
fini sur $\kappa(z)$,
en particulier il est le spectre d'un anneau artinien (par
Algèbre, Lemme \ref{algebra-lemma-finite-dimensional-algebra}). Par le Lemme \ref{lemma-Artinian-affine} le morphisme $X_z \to Y_z$
est affine, en particulier quasi-compact. Il
en découle du Lemme \ref{lemma-fpqc-quotient-topology} que $Y_z$
est un discret fini car cela vaut pour $X_z$.
Par conséquent, $Y_z$ est un schéma affine (c'est un bel exercice ; cela
découle également,
par exemple, de Propriétés, Lemme \ref{properties-lemma-maximal-points-affine} appliqué à
l'ensemble de tous les points de $Y_z$).
Écrivez $Y_z = \Spec(B)$ et $X_z = \Spec(A)$. Alors $A$
est fidèlement plat sur $B$, donc $B \subset A$.
Par conséquent, $\dim_k(B) \leq \dim_k(A) \leq n$ comme souhaité.
\end{proof}





\section{Divers}
\label{section-misc}

\noindent
Résultats qui ne trouvent pas leur place ailleurs.

\begin{lemma}
\label{lemma-factor-through-point}
Soit $f : Y \to X$ un morphisme de schémas. Soit $x \in X$ un point. Supposons
que $Y$ soit réduit et que $f(Y)$ soit contenu au sens d'ensemble
dans $\{x\}$. Alors $f$ se factorise par le morphisme
canonique $x = \Spec(\kappa(x)) \to X$.
\end{lemma}

\begin{proof}
Omis. Indices : en travaillant localement sur les schémas affines, on se ramène à
un lemme d'algèbre commutative. Étant donné un homomorphisme d'anneaux $A \to B$
avec $B$ réduit tel qu'il existe un unique idéal premier
$\mathfrak p \subset A$ dans l'image de $\Spec(B) \to \Spec(A)$, alors $A \to B$ se factorise par $\kappa(\mathfrak p)$.
C'est un bel exercice.
\end{proof}

\begin{lemma}
\label{lemma-factor-through-set-of-points}
Soit $f : Y \to X$ un morphisme de schémas. Soit $E \subset X$. On suppose que
$X$ est localement noethérien, qu'il n'y a pas de spécialisations non triviales
parmi les éléments de $E$, que $Y$ est réduit, et
$f(Y) \subset E$. Alors $f$ se factorise par $\coprod_{x \in E} x \to X$.
\end{lemma}

\begin{proof}
Lorsque $E$ est un singleton, cela
résulte du Lemme \ref{lemma-factor-through-point}. Si $E$ est fini, alors $E$
(avec la topologie induite de $X$) est un espace discret
fini selon notre hypothèse sur les spécialisations. Ainsi,
ce cas se réduit au cas du singleton. En général, il
y a une réduction au cas où $X$ et $Y$ sont des schémas affines.
Disons que $f : Y \to X$ correspond à l'homomorphisme d'anneaux $\varphi : A \to B$.
On note $A' \subset B$ l'image de $\varphi$. Soit $E' \subset \Spec(A') \subset \Spec(A)$ l'ensemble des
idéaux premiers minimaux de $A'$. Par l'Algèbre, Lemme \ref{algebra-lemma-injective-minimal-primes-in-image}, l'ensemble
$E'$ est contenu dans l'image
de $\Spec(B) \to \Spec(A') \subset \Spec(A)$. Nous concluons que $E' \subset E$. Comme $A'$ est noethérien,
nous avons que $E'$ est fini par l'Algèbre, Lemme
\ref{algebra-lemma-Noetherian-irreducible-components}. Puisque tout autre point de l'image de $\Spec(B) \to \Spec(A)$
est une spécialisation d'un élément de $E'$ et dans $E$,
nous concluons que l'image est contenue dans $E'$ (par
notre hypothèse sur les spécialisations entre les points de $E$).
Ainsi, nous réduisons au cas où $E$ est fini,
ce que nous avons traité ci-dessus.
\end{proof}




\begin{multicols}{2}[\section{Autres chapitres}]
\noindent
Préliminaires
\begin{enumerate}
\item \hyperref[introduction-section-phantom]{Introduction}
\item \hyperref[conventions-section-phantom]{Conventions}
\item \hyperref[sets-section-phantom]{Théorie des ensembles}
\item \hyperref[categories-section-phantom]{Catégories}
\item \hyperref[topology-section-phantom]{Topologie}
\item \hyperref[sheaves-section-phantom]{Faisceaux sur les espaces}
\item \hyperref[sites-section-phantom]{Sites et faisceaux}
\item \hyperref[stacks-section-phantom]{Champs}
\item \hyperref[fields-section-phantom]{Corps}
\item \hyperref[algebra-section-phantom]{Algèbre commutative}
\item \hyperref[brauer-section-phantom]{Groupes de Brauer}
\item \hyperref[homology-section-phantom]{Algèbre homologique}
\item \hyperref[derived-section-phantom]{Catégories dérivées}
\item \hyperref[simplicial-section-phantom]{Méthodes simpliciales}
\item \hyperref[more-algebra-section-phantom]{Compléments d'algèbre}
\item \hyperref[smoothing-section-phantom]{Lissage des morphismes d'anneaux}
\item \hyperref[modules-section-phantom]{Faisceaux de modules}
\item \hyperref[sites-modules-section-phantom]{Modules sur les sites}
\item \hyperref[injectives-section-phantom]{Objets injectifs}
\item \hyperref[cohomology-section-phantom]{Cohomologie des faisceaux}
\item \hyperref[sites-cohomology-section-phantom]{Cohomologie sur les sites}
\item \hyperref[dga-section-phantom]{Algèbre différentielle graduée}
\item \hyperref[dpa-section-phantom]{Algèbre à puissances divisées}
\item \hyperref[sdga-section-phantom]{Faisceaux différentiels gradués}
\item \hyperref[hypercovering-section-phantom]{Hyperrecouvrements}
\end{enumerate}
Schémas
\begin{enumerate}
\setcounter{enumi}{25}
\item \hyperref[schemes-section-phantom]{Schémas}
\item \hyperref[constructions-section-phantom]{Constructions de schémas}
\item \hyperref[properties-section-phantom]{Propriétés des schémas}
\item \hyperref[morphisms-section-phantom]{Morphismes de schémas}
\item \hyperref[coherent-section-phantom]{Cohomologie des schémas}
\item \hyperref[divisors-section-phantom]{Diviseurs}
\item \hyperref[limits-section-phantom]{Limites de schémas}
\item \hyperref[varieties-section-phantom]{Variétés}
\item \hyperref[topologies-section-phantom]{Topologies sur les schémas}
\item \hyperref[descent-section-phantom]{Descente}
\item \hyperref[perfect-section-phantom]{Catégories dérivées de schémas}
\item \hyperref[more-morphisms-section-phantom]{Compléments sur les morphismes}
\item \hyperref[flat-section-phantom]{Compléments sur la platitude}
\item \hyperref[groupoids-section-phantom]{Schémas en groupoïdes}
\item \hyperref[more-groupoids-section-phantom]{Compléments sur les schémas en groupoïdes}
\item \hyperref[etale-section-phantom]{Morphismes étales de schémas}
\end{enumerate}
Thèmes de théorie des schémas
\begin{enumerate}
\setcounter{enumi}{41}
\item \hyperref[chow-section-phantom]{Homologie de Chow}
\item \hyperref[intersection-section-phantom]{Théorie de l'intersection}
\item \hyperref[pic-section-phantom]{Schémas de Picard des courbes}
\item \hyperref[weil-section-phantom]{Théories cohomologiques de Weil}
\item \hyperref[adequate-section-phantom]{Modules adéquats}
\item \hyperref[dualizing-section-phantom]{Complexes dualisants}
\item \hyperref[duality-section-phantom]{Dualité pour les schémas}
\item \hyperref[discriminant-section-phantom]{Discriminants et différentes}
\item \hyperref[derham-section-phantom]{Cohomologie de de Rham}
\item \hyperref[local-cohomology-section-phantom]{Cohomologie locale}
\item \hyperref[algebraization-section-phantom]{Géométrie algébrique et formelle}
\item \hyperref[curves-section-phantom]{Courbes algébriques}
\item \hyperref[resolve-section-phantom]{Résolution des surfaces}
\item \hyperref[models-section-phantom]{Réduction semi-stable}
\item \hyperref[functors-section-phantom]{Foncteurs et morphismes}
\item \hyperref[equiv-section-phantom]{Catégories dérivées de variétés}
\item \hyperref[pione-section-phantom]{Groupes fondamentaux de schémas}
\item \hyperref[etale-cohomology-section-phantom]{Cohomologie étale}
\item \hyperref[crystalline-section-phantom]{Cohomologie cristalline}
\item \hyperref[proetale-section-phantom]{Cohomologie pro-étale}
\item \hyperref[relative-cycles-section-phantom]{Cycles relatifs}
\item \hyperref[more-etale-section-phantom]{Compléments de cohomologie étale}
\item \hyperref[trace-section-phantom]{La formule des traces}
\end{enumerate}
Espaces algébriques
\begin{enumerate}
\setcounter{enumi}{64}
\item \hyperref[spaces-section-phantom]{Espaces algébriques}
\item \hyperref[spaces-properties-section-phantom]{Propriétés des espaces algébriques}
\item \hyperref[spaces-morphisms-section-phantom]{Morphismes d'espaces algébriques}
\item \hyperref[decent-spaces-section-phantom]{Espaces algébriques décents}
\item \hyperref[spaces-cohomology-section-phantom]{Cohomologie des espaces algébriques}
\item \hyperref[spaces-limits-section-phantom]{Limites d'espaces algébriques}
\item \hyperref[spaces-divisors-section-phantom]{Diviseurs sur les espaces algébriques}
\item \hyperref[spaces-over-fields-section-phantom]{Espaces algébriques sur des corps}
\item \hyperref[spaces-topologies-section-phantom]{Topologies sur les espaces algébriques}
\item \hyperref[spaces-descent-section-phantom]{Descente et espaces algébriques}
\item \hyperref[spaces-perfect-section-phantom]{Catégories dérivées d'espaces}
\item \hyperref[spaces-more-morphisms-section-phantom]{Compléments sur les morphismes d'espaces}
\item \hyperref[spaces-flat-section-phantom]{Platitude sur les espaces algébriques}
\item \hyperref[spaces-groupoids-section-phantom]{Groupoïdes dans les espaces algébriques}
\item \hyperref[spaces-more-groupoids-section-phantom]{Compléments sur les groupoïdes dans les espaces}
\item \hyperref[bootstrap-section-phantom]{Amorçage}
\item \hyperref[spaces-pushouts-section-phantom]{Sommes amalgamées d'espaces algébriques}
\end{enumerate}
Thèmes de géométrie
\begin{enumerate}
\setcounter{enumi}{81}
\item \hyperref[spaces-chow-section-phantom]{Groupes de Chow des espaces}
\item \hyperref[groupoids-quotients-section-phantom]{Quotients de groupoïdes}
\item \hyperref[spaces-more-cohomology-section-phantom]{Compléments sur la cohomologie des espaces}
\item \hyperref[spaces-simplicial-section-phantom]{Espaces simpliciaux}
\item \hyperref[spaces-duality-section-phantom]{Dualité pour les espaces}
\item \hyperref[formal-spaces-section-phantom]{Espaces algébriques formels}
\item \hyperref[restricted-section-phantom]{Algébrisation des espaces formels}
\item \hyperref[spaces-resolve-section-phantom]{Résolution des surfaces revisitée}
\end{enumerate}
Théorie des déformations
\begin{enumerate}
\setcounter{enumi}{89}
\item \hyperref[formal-defos-section-phantom]{Théorie formelle des déformations}
\item \hyperref[defos-section-phantom]{Théorie des déformations}
\item \hyperref[cotangent-section-phantom]{Le complexe cotangent}
\item \hyperref[examples-defos-section-phantom]{Problèmes de déformation}
\end{enumerate}
Champs algébriques
\begin{enumerate}
\setcounter{enumi}{93}
\item \hyperref[algebraic-section-phantom]{Champs algébriques}
\item \hyperref[examples-stacks-section-phantom]{Exemples de champs}
\item \hyperref[stacks-sheaves-section-phantom]{Faisceaux sur les champs algébriques}
\item \hyperref[criteria-section-phantom]{Critères de représentabilité}
\item \hyperref[artin-section-phantom]{Axiomes d'Artin}
\item \hyperref[quot-section-phantom]{Espaces de Quot et de Hilbert}
\item \hyperref[stacks-properties-section-phantom]{Propriétés des champs algébriques}
\item \hyperref[stacks-morphisms-section-phantom]{Morphismes de champs algébriques}
\item \hyperref[stacks-limits-section-phantom]{Limites de champs algébriques}
\item \hyperref[stacks-cohomology-section-phantom]{Cohomologie des champs algébriques}
\item \hyperref[stacks-perfect-section-phantom]{Catégories dérivées de champs}
\item \hyperref[stacks-introduction-section-phantom]{Introduction aux champs algébriques}
\item \hyperref[stacks-more-morphisms-section-phantom]{Compléments sur les morphismes de champs}
\item \hyperref[stacks-geometry-section-phantom]{La géométrie des champs}
\end{enumerate}
Thèmes de théorie des espaces de modules
\begin{enumerate}
\setcounter{enumi}{107}
\item \hyperref[moduli-section-phantom]{Champs de modules}
\item \hyperref[moduli-curves-section-phantom]{Espaces de modules de courbes}
\end{enumerate}
Divers
\begin{enumerate}
\setcounter{enumi}{109}
\item \hyperref[examples-section-phantom]{Exemples}
\item \hyperref[exercises-section-phantom]{Exercices}
\item \hyperref[guide-section-phantom]{Guide bibliographique}
\item \hyperref[desirables-section-phantom]{Desiderata}
\item \hyperref[coding-section-phantom]{Style de programmation}
\item \hyperref[obsolete-section-phantom]{Obsolète}
\item \hyperref[fdl-section-phantom]{Licence de documentation libre GNU}
\item \hyperref[index-section-phantom]{Index généré automatiquement}
\end{enumerate}
\end{multicols}


\bibliography{my}
\bibliographystyle{amsalpha}

\end{document}
